\documentclass[11pt]{article}
\usepackage[a4paper,margin=26mm]{geometry}
\usepackage[T1]{fontenc}
\usepackage{lmodern,amsmath,amssymb,amsthm,mathtools}
\usepackage{booktabs,longtable,array,microtype}
\usepackage{needspace,etoolbox}
\usepackage[round,authoryear]{natbib}
\usepackage[hidelinks]{hyperref}
\newcommand{\workbenchurl}{https://github.com/KokunoYumeto/zeta-and-l-functions-workbench}
\newcommand{\recordurl}{https://github.com/KokunoYumeto/mathematics-workbenches}
\newcommand{\Q}{\mathbb Q}
\newcommand{\F}{\mathbb F}
\newcommand{\Z}{\mathbb Z}
\newcommand{\C}{\mathbb C}
\newcommand{\ord}{\operatorname{ord}}
\newcommand{\Gal}{\operatorname{Gal}}
\newcommand{\Spec}{\operatorname{Spec}}
\newcommand{\RH}{\operatorname{RH}}
\newcommand{\poly}{\operatorname{poly}}
\newtheorem{theorem}{Theorem}[section]
\newtheorem{proposition}[theorem]{Proposition}
\newtheorem{lemma}[theorem]{Lemma}
\newtheorem{corollary}[theorem]{Corollary}
\theoremstyle{definition}
\newtheorem{definition}[theorem]{Definition}
\newtheorem{example}[theorem]{Example}
\theoremstyle{remark}
\newtheorem{remark}[theorem]{Remark}
\BeforeBeginEnvironment{theorem}{\Needspace{8\baselineskip}}
\BeforeBeginEnvironment{proposition}{\Needspace{8\baselineskip}}
\BeforeBeginEnvironment{lemma}{\Needspace{8\baselineskip}}
\BeforeBeginEnvironment{corollary}{\Needspace{8\baselineskip}}
\BeforeBeginEnvironment{example}{\Needspace{8\baselineskip}}
\lefthyphenmin=3
\righthyphenmin=3
\setlength{\emergencystretch}{2em}
\title{Deterministic finite-field constructions\\
after Alp\"oge, Bach, von zur Gathen,\\
Lenstra, Mumford, and Pila}
\author{}
\date{23 November 2026}
\begin{document}
\maketitle
\begin{abstract}
Bach, von zur Gathen and Lenstra constructed nonresidues from
smooth cyclotomic values, and Pila used Frobenius polynomials of
abelian varieties to find roots of unity. We prove an unconditional
deterministic construction of \(\mathbb F_{p^r}\) with bit cost
polynomial in the odd prime \(r\), \(\log p\), and the part of
\(\Phi_{\operatorname{ord}_r(p)}(p)\) coprime to \(2r\).
Combining Pila's method with Mumford's theta equations and granting
Alp\"oge's Dirichlet GRH also gives a bound polynomial in
\(\log p\) for every fixed prime degree and fixed valuation
\(v_r(p^{\operatorname{ord}_r(p)}-1)\). Exact CM height bounds,
algebraic recognition, marked good reduction and arithmetic over the
full residue algebra supply the geometric construction. The first
bound is uniform in its displayed numerical parameters; the second
has an exponent depending on the fixed degree and valuation. A
bound polynomial uniformly in \(r\) and \(\log p\) remains
open by these methods. We also trace eighteen classical algorithmic
statements to their exact Riemann hypotheses, prove the pure-field
zeta equivalence and Lenstra's Teichm\"uller reduction, and separate
Tang's polynomial-time interfaces from their zero-free inputs.
Character-product and quadratic-value estimates give short auxiliary
searches outside \(O(\sqrt P\log\log P)\) primes at fixed
degree. The mathematics and text were produced by GPT-6 Astra in
ChatGPT Pro mode; the cofactor and fixed-parameter constructions, principal
reductions and average theorems received statement-alone re-proofs
by the same model, earlier proofs received targeted checks by
Claude Opus 5.5 at effort high, and no human has checked this paper.
\end{abstract}

\section{Statements, credit, and scope}
\label{sec:main-statements}
This paper separates the Riemann hypotheses used by deterministic
number-theoretic algorithms and makes their finite-field interfaces explicit.
Write $A$ for the assertion that every nontrivial zero of every Dirichlet
$L$-function has real part $1/2$, including the Riemann zeta function.
We take Alp\"oge's theorem in precisely this form
\citep[p.~1]{Alpoge2026}.  Statements described as consequences of $A$
require no additional Riemann hypothesis.  Alp\"oge's announced theorem
is an input to this paper.  The distinction between this
input and RH for Hecke $L$-functions over other number fields is essential.

All time bounds count bit operations.  A presented field means either
$\F_p[X]/(h)$ with $h$ irreducible, or a basis and its multiplication
table.  These representations are interconvertible in polynomial time
\citep[Theorem~(1.1), pp.~329--330]{Lenstra1991}.
Constructing a presentation from an integer promised to be a prime power
is a different input problem from factoring a polynomial whose coefficient
field is already presented.  Throughout, $\log$ denotes the natural logarithm.


For a positive integer \(n\), \(\Phi_n\) is its cyclotomic polynomial,
and \(v_\ell(n)\) is the exponent of the prime \(\ell\) in \(n\).
Two complementary constructions address the remaining inputs.
Bach, von zur Gathen and Lenstra's nonresidue algorithm uses
smoothness of a cyclotomic value and primitive roots for all its
prime divisors \citep[Theorem~3, author p.~3; Section~5,
pp.~17--20]{BGL2001}. Projecting their separating elements to
one odd-primary part gives the following explicit cofactor bound.
The rational-function argument and odd-prime commutator are theirs;
the primary projection and descent over the whole formal algebra
are proved in Section~\ref{sec:small-cofactor-lifting}.

\begin{theorem}\label{cor:two-primary-cyclotomic-cofactor}
Let \(p\ne r\) be primes, with \(r\) odd, and put
\[
 d=\operatorname{ord}_r(p),\qquad
 a=v_r(p^d-1),\qquad
 M_2=\frac{\Phi_d(p)}{2^{v_2(\Phi_d(p))}r^a}.
\]
There is an unconditional deterministic algorithm constructing
a monic irreducible polynomial of degree \(r\) over
\(\mathbb F_p\) in bit time polynomial in \(r,M_2,\log p\),
with one absolute exponent. It also constructs, in the formal
algebra \(\mathcal D=\mathbb F_p[Z]/\Phi_r(Z)\), an element
\(U\) satisfying
\[
 U^{r^{a-1}}=Z,\qquad
 \sigma_c(U)=U^{\omega_a(c)}
           \quad(c\in(\mathbb Z/r\mathbb Z)^\times),
\]
where \(\sigma_c(Z)=Z^c\) and \(\omega_a(c)\) is the tame
lift modulo \(r^a\): it is congruent to \(c\) modulo \(r\)
and has order dividing \(r-1\). Every component of \(U\) has exact
order \(r^a\). No factorization of \(M_2\), presentation of
\(\mathbb F_{p^d}\), or roots of unity for prime divisors of
\(M_2\) are inputs.
\end{theorem}

For a fixed constant \(C\), the promise
\(M_2\leq(r+1+\log p)^C\) therefore gives polynomial time
in \(r,\log p\), with no Riemann hypothesis and with no bound
on the valuation \(a\). In particular, for cubic extensions it
is the size of \((p+1)/(2^{v_2(p+1)}3^{v_3(p+1)})\) that enters this bound.
The formal algebra is part of the input-free construction: one need
not first choose a factor of its cyclotomic polynomial.

The second result settles every fixed-parameter stratum of the
remaining field-construction problem. Pila used Frobenius polynomials
of abelian varieties to find roots of unity of fixed prime order in
prime fields \citep[Theorem~D, author p.~2; Section~5,
pp.~14--18]{Pila1990}. We combine that idea with Mumford's theta
equations \citep[Section~6, pp.~83--84]{Mumford1967} and
Lenstra's cyclotomic descent \citep[Theorems~(3.1), (5.1),
and (5.2), pp.~334--340]{Lenstra1991}. The root now has a
prescribed prime-power order and lies in an explicit finite algebra.
The prime degree and its valuation are fixed parameters; the
characteristic varies without exception.

\begin{theorem}\label{thm:fixed-wild-cm}
Assume \(\mathrm A\), Dirichlet GRH over \(\mathbb Q\).
Fix an odd prime \(r\) and an integer \(b\geq2\). Put
\[
 g=r^{b-1},\qquad K_b=\mathbb Q(\zeta_{r^b}),\qquad
 T_r=\mathbb Q(\zeta_r),
\]
and let \(R_b\) be the unique degree-\(g\) subfield of \(K_b\).
Fix an integral primitive element \(\theta\) of \(R_b\), with
minimal polynomial \(f\in\mathbb Z[X]\).
There is a deterministic algorithm, uniform in the prime \(p\),
with bit cost \(C_{r,b}(1+\log p)^{c_{r,b}}\), which applies when
\[
 p\ne r,\qquad d=\operatorname{ord}_r(p),\qquad
                         v_r(p^d-1)\geq b.
\]
It computes a unital \(\mathbb F_p\)-algebra homomorphism
\(\mathcal O_{R_b}/p\mathcal O_{R_b}\longrightarrow\mathbb F_p\),
and an element
\[
 U\in\mathcal D_p:=\mathbb F_p[Z]/(\Phi_r(Z))
\]
having order exactly \(r^b\) in every field component of
\(\mathcal D_p\), with \(U^{r^{b-1}}=Z\).
Moreover, for each
\(c\in(\mathbb Z/r\mathbb Z)^\times\),
\begin{equation}
 \sigma_c(U)=U^{\omega_b(c)},\qquad \sigma_c(Z)=Z^c,
 \label{eq:fixed-wild-equivariance}
\end{equation}
where \(\omega_b(c)\) is the unique lift modulo \(r^b\) of
\(c\) whose order divides \(r-1\).
The constants are effective for each fixed pair \((r,b)\).
If \(v_r(p^d-1)=b\), each component of \(U\) generates the
entire \(r\)-primary subgroup of its multiplicative group.
The algorithm does not require a factor of \(\Phi_r\) over
\(\mathbb F_p\).
\end{theorem}

\begin{corollary}\label{cor:fixed-wild-prime-degree-field}
Assume \(\mathrm A\), and fix an odd prime \(r\) and
\(b\geq2\). There is a deterministic algorithm which, for
every prime \(p\ne r\) satisfying
\[
                  v_r(p^{\operatorname{ord}_r(p)}-1)=b,
\]
constructs a monic irreducible polynomial of degree \(r\)
over \(\mathbb F_p\), and hence a presentation of
\(\mathbb F_{p^r}\), in bit time
\(C_{r,b}(1+\log p)^{c_{r,b}}\).
\end{corollary}

In words, fixing the prime degree and the valuation removes the
exceptional characteristics from the construction. This answers that
stratum of Lenstra's finite-field construction problem using only
the rational Dirichlet Riemann hypotheses. The earlier small-order
cases are covered by Theorem~\ref{thm:nonresidue} and
Theorem~\ref{thm:teich}. A bound polynomial in \(r\) and
\(\log p\), with an exponent independent of the valuation,
is not proved here. The explicit CM dimension is \(r^{b-1}\),
and its theta coordinates grow with that dimension.

The proof chooses an imaginary quadratic field with conductor
\(O(\log^2p)\) in which \(p\) splits. Its compositum with
the fixed real wild cyclotomic field has a primitive CM type and an
explicit principal polarization. Exact theta coordinates provide
equations and a complete addition atlas. Reducing the marked model
over all residue fields simultaneously lets Pila's algorithm compute
a Frobenius polynomial without selecting a residue-field factor.
The translated CM supports separate the real embeddings, and the
cyclotomic tensor product turns one such embedding into the required
equivariant root. The proof and its dimension-dependent bit bound
are in Section~\ref{sec:fixed-wild-cm}. The CM algebraicity and
reciprocity theorems, the Colmez and Pazuki height comparisons, and
deterministic algebraic-number recognition carry the intervening
steps; their hypotheses are stated at the point of use.

The smallest hard case gives a particularly short output polynomial.
It is useful both as an example of the general result and as an
explicit construction in any presented quadratic field.

\begin{theorem}[A cubic construction after Mumford and Pila]\label{thm:dirichlet-ninth-root-cubic}
Assume \(\mathrm A\), rational Dirichlet GRH. There is a uniform
deterministic algorithm which, on input a prime
\(p\equiv-1\pmod9\), computes:
\begin{enumerate}
\item a root \(u\in\mathbb F_p\) of
\(f(X)=X^3-3X+1\);
\item an element \(z\) of exact order \(9\) in
\(\mathbb F_p[W]/(W^2+3)\), or in any given presentation of
\(\mathbb F_{p^2}\).
\end{enumerate}
Its deterministic bit cost is \(O((1+\log p)^{C_*})\), for an
absolute effectively computable constant \(C_*\), independent of
\(p\). If \(v_3(p+1)=2\), it also outputs the explicit monic
irreducible polynomial
\[
                 H(Y)=Y^3-3Y-u\in\mathbb F_p[Y],
\]
and hence a presentation of \(\mathbb F_{p^3}\), within the same
bound. The only zero hypotheses used are for Dirichlet \(L\)-functions.
The auxiliary-prime step uses conductors dividing \(4p\). One may
use the sharper CM-height bound with conductors at most \(9q\),
where \(q=O((1+\log p)^2)\), or the unconditional polynomial
height bound of Proposition~\ref{prop:unconditional-cm-recognition}.
\end{theorem}

Thus a prime in the two residue classes \(8,17\pmod{27}\) admits a
deterministic cubic-extension construction from \(A\) alone. These
include \(p=17\), the smallest example emphasized below. The theorem
settles this branch of Lenstra's finite-field construction problem.
The preceding fixed-parameter theorem covers the other fixed
valuations and prime degrees. These are complexity theorems with
large fixed exponents.

For \(r=3\), the general theorem gives every fixed three-adic
valuation. The explicit trace polynomial and the distinction between
a nonmaximal and a maximal root of unity are developed in
Section~\ref{sec:kummer-higher-three-powers}.

\begingroup
\small
\setlength{\tabcolsep}{4pt}
\renewcommand{\arraystretch}{1.18}
\begin{longtable}{@{}p{.15\textwidth}p{.29\textwidth}p{.28\textwidth}p{.22\textwidth}@{}}
\caption{Classical consequences, with the analytic family and the strength of the conclusion kept explicit. A statement marked ``from $A$'' retains any displayed algebraic input promises.}\label{tab:classical}\\
\toprule
Problem & Result and algorithmic meaning & Required zero statement & Source and status \\
\midrule
\endfirsthead
\toprule
Problem & Result and algorithmic meaning & Required zero statement & Source and status \\
\midrule
\endhead
\bottomrule
\endfoot
Miller test & All bases below $2(\log n)^2$ suffice; deterministic polynomial bit time. A perfect-power precheck permits the short proof given below. & Quadratic Dirichlet characters suffice. & From $A$. \citet[Thm.~2.2, pp.~105--106]{Schoof2008}; \citet[Thm.~2, p.~301]{Miller1976}. \\
Prime-field nonresidue & Given prime $r\mid p-1$, find an $r$th nonresidue by a scan of $O(\log^2p)$ candidates. & Dirichlet characters of order $r$ modulo $p$. & From $A$, uniformly in $r$. \citet[pp.~1--2]{LLS2015}. \\
Primitive root with factored $p-1$ & Construct a generator of $\F_p^*$ in polynomial time from the complete factorization of $p-1$. & The preceding prime-order Dirichlet characters. & From $A$ with the given factorization; the construction below proves the transfer. \\
Least primitive root & $g(p)\ll t^4(1+\log t)^4\log^2p\ll\log^6p$, $t=\omega(p-1)$. This gives a short candidate list. & Rational Dirichlet characters modulo $p$. & From $A$. \citet[Thm.~3, author p.~4; Prop.~9, p.~10; proof p.~13]{Shoup1992}. \\
Solovay--Strassen & One Euler--Jacobi trial has composite-input error at most $1/2$; independent repetitions amplify it. & None for the randomized test. A separate derandomization uses a Dirichlet character annihilating its proper liar subgroup. & Randomized result already unconditional. \citet[pp.~84--85]{SolovayStrassen1977}; correction p.~118. \\
Order from factored exponent & Given a factored multiple $M$ of an element's order, repeated prime removal computes the order. & None. & Already unconditional; neither $M$ nor its factorization is provided by this statement. \\
Integer-modulus order & Order computation and integer factoring are polynomial-time Turing equivalent under the stated RH. & Rational Dirichlet $L$-functions in the small-nonresidue step. & Equivalence from $A$; no deterministic polynomial-time factoring algorithm follows. \citet[Thm.~4, pp.~311--312; appendix pp.~315--316]{Miller1976}. \\
Artin primitive roots & Hooley's positive-density asymptotic for a fixed integer $a\ne-1$ that is not a square. & $\RH(\zeta_{\Q(\zeta_k,a^{1/k})})$ for squarefree $k$; these fields are generally nonabelian. & The known proof retains this family beyond $A$. \citet[translation pp.~40--45, 47]{Hooley1967}. \\
Least prime in a progression & For $q>3$, $P(a,q)\le(\varphi(q)\log q)^2$. This is a numerical bound, not polynomial enumeration in $\log q$. & Rational Dirichlet characters modulo $q$. & From $A$. \citet[Cor.~1.2, corrected author p.~3]{LLS2015}. \\
Chebyshev limiting distribution & Normalized progression errors have a limiting logarithmic distribution against bounded continuous test functions. & Rational Dirichlet characters modulo $q$. & From $A$. \citet[Thm.~1.1, p.~174]{RubinsteinSarnak1994}. \\
Strict prime races & The standard strict ordering densities, symmetry formulas, and numerical biases require additional control of zero ordinates. & Dirichlet RH plus the grand simplicity hypothesis, a rational independence condition. & $A$ does not supply that independence. \citet[pp.~176--177]{RubinsteinSarnak1994}. \\
Ideal-class generators & Prime ideals of norm at most $12(\log|D_K|)^2$ generate $\mathrm{Cl}(K)$; this alone supplies no relation-search runtime. & Ideal-class Hecke $L_K(s,\chi)$, including $\zeta_K$ in Bach's argument. RH for the Hilbert class field's zeta is sufficient. & The quoted sharp bound retains Hecke RH. CT separately gives exponent $32$ in Proposition~\ref{prop:ray-class-strip}. \citet[author pp.~2--3]{BDF2008} identify Bach's zero sums. \\
Class-group mixing in relation search & The mixing lemma used in the Hafner--McCurley runtime analysis gives a hitting probability for a class-group walk. & Ideal-class Hecke characters of an imaginary quadratic field, including the trivial character in the spectral estimate. & $A$ covers $\zeta_K$; generally it leaves nontrivial class characters. \citet[Prop.~1, author pp.~8--9]{BiasseErukulangara2023}; \citet[Cor.~1.3, proof p.~9]{JMV2009}. \\
Schoof point counting & The original theorem gives deterministic $O(\log^9q)$ point counting for presented fields of characteristic other than $2,3$. & No number-field RH; the finite-field estimate used is proved. & Already unconditional. \citet[pp.~483, 485--490]{Schoof1985}. \\
ECPP & Certificate correctness is unconditional. The cited practical algorithm's runtime analysis is explicitly heuristic. & None for verification; prime and smoothness distributions enter the heuristic analysis. & No general proved runtime follows just by supplying $A$. \citet[Thm.~5.2, p.~121; pp.~122, 124--125]{Schoof2008}. \\
Isogeny expanders & Spectral gap and rapid mixing on ray-class Cayley graphs; the elliptic application concerns ordinary curves at the relevant endomorphism level. & Finite-order Hecke characters of the narrow ray class group over $K$; imaginary quadratic $K$ in the elliptic application. & $K=\Q$ follows from $A$; the elliptic case generally retains Hecke RH. \citet[Thms.~1.1, 1.5--1.6; proof p.~9]{JMV2009}. \\
CM/CRT split-prime selection & Equation (21) and Lemma 2 control rational primes splitting completely in the ring class field. Only this stage is classified. & RH for the ring class field's Dedekind zeta function. & Generally remains conditional. \citet[corrected author pp.~21--22]{Sutherland2011}. \\
Hilbert Nullstellensatz & Complex feasibility belongs to $\mathrm{AM}$ under the number-field prime-density input; this is not a deterministic-time conclusion. & $\zeta_K$ for a primitive-element field containing one algebraic solution; no normal closure is needed for this estimate. & The A-only transfer leaves this input. CT's fixed strip separately gives AM; Proposition~\ref{prop:fixed-strip-hn-am}. \citet[Thm.~9, p.~282]{Koiran1996}, with the erratum \citep{KoiranErratum1997}. \\
\end{longtable}
\endgroup


\begin{theorem}[The pure-field hypothesis for the classical construction]
\label{thm:kummer}
Let $p$ and $r$ be distinct primes, with $r$ odd, and put
$F=\Q(p^{1/r})$, $K=\Q(\zeta_r)$, and $E=K(p^{1/r})$.
Under $A$, $\RH(\zeta_F)$ and $\RH(\zeta_E)$ are equivalent.
Either assertion implies a prime $\ell\ne p,r$ such that
\[
 \ell\equiv1\pmod r,\qquad p^{(\ell-1)/r}\not\equiv1\pmod\ell,
\]
with $\ell\ll r^2\log^2(rp)$. A completely explicit, less sharp
bound is
\begin{equation}\label{eq:bound}
 \ell\le \left(4r(2r-3)\log r+4(r-1)^2\log p
             +\frac52r(r-1)+5\right)^2.
\end{equation}
Given such a prime, one can construct $\F_{p^r}$ in deterministic time
polynomial in $\ell$ and $\log p$.  Thus the displayed analytic assumption
gives time polynomial in $r$ and $\log p$.
\end{theorem}

The construction is Adleman and Lenstra's Gaussian-period construction
\citep[Algorithm A, pp.~350--351]{AdlemanLenstra1986}; the numerical
prime bound uses Bach and Sorenson \citep[Theorem~5.1, p.~1729]{BachSorenson1996}.
The sharper order of magnitude is the classical least-nonsplit-prime
bound of Murty, as recorded by \citet[p.~1, equation~(1.2)]{Zaman2018};
Theorem~\ref{thm:fixedstrip} proves the required specialization.
What is made explicit here is the single pure-field zeta function remaining
after the abelian factors have been covered by $A$.
This is a sufficient hypothesis for this construction, not a necessary
condition for every algorithm constructing the same field.

\begin{theorem}[Lenstra's Teichm\"uller construction, in prime degree]
\label{thm:teich}
Let $p\ne r$ be primes, with $r$ odd, and set
\[
 s=\ord_r(p),\quad a=v_r(p^s-1),\quad R=r^a,
 \quad D=\F_p[T]/\Phi_r(T),\quad \zeta=T\bmod\Phi_r.
\]
For $b\in\F_r^\times$, let $\rho_b$ be the automorphism
$\rho_b(\zeta)=\zeta^b$, and let $\omega_a(b)$ denote the unique
lift of $b$ modulo $r^a$ whose order divides $r-1$.
The following two inputs are interconvertible, without an analytic
assumption, in deterministic time polynomial in $r$ and $\log p$:
\begin{enumerate}
\item a presentation of $\F_{p^r}$;
\item an element $\eta\in D^\times$ of order $R$ such that
$\rho_b(\eta)=\eta^{\omega_a(b)}$ for every $b\in\F_r^\times$.
\end{enumerate}
If a presentation of $\F_{p^m}$ is given and $r\mid m$, it also
determines such an $\eta$ in time polynomial in $m$ and $\log p$.
\end{theorem}

This is a prime-degree form of Lenstra's constructions
\citep[Theorems~(3.1), (5.1), and (5.2), pp.~334--340]{Lenstra1991}.
The point of writing the reductions out is to identify the input which
is absent when only $p$ and $r$ are known.  Reading a nonresidue from an
already constructed field does not construct that field from its order.

\begin{theorem}[Nonresidues in a presented extension]
\label{thm:nonresidue}
Let $p$ be prime, let $r$ be a prime divisor of $p-1$, and let a
presentation of $\F_{p^m}$ be given.  Under $A$, an $r$th nonresidue
can be found deterministically in time polynomial in $m$ and $\log p$.
For $r=2$ and even $m$ no analytic assumption is needed.
\end{theorem}

Here an $r$th nonresidue means a nonzero element outside
$(\F_{p^m}^{\times})^r$.  The prime-field input is the small-character
nonresidue theorem.  The extension step is linear algebra in Frobenius
eigenspaces, as in Bhargava, Ivanyos, Mittal, and Saxena's work.
The proof below includes a nonvanishing choice of eigenspace vector.

\begin{theorem}[The two polynomial-time interfaces after Tang]
\label{thm:tangfront}
Assume the algebraic prime-field theorem and bit bound of OpenAI
\citep[Theorem~1.2, p.~3; Proposition~10.1, pp.~41--44]{OpenAI2026},
and the uniform Dedekind zero-free theorem of Cureton and Thorner
\citep[Theorem~1.1, p.~3]{CT2026}.
There are deterministic algorithms of polynomial bit complexity for
both of the following tasks:
\begin{enumerate}
\item completely factor a nonzero dense polynomial over a presented
finite field, including its multiplicities;
\item from a binary integer $q>1$ promised to be a prime power,
construct a power-basis presentation of the field of order $q$.
\end{enumerate}
The first interface reduces a degree-$n$ polynomial to at most
$\max(0,n-1)$ split squarefree prime-field queries of degree at most $n$.
For $q=p^m$, the second uses $O(m^2(1+\log_2p))$ prime-field
queries of degree at most $m$ and polynomial additional bit cost.
\end{theorem}

These algorithms and the polynomial conclusion are Tang's
\citep[Theorem~1.1 and Propositions~3.1, 4.1, pp.~2, 4--7]{Tang2026}.
Section~\ref{sec:tang} gives the interface reductions in full and
traces each auxiliary prime to just two Dedekind zeta functions.
The cited external theorems state unconditional results; accepting
them therefore settles the output problem without any unproved RH
assumption. Their proof uses a different analytic input and does not give
a derivation from Dirichlet-GRH estimates alone.  The question pursued
here is a derivation with $A$ as the sole new analytic input, alongside
the classical algebraic reductions.  Since the separate unconditional
zero-free theorem also proves the auxiliary-prime bound, no assertion
of logical independence or of unproved existence is intended.

\begin{theorem}[A finite family for quasipolynomial factorization]
\label{thm:quasipoly}
Let a field $\F_{p^m}$ be presented, and let
$0\ne f\in\F_{p^m}[X]$ have degree $d\geq2$.
If $p\leq d$, its complete factorization can be
computed deterministically in polynomial time in
$d$ and $\log(p^m)$.  If $p>d$, put
\[
 \mathcal B=\{r\leq d: r\text{ prime},\quad
 r\nmid(p-1)m,\quad v_r(p^{\ord_r(p)}-1)>1\}.
\]
Under $A$ and RH for $\zeta_{\Q(p^{1/r})}$ for
$r\in\mathcal B$, complete factorization has
deterministic bit complexity
$(d^{\log d}\log(p^m))^{O(1)}$.
Only primes encountered by the algorithm need their
additional zeta hypothesis.  If every prime $r\leq d$
either divides $m$ or satisfies
$v_r(p^{\ord_r(p)}-1)=1$, the latter bound holds
without an analytic assumption.
\end{theorem}

This finite family refines the earlier quasipolynomial algorithm; it is
not the analytic input to Tang's polynomial theorem.

The next average statements use character orthogonality over the varying input prime. Their lower-search comparisons specialize Montgomery's construction, presented in \citet[p.~432, Exercise~9(a)--(e)]{MontgomeryVaughan2007}. That source treats large least quadratic nonresidues; here the characters encode the auxiliary-prime and Tate-cubic searches, with their exact input restrictions.

\begin{theorem}[Short auxiliary primes for almost all inputs]\label{thm:fixed-r-almost-all-auxiliary}
Assume A and fix an odd prime $r$.  Let
\[
 C_r>\frac{r-1}{2\log r},\qquad
 B=C_r\log P\log\log P.
\]
For all sufficiently large $P$, depending on $r$ and $C_r$, all but
\[
 O_{r,C_r}\!\left(\sqrt P\log\log P\right)
\]
primes $p\in(P/2,P]$ admit a prime $\ell\le B$, with $\ell\ne p,r$, such that $\ell\equiv1\pmod r$ and $p^{(\ell-1)/r}\not\equiv1\pmod\ell$.
In particular this bounds the exceptions among inputs with $r\nmid p-1$ and $r^2\mid p^{\operatorname{ord}_r(p)}-1$.
For $r=3$ one may take $C_3=1$; the assertion remains asymptotic,
and no explicit numerical threshold for $P$ is asserted here.

On the successful inputs, deterministic enumeration and certification
of such an $\ell$, followed by the established auxiliary-prime
reduction, constructs $\mathbb F_{p^r}$ in time polynomial in
$\log p$ for the fixed $r$ and $C_r$.  An algorithm stopping at the displayed bound
can report failure on the exceptional inputs; the theorem is not an
all-input construction.
\end{theorem}
The averaging is over input primes. It makes the classical construction polynomial on a density-one set of inputs of each fixed degree, while leaving the prescribed exceptional input unresolved. The proof uses only rational Dirichlet characters.

For the cubic family, a second classical ingredient is the quadratic-value product estimate in \citet[Lemmas~1, 3, and 4 and \S2.4, pp.~1131--1134]{Cilleruelo2011}. The specialization below turns its large prime divisors into independent cyclic-cubic conductor primes. It gives the following explicit output bound for all prime congruence classes.

\begin{theorem}[A logarithmic cubic search for almost all inputs]
\label{thm:almost-all-small-tate}
Assume A.  Fix $C>1/\log3$ and, for sufficiently large $P$, put
\[
 K=\left\lfloor C\log P\right\rfloor.
\]
All but $O_C(\sqrt P\log\log P)$ primes $p\in(P/2,P]$
admit an integer $1\le k\le K$, $k\not\equiv3\pmod p$,
for which
\[
 Y^3-(k+3)Y^2+kY+1
\]
is irreducible over $\mathbb F_p$.  In particular, the same
exception bound holds among the hard cubic inputs
$p\equiv-1\pmod9$.

Testing these $O_C(\log p)$ cubics deterministically constructs
$\mathbb F_{p^3}$ in time polynomial in $\log p$ on the
successful inputs.  The search can report failure on the
exceptional inputs; the assertion is not an all-input parameter
bound.
\end{theorem}

\begin{theorem}[A simultaneous auxiliary table for almost all characteristics]\label{thm:aa-12r}
Assume \(A\), and fix \(a_0>0\). There is an absolute constant \(C_0\)
such that the following holds for all sufficiently large \(P\)
(the threshold may depend on \(a_0\)).
Let \(1\leq D\leq(\log P)^{a_0}\) be an integer.
Except for
\begin{equation}\label{eq:aa-12r-exceptions}
 O\!\left(
       \frac{\sqrt P\,D^3\log^2(PD)}{\log P}
    \right)
\end{equation}
primes \(p\in(P/2,P]\), every prime \(r\leq D\) simultaneously has
a rational prime \(\ell\) satisfying
\begin{equation}\label{eq:aa-12r-aux}
 \ell\leq C_0r^2\log^2(Pr),\qquad
 \ell\equiv1\pmod{12r},\qquad
 p^{(\ell-1)/r}\not\equiv1\pmod\ell.
\end{equation}
All such primes automatically satisfy
\(\ell\notin\{2,3,p,r\}\).
The constants in the \(O\)-bound are absolute; only the lower
threshold for \(P\) depends on \(a_0\).
\end{theorem}

\begin{corollary}[Average polynomial-time interfaces after Tang]\label{cor:aa-finite-field-algorithms}
Assume \(A\) and the purely algebraic auxiliary-prime algorithm
of \citet[Theorem~1.2, p.~3; Proposition~10.1, pp.~41--44]{OpenAI2026}.
For every prime outside the exceptional set in
Theorem~\ref{thm:aa-12r}, a single uniform deterministic procedure
has the following guarantees.
\begin{enumerate}
 \item Over any explicitly presented finite extension of characteristic
 \(p\), it completely factors every nonzero dense polynomial of degree
 \(n\leq D\), with leading coefficient and multiplicities, in time
 polynomial in \(n+1\) and the bit length of the field order.
 \item Given an integer promised to equal \(q=p^m\), with
 \(1\leq m\leq D\), it constructs a presentation of the field of order
 \(q\) in time polynomial in \(\log q\).
\end{enumerate}
The polynomial exponents are absolute and the algorithms do not
receive \(P\) or \(D\) as input. Their running time uses the input degree, not the global degree cutoff.
\end{corollary}

\paragraph{Why the proofs work.}
The normal closure of $X^r-p$ has affine permutation group
$\F_r\rtimes\F_r^\times$.  Its regular representation separates into
an abelian part and copies of one nonabelian representation.  Artin
formalism then identifies the one extra $L$-function.  The finite-field
construction has a parallel algebraic separation: the cyclotomic
coefficient algebra is explicit, while a maximal $r$-primary unit is
the missing Kummer input.  Frobenius eigenspaces recover that unit when
a field of degree divisible by $r$ is already known.

\section{Classical consequences and their exact hypotheses}
\label{sec:classical-consequences}

Table~\ref{tab:classical} distinguishes a zero-location hypothesis from the
algorithmic conclusion to which it contributes. Three distinctions recur.
A short guaranteed search interval is not automatically a polynomial-time
search in the bit length of the modulus. A set of class-group generators is
not a presentation of that group. And a randomized or interactive complexity
bound remains randomized or interactive when its analytic hypothesis is
proved. The elementary arguments below make the positive transfers explicit;
the remaining paragraphs identify the precise analytic families that survive.

\subsection{Quadratic characters suffice for Miller's cutoff}

The quadratic-character formulation is already explicit in
\citet[Theorem~2.2, pp.~105--106]{Schoof2008}; it is not a new strengthening
claimed here. Handle $n\le16$ by trial division; for larger $n$ the cutoff
is less than $n$, so it does not test a prime against its own zero residue.
The following proof gives a short transfer with an elementary
perfect-power precheck. Suppose $n$ is odd, composite, and not a perfect power,
and write $n-1=2^s d$ with $d$ odd. It has two distinct prime divisors.
For every prime $p\mid n$, write $p-1=2^{t_p}u_p$ with $u_p$ odd.
If a unit $a$ passes the strong probable-prime test, either $a^d=1$ at every
prime divisor of $n$, or there is one $j<s$ for which
$a^{2^j d}=-1$ at every such prime. In the first case the order of $a$ at
each prime has odd order. In the second, its $2$-primary order is exactly
$2^{j+1}$ at every prime.

If the $t_p$ are not all equal, choose $q\mid n$ with
$t_q>\min_{p\mid n}t_p$. Every passing unit then has
$\left(\frac aq\right)=1$. If the $t_p$ are all equal, choose two distinct
prime divisors $p,q$. Their Legendre symbols agree for every passing unit,
so every such unit lies in the kernel of
$\chi(a)=\left(\frac ap\right)\left(\frac aq\right)$.
Thus in either case all passing units lie in the kernel of a nonprincipal
primitive quadratic character of conductor $f\le n$.
The small-character-witness theorem supplies
$a\le2(\log f)^2\le2(\log n)^2$ with $\chi(a)\ne1$.
If $a$ is not a unit modulo $n$, its gcd detects compositeness; otherwise
it fails the strong test. This proves the cutoff after the precheck.

The distinction between $\chi(a)\ne1$ and $\chi(a)=-1$ matters for the
constant $2$: a small nonunit is permitted by the former statement and is
handled by the gcd test. The published formulation without a perfect-power
precheck first excludes repeated prime factors using a smooth-number count,
then applies quadratic characters \citep[pp.~105--106]{Schoof2008}.
All these operations take polynomial bit time. This conclusion concerns
Miller's particular test; primality already has an unconditional deterministic
polynomial-time algorithm, also described in \citet[Section~3]{Schoof2008}.

\subsection{Nonresidues, primitive roots, and order information}

For a given prime $r\mid p-1$, the kernel of a character of order $r$
modulo $p$ is precisely the subgroup of $r$th powers. The Dirichlet
small-witness bound therefore gives an $r$th nonresidue among
$O(\log^2p)$ small integers. A candidate is certified by
$a^{(p-1)/r}\ne1$. The finite number of small primes for which the cutoff
reaches $p$ can be treated directly; for $p\ge17$ it lies below $p$.
This uses rational Dirichlet $L$-functions, uniformly in the variable $r$
\citep[pp.~1--2]{LLS2015}.

Suppose the full factorization $p-1=\prod_i r_i^{e_i}$ is given.
Choose an $r_i$th nonresidue $a_i$ and put
\[
 b_i=a_i^{(p-1)/r_i^{e_i}}.
\]
The $r_i$-part of the order of $a_i$ is $r_i^{e_i}$, so $b_i$ has exact
order $r_i^{e_i}$. Their product has order $p-1$ because these orders are
pairwise coprime. This constructs a primitive root in polynomial time and
does not factor $p-1$.

The least-root theorem gives a different kind of information. Shoup proves
$g(p)\ll t^4(1+\log t)^4\log^2p$, where $t=\omega(p-1)$.
His proof invokes a weighted sum of rational Dirichlet characters, not
number-field Hecke characters \citep[Theorem~3, author p.~4;
Proposition~9, p.~10; Section~6, p.~13]{Shoup1992}.
Since $t!\le p-1$, we have $t(1+\log t)=O(\log p)$ and hence
$g(p)=O(\log^6p)$. A short interval containing a generator is useful,
but identifying a generator requires an order certificate or an adequate
order-computation procedure.

When a factored multiple $M$ of an element's order is given, order
computation is elementary and unconditional. Start with $d=M$ and remove
a prime factor $\ell$ whenever $a^{d/\ell}=1$. The invariant
$\ord(a)\mid d$ persists; at termination no prime factor can divide
$d/\ord(a)$, so equality holds. In contrast, Miller's general theorem
states a polynomial-time Turing equivalence between integer factorization
and rational periods, hence integer-modulus order, under the rational
Dirichlet RH used in his appendix
\citep[Theorem~4, pp.~311--312; pp.~315--316]{Miller1976}.
Discharging that hypothesis establishes the equivalence, not a
polynomial-time algorithm for either problem.

\subsection{Randomized primality tests and their separate derandomization}

Solovay--Strassen is unconditional as a randomized algorithm
\citep{SolovayStrassen1977}. For completeness, let $n$ be odd and
$e=(n-1)/2$. The Euler--Jacobi liars in $(\Z/n\Z)^*$ form the kernel of
the homomorphism $a\mapsto a^e\left(\frac an\right)^{-1}$.
This subgroup is proper. If $p^v\Vert n$ with $v\ge2$, CRT gives a unit
that is $1+p^{v-1}$ modulo $p^v$ and $1$ at the other prime-power factors.
Its Jacobi symbol is $1$, whereas its $e$th power is not $1$ modulo $p^v$.
If $n$ is squarefree composite, choose a nonresidue at one prime divisor
and $1$ at every other prime divisor. Its Jacobi symbol is $-1$, but its
$e$th power is $1$ at another prime divisor. At most half the units,
and therefore at most half the candidate bases, pass. Independent trials
give error at most $2^{-k}$ after $k$ repetitions.
The Dirichlet small-witness theorem can separately make this a
deterministic small-base test by detecting an element outside the proper
liar subgroup. It is this latter statement, not the original randomized
error estimate, that consumes a Riemann hypothesis.

\subsection{Distribution theorems: progression primes, Artin, and races}

The corrected explicit bound
$P(a,q)\le(\varphi(q)\log q)^2$ for $q>3$ uses only rational Dirichlet
characters and therefore follows from $A$
\citep[Corollary~1.2, corrected author p.~3]{LLS2015}.
Enumerating an interval of length comparable to $q^2$ is nevertheless
exponential in the binary length of $q$; this explicit theorem is not by
itself a deterministic polynomial-time prime-search algorithm.

Hooley's proof of the positive-density version of Artin's primitive-root
conjecture uses a different family. For a fixed integer $a\ne-1$ that is not a square, it counts
complete splitting in the Kummer fields
$\Q(\zeta_k,a^{1/k})$, for squarefree $k$, under RH for their Dedekind
zeta functions \citep[translation pp.~40--45, 47]{Hooley1967}.
These fields are generally nonabelian over $\Q$. Factoring out the abelian
pieces covered by $A$ does not remove the remaining pieces from this proof.
This identifies the surviving hypothesis of that proof; it is not a claim
of logical independence from $A$, nor does it conflate an infinitude theorem
with Hooley's full asymptotic.

Rubinstein--Sarnak's first limiting-distribution theorem does follow from
Dirichlet RH \citep[Theorem~1.1, p.~174]{RubinsteinSarnak1994}.
The normalized prime-count errors have a limiting measure for bounded
continuous test functions. A strict race is an indicator of an ordering
region, so possible mass on its boundary must still be controlled.
Their grand simplicity hypothesis is a separate rational-independence
condition on the ordinates of Dirichlet zeros. It enters the explicit
characteristic function and the familiar strictly ordered race densities
\citep[pp.~176--177]{RubinsteinSarnak1994}.
Proving the zeros' real parts are $1/2$ does not prove that independence.

\paragraph{The separate 2026 infinitude result.}
Theorem~1.1 of \citet[p.~2]{OpenAIArtin2026} states that, for each
integer \(a\) that is neither \(-1\) nor a square, there are constants
\(c_a>0\) and \(x_a\) such that every \(x\geq x_a\) satisfies
\[
 \#\{\ell\text{ prime}:x<\ell<2x,\ \ell\nmid a,\
                    \operatorname{ord}_\ell(a)=\ell-1\}
 \geq c_a\,\frac{x}{(\log x)^2}.
\]
This is an infinitude theorem with a quantitative dyadic-interval
lower bound. Hooley's conclusion is the stronger full asymptotic
\(\pi_a(x)\sim\delta_a\operatorname{Li}(x)\), with the predicted
positive Artin density \(\delta_a\)
\citep[translation p.~47]{Hooley1967}.
The displayed \(x/(\log x)^2\) lower bound does not itself imply that
asymptotic or a positive density among the primes.
The 2026 manuscript also states the separate zero-free strip
\(\Re s>1-10^{-6}\) for finite-order Hecke \(L\)-functions over
cyclotomic fields containing \(\mu_{12}\), uniformly in conductor
and height \citep[Theorem~1.2, p.~3]{OpenAIArtin2026}.
Its analytic input is therefore distinct from estimates for rational
Dirichlet \(L\)-functions alone.

\subsection{Class groups and elliptic algorithms}

The bound $12(\log|D_K|)^2$ for prime-ideal generators of
$\mathrm{Cl}(K)$ is often cited under the label GRH. Its analytic input
includes ideal-class Hecke characters over $K$ and $\zeta_K$.
The primary refinement by \citet[author pp.~2--3]{BDF2008} identifies
these two zero sums explicitly; their own field-specific bound removes
the need to assume RH for the trivial-character zero sum, but still
requires the nontrivial ideal-class characters.
The factorization of the Hilbert class field's Dedekind zeta function
into these Hecke functions gives a convenient sufficient single-field
package. Even when $K/\Q$ is quadratic, its Hilbert class field need
not be abelian over $\Q$. Nor does a small generating set, by itself,
bound the time needed to find relations or units.

A precise runtime-related example is the mixing lemma used by
\citet[Proposition~1, author pp.~8--9]{BiasseErukulangara2023}.
It invokes the class-group expansion of \citet[Corollary~1.3]{JMV2009}.
The latter's proof displays the Hecke functions of the relevant class
characters in equations (2.4)--(2.6), author p.~9. This identifies an
exact surviving input without claiming that a generator theorem proves
the entire Hafner--McCurley runtime.

The same spectral theorem governs the ordinary elliptic isogeny
application over imaginary quadratic orders
\citep[Theorems~1.1, 1.5--1.6]{JMV2009}. For base field $K=\Q$ its
ray-class characters are Dirichlet characters and $A$ suffices; the
imaginary-quadratic class characters generally require more.
The conclusion is expansion and rapid mixing at the relevant
endomorphism level, not a polynomial-time discrete-logarithm algorithm.
Similarly, the split-prime selection estimate in the CM/CRT computation
of Hilbert class polynomials uses RH for the ring class field's
Dedekind zeta function
\citep[equation (21) and Lemma~2, corrected author pp.~21--22]{Sutherland2011}.
This is a statement about that stage of the computation.

Schoof point counting occupies the unconditional side of the table.
The original algorithm has deterministic $O(\log^9q)$ complexity in
characteristic different from $2,3$
\citep[pp.~483, 485--490]{Schoof1985}; Hasse's estimate is a proved
finite-field theorem. ECPP certificate correctness is also unconditional.
The practical Atkin-style algorithm described by Schoof has an explicitly
heuristic runtime analysis involving prime and smoothness distributions
\citep[Theorem~5.2, p.~121; pp.~122, 124--125]{Schoof2008}.
Supplying $A$ does not silently supply those additional assumptions.

\subsection{What the separate Dedekind strip gives for class generators}
A fixed zero-free half-plane already gives a polynomial-logarithmic
ray-class generator bound, although it does not give Bach's sharper
quadratic logarithmic bound.  The weighted discriminant argument is
the same one that appears in the auxiliary-prime construction.

\begin{proposition}\label{prop:ray-class-strip}
Fix $1/2\le\theta<1$.  Suppose that for every number field $J$,
$\zeta_J$ has no nontrivial zero with real part greater than $\theta$.
Let $K$ be a number field of degree $n$ and absolute discriminant $D$,
and let $\mathfrak m$ be a modulus with finite part
$\mathfrak m_0$.  Put
\[
 B=1+n+\log D+\log N\mathfrak m_0.
\]
There is a constant $C_\theta$, independent of $K$ and
$\mathfrak m$, such that the ray class group modulo $\mathfrak m$
is generated by prime ideals coprime to $\mathfrak m_0$ of norm at most
\[
 C_\theta B^{1/(1-\theta)}.
\]
In particular, Cureton--Thorner's Theorem~1.1 gives the exponent $32$.
Only its Dedekind-zeta statement is needed for this conclusion.
\end{proposition}

\begin{proof}
Fix once and for all a nonnegative smooth function $w$, supported
in $(1,2)$, with integral $c_w>0$, and use the smoothed prime-ideal
sum $\Psi_J(x)$ defined in Lemma~\ref{lem:smooth}.
The stated zero-free half-plane and the explicit formula give
\[
 \Psi_J(x)=c_wx+
 O_{w,\theta}\bigl(x^\theta(\log D_J+[J:\mathbb Q])\bigr).
\]
Here the constant is independent of $J$; the rapid decay of the
Mellin transform makes the zero-counting sum bounded by a constant
times $\log D_J+[J:\mathbb Q]$.

Let $H$ be the subgroup generated by the prime ideals coprime to
$\mathfrak m_0$ of norm at most $2x$.
Suppose that $H$ is proper.  A maximal subgroup containing it
has prime index $\ell$ in the ray class group.
Class field theory gives a cyclic extension $L/K$ of degree $\ell$
for that quotient, in which every prime counted in $H$ splits
completely.  Every character of this quotient has conductor dividing
$\mathfrak m$.  The conductor--discriminant formula therefore gives
\[
 \log D_L\le\ell\log D+(\ell-1)\log N\mathfrak m_0,
 \qquad [L:\mathbb Q]=\ell n.
\]
This formula follows equally by factoring the completed zeta function
of $L$ into the completed Hecke functions of the characters and
comparing their finite conductor factors.

Remove the Euler factors over prime ideals dividing
$\mathfrak m_0$ from both smoothed sums, and denote the results
by $\Psi_K^{(\mathfrak m)}$ and $\Psi_L^{(\mathfrak m)}$.
Their weighted difference is
\[
 \Psi_K^{(\mathfrak m)}(x)
       -\frac1\ell\Psi_L^{(\mathfrak m)}(x)
 =\left(1-\frac1\ell\right)c_wx+O_{w,\theta}(Bx^\theta).
\]
To justify removing the Euler factors uniformly, for a fixed
prime ideal $\mathfrak p$ of $K$ the interval $(x,2x)$ contains
at most one power of $N\mathfrak p$.
The sum of the residue degrees of primes above $\mathfrak p$
is at most $\ell$.  After the weight $1/\ell$, the removed
contribution is consequently $O_w(\log N\mathfrak p)$ on
either side.  Summing gives $O_w(\log N\mathfrak m_0)$,
which is absorbed in the displayed error.

Every unremoved prime ideal that can contribute has norm at most
$2x$, so its class belongs to $H$ and it splits completely in $L$.
Its entire prime-power contribution to the two weighted sums
cancels.  The left side is therefore zero.
On the other hand, $1-1/\ell\ge1/2$, and choosing
$x=T_\theta B^{1/(1-\theta)}$ with a sufficiently large fixed
$T_\theta$ makes the right side positive.  This is a contradiction.
Thus $H$ is the whole ray class group, and
$C_\theta=2T_\theta$ proves the assertion.
\end{proof}

For the ordinary ideal class group take $\mathfrak m=1$.
If an integral-basis presentation of $O_K$ is given, the resulting
prime ideals can be enumerated by considering rational primes up to
the stated norm bound and factoring the finite algebras
$O_K/pO_K$.  Unconditional algorithms polynomial in the numerical
prime $p$ suffice here.  This constructs a generating list;
it makes no assertion about the time needed to find all relations
or compute the class-group structure.

\subsection{Koiran's AM inclusion needs one solution-field zeta function}

Koiran's extended report states the inclusion of complex polynomial-system
feasibility in $\mathrm{AM}$, with its probabilistic and interactive
meaning explicit \citep[Theorem~2, pp.~5--6]{KoiranExtended1996}.
The journal's prime-density step explicitly notes that only the effective
prime ideal theorem is used \citep[Theorem~9, p.~282]{Koiran1996}.
To see the precise field, choose one algebraic solution, an integral
primitive element $\theta$ for its coordinate field $K$, and the monic
minimal polynomial $f$ of $\theta$. Exclude primes dividing its
discriminant and the denominators of the coordinate formulas.
For the remaining primes, the number $W(p)$ of roots of $f$ modulo $p$
is the number of degree-one prime ideals above $p$.
Higher-degree prime ideals contribute only $O([K:\Q]\sqrt{x})$ to
the prime-ideal count below $x$. RH for $\zeta_K$ therefore estimates
$\sum_{p\le x}W(p)$, and
\[
 W(p)\le[K:\Q]\,\mathbf1_{\{f\text{ has a root modulo }p\}}
\]
gives the necessary lower bound for primes at which the original system
has a solution.

No normal closure is required for this argument.  For example, every
solution of $X^3-2=0$ has a nonnormal degree-three coordinate field.
No such root lies in an abelian extension of $\mathbb Q$: every
subfield of an abelian Galois extension is normal over $\mathbb Q$,
whereas $\mathbb Q(\sqrt[3]{2})$ and its conjugates are not.
Thus $A$ does not in general discharge the solution-field zeta
hypothesis. Choosing an integral primitive element also avoids silently
omitting leading-coefficient bad primes in a nonmonic formulation.
The corrected quantitative bounds in the author's erratum must be used
in the subsequent bit-size analysis \citep{KoiranErratum1997}.
None of these refinements changes the classification of the conclusion
as an $\mathrm{AM}$ inclusion rather than a deterministic polynomial-time
decision algorithm.


\subsection{A fixed uniform zero-free half-plane suffices for the AM protocol}

The AM inclusion also follows from the separate uniform Dedekind
zero-free half-plane.  The following proof spells out the transfer,
including the polynomial bit length of the prime range.

\begin{proposition}\label{prop:fixed-strip-hn-am}
Suppose that a fixed $\theta<1$ has the following property:
for every number field $K$, the Dedekind zeta function $\zeta_K$
has no nontrivial zero with real part greater than $\theta$.
Assume the standard uniform functional equations and local
zero-counting estimates for these zeta functions.
Then feasibility over $\mathbb C$ of a finite system of integer
polynomials belongs to $\mathrm{AM}$, for the usual dense or
sparse binary encoding.
\end{proposition}

\begin{proof}
We identify explicitly the effective algebraic input used below.
For an input of bit length $N$, the standard effective
algebraic-point and arithmetic Nullstellensatz bounds supply
a computable polynomial $Q(N)$ with these two properties.
Write $M=2^{Q(N)}$.
If the system is feasible, it has a point with coordinates in
a number field $K$ and with denominators dividing an integer
$B\ne0$, such that
\[
 n_K\le M,\qquad \log|D_K|\le M,\qquad \log|B|\le M.
\]
More precisely, the coordinates belong to
$O_K[1/B]$.  If the system is infeasible, there is a nonzero
integer $A$ and integer polynomials $G_i$ with
\[
 \sum_iG_iF_i=A,\qquad \log|A|\le M.
\]
These effective algebraic bounds are the inputs in
\citet[Theorems~4, 5, and 7, author pp.~8--10]{KoiranExtended1996},
with the corrected quantitative bounds of \citet{KoiranErratum1997}.
They are external algebraic inputs to the proposition.
The polynomial $Q$ may be increased by a fixed constant below.
For sparse binary exponents the total degree can be exponential
in $N$; the permitted bounds are still exponential in a polynomial
in $N$.

Choose once and for all
\[
 \max(\theta,1/2)<\sigma<1.
\]
We first prove the required uniform prime-density estimate.
For $x\ge2$, define
\[
 W_K(x)=\sum_{N\mathfrak a\le x}
       \Lambda_K(\mathfrak a)(1-N\mathfrak a/x).
\]
Mellin inversion and a contour shift give
\begin{equation}\label{eq:fixed-strip-weighted-primes}
 W_K(x)=\frac{x}{2}
       +O_{\theta,\sigma}\bigl(x^\sigma
                         (\log|D_K|+n_K)\bigr).
\end{equation}
Here the implicit constant is independent of $K$.
For completeness, start with
\[
 \frac{1}{2\pi i}\int_{(2)}
 -\frac{\zeta'_K}{\zeta_K}(s)
 \frac{x^s}{s(s+1)}\,ds.
\]
On the line $\Re s=\sigma$, the standard Hadamard formula
and local zero count imply
\[
 \left|\frac{\zeta'_K}{\zeta_K}(\sigma+it)\right|
 \ll_{\theta,\sigma}\log|D_K|+n_K\log(|t|+3).
\]
One way to see the uniformity is to subtract the Hadamard
formula at $2+it$.  Zeros with $|\Im\rho-t|\le1$ are
separated from the new line by $\sigma-\theta$ and number
$O(\log|D_K|+n_K\log(|t|+3))$.  The difference of the
reciprocal terms for the remaining zeros is summable with
quadratic decay in $|\Im\rho-t|$.  The gamma terms satisfy
the same bound, the pole at $1$ is a fixed distance from the
new line, and $\zeta'_K/\zeta_K(2+it)=O(n_K)$ by its
absolutely convergent Dirichlet series.
The horizontal integrals vanish because the kernel is
$O(|t|^{-2})$.  Only the pole at $s=1$ is crossed, with
residue $x/2$.  The integral on the new line is bounded
by the error in \eqref{eq:fixed-strip-weighted-primes},
because the displayed logarithmic factor is integrable
against $|(\sigma+it)(\sigma+1+it)|^{-1}$.

For each rational prime power $p^j$, the coefficient of
$p^{-js}$ in $-\zeta'_K/\zeta_K(s)$ is at most
$n_K\log p$.  Therefore all terms with $j\ge2$ contribute
at most
\[
 n_K\sum_{p\le\sqrt x}\log x
 \le n_K\sqrt x\log x.
\]
The remaining terms come from prime ideals of norm a
rational prime.  A given rational prime contributes at
most $n_K\log x$.  Thus, when the error in
\eqref{eq:fixed-strip-weighted-primes} is at most $x/4$
and $n_K\sqrt x\log x\le x/8$, there are at least
\[
 \frac{x}{8n_K\log x}
\]
rational primes at most $x$ having a residue-degree-one
prime in $K$.

Choose a fixed integer
\[
 C\ge\max\{8,4/(1-\sigma)\},\qquad
 x=M^C=2^{C Q(N)}.
\]
After increasing the constant term of $Q$, both preceding
error inequalities hold for every input.  If the system
is feasible, a degree-one prime of $K$ over a rational
prime $p\nmid B$ reduces the given algebraic point to a
solution in $\mathbb F_p$.  There are at most
$\log|B|/\log2\le2M$ excluded primes.  Consequently
the set
\[
 S=\{p\le x:p\text{ prime and the input system has
                         a solution in }\mathbb F_p\}
\]
has size at least $256M$, again by a fixed increase of
the constant term of $Q$.  If the system is infeasible,
the Nullstellensatz identity implies $p\mid A$ for
every $p\in S$, so $|S|\le2M$.
No normal closure of $K$ was taken: this is why the
degree and discriminant parameters stay within the
stated exponential bounds.

It remains to realize the cardinality gap by an AM
protocol.  Let $m=CQ(N)+1$, so all integers up to $x$
have fixed-length $m$-bit encodings, and let
$k=Q(N)+4$, so $R=2^k=16M$.
Arthur samples a uniformly random binary matrix
$U\in\mathbb F_2^{k\times m}$ and vector
$b\in\mathbb F_2^k$, and sends them to Merlin.
The affine hash is $h(z)=Uz+b$.
Merlin replies with a prime $p\le x$ and a solution
of the input system in $\mathbb F_p$, satisfying
$h(p)=0$.  Arthur verifies the hash, primality, and all
equations.  Sparse powers are evaluated by binary
exponentiation, so the verification is polynomial in $N$.
The hash is applied only to the prime, not to its possibly
many solution witnesses.

For an infeasible instance, a union bound gives
\[
 \Pr[\exists p\in S:h(p)=0]
 \le |S|/R\le1/8.
\]
For a feasible instance, let $Z$ count the primes of $S$
in the zero bucket.  Pairwise independence gives
$\mathbb EZ=|S|/R\ge16$ and
$\operatorname{Var}Z\le\mathbb EZ$.  Chebyshev's
inequality therefore gives
\[
 \Pr[Z=0]\le1/\mathbb EZ\le1/16.
\]
On the complementary event Merlin can provide an accepted
witness.  Both messages and Arthur's computation have
polynomial length and time, since $C$ is fixed and $Q$
is polynomial.  This is a two-message public-coin AM
protocol with completeness at least $15/16$ and
soundness at most $1/8$.
\end{proof}

Consequently Cureton--Thorner's Theorem~1.1 supplies the hypothesis
of Proposition~\ref{prop:fixed-strip-hn-am}.  The AM conclusion is
available through that separate input, while the original
solution-field RH hypothesis is not discharged by the rational
Dirichlet transfer alone.  The proof above uses one solution-field
zeta function for each feasible instance and no normal closure.

\subsection{Prime-field nonresidues and extension-field hypotheses}
The distinction between a rational Dirichlet character and a ray-class
character over another number field also matters in statements about
finite-field nonresidues.  The weak-GRH estimate of
\citet[Theorem~6.7, p.~8]{BhargavaEtAl2017} concerns
$\mathbb F_p^\times$: its character is explicitly
$x\mapsto x^{(p-1)/r}$.  Theorem~\ref{thm:nonresidue} proves the
extension transfer when $r\mid p-1$.  Neither statement by itself
covers $r\nmid p-1$.

The wider generating-set theorem of
\citet[author pp.~2--3]{BuchmannShoup1996} has complexity
$(n\log p)^{O(n)}$ and hence polynomial complexity for fixed $n$.
Its definition of ERH explicitly covers Dedekind zeta functions of
all number fields \citep[author p.~5]{BuchmannShoup1996}.
There are two relevant uses: Fact~1 on that page invokes the
Adleman--Lenstra auxiliary prime to construct a cyclic degree-$n$
number field in which $p$ is inert; Fact~6 on author p.~9 invokes
Bach's ray-class generator bound over that field, with modulus $p$.
Thus this theorem does not discharge the remaining extension-field
nonresidue step merely by assuming $A$.

One can state the algorithmic input without any convention about the
letters GRH: require deterministic polynomial-time construction of
$\mathbb F_{p^d}$ and deterministic construction of a $k$th nonresidue
in a presented such field in time polynomial in $k$ and $d\log p$.
This is exactly Hypothesis~$(*)$ in
\citet[\S1.1, p.~3]{Guo2017Thesis}.  Its two clauses must both be
justified in an application.  In particular, the prime-field result
in Theorem~6.7 just cited does not alone prove that full hypothesis.
The finite-family formulation in Theorem~\ref{thm:quasipoly}
records the additional input needed by the particular factorization
route analyzed here.

\section{The zeta function of a pure field}
This section proves the analytic equivalence and then identifies the
finite computation which an auxiliary prime makes possible.

\begin{lemma}[The affine permutation identity]\label{lem:affine}
For $F,K,E$ in Theorem~\ref{thm:kummer},
\begin{equation}\label{eq:zetaidentity}
 \zeta_E(z)=\zeta_K(z)
       \left(\frac{\zeta_F(z)}{\zeta_{\Q}(z)}\right)^{r-1}
\end{equation}
as meromorphic functions, including their ramified Euler factors.
Moreover, for a nontrivial character $\chi$ of $\Gal(E/K)$,
\begin{equation}\label{eq:onehecke}
 L_K(z,\chi)=\frac{\zeta_F(z)}{\zeta_{\Q}(z)}.
\end{equation}
Here $L_K$ is the finite-order Hecke $L$-function identified with
the Artin $L$-function by abelian class field theory over $K$.
\end{lemma}
\begin{proof}
Only $r$ ramifies in $K/\Q$, so $p$ is unramified there.
At every prime above $p$, the polynomial $X^r-p$ is Eisenstein.
Thus $[E:K]=r$, and the splitting group over $\Q$ is
$G=\F_r\rtimes\F_r^\times$, acting by $t\mapsto bt+c$ on the
$r$ roots.  Let $N=\F_r$ be its translation subgroup, let $P$ be
the permutation character of that action, and let $Q_0$ be the
regular character of $G/N$, inflated to $G$.  Then
\[
 \operatorname{reg}_G=(r-1)(P-1)+Q_0.
\]
At the identity both sides equal $r(r-1)$.  At a nonzero translation
the right side is $(r-1)(-1)+(r-1)=0$.  At an element with
$b\ne1$, the affine action has exactly one fixed point and
$Q_0$ vanishes.  These cases prove the identity at every element.
Artin formalism applied to this identity proves
\eqref{eq:zetaidentity}.  It applies at a ramified prime as well:
an equality of representations remains an equality after taking
inertia invariants, and therefore gives the same local determinants.

The induced character $\operatorname{Ind}_N^G\chi$ is $r-1$
at the identity, $-1$ at a nonzero translation, and zero off $N$.
Indeed conjugation runs through all nontrivial characters of $N$,
whose sum is $-1$ on a nonidentity element.  Consequently
$\operatorname{Ind}_N^G\chi=P-1$.
Induction of Artin $L$-functions and class field theory give
\eqref{eq:onehecke}.  No general Artin-holomorphy conjecture is used.
\end{proof}

\begin{proof}[Proof of Theorem~\ref{thm:kummer}]
Under $A$, both $\zeta_{\Q}$ and $\zeta_K$ are nonzero at
every point of the open critical strip off the line $\Re z=1/2$.
For $\zeta_K$ this follows from its factorization into Dirichlet
$L$-functions, since $K/\Q$ is abelian.  Equation
\eqref{eq:zetaidentity} therefore says that $\zeta_E$ vanishes at
such a point if and only if $\zeta_F$ does.  This proves the RH
equivalence and also identifies its extra hypothesis as RH for
the single function in \eqref{eq:onehecke}.

The cyclotomic discriminant is $|D_K|=r^{r-2}$.
The relative discriminant of the order $\mathcal O_K[p^{1/r}]$
is generated, up to a unit, by $r^r p^{r-1}$.  The relative
field discriminant divides this ideal.  The discriminant tower
formula consequently gives
\[
 |D_E|\le r^{r(2r-3)}p^{(r-1)^2},\qquad [E:\Q]=r(r-1).
\]
Apply Bach--Sorenson's theorem to the conjugacy class of a
nonidentity translation.  Under the zeta hypotheses just established
it gives an unramified rational prime of size at most
$(4\log|D_E|+\tfrac52[E:\Q]+5)^2$ with that Frobenius class.
Substitution proves \eqref{eq:bound}.

Its restriction to $K$ is trivial, so $\ell\equiv1\pmod r$.
At a prime above $\ell$, write its action on $p^{1/r}$ as
$p^{1/r}\mapsto\zeta_r^c p^{1/r}$, where $c\ne0$.
The defining Frobenius congruence gives
\[
 p^{(\ell-1)/r}\equiv\zeta_r^c\not\equiv1\pmod\ell.
\]
Both $p$ and $r$ ramify, so neither is $\ell$.
The remainder follows from Lemma~\ref{lem:period}, proved next.
Enumerating positive integers up to \eqref{eq:bound}, testing primality,
and testing the two displayed conditions has polynomial bit cost.
\end{proof}

\begin{lemma}[A Gaussian-period construction]\label{lem:period}
Let $p\ne\ell$ be primes, let the prime $r$ divide $\ell-1$, and
suppose $p^{(\ell-1)/r}\not\equiv1\pmod\ell$.
A presentation of $\F_{p^r}$ can be constructed deterministically
in time polynomial in $\ell$ and $\log p$.
\end{lemma}
\begin{proof}
Put $G=(\Z/\ell\Z)^\times$ and let $H$ be its subgroup of
order $(\ell-1)/r$.  We can enumerate $G$, its multiplication
table, and $H$ in polynomial time in $\ell$; factoring integers
at most $\ell$ by trial division also has this cost.  Form the
separable algebra $B=\F_p[Z]/\Phi_\ell(Z)$ and let $G$ act by
$Z\mapsto Z^b$.  Compute $B^H$ by linear algebra.
Over an algebraic closure, the points of $B$ are indexed by
$G$, and those of $B^H$ by $G/H$.  Frobenius acts on these
cosets by multiplication by $p$.  The displayed nonresidue
condition is exactly $p\notin H$.  Since $G/H$ has prime order
$r$, Frobenius is a single $r$-cycle.  A finite reduced algebra
over a finite field is a field precisely when Frobenius is
transitive on its geometric points.  Hence $B^H\cong\F_{p^r}$.
All matrices have dimension at most $\ell-1$, and modular
arithmetic has polynomial bit cost.

There is also an explicit generator.  The elements
$Z,Z^2,\ldots,Z^{\ell-1}$ form a basis of $B$: a relation
would be a scalar multiple of $\Phi_\ell$ with zero constant
coefficient, hence would vanish.  The sums
$\theta_{bH}=\sum_{h\in H}Z^{bh}$ have disjoint supports in
this basis, so they are linearly independent and distinct.
They form a basis of $B^H$.  Frobenius cyclically permutes them;
therefore $\theta_H$ has degree $r$.  Its minimal polynomial,
found by linear dependence among its powers, gives a polynomial
presentation of the field.
\end{proof}

\subsection{The sharper classical bound and a one-sided criterion}
The least-nonsplit-prime theorem gives a better dependence on $r$ than
the explicit general conjugacy-class bound.  Murty's result, in the
form recorded by \citet[p.~1, equation~(1.2)]{Zaman2018}, gives
$P(E/K)\ll(r^{-1}\log|D_E|)^2$ under RH for $\zeta_E$.
Here $P(E/K)$ is the least norm of a degree-one prime of $K$
which is unramified and does not split completely in $E$.
It gives the scale $r^2\log^2(rp)$ in our situation.
We prove the needed specialization directly, including the exclusion
of the two rational primes which ramify, and then state its weaker
one-sided sufficient condition.

Fix once and for all a nonzero nonnegative smooth function $w$
with compact support in $(1,2)$, and write
\[
 c_w=\int_1^2 w(t)\,dt,\qquad
 \Psi_J(x)=\sum_{\mathfrak l}\sum_{j\geq1}
       \log N\mathfrak l\;w((N\mathfrak l)^j/x)
\]
for a number field $J$.  Its absolute degree and discriminant
are denoted by $n_J$ and $D_J$.

\begin{lemma}[Uniform smoothed estimate]\label{lem:smooth}
Let $1/2\leq\theta<1$ be fixed.  If $\zeta_J$ has no
nontrivial zero with real part greater than $\theta$, then,
uniformly for number fields $J$ and real $x\geq2$,
\[
 \Psi_J(x)=c_wx+
 O_w\bigl(x^\theta(\log|D_J|+n_J)\bigr).
\]
\end{lemma}
\begin{proof}
Put $\widehat w(z)=\int_1^2w(t)t^{z-1}\,dt$.
The smoothed explicit formula gives
\[
 \Psi_J(x)=c_wx-\sum_\rho x^\rho\widehat w(\rho)
          +O_w(\log|D_J|+n_J),
\]
where $\rho$ runs over nontrivial zeros with multiplicity.
This is the uniform formula used in
\citet[Lemma~11.2, p.~45]{OpenAI2026}.
For completeness, the uniformity of its zero term follows as follows.
Three integrations by parts give
$|\widehat w(\sigma+it)|\ll_w(1+|t|)^{-3}$,
uniformly for $0\leq\sigma\leq1$.
The standard unit-height zero count is
\[
 N_J(T+1)-N_J(T)\ll\log|D_J|+n_J\log(T+3).
\]
Summing over integer height intervals therefore bounds
$\sum_\rho|\widehat w(\rho)|$ by
$O_w(\log|D_J|+n_J)$.  The zero-free assumption bounds
$|x^\rho|$ by $x^\theta$.
The pole at one gives the displayed main term, and the gamma
factors and trivial zeros give the stated remainder.
These estimates prove the lemma.  Thus the proof is valid
for any fixed $\theta$, even though OpenAI's stated lemma
selects one particular zero-free width.
\end{proof}

\begin{proposition}[The pure-field prime detector]\label{prop:detector}
With $F,K,E$ as in Theorem~\ref{thm:kummer}, put
$\mathcal H(x)=\Psi_K(x)+\Psi_{\Q}(x)-\Psi_F(x)$.
Its prime-power coefficients are nonnegative, including at
ramified primes.  At an unramified rational prime $\ell$, its
coefficient at $\ell^j$ is
\[
 r\log\ell\;\mathbf 1_{\mathcal C}(g_\ell^j),
\]
where $g_\ell$ is a Frobenius element and
$\mathcal C=N\setminus\{1\}$ is the class of nontrivial
translations.  In particular, if no auxiliary prime occurs in
$(x,2x)$, then $\mathcal H(x)\ll_w r\sqrt x$.
\end{proposition}
\begin{proof}
The virtual character $Q_0+1-P$ in Lemma~\ref{lem:affine}
equals $r$ on $\mathcal C$ and zero elsewhere.
The coefficient assertion follows from the local logarithmic
derivative of the corresponding Artin Euler factors.
At a ramified prime, the coefficient is the trace of a power
of Frobenius on inertia invariants.  If $I$ is inertia and
$g$ a lift of Frobenius, that trace for this virtual
representation is
\[
 \frac1{|I|}\sum_{h\in I}(Q_0+1-P)(g^jh),
\]
which is nonnegative and at most $r$.

If no auxiliary prime occurs in $(x,2x)$, all unramified
first-power terms in the support of $w$ vanish.
The other unramified terms have $j\geq2$, so their total is
at most
\[
 r\|w\|_\infty\sum_{j\geq2}\vartheta((2x)^{1/j})
       \ll_w r\sqrt x.
\]
The last inequality follows from the elementary bound
$\vartheta(y)\ll y$: the square term costs $O(\sqrt x)$
and the remaining $O(\log x)$ terms cost
$O(x^{1/3}\log x)=O(\sqrt x)$.
Only $p$ and $r$ can ramify in $E$.
For each, summing its powers in the support of $w$ costs
$O_w(r\log(2x))$, also absorbed by $O_w(r\sqrt x)$.
This proves the upper bound.
\end{proof}

\begin{theorem}[Polynomial search from one zero-free half-plane]
\label{thm:fixedstrip}
Assume $A$ and let $1/2\leq\theta<1$ be fixed.
If $\zeta_{\Q(p^{1/r})}$ has no nontrivial zero with real
part greater than $\theta$, then there is an auxiliary prime
\[
 \ell\ll_\theta \bigl(r\log(rp)\bigr)^{1/(1-\theta)}.
\]
In particular, RH for that pure-field zeta function gives
$\ell\ll r^2\log^2(rp)$ and deterministic field construction
in time polynomial in $r$ and $\log p$.
\end{theorem}
\begin{proof}
The discriminant of $X^r-p$ is, up to sign,
$r^r p^{r-1}$, so
$\log|D_F|\leq r\log r+(r-1)\log p$.
Under $A$, Lemma~\ref{lem:smooth} applies to $K$ and $\Q$
with $\theta=1/2$, and by hypothesis it applies to $F$
with the stated $\theta$.  Thus, with $B=r\log(rp)$,
\[
 \mathcal H(x)=c_wx+O_w(Bx^\theta).
\]
If there is no auxiliary prime in $(x,2x)$,
Proposition~\ref{prop:detector} bounds the left side by
$O_w(r\sqrt x)$.  Set $x=T B^{1/(1-\theta)}$.
After division by $x$, the analytic error is
$O_w(T^{-(1-\theta)})$, and $r/\sqrt x\leq T^{-1/2}$
up to an absolute constant.  A sufficiently large fixed
$T$, depending only on $w$ and $\theta$, makes both
errors smaller than $c_w/3$, a contradiction.
Hence a suitable prime is less than $2x$.
\end{proof}

\begin{proposition}[A finite one-sided sufficient condition]
\label{prop:window}
Fix $0<\varepsilon<1$.  Under $A$, if
\[
 x\geq C_{w,\varepsilon}r^2\log^2r,
 \qquad \Psi_F(x)\leq(2-\varepsilon)c_wx,
\]
for a sufficiently large constant $C_{w,\varepsilon}$,
then an auxiliary prime occurs in $(x,2x)$.
Consequently a uniform polynomial upper bound on some
such $x=x(p,r)$ would suffice for the classical
field-construction algorithm.
\end{proposition}
\begin{proof}
Under $A$, the two abelian fields give
\[
 \Psi_K(x)+\Psi_{\Q}(x)
       =2c_wx+O_w(r\log r\sqrt x).
\]
The assumed upper bound for $\Psi_F$ therefore gives
$\mathcal H(x)\geq\varepsilon c_wx
-O_w(r\log r\sqrt x)$.
If no auxiliary prime occurs, Proposition~\ref{prop:detector}
gives the contradictory upper bound $O_w(r\sqrt x)$
once $C_{w,\varepsilon}$ is large enough.
\end{proof}

The positivity in Proposition~\ref{prop:detector} gives only
$\Psi_F(x)\leq2c_wx+O_w(r\log r\sqrt x)$ under $A$.
It reaches the threshold two without a fixed gap.
No argument here proves the strict one-sided estimate from
$A$.  This is the precise remaining step for this route;
a finite-scale estimate would be enough, so global RH
for $\zeta_F$ is a stronger sufficient premise than the
construction itself requires.


\subsection{An averaged window estimate and a limitation of positive tests}
\label{sec:positive-window-tests}

The one-sided condition in Proposition~\ref{prop:window} can be
forced by an unconditional averaged estimate, but the resulting range
is a power of the discriminant.  We give the estimate and isolate the
conductor term responsible for that scale.  The limitation proved below
concerns a specified use of positive transforms; it does not exclude
other deductions from $A$.

Let $J$ be a number field of degree $n_J$, with $r_{1,J}$ real
embeddings and absolute discriminant $D_J$.  For an even real compactly
supported test function $f$, define
\begin{align*}
 \Phi_f(s)&=\int_{\mathbb R}f(t)e^{(s-1/2)t}\,dt,\\
 T_J(f)&=\sum_{\mathfrak p}\sum_{k\geq1}
 \frac{\log N\mathfrak p}{(N\mathfrak p)^{k/2}}
 f(k\log N\mathfrak p),\\
 Z_J(f)&=\sum_{\rho_J}\Phi_f(\rho_J).
\end{align*}
The zeros in $Z_J$ are nontrivial zeros, counted with multiplicity.
We take $f$ to satisfy the admissibility conditions in
\citet[author pp.~4--5, equations~(3)--(4)]{BDF2008}:
$f$ is continuous, $f(t)e^{(1/2+\eta)t}$ has bounded variation
and is integrable on $[0,\infty)$ for some $\eta>0$, and
$(f(0)-f(t))/t$ has bounded variation there.  Their unconditional
explicit formula, specialized to the trivial character, is
\begin{equation}\label{eq:positive-test-explicit}
 2T_J(f)=2\Phi_f(1)+f(0)\log D_J+\mathcal A_J(f)-Z_J(f),
\end{equation}
where, with $\gamma_{\!E}$ denoting Euler's constant,
\begin{align*}
 \mathcal A_J(f)={}&
 f(0)\left(-n_J\gamma_{\!E}-n_J\log(8\pi)
                   -\frac{r_{1,J}\pi}{2}\right)\\
 &+r_{1,J}\int_0^\infty
       \frac{f(0)-f(t)}{2\cosh(t/2)}\,dt
  +n_J\int_0^\infty
       \frac{f(0)-f(t)}{2\sinh(t/2)}\,dt .
\end{align*}
The factor $2$ in the pole term is important: for an even test,
$\Phi_f(0)=\Phi_f(1)$, so the two poles contribute
$\Phi_f(0)+\Phi_f(1)=2\Phi_f(1)$.

\begin{lemma}[A compact test positive throughout the strip]
\label{lem:strip-positive-triangle}
For $S\geq1$, put
\[
 h_S(t)=\left(1-\frac{|t|}{S}\right)_+,
 \qquad f_S(t)=e^{-|t|/2}h_S(t).
\]
Then $\operatorname{Re}\Phi_{f_S}(\beta+i\gamma)\geq0$
for every $0\leq\beta\leq1$ and $\gamma\in\mathbb R$.
Moreover,
\begin{equation}\label{eq:triangle-pole}
 \Phi_{f_S}(1)=\frac S2+1-\frac{1-e^{-S}}{S}.
\end{equation}
\end{lemma}
\begin{proof}
Use the Fourier convention
$\widehat g(\gamma)=\int_{\mathbb R}g(t)e^{i\gamma t}\,dt$.
The real part in question is the Fourier transform of
\[
 f_S(t)\cosh((\beta-1/2)t)
 =\frac{h_S(t)}2
   \left(e^{-\beta|t|}+e^{-(1-\beta)|t|}\right).
\]
The two functions
\[
 \widehat h_S(\gamma)
 =\frac{2(1-\cos(S\gamma))}{S\gamma^2},
 \qquad
 \widehat{e^{-a|t|}}(\gamma)=\frac{2a}{a^2+\gamma^2}
 \quad(a>0)
\]
are nonnegative, with the first expression interpreted continuously
at zero.  A product of these original functions has a Fourier transform
equal to their convolution divided by $2\pi$, and is therefore
nonnegative.  The cases $\beta=0,1$ follow by continuity.
Finally,
\[
 \Phi_{f_S}(1)=\int_0^S
       (1+e^{-t})\left(1-\frac tS\right)dt,
\]
which gives \eqref{eq:triangle-pole}.
The test is admissible in \eqref{eq:positive-test-explicit}.
It is continuous and piecewise smooth, and its transform is
$O_S((1+|\gamma|)^{-2})$ uniformly in the closed strip.
Thus its zero sum converges absolutely by the usual unit-height
zero count $O(\log D_J+n_J\log(|\gamma|+3))$.
\end{proof}

\begin{proposition}[An averaged one-sided estimate]
\label{prop:averaged-window}
For the fixed nonnegative smooth function $w$ in
Proposition~\ref{prop:window}, and every number field $J$,
\begin{equation}\label{eq:averaged-window-bound}
 \int_1^{e^S}\frac{\Psi_J(x)}{x^2}
       \left(1-\frac{\log x}{S}\right)dx
 \leq\frac{c_w}{2}(S+\log D_J)+O_w(n_J)
 \qquad(S\geq1).
\end{equation}
This estimate is unconditional and uniform in $J$ and $S$.
\end{proposition}
\begin{proof}
Apply \eqref{eq:positive-test-explicit} to $f_S$.
Its zero sum is real and nonnegative by
Lemma~\ref{lem:strip-positive-triangle}.
The archimedean term is uniformly $O(n_J)$, and in fact is
nonpositive for $S\geq1$.  To verify the latter assertion directly,
the real-place integral is at most
$\int_0^\infty(2\cosh(t/2))^{-1}dt=\pi/2$, and
\begin{align*}
 \int_0^\infty\frac{1-f_S(t)}{2\sinh(t/2)}dt
 &=2\log2+\int_0^\infty
       \frac{\min(t/S,1)}{e^t-1}\,dt\\
 &\leq2\log2+\frac{\pi^2}{6S}.
\end{align*}
Consequently
\[
 \mathcal A_J(f_S)
 \leq-n_J\left(\gamma_{\!E}+\log(2\pi)
                         -\frac{\pi^2}{6}\right)<0.
\]
If $\Lambda_J(\mathfrak a)$ denotes the ideal von Mangoldt
function, then
\begin{equation}\label{eq:harmonic-prime-bound}
 \begin{split}
 T_J(f_S)
 &=\sum_{N\mathfrak a\leq e^S}
   \frac{\Lambda_J(\mathfrak a)}{N\mathfrak a}
       \left(1-\frac{\log N\mathfrak a}{S}\right)\\
 &\leq\frac S2+\frac12\log D_J+1.
 \end{split}
\end{equation}

Write $X=e^S$ and $c_{w,1}=\int_1^2w(v)\log v\,dv$.
For $n=N\mathfrak a\leq X$, substitution $v=n/x$ gives
\[
 \int_1^X\frac{w(n/x)}{x^2}
       \left(1-\frac{\log x}{S}\right)dx
 =\frac1n\left[c_w\left(1-\frac{\log n}{S}\right)
                         +\frac{c_{w,1}}S\right].
\]
The terms with $X<n<2X$ contribute $O_w(n_J/S)$.
Indeed, in that interval the integrand after substitution has
logarithmic weight at most $(\log2)/S$, and
\[
 \sum_{X<N\mathfrak a<2X}
       \frac{\Lambda_J(\mathfrak a)}{N\mathfrak a}\ll n_J.
\]
Also
\[
 \sum_{N\mathfrak a\leq X}
       \frac{\Lambda_J(\mathfrak a)}{N\mathfrak a}
       \ll n_J(1+\log X).
\]
For both bounds, the coefficient at a rational prime power in
$-\zeta'_J/\zeta_J$ is at most $n_J$ times the corresponding
rational von Mangoldt coefficient; Chebyshev's bound
$\psi_{\mathbb Q}(x)\ll x$ and partial summation then suffice.
It follows that the integral in \eqref{eq:averaged-window-bound}
equals $c_wT_J(f_S)+O_w(n_J)$.  Equation
\eqref{eq:harmonic-prime-bound} proves the proposition.
\end{proof}

For $F=\mathbb Q(p^{1/r})$, fix $0<\varepsilon<1$ and $u_0\geq1$.
If $\Psi_F(x)\geq(2-\varepsilon)c_wx$ for every
$e^{u_0}\leq x\leq e^S$, the left side of
\eqref{eq:averaged-window-bound} is at least
\[
 \frac{(2-\varepsilon)c_w(S-u_0)^2}{2S}.
\]
Comparison gives the necessary inequality
\begin{equation}\label{eq:window-scale-obstruction}
 (1-\varepsilon)S
 \leq\log D_F+O_w(r)+2(2-\varepsilon)u_0.
\end{equation}
Thus, for a suitable effective constant $C_w$, choosing
\[
 S>\frac{\log D_F+C_wr+2(2-\varepsilon)u_0}{1-\varepsilon}
\]
forces at least one $x\in[e^{u_0},e^S]$ satisfying the strict
one-sided bound.  Under $A$, take
$e^{u_0}\geq C_{w,\varepsilon}r^2\log^2r$ and apply
Proposition~\ref{prop:window}.  This gives an auxiliary prime with
\begin{equation}\label{eq:averaged-auxiliary-bound}
 \ell\leq D_F^{1/(1-\varepsilon)}
                  \exp(O_{w,\varepsilon}(r)).
\end{equation}
This conclusion is useful for locating the loss in the averaging
argument, even though sharper classical discriminant-power bounds are
available.  Its logarithmic search length is proportional to
$\log D_F$, rather than to $\log r+\log\log p$.

The next estimate uses the classical interval case of Tur\'an's
extremal problem for positive definite functions; see
\citet[author version, Theorem~6, equation~(29), p.~14,
and Remark~4, p.~15]{KR2006Turan}.  We include a periodization
proof to fix the normalization and the support convention.

\begin{lemma}[Tur\'an bound for tests positive throughout the strip]
\label{lem:strip-positive-barrier}
Let $f$ be a continuous real even function supported on $[-S,S]$.
Suppose
\[
 \operatorname{Re}\Phi_f(\beta+i\gamma)\geq0
 \quad(0<\beta<1,\ \gamma\in\mathbb R).
\]
Then $f(0)\geq0$, and
\begin{equation}\label{eq:strip-positive-pole-bound}
 0\leq\Phi_f(1)\leq Sf(0),
 \qquad
 \Phi_f(0)+\Phi_f(1)\leq2Sf(0).
\end{equation}
If $f(0)=0$, then $f$ is identically zero.
\end{lemma}
\begin{proof}
Continuity in $\beta$ gives positivity also at $\beta=1$.
The function $g(t)=f(t)\cosh(t/2)$ therefore has nonnegative
Fourier transform.  Here no smoothness assumption or informal
Fourier inversion is needed: Gaussian regularization gives
\[
 \frac1{2\pi}\int_{\mathbb R}\widehat g(\gamma)
                           e^{-a\gamma^2}\,d\gamma
      \longrightarrow g(0)\qquad(a\downarrow0).
\]
Monotone convergence shows that $\widehat g$ is integrable,
with integral $2\pi g(0)$.  Fourier inversion now implies
\[
 |g(t)|\leq\frac1{2\pi}\int_{\mathbb R}
                   \widehat g(\gamma)\,d\gamma=g(0)=f(0).
\]
The integral of $g$ is $\Phi_f(1)=\widehat g(0)\geq0$.
To obtain the stronger interval bound, fix $T>S$ and periodize:
\[
 G_T(t)=\sum_{k\in\mathbb Z}g(t+kT).
\]
This is a continuous $T$-periodic function, and
$G_T(0)=g(0)$ because $g$ is supported on $[-S,S]$.
Its Fourier coefficients are
\[
 c_n=\frac1T\widehat g(2\pi n/T)\geq0.
\]
The Fej\'er sums at zero therefore satisfy
\[
 \sum_{|n|\leq N}\left(1-\frac{|n|}{N+1}\right)c_n
       \geq c_0=\frac{\widehat g(0)}T.
\]
By Fej\'er's theorem, the left side converges to
$G_T(0)=g(0)$.  Hence $\Phi_f(1)\leq Tf(0)$.
Letting $T\downarrow S$ proves the first bound.
Evenness gives $\Phi_f(0)=\Phi_f(1)$ and the second.
If $f(0)=0$, the preceding pointwise bound forces $g=f=0$.
\end{proof}

The exact conductor for the detector is smaller than $D_F$,
but has the same dependence on $p$.  Put
\[
 K=\mathbb Q(\zeta_r),\qquad
 L(s)=\frac{\zeta_F(s)}{\zeta_{\mathbb Q}(s)},\qquad
 Q_\chi=\frac{D_F}{D_K}.
\]
The Artin induction identities identify $L(s)$ with a nontrivial
finite-order Hecke $L$-function over $K$, of finite conductor
$\mathfrak f_\chi$, and give
$D_KN\mathfrak f_\chi=D_F$.
Thus $Q_\chi=N\mathfrak f_\chi$ is a positive integer.
For $T_{\mathcal H}=T_K+T_{\mathbb Q}-T_F$,
subtracting the three explicit formulas gives the exact identity
\begin{equation}\label{eq:detector-positive-test}
 \boxed{
 2T_{\mathcal H}(f)
   =2\Phi_f(1)-f(0)\log Q_\chi-Z_K(f)+Z_L(f).
 }
\end{equation}
Here $Z_L$ is the zero sum of the entire nontrivial Hecke function
$L$, so $Z_F=Z_{\mathbb Q}+Z_L$.
There is no archimedean remainder in
\eqref{eq:detector-positive-test}:
\[
 n_K+n_{\mathbb Q}-n_F=(r-1)+1-r=0,
 \qquad
 r_{1,K}+r_{1,\mathbb Q}-r_{1,F}=0+1-1=0.
\]

\begin{proposition}[A limitation of discarding the unknown zero sum]
\label{prop:positive-discard-limitation}
Let $f$ be a nonzero admissible test satisfying the hypotheses of
Lemma~\ref{lem:strip-positive-barrier}.  The lower bound obtained
from \eqref{eq:detector-positive-test} by using only
$Z_L(f)\geq0$ cannot be strictly positive if
\[
 S\leq\frac12\log Q_\chi.
\]
\end{proposition}
\begin{proof}
Both $Z_L(f)$ and $Z_K(f)$ are nonnegative by the transform
hypothesis.  The stated use of positivity gives only
\[
 2T_{\mathcal H}(f)
 \geq 2\Phi_f(1)-f(0)\log Q_\chi-Z_K(f).
\]
The expression on the right is at most
$f(0)(2S-\log Q_\chi)$ by
Lemma~\ref{lem:strip-positive-barrier}, and $f(0)>0$.
It therefore cannot certify a positive value in the stated range,
even if the abelian zero sum $Z_K(f)$ is known exactly.
\end{proof}

Only $p$ and $r$ ramify in $F$, and $p$ is tamely totally
ramified, so $v_p(D_F)=r-1$.  Since
$D_F\leq r^rp^{r-1}$ and $D_K=r^{r-2}$, the integral conductor
$Q_\chi$ satisfies
\begin{equation}\label{eq:detector-conductor-size}
 p^{r-1}\leq Q_\chi\leq r^2p^{r-1},
 \qquad
 (r-1)\log p\leq\log Q_\chi
                   \leq(r-1)\log p+2\log r.
\end{equation}
In the difficult scalar case, $s=\operatorname{ord}_r(p)$
and $v_r(p^s-1)\geq2$, one has the exact identity
$Q_\chi=p^{r-1}$.  Indeed, then
$p^{r-1}\equiv1\pmod{r^2}$; the decomposition
$\mathbb Z_r^\times=\mu_{r-1}\times(1+r\mathbb Z_r)$
and the $r$-adic logarithm show that $p$ is an $r$th power
in $\mathbb Q_r$.  Since $K$ contains $\mu_r$, the extension
$K(p^{1/r})/K$ splits completely at its prime above $r$.
The character conductor is consequently supported only above $p$,
and its norm has the already determined $p$-adic exponent $r-1$.

For a search bounded by a fixed polynomial in $r$ and $\log p$,
the logarithmic support is $S=O(\log r+\log\log p)$.
Equations \eqref{eq:strip-positive-pole-bound} and
\eqref{eq:detector-conductor-size} explain why the strategy in
Proposition~\ref{prop:positive-discard-limitation} cannot supply
that range uniformly.  This is already visible with $r$ fixed.
It is a limitation of discarding the nonabelian zero sum, not an
impossibility assertion about the one-sided window condition or $A$.

For comparison, suppose transform positivity is required only for
$\delta\leq\beta\leq1-\delta$, where $0<\delta\leq1/2$.
The same argument applied to
$g_\delta(t)=f(t)\cosh((1/2-\delta)t)$ gives
\[
 |\Phi_f(1)|\leq
 4f(0)\frac{e^{\delta S}-1}{\delta}.
\]
Indeed, $|g_\delta(t)|\leq f(0)$ and
$\cosh(t/2)/\cosh((1/2-\delta)t)\leq2e^{\delta t}$
for $t\geq0$.  With fixed $\delta>0$, this upper bound permits
a pole response large enough at $S$ proportional to
$\log\log Q_\chi$; it does not construct such a test.
If instead $\delta\asymp1/\log Q_\chi$ and
$S=O(\log\log Q_\chi)$, the bound remains $O(Sf(0))$.
This comparison does not cover a height-dependent positivity region
with a separately estimated tail, nor a quantitative lower bound
for $Z_L(f)$, nor a sign-changing transform whose zero
terms can be controlled by additional information.

The common multiplicity in the normal closure does not create an
additional conductor discount in these formulas.  The identities
\[
 \zeta_E=\zeta_K L^{r-1},\qquad D_E=D_KD_F^{r-1}
\]
give
\[
 Z_E(f)-Z_K(f)=(r-1)Z_L(f),\qquad
 \frac{\log D_E-\log D_K}{r-1}=\log D_F.
\]
The same factor $r-1$ has therefore already been removed from
both the zero sum and the conductor in passing to $L$.
Combining its explicit formula with that of $K$ leaves exactly
$\log Q_\chi$ in \eqref{eq:detector-positive-test}.
Any further saving would require additional information on the
remaining zero sum, rather than the multiplicity identity alone.

\subsubsection{Rational twists and the full tensor-positivity cone}
\label{sec:tensor-twist-audit}

This calculation determines the constraints imposed by rational twists and positive tensor powers.
Let \(r\geq3\) and \(p\ne r\) be primes, and put
\[
 K=\mathbb Q(\zeta_r),\qquad F=\mathbb Q(p^{1/r}),\qquad
 E=K(p^{1/r}),\qquad n=r-1.
\]
Write \(G=\operatorname{Gal}(E/\mathbb Q)=N\rtimes\mathbb F_r^\times\),
where \(N\simeq C_r\), and let \(\rho\) denote its irreducible
representation of degree \(n\). Let \(Q_0\) be the regular representation
of \(G/N\), inflated to \(G\). Thus
\begin{equation}\label{eq:twist-character-values}
 \rho(g)=
 \begin{cases}n&g=1,\\-1&g\in N\setminus\{1\},\\0&g\notin N,\end{cases}
 \qquad Q_0(g)=n\,1_N(g),
 \qquad
 L(s,\rho)=\frac{\zeta_F(s)}{\zeta_{\mathbb Q}(s)}.
\end{equation}
The same notation denotes a representation and its character.
For an unramified prime \(\ell\), the value
\((Q_0-\rho)(\operatorname{Frob}_\ell)\) is \(r\) exactly when
\(\ell\equiv1\pmod r\) and \(p\) is not an \(r\)-th power modulo
\(\ell\), and is zero otherwise. These are the auxiliary primes.
The assertions below concern particular positivity and averaging
arguments. They do not assert a limitation on every possible deduction
from rational Dirichlet GRH.

\subsubsection{The inequalities from all tensor powers}

Fix nonnegative weights \(w_{\ell,k}\), only finitely many of which
are nonzero. For a class function \(h\) on \(G\), define
\[
 \mathcal P(h)=
 \sum_{\ell,k}w_{\ell,k}\log\ell\,
 \frac1{|I_\ell|}\sum_{u\in I_\ell}
 h(\phi_\ell^k u),
\]
where \(I_\ell\) is an inertia group and \(\phi_\ell\) is a lift of
arithmetic Frobenius. For a representation this is the corresponding
weighted logarithmic-derivative prime-power sum: averaging over inertia
is its projector onto inertia invariants. In particular,
\(\Psi_K=\mathcal P(Q_0)\) and \(\Psi_L=\mathcal P(\rho)\).
Pointwise nonnegative class functions have nonnegative \(\mathcal P\).
This formulation includes ramified primes.

\begin{proposition}\label{prop:all-tensor-powers}\label{prop:kummer-tensor-powers}
For every integer \(m\geq1\),
\begin{equation}\label{eq:all-tensor-powers}
 \rho^{\otimes m}=u_m Q_0+v_m\rho,\qquad
 u_m=\frac{n^{m-1}+(-1)^m}{r},\qquad
 v_m=\frac{n^m-(-1)^m}{r}.
\end{equation}
The coefficients are nonnegative integers. Odd powers are not
pointwise nonnegative. Even powers give the bounds
\[
 \Psi_L\geq-\frac{u_{2k}}{v_{2k}}\Psi_K,\qquad
 \frac{u_{2k}}{v_{2k}}\searrow\frac1n.
\]
Consequently all even powers improve the lower bound from
\(\rho^{\otimes2}=Q_0+(r-2)\rho\), but give no upper bound.
Their limiting lower bound is already the positivity bound from
\(\operatorname{Reg}_G=Q_0+n\rho\), equivalently
\(\zeta_E=\zeta_K L(s,\rho)^n\).
\end{proposition}

\begin{proof}
At \(1\) and at a nonidentity element of \(N\), respectively, the
required equations are
\[
 n(u_m+v_m)=n^m,\qquad nu_m-v_m=(-1)^m.
\]
They have the displayed solution, which also agrees with the zero
values outside \(N\). Integrality follows from \(n\equiv-1\pmod r\);
nonnegativity is immediate, including \(u_1=0\).
At a nonidentity translation an odd power has value \(-1\).
For an even power every value is nonnegative, including each value
in the inertia average. Alternatively, writing
\(S=\rho^{\otimes k}\), its local trace is
\[
 \operatorname{tr}\bigl(\phi_\ell^j\mid
                  (S\otimes S^*)^{I_\ell}\bigr)
 =\frac1{|I_\ell|}\sum_{u\in I_\ell}
       |\operatorname{tr}S(\phi_\ell^j u)|^2\geq0.
\]
Finally
\[
 \frac{u_{2k}}{v_{2k}}
 =\frac1n+\frac{n+1}{n(n^{2k}-1)},
\]
which decreases to \(1/n\). The regular representation has character
\(rn\) at the identity and zero elsewhere, proving the last assertion.
\end{proof}

\begin{proposition}\label{prop:full-tensor-cone}
The pointwise nonnegative functions in
\(\operatorname{span}_{\mathbb R}\{Q_0,\rho\}\) form the cone
generated by \(\operatorname{Reg}_G\) and \(Q_0-\rho\).
All inequalities obtained by applying \(\mathcal P\) to this cone
are exactly
\begin{equation}\label{eq:full-tensor-interval}
 -\frac{\Psi_K}{n}\leq\Psi_L\leq\Psi_K.
\end{equation}
Moreover, after fixing the nonnegative mass outside \(N\), every
polynomial-square moment condition in \(1,Q_0,\rho\), of arbitrary
degree, is already implied by the positive semidefiniteness of a
single \(2\)-by-\(2\) Gram matrix.
In particular these conditions give no strict improvement of the
upper endpoint in \eqref{eq:full-tensor-interval}.
\end{proposition}

\begin{proof}
By \eqref{eq:twist-character-values}, \(aQ_0+b\rho\geq0\) precisely when
\(a+b\geq0\) and \(na-b\geq0\). The exact decomposition is
\[
 aQ_0+b\rho
 =\frac{a+b}{r}\operatorname{Reg}_G
  +\frac{na-b}{r}(Q_0-\rho).
\]
For the moment assertion set \(q=Q_0/n=1_N\). In the algebra of
functions supported on \(N\),
\[
 q^2=q,\qquad q\rho=\rho,\qquad
 \rho^2=nq+(n-1)\rho.
\]
Its two disjoint idempotents are
\[
 e_0=\frac{q+\rho}{r}=1_{\{1\}},
 \qquad
 e_1=\frac{nq-\rho}{r}=1_{N\setminus\{1\}}.
\]
Every polynomial supported on \(N\) reduces to its two values, and
the square of its modulus is a nonnegative linear combination of
\(e_0,e_1\). A general polynomial in \(1,q,\rho\) has the additional
orthogonal idempotent \(1-q\), whose nonnegative mass is part of the
fixed abelian marginal and does not involve \(\mathcal P(\rho)\).
Put \(M_0=\mathcal P(q)\) and
\(M_1=\mathcal P(\rho)\). The Gram matrix for \(q,\rho\) is
\[
 \begin{pmatrix}
 M_0&M_1\\
 M_1&nM_0+(n-1)M_1
 \end{pmatrix},
 \qquad
 \det=(nM_0-M_1)(M_0+M_1).
\]
For \(M_0>0\) it is positive semidefinite exactly when
\(-M_0\leq M_1\leq nM_0\); these inequalities make the other diagonal
entry nonnegative as well. For \(M_0=0\), positivity requires
\(M_1=0\), giving the same conclusion. Since \(\Psi_K=nM_0\),
this is \eqref{eq:full-tensor-interval}.
Conversely, these inequalities say
\(\mathcal P(e_0),\mathcal P(e_1)\geq0\); together with
\(\mathcal P(1-q)\geq0\), they imply every higher
polynomial-square condition. A positive measure supported at \(1\)
attains the upper endpoint. This last statement concerns the moment
constraints; it does not claim that an arbitrary prescribed measure
comes from the primes of a fixed number field.
\end{proof}

For example, the upper comparison between any even power \(m\geq4\)
and the second power is exactly a multiple of the existing detector:
\[
 n^{m-2}\rho^2-\rho^m
 =\frac{n^{m-2}-1}{r}(Q_0-\rho).
\]
Thus the full moment audit is more precise than saying that all powers
give the same inequality: they improve the lower endpoint, while
the upper endpoint is unchanged.

\subsubsection{Rational twists retain the nonabelian part and its conductor}

\begin{proposition}\label{prop:rational-twist-conductor}
Let \(\psi\) be any finite-order rational Dirichlet character.
If \(d=\operatorname{ord}_r(p)>1\), then its twist satisfies
\begin{equation}\label{eq:twist-conductor-floor}
 a_p(\rho\otimes\psi)\geq n,
 \qquad \mathfrak f(\rho\otimes\psi)\geq p^n,
\end{equation}
where \(a_p\) is the local Artin-conductor exponent and
\(\mathfrak f\) the positive integer finite conductor.
This holds even when \(p\) divides the conductor of \(\psi\).
If \(\psi\) is unramified at \(p\), equality holds in the first bound.
If \(d=1\), the weaker uniform bound \(a_p(\rho\otimes\psi)\geq n-1\)
holds.
\end{proposition}

\begin{proof}
At \(p\), the inertia image in \(G\) is \(N\); the decomposition group
is \(N\rtimes C_d\), and a Frobenius lift acts on \(N\) by
\(\tau\mapsto\tau^p\). If \(d>1\), the commutators of a Frobenius lift
with inertia map onto \(N\), because \(p-1\) is invertible modulo \(r\).
They lie in absolute inertia and are killed by every one-dimensional
character \(\psi\) of the absolute local Galois group.
The restriction of \(\rho\otimes\psi\) to these elements therefore
has the same invariant vectors as the restriction of \(\rho\) to
\(N\), which has none. Its inertia-invariant space is zero.
The conductor formula
\[
 a_p(V)=\dim V-\dim V^{I_p}+\operatorname{Swan}_p(V)
\]
gives \eqref{eq:twist-conductor-floor}. An unramified twist preserves
the tame representation at \(p\), giving exponent \(n\).

For \(d=1\), if \(\psi\) is nontrivial on the kernel of the inertia
action in \(N\), there are again no invariants. Otherwise it factors
through \(N\). The restriction
\(\rho|_N\) is the sum of the \(r-1\) distinct nontrivial characters
of \(N\); a scalar twist makes at most one summand trivial.
Thus the invariant dimension is at most one and the claimed
bound follows.
\end{proof}

\begin{proposition}\label{prop:twist-compositum-identity}
Let \(B=\mathbb Q(\zeta_M)\), and use complete Artin local factors
for the characters \(\psi\) of \(\operatorname{Gal}(B/\mathbb Q)\).
Then
\begin{equation}\label{eq:twist-compositum}
 \prod_{\psi}L(s,\rho\otimes\psi)
 =\frac{\zeta_{FB}(s)}{\zeta_B(s)}.
\end{equation}
Every character \(\eta\) of \(G/N\) is a self-twist of \(\rho\):
\(\rho\otimes\eta\simeq\rho\). Every diagonal Rankin--Selberg
product of a rational twist is likewise unchanged:
\[
 (\rho\otimes\psi)\otimes(\rho\otimes\psi)^*
 \simeq\rho\otimes\rho^*.
\]
Consequently these two averaging operations do not replace the
nonabelian factor by rational Dirichlet \(L\)-functions.
\end{proposition}

\begin{proof}
The commutator subgroup of \(G\) is \(N\); hence \(E\cap B\subseteq K\).
The subgroup \(\operatorname{Gal}(EB/B)\), acting on the roots
of \(X^r-p\), contains \(N\) and is transitive. Therefore
\([FB:B]=r\). The permutation representation attached to \(FB\)
is the tensor product of those attached to \(F\) and \(B\).
The former is \(1+\rho\), and the latter is
\(\bigoplus_\psi\psi\). Artin formalism, at every local factor, gives
\[
 \zeta_{FB}(s)
 =\prod_\psi L(s,\psi)L(s,\rho\otimes\psi),
\]
which proves \eqref{eq:twist-compositum}.
The self-twist statement follows immediately from the character
values: \(\eta=1\) on \(N\), and \(\rho=0\) outside \(N\).
The Rankin--Selberg identity follows by cancellation of
\(\psi\overline\psi\). In Hecke notation the general twist is
\[
 L(s,\rho\otimes\psi)
 =L_K\bigl(s,\chi\cdot(\psi\circ N_{K/\mathbb Q})\bigr).
\]
The norm-pullback factor is Galois invariant, so multiplication by it
does not shorten the full Galois orbit of the order-\(r\) character
\(\chi\). It remains an induced factor of degree \(n\).
\end{proof}

There is also an exact finite moment statement. Enlarge \(B\) to
contain \(K\). On an abelian Frobenius fiber whose multiplier in
\(\mathbb F_r^\times\) is \(1\), every rational twist has value
\(\psi(a)\rho(g)\), with the scalar \(\psi(a)\) fixed on that fiber.
The remaining values of \(\rho(g)\) are just \(n,-1\).
All mixed tensor-square conditions therefore reduce, on each such
fiber, to the two-idempotent algebra in
Proposition~\ref{prop:full-tensor-cone}.
On the other fibers \(\rho=0\).
Knowledge of the total mass in every abelian fiber, even exact
knowledge, does not force a positive mass at nonidentity translations:
all mass in a multiplier-\(1\) fiber may satisfy the moment constraints
at the value \(n\). This proves a limitation of these finite moment
conditions, not an impossibility of using additional analytic
relations between the twisted \(L\)-functions.

\subsubsection{A quantitative attempt using twist zero density}

The preceding conductor obstruction can be compared with a uniform zero-density theorem. Cyclic automorphic induction gives a cuspidal
representation \(\pi_\rho\) of \(\mathrm{GL}_n(\mathbb A_{\mathbb Q})\)
with \(L(s,\pi_\rho)=L(s,\rho)\); its cuspidality follows from the full
orbit of \(\chi\)
\citep[Chapter~III, \S6, Theorem~6.2 and Lemma~6.4]{ArthurClozel1989}.
Write \(C_\rho\) for its analytic conductor at height zero, in the
normalization of Lemke Oliver and Thorner.
Their Corollary~1.4 gives absolute constants \(c_1>0\) and \(C\geq1\)
such that
\begin{equation}\label{eq:LOT-twist-density}
 \sum_{\substack{q\leq T\\(q,\mathfrak f(\rho))=1}}
 \ \sum_{\psi\bmod q}^{*}
 N_{\pi_\rho\otimes\psi}(1-\delta,T)
 \leq C(C_\rho T)^{c_1n\delta},
 \qquad T\geq1,\quad 0\leq\delta\leq\tfrac12.
\end{equation}
Here the star restricts to primitive characters; \(N\) counts zeros
with real part at least \(1-\delta\) and imaginary part of absolute
value at most \(T\).
The source is \citet{LemkeOliverThorner2019}, author version dated
19~August~2017, Corollary~1.4 on printed/PDF page~5;
the analytic-conductor definition is equation~(2.3), page~12.
This theorem is unconditional. Applying rational Dirichlet GRH to
the abelian factors does not remove \(C_\rho\) from its stated bound.

In the hard case \(p^{r-1}\equiv1\pmod{r^2}\), the previously computed
discriminant is
\(\mathfrak f(\rho)=D_F=r^{r-2}p^{r-1}\).
The archimedean shifts are \(n/2\) copies of \(0\) and \(n/2\) copies
of \(1\). Thus, in that source's normalization,
\begin{equation}\label{eq:twist-analytic-conductor}
 C_\rho=12^{n/2}r^{r-2}p^{r-1},
 \qquad
 \log C_\rho
 =n\log p+(r-2)\log r+\tfrac n2\log12.
\end{equation}
In particular \(C_\rho\geq p^n\), a feature that persists under the
arbitrary rational twists in
Proposition~\ref{prop:rational-twist-conductor} when \(d>1\).

One possible use of \eqref{eq:LOT-twist-density} is to seek a twist
having no zero in the fixed finite strip by comparing the zero count
with the number of twists. Even allowing the optimistic upper bound
\(T^2\) for the family size, this pigeonhole certificate requires
\[
 C(C_\rho T)^{c_1n\delta}<T^2.
\]
Hence a necessary condition for this particular certificate is
\begin{equation}\label{eq:twist-density-scale}
 (2-c_1n\delta)\log T>
       c_1n\delta\log C_\rho+\log C.
\end{equation}
If \(c_1n\delta\geq2\), this condition cannot hold. Otherwise it
requires
\[
 \log T>
 \frac{c_1n\delta}{\,2-c_1n\delta\,}\log C_\rho.
\]
For fixed \(r,\delta>0\) this is a positive power of \(p\), not a
fixed power of \(\log p\). Conversely, for
\(T\leq(r\log p)^A\), a fixed \(A\), the available strip width for
this certificate is at most
\[
 \delta<
 \frac{2\log T}{c_1n(\log C_\rho+\log T)}
 \leq
 \frac{2A\log(r\log p)}{c_1(r-1)^2\log p}.
\]
These inequalities describe what the quoted bound can certify;
they do not assert the existence of zeros in any twist.
Moreover a zero-free estimate for one complex twist would still need
a separate untwisting argument to yield the positive auxiliary-prime
detector.

For completeness, the elementary finite mean square has the same
absence of a strict saving. Let
\(a_j=\Lambda_\rho(j)w(j/x)\), where \(w\) is bounded and supported
in \([0,2]\), and set
\[
 S_\psi=\sum_{\substack{j\leq2x\\(j,M)=1}}a_j\psi(j).
\]
Orthogonality and Cauchy--Schwarz give the exact identity and bound
\[
 \frac1{\varphi(M)}\sum_{\psi\bmod M}|S_\psi|^2
 =\sum_{a\in(\mathbb Z/M\mathbb Z)^\times}
       \left|\sum_{\substack{j\leq2x\\j\equiv a\bmod M}}a_j\right|^2
 \leq
 \left(1+\frac{2x}{M}\right)\sum_{j\leq2x}|a_j|^2.
\]
Since \(|\Lambda_\rho(j)|\leq n\Lambda(j)\), the elementary bound
\(\sum_{j\leq2x}\Lambda(j)^2\ll x\log x\) yields
\[
 |S_1|\ll_w
 n\sqrt{\varphi(M)\left(1+\frac{2x}{M}\right)x\log x}.
\]
This does not by itself force the strict upper margin needed for the
detector. Adding the exact polynomial-square constraints gives
precisely the endpoint in
Proposition~\ref{prop:full-tensor-cone}.
Thus this rational-twist attempt has an explicit residual requirement:
an analytic estimate must distinguish the two translation classes
inside an abelian Frobenius fiber, with conductor dependence stronger
than \eqref{eq:LOT-twist-density} at the requested polynomial-logarithmic
scale. No such estimate is established here.

The classical conductor-uniform zero-free theorem also preserves the dependence on $p$: applied over $K$ to the nonreal character $\chi$ and the trivial character, it gives width of order
\[
\bigl(\log D_F+\log D_K+r+(r-1)\log(|t|+3)\bigr)^{-1},
\]
with no exceptional real zero in this application
\citep[Appendix A, Theorem A.1, p.~12, and Remark A.4, p.~13]{HumphriesBrumley2019}.
For fixed $r$ its product with $O(\log\log p)$ tends to zero, so this estimate does not supply the fixed width used in Theorem~\ref{thm:fixedstrip}.

\subsection{An exact residual zero criterion in the difficult scalar case}
\label{sec:kummer-residual-zero}

Let $r$ be an odd prime, let $p\ne r$ be prime, and suppose
\[
 d=\operatorname{ord}_r(p),\qquad v_r(p^d-1)\geq2.
\]
Put $F=\mathbb Q(p^{1/r})$, $K=\mathbb Q(\zeta_r)$,
$E=K(p^{1/r})$, and
\[
 L(z)=\frac{\zeta_F(z)}{\zeta_{\mathbb Q}(z)}=L_K(z,\chi),
 \qquad Q_\chi=p^{r-1},
\]
where $\chi$ is any nontrivial character of
$\operatorname{Gal}(E/K)$, interpreted as a primitive Hecke character.
Call a rational prime $\ell$ auxiliary if
\[
 \ell\nmid pr,\qquad \ell\equiv1\pmod r,
 \qquad p^{(\ell-1)/r}\not\equiv1\pmod\ell.
\]
The following criterion concerns the zero sum of $L$.
It is an equivalent finite-scale detector, not a proof that the
required strict inequality has been derived with $A$ as the sole
new analytic input; the separate fixed-strip theorem is excluded from
this particular derivation.

\begin{lemma}[The local factors of the detector]
\label{lem:kummer-detector-local}
For $\operatorname{Re}z>1$,
\begin{equation}\label{eq:kummer-detector-euler}
 \frac{\zeta_K(z)}{L(z)}
 =(1-p^{-dz})^{-(r-1)/d}
   \prod_{\ell\ \mathrm{auxiliary}}
       \frac{1-\ell^{-rz}}{(1-\ell^{-z})^r}.
\end{equation}
In particular, the negative logarithmic derivative on the right has
nonnegative coefficients.  Apart from the displayed factor at $p$,
its support consists precisely of powers $\ell^k$ of auxiliary
primes for which $r\nmid k$.
\end{lemma}
\begin{proof}
The Galois group of $E/\mathbb Q$ is the affine group
$\mathbb F_r\rtimes\mathbb F_r^\times$, acting on the $r$
roots of $X^r-p$.  At an unramified prime, a Frobenius with
nontrivial multiplier of order $e$ has one fixed point and
$(r-1)/e$ cycles of length $e$.  Its local factor in $\zeta_F$
is therefore the product of the local factors in
$\zeta_K\zeta_{\mathbb Q}$.  For multiplier $1$, Frobenius
is either the identity, when the same cancellation occurs, or
a nontrivial translation, which is an $r$-cycle.  The latter
case is exactly the auxiliary-prime condition and gives the
displayed factor.  Equivalently, no power of an affine element
outside the nontrivial translations becomes a nontrivial
translation: if its multiplier $c\ne1$ satisfies $c^k=1$,
then its translation part is multiplied by
$1+c+\cdots+c^{k-1}=0$.

At $p$, the polynomial $X^r-p$ is Eisenstein, so $F$ has
one prime of residue degree $1$.  Its factor cancels that of
$\zeta_{\mathbb Q}$.  The field $K$ has $(r-1)/d$ primes
of residue degree $d$, giving the displayed factor at $p$.
At $r$, the assumption implies $p^{r-1}\equiv1\pmod{r^2}$.
The decomposition
$\mathbb Z_r^\times=\mu_{r-1}\times(1+r\mathbb Z_r)$
and the $r$-adic logarithm imply that $p$ is an $r$th power
in $\mathbb Q_r$.  Consequently
\[
 F\otimes\mathbb Q_r\simeq
 \mathbb Q_r\times\mathbb Q_r(\zeta_r).
\]
Both factors have residue degree $1$, so the local factors
of $\zeta_K\zeta_{\mathbb Q}/\zeta_F$ cancel also at $r$.
Finally, logarithmic differentiation gives
\[
 -\frac{d}{dz}\log\frac{\zeta_K(z)}{L(z)}
 =(r-1)\log p\sum_{m\geq1}p^{-dmz}
 +r\sum_{\ell\ \mathrm{auxiliary}}\log\ell
       \sum_{\substack{k\geq1\\r\nmid k}}\ell^{-kz}.
\]
All operations are justified by absolute convergence in
$\operatorname{Re}z>1$.
\end{proof}

Use the transforms and zero sums of
Section~\ref{sec:positive-window-tests}.  For an admissible
even real test $f$ supported on $[-S,S]$, set
\begin{align*}
 P_p(f)&=(r-1)\log p\sum_{m\geq1}p^{-dm/2}f(dm\log p),\\
 \mathcal E_f&=Z_L(f)-Z_K(f)-f(0)\log Q_\chi
                  +2\Phi_f(1)-2P_p(f).
\end{align*}

\begin{proposition}[Exact zero excess]
\label{prop:kummer-exact-zero-excess}
One has the unconditional identity
\begin{equation}\label{eq:kummer-exact-zero-excess}
 \mathcal E_f
 =2r\sum_{\ell\ \mathrm{auxiliary}}\log\ell
      \sum_{\substack{k\geq1\\r\nmid k}}
            \ell^{-k/2}f(k\log\ell).
\end{equation}
If $f\geq0$ and $f(t)>0$ for $0<t<S$, then
$\mathcal E_f\geq0$, and $\mathcal E_f>0$ if and only if
there is an auxiliary prime $\ell<e^S$.
\end{proposition}
\begin{proof}
The logarithmic derivative in
Lemma~\ref{lem:kummer-detector-local} shows that
$T_K(f)+T_{\mathbb Q}(f)-T_F(f)$ equals $P_p(f)$ plus
the auxiliary-prime sum on the right of
\eqref{eq:kummer-exact-zero-excess}, divided by $2$.
Substitute this in \eqref{eq:detector-positive-test}.
The positivity statement follows term by term.  If the sum
is positive, some power $\ell^k<e^S$ occurs, hence
$\ell<e^S$.  Conversely such a prime contributes a strictly
positive $k=1$ term.  A prime at $e^S$ contributes zero
because continuity and compact support imply $f(S)=0$.
\end{proof}

For the strip-positive triangle $f_S=e^{-|t|/2}(1-|t|/S)_+$,
$S\geq1$, this becomes
\begin{align}
 P_p(S)&=(r-1)\log p\sum_{m\geq1}p^{-dm}
                  \left(1-\frac{dm\log p}{S}\right)_+,
                  \label{eq:kummer-ramified-correction}\\
 \mathcal E(S)&=Z_L(f_S)-Z_K(f_S)-(r-1)\log p
           +S+2-\frac{2(1-e^{-S})}{S}-2P_p(S)
                  \label{eq:kummer-triangle-excess}\\
 &=2r\sum_{\ell\ \mathrm{auxiliary}}\log\ell
       \sum_{\substack{k\geq1\\r\nmid k}}
          \frac1{\ell^k}
          \left(1-\frac{k\log\ell}{S}\right)_+.
                  \nonumber
\end{align}
The ramified correction is uniformly small:
\begin{equation}\label{eq:kummer-ramified-small}
 0\leq P_p(S)\leq\frac{(r-1)\log p}{p^d-1}
                          \leq\frac{3\log r}{r}.
\end{equation}
Indeed $B=p^d\geq r^2+1$, and $\log B/(B-1)$ decreases
for $B>1$.  Thus the middle expression is at most
$(r-1)\log(r^2+1)/r^2\leq3\log r/r$.
There is also a quantitative witness to strictness.  If an
auxiliary prime satisfies $\ell\leq X$, $X\geq2$, its first
power contributes at $S=\log(2X)$ at least
\begin{equation}\label{eq:kummer-positive-gap}
 \mathcal E(\log(2X))
 \geq \frac{2r(\log2)^2}{X\log(2X)}.
\end{equation}
Here $1/\ell\geq1/X$, $\log\ell\geq\log2$, and
$1-\log\ell/\log(2X)\geq\log2/\log(2X)$.

Consequently a fixed polynomial auxiliary-prime bound is
equivalent to a strict inequality
\[
 Z_L(f_S)>
 (r-1)\log p+Z_K(f_S)-2\Phi_{f_S}(1)+2P_p(S)
\]
at some $S=O(\log(r+\log p+3))$, with a uniform implied
constant.  Assumption $A$ concerns the rational Dirichlet
factors of $\zeta_K$ and hence the abelian sum $Z_K$.
The uncompleted step of this particular derivation is a lower bound
for the zero sum $Z_L$ with a strict excess over the displayed
threshold.  Nonnegativity of $Z_L$ alone is much weaker.

\subsubsection{A quantitative attempt using the high zeros}

The high zeros do yield an unconditional lower bound for
$Z_L(f_S)$, with all conductor dependence visible.
For $0\leq a\leq1$, let
\[
 A_{a,S}(\gamma)=\int_0^S(1-t/S)e^{-at}\cos(\gamma t)\,dt.
\]
For $|\gamma|\geq3$ and $S\geq1$,
\begin{equation}\label{eq:kummer-high-zero-kernel}
 \operatorname{Re}\Phi_{f_S}(\beta+i\gamma)
 = A_{\beta,S}(\gamma)+A_{1-\beta,S}(\gamma)
 \geq\frac1{1+\gamma^2}
 \qquad(0\leq\beta\leq1).
\end{equation}
To check this directly, put $z=a-i\gamma$ and integrate:
\[
 A_{a,S}(\gamma)
 =\operatorname{Re}\left(
       \frac1z-\frac{1-e^{-Sz}}{Sz^2}\right).
\]
Writing $v=e^{-aS}$, the real numerator of
$(1-e^{-Sz})/z^2$ is
\[
 (a^2-\gamma^2)(1-v\cos(S\gamma))
                  +2a\gamma v\sin(S\gamma).
\]
For $a>0$ it is at most
$av[-S(\gamma^2-a^2)+2|\gamma|]\leq0$;
use $1-v\geq aSv$ and $|\gamma|\geq3$.
The case $a=0$ is immediate.  Hence
$A_{a,S}(\gamma)\geq a/(a^2+\gamma^2)$.
If $u=\beta(1-\beta)$ and $t=\gamma^2$, the numerator of
\[
 \frac{\beta}{\beta^2+t}
 +\frac{1-\beta}{(1-\beta)^2+t}-\frac1{1+t}
\]
over its positive common denominator is
$u(1-u+3t)\geq0$.  This proves
\eqref{eq:kummer-high-zero-kernel}.

Let $n=r-1$ and $N_L(T)=\#\{\rho_L:|\operatorname{Im}\rho_L|\leq T\}$,
with multiplicity.  The identity $\zeta_F=\zeta_{\mathbb Q}L$
and the explicit zero count of
\citet[Corollary~1.2, p.~278]{HSW2022} give, for $T\geq1$,
\begin{align}
 N_L(T)&=\frac T\pi\log D_F
             +\frac{nT}\pi\log\frac{T}{2\pi e}+E_L(T),
             \label{eq:kummer-zero-count}\\
 |E_L(T)|&\leq B(T):=
 0.228\log D_F+0.228(r+1)\log T
                +23.108(r+1)+9.040.\nonumber
\end{align}
For any $T\geq3$, positivity of the omitted low-zero terms
and partial summation therefore imply
\begin{align}
 Z_L(f_S)\geq{}&
 \frac{\log D_F}{\pi}\arctan(1/T)
 +\frac n\pi\int_T^\infty
          \frac{\log(t/(2\pi))}{1+t^2}\,dt
 \label{eq:kummer-high-zero-bound}\\
 &-\frac{2B(T)}{1+T^2}
 -0.114(r+1)\log(1+T^{-2}).\nonumber
\end{align}
For completeness, with $w(t)=(1+t^2)^{-1}$, the error in
$\int_{(T,\infty)}w\,dN_L$ is bounded by
\[
 w(T)B(T)+\int_T^\infty(-w'(t))B(t)\,dt
 =\frac{2B(T)}{1+T^2}
   +0.114(r+1)\log(1+T^{-2}).
\]
The formula remains valid when a zero has height exactly
$T$, because the integral uses $(T,\infty)$ and $N_L(T)$
includes the endpoint.

This is a positive conductor contribution, but it does not
prove the required excess.  Before its error is subtracted,
the coefficient of $\log D_F$ furnished by
\eqref{eq:kummer-high-zero-bound} is at most
$\pi^{-1}\arctan(1/3)<0.103$.
For bounded $T$, the zero-counting error also has order
$\log D_F$.  In particular, this estimate does not cancel
the coefficient $1$ of $\log Q_\chi$ to an error of size
$O(\log r+\log\log p)$, even with $r$ fixed.
It is weaker than the exact non-strict inequality already
contained in \eqref{eq:kummer-exact-zero-excess}.

One may try to recover more conductor mass with a test
concentrated close to zero in the logarithmic variable.
The following exact observation isolates what that recovery
does.  If $g$ is any admissible even test supported strictly
inside $(-\log2,\log2)$, no prime-ideal power enters its
explicit formula.  In particular $P_p(g)=0$, and
\begin{equation}\label{eq:kummer-small-support-null}
 Z_L(g)-Z_K(g)=g(0)\log Q_\chi-2\Phi_g(1),
 \qquad \mathcal E_g=0.
\end{equation}
Thus adding such a narrow logarithmic test can recover a
known conductor contribution to the zero sums, but changes
the residual excess by exactly zero.  This observation does
not exclude other transforms or a lower bound using
arithmetic information about the low zeros.  It
identifies why conductor recovery by itself does not establish
strictness at the desired finite scale.

\subsection{An almost-all auxiliary-prime bound from rational Dirichlet GRH}

The averaging variable in this subsection is the input prime $p$.
The conclusion does not give a bound for every individual input.
All logarithms are natural.

\begin{lemma}\label{lem:dirichlet-log-weight}
Assume A.  There is an absolute constant $C_0$ such that, for every
Dirichlet character $\chi$ modulo $q\ge1$ and every $x\ge2$,
\[
 \left|\operatorname{Re}\sum_{n\le x}\Lambda(n)\chi(n)
                 \log(x/n)-\delta_\chi x\right|
 \le C_0\sqrt{x}\log(2q),
\]
where $\delta_\chi=1$ for the principal character and is zero otherwise.
\end{lemma}

\begin{proof}
We include the uniformity argument because the modulus below can be
larger than any fixed power of $x$.
For a primitive nonprincipal character of conductor $q\ge3$, use
Lemma~2.2 and the real-part identity in the corrected author version
of \citet[corrected author pp.~5--7, Lemmas~2.1--2.3]{LLS2015}.  Their Lemma~2.3, evaluated at the fixed value
$x_0=100$, gives
\[
 |\operatorname{Re}B(\chi)|\ll\log(2q).
\]
Indeed the denominator in that identity is at least
$1-2/10+1/100=81/100$, and its finite prime sum and elementary
terms are bounded absolutely.  Lemma~2.2 therefore gives the asserted
bound: its remaining terms are
$O(\log(2q)\log x+(\log x)^2+1)$, which are absorbed in
$O(\sqrt{x}\log(2q))$.  The primitive principal character is
the Riemann zeta case, covered by their Lemma~2.1.

If $\chi$ is induced by a primitive character $\chi^*$, the two sums
differ only at prime powers whose prime divides $q$.  For each such
prime $v$ the absolute difference is at most
\[
 \log v\sum_{1\le j\le\log x/\log v}\log(x/v^j)
 \le(\log x)^2.
\]
There are at most $\log(2q)/\log2$ such primes.  Finally
$(\log x)^2\ll\sqrt{x}$ uniformly for $x\ge2$.
This proves the assertion, including imprimitive principal characters,
with an absolute constant and without a restriction on $q$.
\end{proof}

\begin{proposition}\label{prop:auxiliary-character-sieve}
Let $r$ be an odd prime, $B\ge2$, and put
\[
 \mathcal L=\{\ell\le B:\ell\text{ prime},\ \ell\equiv1\pmod r\},
 \qquad m=|\mathcal L|,\qquad M=\prod_{\ell\in\mathcal L}\ell.
\]
For an integer $n$, let $G(n)$ be one if $n$ is a nonzero $r$th
power modulo every $\ell\in\mathcal L$, and zero otherwise.
The empty product is interpreted as one.  Under A, uniformly in
$r,B,x\ge2$,
\begin{equation}\label{eq:auxiliary-sieve-log-sum}
 \sum_{n\le x}\Lambda(n)G(n)\log(x/n)
 =\frac{x}{r^m}+O\!\left(\sqrt{x}\log(2M)\right),
\end{equation}
with an absolute implied constant.

For $P\ge4$, let $E_r(P,B)$ count primes $p\in(P/2,P]$ with
$p\nmid M$ for which no prime $\ell\le B$ satisfies
\[
 \ell\equiv1\pmod r,
 \qquad p^{(\ell-1)/r}\not\equiv1\pmod\ell.
\]
Then
\begin{equation}\label{eq:auxiliary-exception-bound}
 E_r(P,B)\ll
 \frac{P/r^m+\sqrt P\log(2M)}{\log P},
\end{equation}
again with an absolute implied constant.

For the cubic hard congruence, define
$G^-(n)=1_{n\equiv-1\pmod9}G(n)$ with $r=3$, and put $Q=9M$.
The corresponding exact main term is
\begin{equation}\label{eq:cubic-hard-sieve-log-sum}
 \sum_{n\le x}\Lambda(n)G^-(n)\log(x/n)
 =\frac{x}{6\cdot3^m}+O\!\left(\sqrt{x}\log(2Q)\right).
\end{equation}
In particular, for exceptional primes $p\equiv-1\pmod9$ the
$P/r^m$ term in \eqref{eq:auxiliary-exception-bound} can be replaced
by $P/(6\cdot3^m)$.  Such primes never divide $Q$.
\end{proposition}

\begin{proof}
For each $\ell\in\mathcal L$, choose a Dirichlet character
$\eta_\ell$ of exact order $r$ modulo $\ell$.  Extend every character,
including the principal character $\eta_\ell^0$, by zero on
nonunits.  Orthogonality in the cyclic quotient of order $r$ gives
\[
 1_{\{n\ne0\text{ and }n\text{ an }r\text{th power}\bmod\ell\}}
 =\frac1r\sum_{j=0}^{r-1}\eta_\ell^j(n).
\]
Consequently
\[
 G(n)=r^{-m}\sum_{\boldsymbol j\in\{0,\ldots,r-1\}^m}
          \prod_{\ell\in\mathcal L}\eta_\ell^{j_\ell}(n).
\]
These are distinct characters modulo $M$, by the Chinese remainder
theorem, and exactly one is principal.  All coefficients are real and
their absolute values sum to one.  Apply
Lemma~\ref{lem:dirichlet-log-weight} to each character and take real
parts.  The total error is bounded by the largest individual error;
it is not multiplied by $r^m$.  This proves
\eqref{eq:auxiliary-sieve-log-sum}.

Every prime counted by $E_r(P,B)$ has $G(p)=1$ and contributes at
least $\log(P/2)\log2$ to the nonnegative sum with $x=2P$.
The same sum can contain additional primes and prime powers, all
with nonnegative coefficients.  Thus no prime-power subtraction is
needed, and the first bound follows directly from
\eqref{eq:auxiliary-sieve-log-sum}.

For the cubic refinement, expand
\[
 1_{n\equiv-1\pmod9}
   =\frac16\sum_{\psi\bmod9}\overline{\psi(-1)}\psi(n).
\]
Every coefficient $\overline{\psi(-1)}$ is real, equal to $1$ or
$-1$.  Since $(9,M)=1$, multiplying the two expansions gives
$6\cdot3^m$ distinct characters modulo $Q$, exactly one principal,
with coefficient $1/(6\cdot3^m)$.  Their coefficient absolute values
again sum to one.  The same argument proves
\eqref{eq:cubic-hard-sieve-log-sum} and the refined exception bound.
A prime $p\equiv-1\pmod9$ is neither $3$ nor any
$\ell\equiv1\pmod3$, which proves the final assertion.
\end{proof}

\begin{proof}[Proof of Theorem~\ref{thm:fixed-r-almost-all-auxiliary}]
The prime number theorem in the fixed progression $1\pmod r$ gives
\[
 m=\frac{B}{(r-1)\log B}(1+o(1))
   =\left(\frac{C_r}{r-1}+o(1)\right)\log P.
\]
The strict inequality for $C_r$ implies $m\ge\log P/(2\log r)$
for all sufficiently large $P$, so $r^m\ge\sqrt P$.
The elementary Chebyshev bound gives
$\log M\le\sum_{\ell\le B}\log\ell\ll B$.
Also $B<P/2$ eventually, so no prime in the interval divides $M$.
Substitution into \eqref{eq:auxiliary-exception-bound} proves the
asserted bound.  For $r=3$, $1>1/\log3$.

Enumerating primes at most $B$ and testing
$p^{(\ell-1)/r}\not\equiv1\pmod\ell$ use a number of bit
operations polynomial in $B$, $r$ and $\log p$.  A successful test is
a certificate independent of A.  The auxiliary-prime reduction then
uses an explicitly represented cyclotomic algebra of dimension
$\ell-1$, so its established running time is polynomial in these
parameters.  Since $B=O_{r,C_r}(\log p\log\log p)$ on the dyadic
interval, this gives the stated algorithmic conclusion.
\end{proof}

\begin{corollary}\label{cor:uniform-r-almost-all-auxiliary}
There is a sufficiently large absolute constant $C$ such that,
under A, for every odd prime $r$ and every $P\ge4$, setting
\[
 L=\log(Pr),\qquad B=C r^2L^2
\]
gives $m\ge\log P/\log r$.  The number of hard primes
$p\in(P/2,P]$ having no good auxiliary prime at most $B$ is
\[
 O\!\left(\frac{\sqrt P\,r^2\log^2(Pr)}{\log P}\right),
\]
with an absolute implied constant.  Relative to the number of all
prime inputs in the interval, this is negligible in ranges where
$r^2\log^2(Pr)=o(\sqrt P)$.  No uniform relative-density
assertion within the thinner set of hard inputs is made here.
On every successful input the auxiliary-prime search and field
construction take time polynomial in $r$ and $\log p$, uniformly
in these parameters.
\end{corollary}

\begin{proof}
We first obtain a uniform progression count with a positive weight.
Let
\[
 S(y;r,1)=\sum_{\substack{n\le y\\n\equiv1\bmod r}}
                 \Lambda(n)\log(y/n).
\]
Character orthogonality and Lemma~\ref{lem:dirichlet-log-weight}
give $S(y;r,1)=y/(r-1)+O(\sqrt y\log(2r))$.
Form
\[
 T(B)=S(B;r,1)-2S(B/2;r,1)+S(B/4;r,1).
\]
The coefficient of $\Lambda(n)$ is zero for $n\le B/4$, equals
$\log(4n/B)$ for $B/4<n\le B/2$, and equals $\log(B/n)$
for $B/2<n\le B$.  It is therefore nonnegative and at most
$\log2$.  Also
\[
 T(B)=\frac{B}{4(r-1)}+O(\sqrt B\log(2r)).
\]
The total contribution of proper prime powers is $O(\sqrt B)$:
Chebyshev's bound gives
$\sum_{k\ge2}\vartheta(B^{1/k})\ll
\sqrt B+B^{1/3}\log B\ll\sqrt B$.
For an absolute $C_1$, if $B\ge C_1r^2\log^2(2r)$, it follows
that the prime part of $T(B)$ is at least $B/[8(r-1)]$.
As each such prime contributes at most $(\log2)\log B$,
\[
 m\ge\frac{B}{8(r-1)(\log2)\log B}.
\]
Now take $B=C r^2L^2$ with $L=\log(Pr)$.
The preceding size condition holds for all $P\ge4$, $r\ge3$,
when the absolute $C$ is sufficiently large.
Moreover $\log B\le c_C L$, where $c_C=O(1+\log C)$.
Thus, after increasing the same absolute $C$ if needed,
\[
 m\ge\frac{C rL}{8(\log2)c_C}
       \ge\frac{\log P}{\log r}.
\]
The elementary bound $\log M\ll B$ and
\eqref{eq:auxiliary-exception-bound} now prove the conclusion.
Every hard prime has $\operatorname{ord}_r(p)>1$, whereas every
prime factor of $M$ is $1\pmod r$.  Thus the exclusion $p\nmid M$
causes no loss, even when $B\ge P/2$.
\end{proof}

The sieve argument averages fixed rational Dirichlet characters over
the varying input prime.  It proves neither a pointwise upper bound
for the least good auxiliary prime nor the required positive
nonabelian prime-detection gap for an arbitrarily specified $p$.

\section{Shifts, norms, and rational twists at a specified characteristic}
\label{sec:fixed-input-uniformization}

Throughout this section \(r\geq3\) and \(p\ne r\) are primes,
\(K=\mathbb Q(\zeta_r)\), \(E=K(p^{1/r})\), \(F=\mathbb Q(p^{1/r})\),
and \(n=r-1\). An auxiliary prime means a prime
\(\ell\equiv1\pmod r\) for which \(p\) is not an \(r\)-th power modulo
\(\ell\). The stronger progression \(1\bmod 12r\) used in the algebraic
factorization interface is not imposed in the statements below.
The purpose is to examine specific ways of transferring an average
over characteristics to a specified \(p\). None of the following
restrictions rules out every possible use of rational Dirichlet GRH.

\subsection{Which rational substitutions preserve the Kummer extension?}

\begin{proposition}\label{prop:rational-kummer-orbit}
For \(a\in\mathbb Q^\times\),
\[
 K(a^{1/r})=E
 \quad\Longleftrightarrow\quad
 a=p^j c^r\quad
 (1\leq j\leq r-1,\ c\in\mathbb Q^\times).
\]
In particular, the only positive rational prime in this orbit is \(p\).
The number of positive integers in this orbit not exceeding \(X\) is
\[
 \sum_{j=1}^{r-1}
 \left\lfloor (X/p^j)^{1/r}\right\rfloor
 \leq (r-1)(X/p)^{1/r}.
\]
\end{proposition}
\begin{proof}
Put \(t=p^{1/r}\), and choose the generator \(\tau\) of
\(\operatorname{Gal}(E/K)\) with \(\tau(t)=\zeta_r t\).
If \(\alpha^r=a\) generates \(E/K\), then
\(\tau(\alpha)/\alpha=\zeta_r^j\) for some \(1\leq j<r\).
The eigenspaces of \(\tau\), in the basis \(1,t,\ldots,t^{r-1}\),
are \(Kt^j\). Thus \(\alpha=bt^j\), where \(b\in K^\times\), and
\(z:=a/p^j=b^r\in\mathbb Q^\times\).
Taking norms to \(\mathbb Q\) gives
\((N_{K/\mathbb Q}b)^r=z^{r-1}\).
Consequently \(c=z/N_{K/\mathbb Q}b\) is rational and \(c^r=z\).
The reverse implication follows because \(j\) is coprime to \(r\).
If \(a\) is a positive integer, its valuations force \(c\) to be
a positive integer: outside \(p\) this is immediate, and at \(p\)
the inequality \(j+r v_p(c)\geq0\), with \(0<j<r\), implies
\(v_p(c)\geq0\). The representations have distinct \(p\)-adic
valuations modulo \(r\), which proves the count and the assertion
about primes.
\end{proof}

This proposition prevents the prime exceptional-set estimate from
being applied to multiple prime radicands defining this same extension.
It makes no assertion about transformations that change the extension
and provide some additional, separately proved transport map.

\begin{proposition}\label{prop:prime-norm-exclusion}
Let \(q\notin\{p,r\}\) be prime and \(q\equiv1\pmod r\).
If \(v_q(N_{F/\mathbb Q}\alpha)\) is not divisible by \(r\), for
some \(\alpha\in F^\times\), then \(p\) is an \(r\)-th power modulo \(q\).
In particular, a prime \(q=|N_{F/\mathbb Q}\alpha|\) in this
progression is not an auxiliary prime.
\end{proposition}
\begin{proof}
If \(p\) is not an \(r\)-th power modulo \(q\), Frobenius acts
transitively on the \(r\) roots of \(X^r-p\) over \(\mathbb F_q\):
it acts by a nonidentity translation in the group of \(r\)-th roots
of unity. Hence \(X^r-p\) is irreducible modulo \(q\).
Since \(q\nmid r p\), the prime is unramified and inert in \(F\),
with residue degree \(r\). The valuation formula for norms then gives
\(v_q(N\alpha)=r v_{\mathfrak q}(\alpha)\), also for fractional
\(\alpha\). This proves the contrapositive.
\end{proof}

\subsection{The coefficient cost of an additive shift}

The sparse multiplicative character expansion used to count exceptional
characteristics has coefficient norm \(1\). A nonzero additive shift
does not preserve that property.

\begin{proposition}\label{prop:shift-fourier-cost}
Let \(\ell\) be prime, \(r\mid\ell-1\), and \(\ell\geq2r+1\).
Put \(H=(\mathbb F_\ell^\times)^r\), fix \(c\in\mathbb F_\ell^\times\),
and define \(f_c(u)=1_H(u+c)\) for \(u\in\mathbb F_\ell^\times\).
For a multiplicative character \(\psi\), set
\[
 \widehat f_c(\psi)=
 \frac1{\ell-1}\sum_{u\in\mathbb F_\ell^\times}
 f_c(u)\overline{\psi(u)},\qquad
 \mu=\frac1r-\frac{1_H(c)}{\ell-1}.
\]
Then
\begin{align*}
 \widehat f_c(1)&=\mu,\\
 |\widehat f_c(\psi)|
 &\leq\frac{1+(r-1)\sqrt\ell}{r(\ell-1)}
 \quad(\psi\ne1),\\
 \frac{\sqrt\ell}{8r}
 &\leq \sum_\psi|\widehat f_c(\psi)|
 \leq \sqrt{(\ell-1)/r}.
\end{align*}
For a squarefree product \(M=\prod_{\ell\in\mathcal L}\ell\)
of such primes, the coefficient norm of
\[
 u\longmapsto
 \prod_{\ell\in\mathcal L}1_{(\mathbb F_\ell^\times)^r}(u+c)
 \quad\text{on }(\mathbb Z/M\mathbb Z)^\times
\]
is the product of the local coefficient norms. In particular it is
at least
\[
 \prod_{\substack{\ell\in\mathcal L\\\ell\nmid c}}
 \frac{\sqrt\ell}{8r}.
\]
An unshifted local factor has coefficient norm exactly \(1\).
\end{proposition}
\begin{proof}
Let \(\chi\) have order \(r\), and extend every character, including
the principal character, by zero at zero. Then
\[
 1_H(v)=\frac1r\sum_{a=0}^{r-1}\chi^a(v).
\]
There are \((\ell-1)/r-1_H(c)\) allowed values of \(u\), giving
the formula for \(\mu\). If \(\psi\ne1\), the term \(a=0\) in the
Fourier coefficient has numerator of absolute value \(1\).
For \(a\ne0\), substituting \(u=-cv\) gives the numerator
\[
 \chi^a(c)\overline{\psi(-c)}
 J(\overline\psi,\chi^a).
\]
Its absolute value is \(1\) when \(\overline\psi\chi^a=1\), and
\(\sqrt\ell\) otherwise. These Jacobi-sum identities follow from
\(|G(\eta)|^2=\ell\) for nonprincipal \(\eta\),
\(G(\alpha)G(\beta)=J(\alpha,\beta)G(\alpha\beta)\), and
\(J(\alpha,\overline\alpha)=-\alpha(-1)\).
They give the stated maximum coefficient bound.

Parseval and \(f_c^2=f_c\) give
\[
 \sum_{\psi\ne1}|\widehat f_c(\psi)|^2=\mu(1-\mu).
\]
Here \(\mu\geq1/(2r)\) and \(1-\mu\geq1/2\), while
the maximum nonprincipal coefficient is at most
\(\sqrt\ell/(\ell-1)\).
The inequality
\(\sum |z_i|\geq\sum|z_i|^2/\max|z_i|\) now gives the
lower bound. Cauchy--Schwarz and Parseval give the upper bound.
Finally, the Chinese remainder theorem makes both the Fourier
transform and its coefficient norm factor over \(\ell\).
When \(c=0\) locally, the original expansion has \(r\)
coefficients equal to \(1/r\).
\end{proof}

Thus an estimate obtained by applying individual Dirichlet bounds
and the triangle inequality to a shifted exceptional-set indicator
must include this coefficient norm. For example, each changed
local factor with \(\ell\geq256r^2\) contributes a factor at least
\(2\). This is a statement about the bound furnished by that method,
not a lower bound for the prime-counting error.

\subsection{A large family of twists with finite-height zero-free strips}

There is a positive zero-density consequence for characters whose
modulus depends on \(p\). It still requires an untwisting argument
to say anything about the specified untwisted function.

\begin{proposition}\label{prop:p-specific-twist-pool}
Suppose \(d=\operatorname{ord}_r(p)>1\) and
\(v_r(p^d-1)\geq2\). Let
\(\rho=\operatorname{Ind}_{K}^{\mathbb Q}\chi\), where \(\chi\)
is a nontrivial character of \(\operatorname{Gal}(E/K)\).
There is an absolute effective constant \(c_5\geq1\) such that,
with \(\delta=1/(4c_5n)\), all but
\[
 O\bigl((pT)^{1/4}/\sqrt n\bigr)
\]
of the \(p-1\) Dirichlet characters \(\psi\bmod p\) have
\(L(s,\rho\otimes\psi)\) free of zeros in
\[
 \Re s>1-\delta,\qquad |\Im s|\leq T.
\]
This holds unconditionally for every \(T\geq1\).
In particular, for \(T\leq p\) the exceptional number is
\(O(\sqrt p/\sqrt n)\).
\end{proposition}
\begin{proof}
Put \(B=\mathbb Q(\zeta_p)\). Since \(r\nmid p-1\), the extensions
\(E/K\) and \(KB/K\) have coprime degrees. Thus
\(\operatorname{Gal}(EB/K)=C_r\times\mathbb F_p^\times\).
The hard-case congruence implies
\(p^{r-1}\equiv1\pmod{r^2}\), so \(p\) is an \(r\)-th power
in \(\mathbb Q_r^\times\). Hence \(E/K\) is unramified at \(r\);
its only ramification is tame and total at the primes above \(p\).
Each \(\chi^j\), \(j\ne0\), therefore has conductor \(p\mathcal O_K\).

The characters associated with \(EB/K\) are
\(\chi^j(\psi\circ N_{K/\mathbb Q})\).
At a prime above \(p\), the two factors have respective orders
dividing \(r\) and \(p-1\). Their nontrivial inertia characters
cannot cancel when \(j\ne0\). All these conductors divide
\(p\mathcal O_K\), and the maximum conductor norm is \(p^n\).
The \(n\) characters with a fixed \(\psi\) and \(j\ne0\)
have the same \(L\)-function, namely \(L(s,\rho\otimes\psi)\):
the Galois action on \(K/\mathbb Q\) permutes them transitively
and fixes the norm pullback.

The log-free density estimate of Thorner and Zaman
\citep[Theorem~3.2, printed p.~1051; equation~(3-1), p.~1050]
{ThornerZaman2019} gives
\[
 \sum_\eta N(1-\delta,T,\eta)
 \ll
 \bigl(r^{r-2}p^n n^n T^n\bigr)^{c_5\delta}.
\]
The optional exceptional-zero factor in that estimate is at most
\(1\) and has been discarded. A bad \(\psi\) contributes at least
\(n\) to the left-hand side. Substituting the specified \(\delta\)
and using \(r=n+1\) proves
\[
 \#\{\text{bad }\psi\}
 \ll \frac{r^{(r-2)/(4n)}n^{1/4}(pT)^{1/4}}{n}
 \ll \frac{(pT)^{1/4}}{\sqrt n}.
\]
\end{proof}

The conclusion does not identify a good character in deterministic
time polynomial in \(r\) and \(\log p\), and it need not include
the principal character. The next calculation concerns the additional
problem of transferring bounds from such twists to the untwisted sum.

\subsection{An absolute-error untwisting calculation}

\begin{lemma}\label{lem:complex-dirichlet-tent}
Assume rational Dirichlet GRH. There is an absolute constant \(C\)
such that, for every character \(\xi\bmod q\) and \(x\geq2\),
\[
 \left|\sum_{m\leq x}\Lambda(m)\xi(m)\log(x/m)
       -1_{\xi\ {\rm principal}}x\right|
 \leq C\sqrt x\log(2q).
\]
This is an absolute complex estimate, not only a bound on its real
part.
\end{lemma}
\begin{proof}
For primitive nonprincipal \(\xi\), use the Hadamard product
and the full complex identity in
\citet[equation~(2.3) and Lemma~2.2, corrected author p.~5]{LLS2015}.
Writing \(B_\xi\) for the logarithmic derivative at zero of the
completed function, logarithmic differentiation at \(2\) gives
\[
 B_\xi=\frac{\Lambda_\xi'}{\Lambda_\xi}(2)
       -\sum_\varrho\left(\frac1{2-\varrho}+\frac1\varrho\right).
\]
The local zero count
\[
 \#\{\varrho:|\Im\varrho-t|\leq1\}
 \ll \log(2q)+\log(|t|+3)
\]
is the \(K=\mathbb Q\) case of
\citet[Lemma~2.5, printed p.~1050]{ThornerZaman2019}.
Under GRH it implies
\(\sum_\varrho(1+|\Im\varrho|^2)^{-1}\ll\log(2q)\).
Since
\[
 \frac1{2-\varrho}+\frac1\varrho
 =\frac2{\varrho(2-\varrho)}
 \ll \frac1{1+|\Im\varrho|^2},
\]
and the completed logarithmic derivative at \(2\) is
\(O(\log(2q))\), we have \(|B_\xi|\ll\log(2q)\).
Lemma~2.2 cited above now bounds the smoothed sum by
\[
 O\bigl(\sqrt x\log(2q)+
       \log x\log(2q)+(\log x)^2+1\bigr),
\]
which has the claimed size. The principal primitive case follows
from the zeta-function identity in their Lemma~2.1.
Passing from a primitive character to a character modulo \(q\)
removes Euler factors at primes dividing \(q\); their total
contribution is at most
\(\omega(q)(\log x)^2\ll\sqrt x\log(2q)\).
\end{proof}

\begin{proposition}\label{prop:absolute-untwisting-cost}
Assume rational Dirichlet GRH, and let \(2\leq x<p\).
Define nonnegative weights
\[
 w_m=n\Lambda(m)1_{m\equiv1\bmod r}\log(x/m),\qquad
 S_0=\sum_{m\leq x}w_m.
\]
There is an absolute \(C_A>0\) such that
\[
 S_0=x+O(n\sqrt x\log(2r)),\qquad
 \left|\sum_{m\leq x}w_m\psi(m)\right|
 \leq C_A n\sqrt x\log(2pr)
 \quad(\psi\bmod p,\ \psi\ne1).
\]
Suppose a linear combination
\(f(m)=\sum_{\psi\ne1}c_\psi\psi(m)\) satisfies
\[
 \sum_{m\leq x}w_m|1-f(m)|\leq\varepsilon S_0,
 \qquad 0\leq\varepsilon<1.
\]
Then its coefficient norm \(U=\sum_{\psi\ne1}|c_\psi|\) obeys
\begin{equation}\label{eq:untwisting-l1-necessary}
 U\geq
 \frac{(1-\varepsilon)S_0}
 {C_A n\sqrt x\log(2pr)}.
\end{equation}
If the auxiliary-prime window up to \(x\) is empty, these same
sums are exactly the twists of the \(\rho\) prime-power sum,
with the Euler factors at \(rp\) removed.
\end{proposition}
\begin{proof}
The identity
\(n1_{m\equiv1\bmod r}=\sum_{\eta\bmod r}\eta(m)\)
and Lemma~\ref{lem:complex-dirichlet-tent} give the two estimates.
For \(\psi\ne1\), every product \(\eta\psi\bmod rp\) is
nonprincipal by the Chinese remainder theorem.
The approximation assumption gives
\[
 (1-\varepsilon)S_0
 \leq\left|\sum_{m\leq x}w_m f(m)\right|
 \leq U C_A n\sqrt x\log(2pr),
\]
which proves \eqref{eq:untwisting-l1-necessary}.

For the final assertion, the character of \(\rho\) is \(n\) at
the identity, \(-1\) on \(N\setminus\{1\}\), and zero outside
\(N=\operatorname{Gal}(E/K)\). An element of
\(N\rtimes\mathbb F_r^\times\) outside \(N\) has order equal to
the order of its multiplier. Its powers therefore either remain
outside \(N\) or equal the identity. For a prime
\(\ell\equiv1\bmod r\) below \(x\), the empty-window assumption
forces Frobenius to be the identity. Consequently, for all prime
powers \(m\leq x\) coprime to \(rp\), the coefficient of
\(-L'/L(s,\rho)\) is
\(n\Lambda(m)1_{m\equiv1\bmod r}\).
Multiplication by \(\psi(m)\) proves the asserted equality.
\end{proof}

For example, once \(x\) is sufficiently larger than
\(n^2\log^2(2pr)\), a convex combination of nonprincipal
characters cannot satisfy the approximation hypothesis with any
fixed \(\varepsilon<1\). Allowing complex coefficients removes
that obstruction only by increasing their norm.

Here is an exact limitation of one common certificate based on
such coefficients. Suppose one also has, for all selected twists,
a bound of the usual zero-free-strip form
\[
 |T_{\rho\psi}(x)|
 \leq C_\delta n\log(2pr)\,x^{1-\delta},
 \qquad 0<\delta<\tfrac12,
\]
where \(T_{\rho\psi}\) denotes the same sum with the Euler factors
at \(rp\) removed. Combining these bounds only by the triangle
inequality would certify a contradiction to an empty window only if
\[
 U C_\delta n\log(2pr)x^{1-\delta}
 <(1-\varepsilon)S_0.
\]
Together with \eqref{eq:untwisting-l1-necessary}, this requires
\begin{equation}\label{eq:untwisting-budget-conflict}
 x^{1/2-\delta}<C_A/C_\delta.
\end{equation}
Thus this absolute-error transfer, with a fixed
\(\delta<1/2\) and fixed constants, cannot give arbitrarily large
windows, even if the selected twists satisfy those bounds.
For a variable \(\delta\) the exact restriction is
\eqref{eq:untwisting-budget-conflict}; no uniform claim about
the dependence of \(C_\delta\) is needed.
The conclusion concerns this particular use of individual bounds.
It gives neither a lower bound for the twisted sums
nor an obstruction to a new correlation estimate retaining
cancellation among the twists.

\subsection{What a bilinear character calculation adds}

The issue is not that interpolation by good twists is algebraically
impossible. A sufficiently large family always interpolates a short
window. This can be proved without any analytic assumption.

\begin{proposition}\label{prop:good-twist-interpolation}
Let \(\mathcal G\) be any set of characters modulo \(p\), and let
\(b=(p-1)-|\mathcal G|\). For a set
\(\mathcal S\subseteq\mathbb F_p^\times\) of size \(h\), suppose
\(bh<p-1\). Then there are coefficients \(c_\psi\), \(\psi\in\mathcal G\),
such that
\[
 \sum_{\psi\in\mathcal G}c_\psi\psi(u)=1
 \quad(u\in\mathcal S),\qquad
 \sum_{\psi\in\mathcal G}|c_\psi|
 \leq \sqrt{\frac{(p-1)h}{p-1-bh}}.
\]
Thus the good twists in Proposition~\ref{prop:p-specific-twist-pool},
after also excluding the principal character, permit exact
interpolation whenever the displayed size condition holds.
This last assertion is existence of coefficients; it does not
give a polynomial-time procedure to identify the good twists.
\end{proposition}
\begin{proof}
Form the \(h\)-by-\(|\mathcal G|\) matrix \(V\) with
\(V_{u,\psi}=\psi(u)\). Orthogonality for the complete character
family gives
\[
 VV^*=(p-1)I_h-WW^*,
\]
where \(W\) is the matrix of the \(b\) omitted columns.
Since \(WW^*\) is positive semidefinite and its trace is \(bh\),
the least eigenvalue of \(VV^*\) is at least \(p-1-bh>0\).
For \({\bf 1}\in\mathbb C^h\), choose
\(c=V^*(VV^*)^{-1}{\bf 1}\). Then \(Vc={\bf 1}\), and
\[
 \|c\|_2^2={\bf 1}^*(VV^*)^{-1}{\bf 1}
 \leq\frac{h}{p-1-bh}.
\]
Cauchy--Schwarz and \(|\mathcal G|\leq p-1\) finish the proof.
\end{proof}

For comparison, a small selected family is nearly orthogonal to
the constant function in the prime weights under A.
This is a bilinear estimate in distinct rational characters,
not an additional tensor-power identity for \(\rho\).

\begin{proposition}\label{prop:bilinear-small-menu}
Use the weights of Proposition~\ref{prop:absolute-untwisting-cost},
and put \(B=C_A n\sqrt x\log(2pr)\).
For \(k\) distinct nonprincipal characters
\(\psi_1,\ldots,\psi_k\bmod p\), suppose \(S_0>(k-1)B\).
Then every \(c\in\mathbb C^k\) satisfies
\[
 \sum_{m\leq x}w_m
 \left|1-\sum_{i=1}^k c_i\psi_i(m)\right|^2
 \geq
 S_0-\frac{kB^2}{S_0-(k-1)B}.
\]
\end{proposition}
\begin{proof}
The weighted Gram matrix
\(G_{ij}=\sum_m w_m\overline{\psi_i(m)}\psi_j(m)\)
has diagonal entries \(S_0\), and its off-diagonal entries have
absolute value at most \(B\), because
\(\overline{\psi_i}\psi_j\) is nonprincipal.
Consequently \(G\geq(S_0-(k-1)B)I\).
The vector \(v_i=\sum_m w_m\overline{\psi_i(m)}\) satisfies
\(\|v\|_2^2\leq kB^2\).
Minimizing the quadratic expression
\(S_0-2\Re(c^*v)+c^*Gc\) gives
\(S_0-v^*G^{-1}v\), which is bounded as stated.
\end{proof}

For \(kB\leq S_0/2\), the lower bound is
\(S_0-2kB^2/S_0\). The large-family interpolation proposition
and this small-family bound are consistent: they quantify different
sizes of the available character family. Neither yields the new
cancellation estimate for nonabelian twisted sums that would be
needed to advance the uniform auxiliary-prime problem.


\section{The maximal Teichm\"uller unit}
This section proves both directions of Theorem~\ref{thm:teich},
so that a use of the field as an input cannot be mistaken for its construction.

\begin{proof}[Proof of Theorem~\ref{thm:teich}]
Each irreducible factor of $\Phi_r$ over $\F_p$ has degree
$s=\ord_r(p)$.  Thus $D$ is a product of $(r-1)/s$ copies
of $\F_{p^s}$.  This description proves the properties used
below; the algorithm never needs to factor $\Phi_r$.

First let a field $L\cong\F_{p^r}$ be presented.
Compute a normal-basis element $\alpha$ for $L/\F_p$ by
Lenstra's deterministic algorithm, and put $C=L\otimes_{\F_p}D$.
Let $\sigma$ be $p^s$-power Frobenius on $L$ and the identity
on $D$.  Since $s\mid r-1$, $\gcd(s,r)=1$, and $\sigma$
generates $\Gal(L/\F_p)$.  In every field component of $C$,
the elements $\sigma^i(\alpha)$, $0\le i<r$, are a basis over
$\F_{p^s}$: $L$ and $\F_{p^s}$ are linearly disjoint.
It follows that
\[
 \beta=\sum_{i=0}^{r-1}\zeta^i\sigma^i(\alpha)
\]
is nonzero in every component, hence is a unit.  A cyclic
shift gives $\sigma(\beta)=\zeta^{-1}\beta$.
On all of $C$, $\sigma$ agrees with $p^s$-power Frobenius.
Writing $w=(p^s-1)/R$ and $\delta=\beta^w$, we obtain
\[
 \delta^R=\zeta^{-1},\qquad \delta^{rR}=1.
\]

Use the Teichm\"uller lifts modulo $rR$, and define
\[
 y=\prod_{b\in\F_r^\times}
          \rho_b(\delta)^{\omega_{a+1}(b)^{-1}},
\]
where $\rho_b$ acts on $D$ and fixes $L$.
The exponents are read modulo $rR$.  Reindexing this product
by $b\mapsto cb$ shows $\rho_c(y)=y^{\omega_{a+1}(c)}$.
Also
\[
 y^R=\prod_b\zeta^{-b\omega_{a+1}(b)^{-1}}
     =\zeta^{-(r-1)}=\zeta.
\]
Thus $y$ has order $rR$, and $\eta=y^r$ has order $R$.
Since $R\mid p^s-1$, $\eta^{p^s}=\eta$.
The fixed algebra of $\sigma$ in $C$ is $D$, so $\eta\in D$.
Reduction of the equivariance congruence modulo $R$ proves
the required property of $\eta$.

Conversely, let such an $\eta\in D^\times$ be given.
The automorphisms $\rho_b$ are transitive on the field
components of $D$, and exponentiation by $\omega_a(b)$
preserves order.  Hence the image of $\eta$ in every component
has the same order, namely $R$.  It is an $r$th nonresidue
there, because $R$ is the full $r$-part of $p^s-1$.
Over a field containing $\mu_r$, a polynomial $Y^r-c$ for
prime $r\ne p$ is irreducible when $c$ is an $r$th nonresidue:
adjoining one root splits the polynomial, and its Galois group
embeds into the prime-order group $\mu_r$, so a nontrivial
extension has degree $r$.
Consequently
\[
 C'=D[Y]/(Y^r-\eta)
\]
is a product of fields of degree $sr$ over $\F_p$.

Extend the automorphisms to $C'$ by
$\rho_b(Y)=Y^{\omega_{a+1}(b)}$.
They preserve $Y^r=\eta$ and form a group action: $Y$ has
order $rR$, and the lifts are multiplicative modulo $rR$.
The action is transitive on the field components, so its
fixed algebra is a field.  To compute its degree, extend
scalars to an algebraic closure.  There are $r(r-1)$ geometric
points of $C'$, and the action on them is free because its
action on the primitive $r$th roots $\zeta$ is free.
There are therefore $r$ orbits.  Invariants commute with
this scalar extension, being kernels of finite systems of
linear equations.  It follows that $(C')^{\F_r^\times}$
is a degree-$r$ field over $\F_p$.

Both constructions involve algebras of dimension at most
$r(r-1)$.  The integers $R$ and $rR$ have $O(r\log p)$
bits.  The lifts are computed by Hensel lifting the simple
roots of $X^{r-1}-1$ modulo $r$; all exponentiations and
linear systems have polynomial bit cost.  Finally, if a
field of degree $m$ is given with $r\mid m$, its unique
degree-$r$ subfield is the kernel of $z\mapsto z^{p^r}-z$.
This is an $\F_p$-linear computation of polynomial size.
\end{proof}

\begin{corollary}[A presented finite ground field]\label{cor:teich-basefield}
Theorem~\ref{thm:teich} remains valid with the ground field
$\F_p$ replaced by a presented finite field $k=\F_q$, where
$r$ is odd and prime to its characteristic, $s=\ord_r(q)$,
and $R=r^{v_r(q^s-1)}$. The reductions have bit cost polynomial
in $r$ and $\log q$.
\end{corollary}
\begin{proof}
Use $q$th-power Frobenius and linear algebra over $k$ throughout
that proof. The factors of $\Phi_r$ have degree $s$ over $k$;
the degree-$r$ extension of $k$ and its degree-$s$ extension are
linearly disjoint. The same normal-basis resolvent, multiplicative
projection, and free action on geometric points therefore give both
reductions. Each algebra has dimension at most $r(r-1)$ over $k$,
and the exponents have $O(r\log q)$ bits. Arithmetic and linear
algebra over the given presentation of $k$ have polynomial bit cost.
\end{proof}

\begin{corollary}[The easy valuation case]\label{cor:valuationone}
If $v_r(p^{\ord_r(p)}-1)=1$, the degree-$r$ field can be
constructed without an analytic assumption.
\end{corollary}
\begin{proof}
In this case $R=r$, and $\eta=\zeta\in D$ has order $R$
and satisfies the equivariance condition modulo $r$.
Apply Theorem~\ref{thm:teich}.
\end{proof}

The remaining case in this particular approach is therefore
$r\nmid p-1$ and $r^2\mid p^{\ord_r(p)}-1$.
Theorem~\ref{thm:teich} states an equivalence of two output
problems in this case; it supplies neither input from $p,r$ alone.

\section{Arithmetic algorithms over products of finite fields}
\label{sec:cm-etale-interface}

Arithmetic algorithms over fields can be executed simultaneously
over a finite product by splitting only the supports of their zero
tests. This supplies the coefficient interface for the root-extraction
and CM constructions below.

\begin{proposition}\label{prop:etale-arithmetic-simulation}
Let \(h\in\mathbb F_p[T]\) be monic and squarefree of degree \(N\),
and put \(A=\mathbb F_p[T]/(h)\).
Consider a deterministic algorithm over finite fields whose
coefficient-dependent operations are field arithmetic, division by
tested nonzero elements, and equality tests. Suppose it uses at most
\(U\) such operations and at most \(U\) bit operations on
non-field data on every component field of \(A\), when run on
the reductions of a given input over \(A\). Integer operations and
loop bounds must be determined by the field size and previous
test outcomes; they must not inspect a representation-specific
encoding of a field element.
Assume also that \(U\) is at least the number of input field
elements and the bit length of the non-field input.

There is a deterministic simulation, polynomial in
\(N,U,\log p\), that partitions \(h\) into coprime factors and
computes the correct output over every resulting quotient. It
does not require an irreducible factor of \(h\) or a selected root
of \(h\). When the outputs are integers or integer polynomials,
the output on every factor is a single such value, valid for all
of that factor's component fields.
\end{proposition}
\begin{proof}
First use distinct-degree factorization to partition \(h\) into
products of irreducibles of a common degree \(e\). This requires
only modular powering and greatest common divisors, and has a
deterministic polynomial cost. Every component in such a block
then has the same known cardinality \(p^e\).

Run the algorithm modulo a current factor \(v\) of \(h\).
Arithmetic operations are polynomial arithmetic modulo \(v\).
For a zero test on an element represented by \(c(T)\), compute
\(w=\gcd(v,c)\). If \(w=v\), the element vanishes on every
component. If \(w=1\), it vanishes on none, and its inverse,
when requested, is obtained by the extended Euclidean algorithm.
If \(0<\deg w<\deg v\), split the state into its reductions
modulo \(w\) and \(v/w\), and continue the two appropriate
branches. Squarefreeness ensures that these factors are coprime
and that these cases exactly describe the componentwise zero test.

Every irreducible factor follows exactly the path of the original
field algorithm, by induction on its operations. There are at
most \(N\) final branches and at most \(N-1\) genuine splits.
Charging an operation to every final branch below it bounds the
number of simulated operations by \(NU\). All polynomial
arithmetic, greatest common divisors, inversions, and state
reductions have polynomial cost in \(N\) and \(\log p\).
The stored state has polynomial size because the original run
has at most \(U\) operations. This proves the claimed complexity.
If an output is an integer, its computation depends on the
common field size and the common test outcomes along that
branch, so the output is identical in all its components.
\end{proof}

\begin{proposition}\label{prop:structure-etale-simulation}
The preceding conclusion also holds for any finite \(\mathbb F_p\)-algebra
\(A\) that is a product of finite fields and is given by structure
constants in a basis of size \(N\). The output partition is then a
direct-product decomposition into explicitly represented ideals with
their own identities. A monogenic presentation is unnecessary.
\end{proposition}
\begin{proof}
All the required decompositions can be computed by linear algebra over
\(\mathbb F_p\). To separate the component degrees, let \(F:x\mapsto x^p\).
For \(e=1,\ldots,N\), in the as yet unassigned part of the algebra,
form the ideal \(J_e\) generated by \(F^e(b)-b\), where \(b\) runs through
a basis. A component of degree \(d\) survives in \(A/J_e\) exactly when
\(d\mid e\): Frobenius to the \(e\)-th power is the identity precisely
in that case, and otherwise some generator is nonzero in that field.
All proper divisors of \(e\) have already been assigned at earlier
iterations. Thus the new nonzero quotient consists exactly of the
degree-\(e\) components. To compute \(J_e\), span the products of its
generators with the algebra basis. Compute its annihilator by solving
linear systems. Since the algebra is a product of fields,
\(A=J_e\oplus\operatorname{Ann}(J_e)\); decomposing \(1\) gives the two
identities and realizes the quotient as a subalgebra. All Frobenius
computations use repeated powering with \(e\leq N\).

On any current block \(B\), a test on \(c\in B\) is handled by the
multiplication matrix \(M_c\). Its kernel and image are complementary
ideals: the former consists of the components on which \(c=0\), and the
latter of those on which \(c\ne0\). Decomposing \(1_B\) again computes
their identities. On the latter ideal, solve \(cd=1\) by linear algebra.
This gives exactly the same branch simulation and the same \(NU\)
operation charge as before. The structure constants, subspace bases,
state restrictions, and all linear systems have polynomial size in
\(N,U,\log p\).
\end{proof}

\begin{corollary}\label{cor:pmax-etale-interface}
Given an explicitly represented order \(\mathcal R\) in a number field
and a prime \(p\), one can deterministically compute, in polynomial time
in the input length, a \(p\)-maximal overorder \(\mathcal R'\) and the
finite \(\mathbb F_p\)-algebra
\((\mathcal R'/p\mathcal R')_{\mathrm{red}}\), by structure constants.
The latter is a product of the residue fields above \(p\), and the
simulation just proved applies without selecting any of those fields.
\end{corollary}
\begin{proof}
Algorithm~6.1 and Proposition~6.2 of
\citet[pp.~246--247]{BuchmannLenstra1994} compute the
\(p\)-maximal overorder. Each iteration computes a nilradical by
Frobenius linear algebra and then a multiplier order; the proof bounds
the number of strict iterations by
\(\log|\operatorname{disc}(\mathcal R)|/(2\log p)\).
For any finite algebra of dimension \(n\), its nilradical is the kernel
of \(x\mapsto x^{p^t}\) as soon as \(p^t\geq n\). This follows either
from its decomposition into Artinian local rings, or from the descending
chain of images of multiplication by an element. Compute this kernel
and quotient by it. The resulting finite reduced algebra is a product
of finite fields and has the asserted interpretation because the order
is maximal at \(p\). Neither operation uses finite-field factorization.
\end{proof}

\begin{corollary}\label{cor:etale-pila}
Let \(A\) be a finite \(\mathbb F_p\)-algebra as in either simulation
proposition. Suppose an abelian scheme over \(A\), including its identity
section, is explicitly given by projective equations whose
affine-chart ideals are radical, together with a
complete addition law of fixed shape:
the ambient dimension, number of equations, and all their
degrees are bounded by constants. Then one can compute a
direct-product decomposition of \(A\) and the Frobenius characteristic polynomial
of every component fiber in time polynomial in \(N\) and
\(\log p\), without selecting a component field.
The Frobenius on a component of degree \(e\) is the
\(p^e\)-power Frobenius.
\end{corollary}
\begin{proof}
Apply Propositions~\ref{prop:etale-arithmetic-simulation}
and~\ref{prop:structure-etale-simulation} to the
algorithm of \citet[Theorem~A; Sections~3--4,
author/PDF pp.~8--14]{Pila1990}. Its input format, including
radical affine-chart ideals, addition atlas, and identity, is
specified on pp.~8--9. The monomial-dependence tests on
pp.~10--11 are explicitly reduced to inhomogeneous linear
systems over the coefficient field, using Theorem~2.7 on p.~8.
The Frobenius powers are computed by repeated squaring and
reduction on p.~11. Section~4, p.~12, expressly lists the four
activities: polynomial arithmetic, linear-system solution,
enumeration over the small torsion prime fields, and integer
Chinese remaindering. Its chosen irreducible-polynomial
presentation is used to implement coefficient-field arithmetic;
none of these four activities branches on a component's
representation-specific encoding. The enumeration on p.~9
is over \(\mathbb F_\ell\), with \(\ell=O(\log(p^e))\),
not over the coefficient field. Thus the entire algorithm has
the arithmetic form used in the proposition. The operation
bounds on pp.~12--14 are polynomial in
\(\log(p^e)=e\log p\), with constants depending on the fixed
shape. Since \(e\leq N\), the simulation has the stated cost.
\end{proof}

This interface assumes that the equations and addition law
already define smooth abelian varieties in every component.
It neither proves a polynomial cost for constructing a varying
CM model nor provides the characteristic-zero identification of
the resulting Frobenius with a specified CM element. Those requirements are established in
Sections~\ref{sec:sextic-cm-periods}--\ref{sec:composed-cubic-cm}.


\section{A uniform construction when the cyclotomic cofactor is small}
\label{sec:small-cofactor-lifting}

The nonresidue algorithm of Bach, von zur Gathen, and Lenstra
uses smoothness of a cyclotomic value, together with roots of
unity for all its prime divisors. Their proof also permits a
different specialization: project to the primary part of one
odd prime before choosing the separating element. This costs
the numerical cofactor, but requires no roots of unity for its
prime divisors. We give the specialization explicitly, following
the rational-function and commutator arguments in
\citet[Section~5, author pp.~17--20]{BGL2001}.


Sze's root-extraction theorem has a related small-cofactor parameter,
but it belongs to a partial factorization of the full multiplicative
order of an already presented field
\citep[Definition~2.1 and Theorem~2.2, p.~4]{Sze2011}.
At field size \(p^d\), that residual parameter also contains the
proper cyclotomic factors of \(p^d-1\). The projection below
removes those factors before the bounded search. Sze's selected
prime-two case requires a primitive fourth root; the present
odd-degree construction uses only the comparison with \(1,-1\).
Thus neither of these earlier statements is invoked with a
silently changed cofactor or a missing field presentation.

\begin{theorem}\label{thm:small-cyclotomic-cofactor}
Let \(p\ne r\) be primes, with \(r\) odd, and put
\[
 d=\operatorname{ord}_r(p),\qquad
 a=v_r(p^d-1),\qquad
 M=\frac{\Phi_d(p)}{r^a}.
\]
There is an unconditional deterministic algorithm constructing
a monic irreducible polynomial of degree \(r\) over
\(\mathbb F_p\) in bit time polynomial in \(r,M,\log p\),
with one absolute exponent. It also constructs, in the formal
algebra \(\mathcal D=\mathbb F_p[Z]/\Phi_r(Z)\), an element
\(U\) satisfying
\[
 U^{r^{a-1}}=Z,\qquad
 \sigma_c(U)=U^{\omega_a(c)}
           \quad(c\in(\mathbb Z/r\mathbb Z)^\times),
\]
where \(\sigma_c(Z)=Z^c\) and \(\omega_a(c)\) is the tame
lift modulo \(r^a\). Every component of \(U\) has exact
order \(r^a\). No factorization of \(M\), presentation of
\(\mathbb F_{p^d}\), or roots of unity for prime divisors of
\(M\) are inputs.
\end{theorem}

We first prove the field calculation that will be simulated
over \(\mathcal D\).

\begin{lemma}\label{lem:projected-primary-root-extraction}
Use \(p,r,d,a,M\) as above, let \(K\) be a presented field
of cardinality \(p^d\), and suppose a primitive \(r\)-th
root \(\eta\in K\) is supplied. Put
\(t=2^{\omega(d)}\), where \(\omega(d)\) counts distinct
prime divisors and \(t=1\) if \(d=1\).
If \(p\geq trM\), an \(r\)-th root in \(K\) of any
promised nonzero \(r\)-th power \(c\in K\) can be computed
using a number of field operations and equality tests
polynomial in \(r,d,M,\log p\). The same bound holds for
the non-field bit operations. The algorithm uses no other
operations on the representation of field elements.
\end{lemma}
\begin{proof}
Since \(d=\operatorname{ord}_r(p)\), the prime \(r\) divides
\(\Phi_d(p)\) and none of the other factors of \(p^d-1\).
Thus \(M\) is an integer prime to \(r\). This includes
\(d=1\), when \(\Phi_1(p)=p-1\).

Form the split algebra
\[
 R=K[X]/(X^r-c),\qquad \alpha=X\bmod(X^r-c),
 \qquad \tau(\alpha)=\eta\alpha,
\]
where \(\tau\) fixes \(K\). The power-residue promise implies
\(R\cong K^r\). The absolute Frobenius \(\sigma(x)=x^p\)
satisfies \(\sigma^d=1\) and commutes with \(\tau\).
For a squarefree divisor \(e\mid d\), define
\[
 \sigma_e=\prod_{\ell\mid e}\sigma^{d/\ell},\qquad
 E_d=\prod_{\ell\mid d}(1-p^{d/\ell}),
\]
with empty products equal to the identity and to one,
respectively. For any \(v\in R^\times\), put
\[
 P(v)=\prod_{e\mid\operatorname{rad}(d)}
                          \sigma_e(v)^{\mu(e)}=v^{E_d}.
\]
The integer \(\Phi_d(p)E_d\) is divisible by \(p^d-1\):
each proper cyclotomic factor of \(X^d-1\) divides some
\(X^{d/\ell}-1\). Consequently
\begin{equation}\label{eq:cofactor-cyclotomic-projection}
 P(v)^{\Phi_d(p)}=1,
 \qquad P(v)^M\text{ has }r\text{-power order dividing }r^a.
\end{equation}
Also \(r\nmid E_d\), because \(p^{d/\ell}\not\equiv1\pmod r\)
for each prime \(\ell\mid d\).

\paragraph{A bounded search after primary projection.}
For \(i=1,\ldots,trM-1\), first test whether \(\alpha+i\)
is a unit. If it is not, then \((-i)^r=c\), already giving
the requested root. Otherwise calculate
\[
 \beta_i=P(\alpha+i),\qquad \gamma_i=\beta_i^M,
 \qquad \delta_i=\tau(\gamma_i)/\gamma_i.
\]
For at least one such \(i\), the element \(\delta_i\) is
outside the scalar subalgebra \(K\subset R\).

Suppose to the contrary that every \(\delta_i\) is scalar.
For any scalar coboundary \(\delta=\tau(\gamma)/\gamma\),
the identity \(\tau^r=1\) gives \(\delta^r=1\).
Choose, solely for this proof, an \(r\)-th root \(b\) of
\(c\) in \(K\), and specialize \(\alpha\mapsto b\).
Define
\[
 H(Y)=\prod_{e\mid\operatorname{rad}(d)}
  \left(\frac{\sigma_e(\eta b)+Y}{\sigma_e(b)+Y}\right)^{\mu(e)}
                         \in K(Y).
\]
Its numerator and denominator are monic of degree \(t\)
before possible cancellation. The scalar assumption implies
\(H(i)^{rM}=1\) at all the searched integers. The same
identity holds at zero, because
\[
                   H(0)=\eta^{E_d}\ne1,
                   \qquad H(\infty)=1.
\]
All these evaluations are defined: the unit tests ensure the
nonzero denominators at the searched integers, and \(b\ne0\).
The difference between the numerator and denominator of
\(H^{rM}\) has degree at most \(trM-1\), since their leading
coefficients agree. It has \(trM\) distinct zeros in
\(\mathbb F_p\), as \(p\geq trM\). It must vanish
identically. Hence \(H^{rM}=1\) in \(K(Y)\), so \(H\)
is constant, contradicting its distinct values at zero and
infinity. The search therefore succeeds. Scalar membership
is tested by the nonconstant coefficients in the power basis
\(1,\alpha,\ldots,\alpha^{r-1}\).

\paragraph{An odd-primary commutator gives a zero divisor.}
Write \(\gamma=\gamma_i\), \(\delta=\delta_i\) for the
successful search. The element \(\delta\notin K\) has order
\(r^m\) for some \(1\leq m\leq a\). Its order is found
by repeated \(r\)-th powering, without factoring any integer.
Whenever an element \(v\notin K\) satisfies \(v^r=1\),
comparison with the known scalars
\(1,\eta,\ldots,\eta^{r-1}\) finds a nonzero zero divisor:
the components of \(v\) are these scalars, and are not all
equal. This settles \(m=1\).

If \(m\geq2\), put
\[
 \xi=\delta^{r^{m-2}},\qquad
 \nu=\gamma^{r^{m-2}},\qquad
 \xi=\tau(\nu)/\nu.
\]
Then \(\xi\) has order \(r^2\). If \(\xi^r\notin K\),
the preceding comparison gives a zero divisor. Otherwise
\(\xi^r\) is a nontrivial scalar \(r\)-th root of unity.
Now \(\epsilon=\tau(\xi)/\xi\) satisfies \(\epsilon^r=1\).
If \(\epsilon\notin K\), again use the comparison. In the
remaining case \(\epsilon\in\langle\eta\rangle\subset K\),
and induction gives
\[
       \tau^j(\nu)=\epsilon^{j(j-1)/2}\xi^j\nu.
\]
Since \(r\) is odd and \(\epsilon^r=1\), setting \(j=r\)
would give \(\nu=\xi^r\nu\), contradicting \(\xi^r\ne1\).
Thus the remaining case is impossible and a nonzero zero
divisor is always obtained. This is the odd-prime part of
\citet[Section~5.10, author p.~20]{BGL2001}; its exceptional
quadratic branch is not used.

Take the polynomial gcd of the zero divisor with \(X^r-c\),
obtaining a monic factor \(h\) of degree \(j\) with
\(0<j<r\). The nonzero product of its roots,
\(w=(-1)^jh(0)\), satisfies \(w^r=c^j\).
Choose integers \(u,v\) with \(uj+vr=1\). Then
\(w^u c^v\) is an \(r\)-th root of \(c\).

There are fewer than \(drM\) shifts because \(t\leq d\).
All polynomials over \(K\) have degree at most \(r\).
The exponent lengths, including that of \(E_d\), are
polynomial in \(d,\log p\); and
\(a\log r\leq d\log p\). All inverses, gcds, unit tests,
and scalar tests use coefficient arithmetic and equality.
These observations prove every asserted complexity bound.
\end{proof}

\begin{proof}[Proof of Theorem~\ref{thm:small-cyclotomic-cofactor}]
The integers \(d,a,\Phi_d(p),M\) are computed in polynomial
time in \(r,\log p\); their bit lengths are at most a
polynomial in these parameters. Factoring \(d\leq r-1\)
by trial division gives \(t=2^{\omega(d)}\leq d\).
If \(a=1\), take the tautological maximal unit \(U=Z\)
and apply Theorem~\ref{thm:teich}.

If \(p<trM\), use the unconditional prime-field
factorization-to-presentation reduction of
Proposition~\ref{prop:tangpresentation}, with Berlekamp's
deterministic factorization algorithm. Its bit cost is
polynomial in \(p,r,\log p\), and therefore polynomial in
\(r,M,\log p\), because \(p<drM\leq r^2M\).
Theorem~\ref{thm:teich} gives the maximal equivariant unit
from that field. Its \(r^{a-1}\)-th power is \(Z^j\) for
some \(j\not\equiv0\pmod r\): every equivariant strict
\(r\)-th root in \(\mathcal D\) is a power of \(Z\).
Testing the \(r-1\) possibilities determines \(j\);
raising the unit to an inverse of \(j\) modulo \(r^a\)
gives the stated normalization. This branch includes \(p=2\).

We may now assume \(p\geq trM\). Every component of
\(\mathcal D\) is a field of cardinality \(p^d\), and
\(Z\) is a primitive \(r\)-th root in each. In a single
component, repeatedly apply
Lemma~\ref{lem:projected-primary-root-extraction}, starting
with \(Z\), exactly \(a-1\) times. At the \(j\)-th step
the current element has order \(r^j\), so it is an
\(r\)-th power while \(j<a\); every returned root has
order \(r^{j+1}\).

This coefficient-field algorithm uses only arithmetic,
equality tests, and non-field integers, with the same known
field size \(p^d\) in every component. It can therefore
be run over the explicit product algebra \(\mathcal D\)
using Proposition~\ref{prop:structure-etale-simulation}.
Tests and divisions split its support only when necessary.
There are at most \(r-1\) final blocks, and their returned
elements are recombined using their idempotents. No field
component or irreducible factor of \(\Phi_r\) is selected.
The resulting \(u\in\mathcal D\) satisfies
\[
                        u^{r^{a-1}}=Z,
                        \qquad u^{r^a}=1.
\]
Since \(\dim_{\mathbb F_p}\mathcal D=r-1\), the support
operations and recombination add only a polynomial overhead
in \(r,\log p\).

Finally compute the tame projection
\[
 V=\prod_{c\in(\mathbb Z/r\mathbb Z)^\times}
           \sigma_c(u)^{\omega_a(c)^{-1}},\qquad U=V^{-1},
\]
reading exponents modulo \(r^a\). Reindexing by multiplication
in \((\mathbb Z/r\mathbb Z)^\times\) gives
\(\sigma_c(V)=V^{\omega_a(c)}\). Moreover,
\[
 V^{r^{a-1}}
   =\prod_c Z^{c\omega_a(c)^{-1}}
   =Z^{r-1}=Z^{-1}.
\]
Thus \(U\) has the required normalization, equivariance,
and componentwise order. Hensel lifting computes the tame
exponents in polynomial time, since \(a\log r\leq d\log p\).
Theorem~\ref{thm:teich} constructs the degree-\(r\) field.
Taking the minimal polynomial of any element outside its
scalar subspace produces a monic irreducible polynomial of
degree \(r\). This proves the theorem with an absolute
polynomial exponent.
\end{proof}

\begin{proof}[Proof of Theorem~\ref{cor:two-primary-cyclotomic-cofactor}]
Write \(u=v_2(\Phi_d(p))\). In the field lemma, replace
\(\gamma_i=\beta_i^M\) by \(\gamma_i=\beta_i^{M_2}\).
These elements and their coboundaries have order dividing
\(2^u r^a\). The scalar-coboundary implication
\(\delta_i^r=1\) is unchanged. The same rational-function
proof, now using \(H^{rM_2}\), finds a nonscalar coboundary
after fewer than \(trM_2\) shifts when \(p\geq trM_2\).
Its exact order is computed using only the known primes
\(2,r\) and their exponents.

If the coboundary's order is even, a suitable power is an
element \(v\) of order two which is still a coboundary
\(v=\tau(w)/w\). It cannot be scalar: scalarity would imply
\(v^r=1\), contrary to its order two and the oddness of
\(r\). Comparing \(v\) with the known scalars \(1,-1\)
therefore gives a nonzero zero divisor. This uses no square
root or quadratic nonresidue. If its order is odd, it is an
\(r\)-power and the preceding odd-prime commutator proof
applies unchanged.

The small-prime branch uses \(p<trM_2\leq r^2M_2\), so
Berlekamp's factorization-to-presentation bound has the required
size. The root-lifting iteration, simulation on every component,
and final Teichm\"uller projection are unchanged. This proves
the strengthened bound.
\end{proof}

\paragraph{The cubic case and unbounded valuation.}
For \(p\equiv-1\pmod3\), the theorem uses
\[
             M=(p+1)/3^{v_3(p+1)}.
\]
Hence every fixed bound \(M\leq(1+\log p)^C\) gives a
deterministic polynomial-time construction, uniformly in the
unbounded valuation \(v_3(p+1)\). The conclusion is
unconditional. The stronger corollary replaces this cofactor by
\[
 M_2=\frac{p+1}{2^{v_2(p+1)}3^{v_3(p+1)}}.
\]
Both promises are tested by divisions by two and three;
no factorization of the remaining cofactor is required.
This is an additional input class, rather than a solution
for arbitrary cofactors.

\begin{corollary}\label{cor:fermat-supported-cubic}
Let \(p\equiv-1\pmod3\) be prime. Suppose every odd prime
\(\ell\mid p+1\) satisfies that \(\ell-1\) is a power of
two, and let \(S\) be the largest prime divisor of \(p+1\).
There is an unconditional deterministic algorithm producing
an element of maximal three-power order in \(\mathbb F_{p^2}\)
and a monic irreducible cubic over \(\mathbb F_p\), with bit
cost polynomial in \(S,\log p\), with an absolute exponent.
In particular, this cost is polynomial in \(\log p\),
uniformly in both exponents, when \(p+1=2^u3^a\).
\end{corollary}
\begin{proof}
For \(p=2\), use \(\mathbb F_4=\mathbb F_2[s]/(s^2+s+1)\)
and the irreducible cubic \(Y^3+Y+1\). Otherwise
\(-3\) is a nonsquare modulo \(p\), so
\(K=\mathbb F_p[s]/(s^2+3)\) is an explicit quadratic field.
Trial division factors \(p+1\) in time polynomial in
\(S,\log p\), and verifies the support promise.
For \(\ell=2\), the primitive root is \(-1\); for
\(\ell=3\), it is \((-1+s)/2\).
For any other prime divisor \(\ell\), apply
\citet[Theorem~2, author p.~3; proof, pp.~13--16]{BGL2001}
with \(h=2\). Its additional inputs are irreducible
polynomials of prime degrees dividing \(\ell-1\) but not
dividing \(h\). There are no such degrees under this
support promise. The theorem therefore constructs a primitive
\(\ell\)-th root in \(K\) unconditionally in time polynomial
in \(\ell,\log p\).

Now \citet[Theorem~3, author p.~3; proof, pp.~17--20]{BGL2001}
applies with \(k=2\), since \(\Phi_2(p)=p+1\). All of its
required primitive roots have just been constructed; it
returns a noncube \(c\in K\) in time polynomial in
\(S,\log p\). Put \(a=v_3(p+1)\) and
\(z=c^{(p^2-1)/3^a}\). Then \(z\) has exact order \(3^a\)
and norm one, and \(u=z+z^{-1}\) belongs to \(\mathbb F_p\).
The polynomial
\[
                         Y^3-3Y-u
\]
is irreducible: a root \(\beta\in\mathbb F_p\) would give
a root \(w\in\mathbb F_{p^2}\) of
\(W^2-\beta W+1\), and then \(w^3\in\{z,z^{-1}\}\),
contradicting that these maximal three-primary elements are
noncubes. The number of prime divisors is at most \(\log_2(p+1)\),
so constructing all prerequisite roots preserves the stated
uniform polynomial bound. No analytic assumption occurs.
\end{proof}

The last corollary is a direct specialized transfer of the
classical smoothness theorem, with its missing-root inputs
discharged explicitly. It also covers cofactors that are
numerically large but have the indicated prime support;
Theorem~\ref{thm:small-cyclotomic-cofactor} instead permits
an arbitrary cofactor and pays for its numerical size.


\section{Nonresidues from Frobenius eigenspaces}
This section checks the extension-field step and gives a direct proof
of its unconditional even-degree quadratic case.

\begin{proof}[Proof of Theorem~\ref{thm:nonresidue} for odd $r$]
Under $A$, the Dirichlet character of order $r$ modulo $p$
has a nontrivial value at an integer of size $O((\log p)^2)$,
by the small-character bound.  Testing these integers gives
$b\in\F_p^\times$ with $\zeta=b^{(p-1)/r}\ne1$.
As $r$ is prime, $\zeta$ has order $r$.
If $r\nmid m$, then
\[
 b^{(p^m-1)/r}=\zeta^{1+p+\cdots+p^{m-1}}=\zeta^m\ne1,
\]
so the constant $b$ is already a nonresidue in the given field.

Suppose $r\mid m$ and put $Q=p^{m/r}$.
In the cyclic group $\F_{p^m}^\times$, the image of the map
$x\mapsto x^{Q-1}$ has order $S=1+Q+\cdots+Q^{r-1}$.
Because $Q\equiv1\pmod r$, $r\mid S$, so this image contains
$\zeta$.  The $\F_p$-linear equation $u^Q=\zeta u$ therefore
has a nonzero solution.  Compute one by linear algebra.
Writing $Q=1+rc$ and expanding modulo $r^2$ gives
\[
 S\equiv r+rc\frac{r(r-1)}2\equiv r\pmod{r^2}.
\]
Thus $u^{(p^m-1)/r}=\zeta^{S/r}=\zeta\ne1$.
All exponents have $O(m\log p)$ bits; in this branch $r\le m$.
The claimed uniform polynomial bound follows.
\end{proof}

\begin{proof}[The quadratic case]
For odd $m$, a prime-field nonsquare remains a nonsquare,
so the preceding small-character bound applies.
Let now $m$ be even and set $Q=p^{m/2}$.
If $p\equiv1\pmod4$ or $4\mid m$, then $Q\equiv1\pmod4$.
For a polynomial generator $\alpha$ of the given field,
$u=\alpha^Q-\alpha\ne0$ and $u^Q=-u$.  Therefore
\[
 u^{(p^m-1)/2}=(-1)^{(Q+1)/2}=-1.
\]

It remains that $p\equiv3\pmod4$ and $m\equiv2\pmod4$.
Compute $K=\ker(z\mapsto z^{p^2}-z)$ and choose
$v\in K\setminus\F_p$.  Then $u=v^p-v\ne0$,
$u^p=-u$, and $c=u^2$ is a nonsquare in $\F_p$.
Since $-c$ is a square, $t=(-c)^{(p+1)/4}$ satisfies
$t^2=-c$.  Put $i=u/t$; thus $i^2=-1$ and $K=\F_p(i)$.

Start with $z=i$, which has order four and norm one.
While $N_{K/\F_p}(z)=1$, write $z=a+bi$.
Its order is a power of two at least four; since the
$2$-part of $\F_p^\times$ has order two, $b\ne0$.
The nonzero numbers $(a+1)/2$ and $(a-1)/2$ have product
$-b^2/4$, a nonsquare.  Exactly one, say $\delta$, is a square.
Test by Euler's criterion, put $x=\delta^{(p+1)/4}$,
and set $y=b/(2x)$.  The relation $a^2+b^2=1$ verifies
$(x+yi)^2=z$.  The norm of this root is $1$ when
$\delta=(a+1)/2$, and $-1$ when $\delta=(a-1)/2$.
Replace $z$ by the root.

Each step doubles the order.  The norm-one group has
order $p+1$, so the process must reach norm $-1$ within
$O(\log p)$ steps.  At that point
\[
 z^{(p^2-1)/2}=N_{K/\F_p}(z)^{(p-1)/2}=-1.
\]
The extension degree $[\F_{p^m}:K]=m/2$ is odd, so this
element remains a nonsquare in the given field.  The entire
construction uses polynomially many field operations and
linear-algebra operations, with no analytic hypothesis.
In characteristic two every element is a square, so no
quadratic-nonresidue assertion is made there.
\end{proof}

\section{Zero divisors and the earlier quasipolynomial route}
This section proves the coefficient-unit step explicitly and shows
how it completes the earlier factorization algorithm.
The final polynomial-time algorithm of Tang uses a different
prime-field theorem; the two degree dependences must be kept separate.

\begin{proposition}[Splitting from a maximal coefficient unit]
\label{prop:zerodivisor}
Let $A_0$ be an explicitly presented split algebra
$\F_p^t$, with $t>1$, and let $\tau$ be an automorphism
of prime order $r\ne p$.  Put
$s=\ord_r(p)$, $R=r^{v_r(p^s-1)}$, and
$D=\F_p[T]/\Phi_r(T)$.
Suppose a maximal Teichm\"uller unit $\eta\in D^\times$
is known.  A nonzero zero divisor in $A_0$ can be
constructed in time polynomial in $t,r,\log p$.
The same assertion holds for $r=2$, with $D=\F_p$
and a scalar unit of the full $2$-primary order.
\end{proposition}
\begin{proof}
Compute the ideal generated by $\tau(b)-b$ as $b$
runs through a basis of $A_0$.  In split coordinates,
its support is precisely the set of points moved by
$\tau$.  It is nonzero.  If it is proper, its identity
element is the required zero divisor.  Otherwise all
cycles have length $r$; assume this from now on.

Work in $C=A_0\otimes_{\F_p}D$ and extend $\tau$
trivially on $D$.  Write $\zeta=T\bmod\Phi_r$.
We first construct a unit $y\in C$ with
$\tau(y)=\zeta y$.
On any nonzero $\tau$-invariant component ideal $eC$,
the projection
\[
 P_e(z)=\sum_{i=0}^{r-1}\zeta^{-i}\tau^i(z)
\]
is not the zero map.  Over an algebraic closure this
follows by applying it to a point idempotent in a
length-$r$ orbit: the $r$ terms have disjoint supports.
Since the map is $\F_p$-linear, a basis search finds
$x=P_e(z)\ne0$.  The ideal $xC$ has an identity $f$,
computable by solving its multiplication equations.
This identity is fixed by $\tau$, because
$\tau(x)=\zeta x$ preserves the ideal.  The element
$x$ is a unit on $fC$.  Replace $e$ by $e-f$ and
repeat.  At each step the remaining dimension decreases.
Adding the vectors on their disjoint supports gives
the required unit $y$ in at most $\dim C$ steps.

Every field factor of $C$ has cardinality $p^s$.
Write $p^s-1=Rw$, and choose a positive $h$ with
$wh\equiv-1\pmod r$.  Then $z=y^{wh}$ has
$r$-power order dividing $R$ and
$\tau(z)=\zeta^{-1}z$.
For odd $r$, use the lifts $\omega$ modulo $R$ and put
\[
 x=\prod_{b\in\F_r^\times}
          \rho_b^{-1}\bigl(z^{\omega(b)}\bigr).
\]
Reindexing proves $\rho_c(x)=x^{\omega(c)}$.
Because $\tau$ commutes with $\rho_b$ and
$b^{-1}\omega(b)\equiv1\pmod r$, it also gives
$\tau(x)=\zeta^{-(r-1)}x=\zeta x$.
For $r=2$ take $x=z$; the same properties hold.

The two Teichm\"uller elements $x$ and $\eta$ generate
a noncyclic $r$-group.  Indeed its exponent divides $R$,
and a cyclic such group containing $\eta$ of order
$R$ would be generated by $\eta$.  Then $x$ would be
a power of $\eta$ and fixed by $\tau$, contradicting
$\tau(x)=\zeta x$.

Here is a concrete way to turn this discrepancy into
a zero divisor.  In every field component of $C$,
$\eta$ generates the entire $r$-primary group, so
write $x=\eta^{c_i}$ in that component for the
correctness argument.  The algorithm does not know
these exponents.  Let $R=r^a$ and
$\xi=\eta^{R/r}$.  Starting with $x_0=x$,
recover common base-$r$ digits successively.  At step
$k$, $0\leq k<a$, the previous $k$ common digits
have been removed, and compute
$v=x_k^{r^{a-1-k}}$.
Each component of $v$ is one of $\xi^j$, $0\leq j<r$.
Test the multiplication matrix of every $v-\xi^j$.
If one is zero, the $k$th digit is common to all
components; replace $x_k$ by $x_k\eta^{-jr^k}$ and
continue.  If none is zero, at least one is nonzero
and singular, because some coordinate value occurs.
This gives a nonzero zero divisor $z_0$ in $C$.
All $a$ steps cannot finish with a common digit,
since that would give $x\in\langle\eta\rangle$.

The vanishing support of $z_0=v-\xi^j$ is invariant
under every $\rho_b$.  Indeed $v$ and $\xi^j$ are
Teichm\"uller elements, and raising two $r$-primary
units to the common exponent $\omega(b)$ preserves
their equality, since that exponent is prime to $r$.
The identity of $\operatorname{Ann}_C(z_0)$ is therefore
a nonzero proper idempotent fixed by all $\rho_b$.
It lies in $C^{\F_r^\times}=A_0$, and is the
required zero divisor.
All ideals, identities, and fixed spaces are computed
by linear algebra in dimension at most $t(r-1)$.
Also $a\log r\leq s\log p$.  This proves the
polynomial bit bound.
\end{proof}

This proof specializes the constructions of
\citet[Lemma~4.2, author p.~16; Proposition~4.6, pp.~18--19]{IKRS2012}.
It makes explicit both the nonvanishing eigenvector and
the descent of the idempotent from the cyclotomic algebra.
For an arbitrary nonidentity automorphism, its order
can first be determined using the known exponent $t!$
of the permutation group on $t$ points.  The factorization
of $t!$ is obtained by trial division up to $t$;
successively remove prime factors while a matrix power
is the identity.  Taking a suitable power then gives
an automorphism of prime order.

\subsection{A complete transfer from a factor-count witness}
The next proof gives the form of the factor-count theorem needed for
the cyclotomic and elliptic constructions below. It is the prime-degree
consequence of \citet[Theorems~1.2 and 1.4, p.~2]{BhargavaEtAl2017}.
We include the nonvanishing construction of the resolvent and its
relative norm, so that the conclusion uses a deterministically chosen
unit in every component.

\begin{proposition}[A factor-count witness gives the coefficient unit]
\label{prop:factorcount-transfer}
Let $k$ be a presented finite field of cardinality $q$ and
characteristic $p$, let $r\ne p$ be an odd prime, and let
$f\in k[X]$ be squarefree of degree $N$.
Suppose that, for some $d$ divisible by $r$, the number $n$
of irreducible degree-$d$ factors of $f$ is not divisible by $r$.
In deterministic time polynomial in $N$ and $\log q$, one can
construct the maximal Teichm\"uller unit in
$D=k[T]/\Phi_r(T)$, and hence a presentation of $\F_{q^r}$.
Neither the degree-$d$ factors nor the factorization of $\Phi_r$
is an input to this algorithm.
\end{proposition}
\begin{proof}
Distinct-degree factorization uses only the gcds of successive
$X^{q^j}-X$ with the remaining polynomial, for $1\leq j\leq N$.
It determines the counts and gives the product $g$ of the degree-$d$
factors in polynomial time. Set $A=k[X]/(g)$, so that,
for the correctness argument, $A$ is a product of $n$ fields of
degree $d$. Put
\[
 s=\ord_r(q),\quad h=\gcd(d,s),\quad Q=q^s,\quad
 R=r^{v_r(Q-1)},\quad C=A\otimes_kD.
\]
Above each component $\F_Q$ of $D$, the algebra $C$ has
$nh$ components, each of degree $d/h$ over $\F_Q$.
This follows by decomposing the tensor product of finite fields:
their common subfield has degree $h$, and their compositum has
degree $ds/h$ over $k$.

Let $F$ denote $Q$th-power Frobenius on $C$. It fixes $D$.
The automorphism $\sigma=F^{d/(rh)}$ has order $r$ and
has only length-$r$ cycles on the geometric points of $C$.
Here $r\nmid h$, since $h\mid s\mid r-1$.
With $\zeta=T\bmod\Phi_r$, construct a unit
$\beta\in C$ satisfying $\sigma(\beta)=\zeta\beta$.
The support construction in Proposition~\ref{prop:zerodivisor}
does this directly: on each remaining invariant ideal apply
$\sum_{i=0}^{r-1}\zeta^{-i}\sigma^i$ to basis vectors,
take the identity of the nonzero image ideal, and continue on
its complementary support. The projection is nonzero on
each length-$r$ orbit, and the sum of the vectors on disjoint
supports is a unit. This uses no factorization or division
by the order of an auxiliary group.

Set $M=Q^{d/(rh)}$. On each field component,
$\beta^{M-1}=\zeta$. Since $M\equiv1\pmod r$ and $r$
is odd,
\[
 1+M+\cdots+M^{r-1}\equiv r\pmod{r^2}.
\]
Consequently
\[
 \beta^{(Q^{d/h}-1)/r}=\zeta
\]
in every component. Take the algebra norm $b=N_{C/D}(\beta)$.
For each of the $nh$ components over $\F_Q$, the norm to
$\F_Q$ has $r$th-power character $\zeta$. Their product
therefore gives
\[
 b^{(Q-1)/r}=\zeta^{nh}.
\]
Since $r\nmid nh$, this is a primitive $r$th root in every
component of $D$.

Write $Q-1=Rw$ and set $\delta=b^w$. Then
$\delta^{R/r}=\zeta^{nh}$, so $\delta$ has full order $R$.
Using the Teichm\"uller lifts modulo $R$, put
\[
 \eta=\prod_{c\in\F_r^\times}
          \rho_c(\delta)^{\omega_a(c)^{-1}},
 \qquad a=v_r(Q-1).
\]
Reindexing proves $\rho_c(\eta)=\eta^{\omega_a(c)}$,
whereas
$\eta^{R/r}=\zeta^{nh(r-1)}=\zeta^{-nh}$ proves
that its order is $R$. Corollary~\ref{cor:teich-basefield} now constructs
the required field.

All algebras have dimension at most $N(r-1)$, and $r\leq N$.
The norm is the determinant of a multiplication matrix over $D$;
a division-free determinant algorithm computes it over this
commutative algebra in polynomial time. The exponents have
polynomially many bits. This proves the uniform bit bound.
\end{proof}

The nonvanishing choice is essential. In
$\F_7[\alpha]/(\alpha^3-2)$, with $\zeta=4$, one has
$\alpha^7=4\alpha$, $\alpha^{49}=2\alpha$, and
$\alpha+4\alpha^7+2\alpha^{49}=0$.
Thus a resolvent of a fixed polynomial generator can vanish.
The support construction above, or a basis search in a single
field eigenspace, proves the form used here without assuming
that a prescribed resolvent is nonzero.


\subsection{Proof of Theorem~\ref{thm:quasipoly}}

\begin{proof}
Squarefree decomposition and inverse Frobenius retain
all multiplicities with polynomial cost.  For a
squarefree factor $g$, set
$A=\F_{p^m}[X]/(g)$ and compute
$B=\ker(z\mapsto z^p-z)$ by $\F_p$-linear algebra.
If $g$ has $t$ irreducible factors, then $B\cong\F_p^t$;
this follows from the Chinese remainder theorem and
does not require those factors as input.
When $t=1$, $g$ is irreducible.

If $p\leq d$ and $t>1$, choose $b\in B\setminus\F_p1$
and test $b-c$ for $c\in\F_p$.  One of these is a
nonzero zero divisor, so at most $p\leq d$ tests
split the polynomial.  Repeating gives the first claim.

For $p>d$, apply the unconditional algebra theorem of
\citet[Theorem~5.8, author pp.~28--29]{IKRS2012} to
$B$ over its scalar subalgebra $\F_p$.
It gives either a nonzero zero divisor or a semiregular
automorphism of order $t$.  Choose a prime $r\mid t$
and take its order-$r$ power.  The moved-point ideal
in Proposition~\ref{prop:zerodivisor} either splits
$B$ or reduces to its semiregular case.
Here $r\leq d<p$.

There are four ways to obtain the maximal coefficient
unit for that prime.  If $r$ is odd and $r\mid m$,
extract the degree-$r$ subfield of the given coefficient
field and apply Theorem~\ref{thm:teich}.
If $r=2$ and $m$ is even, use the unconditional part of
Theorem~\ref{thm:nonresidue} to find a nonsquare $z$ in
that field. Its norm $b$ to $\F_p$ is a nonsquare because
$b^{(p-1)/2}=z^{(p^m-1)/2}=-1$.
Raising $b$ to the odd part of $p-1$ gives the required
scalar unit.
If $r\mid p-1$, find a prime-field nonresidue under
$A$. For odd $r$, construct the irreducible binomial
$X^r-c$ and apply Theorem~\ref{thm:teich}; for $r=2$,
raise the nonsquare to the odd part of $p-1$.
If $v_r(p^{\ord_r(p)}-1)=1$, use the tautological
unit $\zeta$ directly.
Every remaining $r$ lies in $\mathcal B$; its one
pure-field zeta hypothesis gives the field by
Theorem~\ref{thm:kummer}, and then the unit by
Theorem~\ref{thm:teich}.
Proposition~\ref{prop:zerodivisor} now splits $B$.

For a zero divisor $b\in B\subset A$, represented by
a polynomial of degree below $\deg g$, the gcd
$\gcd(g,b)$ is a proper nonconstant factor.
There are at most $2d-1$ nodes in the splitting tree.
The IKRS cost at each node is
$d^{O(\log d)}(\log(p^m))^{O(1)}$, and every other
step has polynomial cost.  This proves the stated
bound.  Under the final promise, select the
given-subfield or valuation-one branch before any
analytic branch; no analytic hypothesis is then used.
\end{proof}

The earlier construction through Adleman--Lenstra and
\citet[Theorem~2, author p.~3]{BGL2001} uses the sufficient
family $\zeta_{\Q(\zeta_{2\nu},p^{1/\nu})}$, with
construction degrees $\nu\leq d^2$, together with
$\zeta_{\Q}$.  Theorem~\ref{thm:quasipoly} states
the smaller prime-degree family obtained by applying
Lenstra's equivalence directly.  Both are sufficient
hypotheses for the earlier quasipolynomial algorithm;
neither is a necessity claim, and neither is Tang's
polynomial-time theorem.


\section{Tang's polynomial algorithm and its analytic input}
\label{sec:tang}

This section separates Tang's two input interfaces and identifies the
zero-free information that makes the prime-field calls polynomial.
Tang's Theorem~1.1 \cite[pp.~1--2]{Tang2026} factors a nonzero dense
polynomial over a presented finite field and, as a separate operation,
constructs a presentation from an integer promised to be a prime power.
In the second interface the polynomial coefficients are entered after
the power basis has been constructed. A list of integers called
coefficients, without a field presentation, does not specify the first
problem.

The algebraic input is OpenAI's Theorem~1.2 (p.~3) and its explicit bit
bound, Proposition~10.1 (pp.~41--44) \cite{OpenAI2026}. For a degree-$n$ prime-field input put
\[
 B=20+(n+1)(\lceil\log_2p\rceil+1).
\]
When $n\geq2$ and $p>B^{200000}$, the input table contains, for each
prime $r\leq n$, a prime $\ell$ such that
\begin{equation}\label{eq:tangaux}
 \ell\notin\{2,3,p,r\},\qquad
 \ell\equiv1\pmod{12r},\qquad
 p^{(\ell-1)/r}\not\equiv1\pmod\ell.
\end{equation}
The other branch uses the elementary prime-field algorithm and an
empty table. If $E$ is the largest table entry, or $2$ for an empty
table, the bound is
\begin{equation}\label{eq:tangbit}
 O(B^{500000}+B^{100}E^{20}).
\end{equation}
This charges the numerical value of $E$, rather than its bit length.
The algebraic theorem contains no zero-free hypothesis. Its proof and
bit bound are external inputs here.

\subsection{Exactly two Dedekind zeta functions for an auxiliary entry}

The following proof records why a fixed half-plane suffices, including
the prime $r=2$. It is Tang's Lemma~2.1 \cite[pp.~2--4]{Tang2026},
using the smoothed explicit formula of
\cite[Lemma~11.2, p.~45]{OpenAI2026} and
\cite[Theorem~1.1, p.~3]{CT2026}.

\begin{proposition}\label{prop:tangaux}
Let $n\geq2$, let $p>\max(n,3)$ be prime, and let $r\leq n$ be prime.
Set
\[
 K=\mathbb Q(\mu_{12r}),\qquad M=K(p^{1/r}).
\]
If both $\zeta_K$ and $\zeta_M$ have no nontrivial zero with real part
greater than $31/32$, then a prime satisfying \eqref{eq:tangaux}
exists with $\ell=O(B^{32})$. The implied constant is absolute.
\end{proposition}

\begin{proof}
The prime $p$ is unramified in $K$. At a prime above $p$, the valuation
of $p$ is one, so $X^r-p$ is Eisenstein. Thus $[M:K]=r$; since
$\mu_r\subset K$, the extension is cyclic. Write $d_F=[F:\mathbb Q]$
and $D_F$ for the absolute discriminant of a number field $F$.
The cyclotomic degree and the power-basis discriminants give
\[
 d_K\leq4r,\quad \log D_K\leq d_K\log(12r),\quad d_M=rd_K,
\]
and
\[
 \log D_M\leq r\log D_K+
 d_K\bigl(r\log r+(r-1)\log p\bigr).
\]
Consequently
\begin{equation}\label{eq:tangdisc}
 (\log D_K+d_K)+r^{-1}(\log D_M+d_M)
   =O(r(1+\log p))=O(B).
\end{equation}
Here $r<p$, and the constants do not depend on the two fields.

Fix a nonzero nonnegative $w\in C_c^\infty((1,2))$, and let
$c_w=\int_1^2w(t)\,dt>0$. Define
\[
 \Psi_F(x)=\sum_{\mathfrak p}\sum_{j\geq1}
       \log N\mathfrak p\,w((N\mathfrak p)^j/x).
\]
The smoothed explicit formula expresses this as $c_wx$, minus
$\sum_\rho x^\rho\widehat w(\rho)$, and an error
$O_w(\log D_F+d_F)$. Uniform zero counting in unit height intervals
and three integrations by parts give
$\sum_\rho|\widehat w(\rho)|=O_w(\log D_F+d_F)$.
The stated zero-free half-plane therefore gives
\[
 \Psi_F(x)=c_wx+
 O_w\bigl(x^{31/32}(\log D_F+d_F)\bigr).
\]
Although the statement of OpenAI's Lemma~11.2 fixes a different width,
its displayed proof establishes precisely this implication for a fixed
width. Applying it to $K$ and $M$, with their indicated weights, yields
\begin{equation}\label{eq:tangdifference}
 \Psi_K(x)-r^{-1}\Psi_M(x)
 =(1-r^{-1})c_wx+O_w(Bx^{31/32}).
\end{equation}

Suppose no prime satisfying \eqref{eq:tangaux} lies in $(x,2x)$.
A degree-one prime of $K$ above
$\ell\notin\{2,3,p,r\}$ has $\ell\equiv1\pmod{12r}$.
The supposition says that $p$ is an $r$th power modulo $\ell$.
Thus $X^r-p$ has $r$ distinct roots in the residue field, and this
prime of $K$ splits completely in $M$. Its first-power contribution
to $\Psi_K$ cancels with the $r$ contributions to $r^{-1}\Psi_M$.
Every uncancelled term from $M$ has a negative sign and can be
discarded in an upper bound. Apart from the four exceptional rational
primes, the remaining positive terms lie over rational primes at most
$\sqrt{2x}$. For any prime ideal, the sum over eligible exponents is
at most $\|w\|_\infty\log(2x)$. Hence
\begin{equation}\label{eq:tangnop}
 \Psi_K(x)-r^{-1}\Psi_M(x)
 \leq C_wr\sqrt{x}\log(2x).
\end{equation}
Put $x=TB^{32}$ with $T$ an absolute constant. After division by $x$,
the error in \eqref{eq:tangdifference} is $O_w(T^{-1/32})$, whereas
the right side of \eqref{eq:tangnop} is
$O_w(T^{-1/2}(1+\log T))$, uniformly in $B$. For sufficiently large
fixed $T$, both are less than $c_w/4$, contradicting the main term
$(1-r^{-1})c_w\geq c_w/2$. Thus $\ell<2TB^{32}$.
\end{proof}

Cureton--Thorner's Theorem~1.1 provides exactly the two zero-free
statements in Proposition~\ref{prop:tangaux}. Their Hecke-character
corollary is not required. Nor does this proof invoke the stronger
analytic input quoted as OpenAI's Theorem~11.1: a half-plane of width
$10^{-6}$ for every finite-order Hecke character of every cyclotomic
field containing $\mu_{12}$. Substituting $E=O(B^{32})$ in
\eqref{eq:tangbit} gives $O(B^{500000})$, since $100+20\cdot32=740$.
Upward search finds the primes without knowing the implied constant.

For $r=2$, the fields are
$K=\mathbb Q(\mu_{24})$ and $M=K(\sqrt p)$. They are abelian over
$\mathbb Q$, so Dirichlet GRH supplies RH for both zeta functions.
For odd $r$, the additional radical extension is generally nonabelian
over $\mathbb Q$. Tang's proof supplies a polynomial search bound by
the fixed Dedekind zero-free theorem; it does not derive that theorem,
or RH for the radical field, from Dirichlet GRH.

\subsection{Factoring in a field that has already been presented}

The absolute Frobenius fixes a much smaller scalar field than the
$q$th-power map. This gives the following explicit reduction, due to
Tang \cite[Proposition~3.1, pp.~4--5]{Tang2026}, in the tradition of
Berlekamp's reduction to prime-field root finding.

\begin{proposition}\label{prop:tangfrobenius}
Given $\mathbb F_q=\mathbb F_p[U]/(h)$, with $h$ monic irreducible
of degree $m$, a nonzero degree-$n$ polynomial over $\mathbb F_q$
can be completely factored using at most $\max(0,n-1)$ prime-field
factorization calls. Each query is split and squarefree of degree
at most $n$. All remaining bit operations have cost polynomial in
$n+1$ and $\log_2q$.
\end{proposition}

\begin{proof}
First let $g$ be monic and squarefree, and put
$A=\mathbb F_q[X]/(g)$. The matrix of
$b\mapsto b^p-b$ over $\mathbb F_p$ is obtained by binary powering
the $m\deg g$ elements of the tensor power basis. Compute its kernel
$B$ by Gaussian elimination. If $g$ has $s$ irreducible factors, the
Chinese remainder theorem identifies $A$ with a product of $s$
finite fields and $B$ with $\mathbb F_p^s$. Thus $\dim B=s$.

When $s=1$, $g$ is irreducible. Otherwise choose a basis element
$b\in B$ outside $\mathbb F_p1$. Its minimal polynomial is found
from the first dependence among $1,b,b^2,\ldots$. It is the product
of $Y-a$ over the distinct coordinate values of $b$, so it is split,
squarefree, and has degree between $2$ and $s$. Factor it by one
prime-field call. For each root $a$, the polynomial
$\gcd(g,b-a)$ selects precisely the components with coordinate $a$.
These gcds partition the factors into at least two nonempty groups.
Recurse on each group. A rooted splitting tree with $s$ leaves and
at least two children at each internal vertex has at most $s-1$
internal vertices, hence at most $s-1$ queries.

For a general monic $f$, the polynomial $g=f/\gcd(f,f')$ is the
product of the irreducible factors whose multiplicities are not
divisible by $p$. Indeed an irreducible $u$ is separable over a
finite field, and differentiating $u^ev$, with $u\nmid v$, shows
that its multiplicity in $\gcd(f,f')$ is $e-1$ if $p\nmid e$ and
$e$ if $p\mid e$. Factor $g$, remove the full multiplicity of each
factor by exact division, and take the $p$th root of the remainder.
This root is obtained by dividing monomial exponents by $p$ and
applying $a\mapsto a^{p^{m-1}}$ to the coefficients. Repeat.
Different squarefree layers contain disjoint sets of irreducibles.
Their splitting trees therefore use at most $n-1$ queries altogether
when $n>0$. Constants use none.

The matrices have dimension at most $mn$, all quotient arithmetic
is polynomial in that dimension, and every exponent has
$O(m\log p)$ bits. There are only polynomially many divisions and
root steps. Restoring the leading coefficient and multiplicities
therefore preserves the claimed bound.
\end{proof}

\subsection{Constructing a presentation from its order}

Here the prime-field oracle also supplies the roots needed to build
the coefficient field. The query count below is Tang's
Proposition~4.1 \cite[pp.~5--7]{Tang2026}; its arithmetic construction
is included to expose the precise point at which a maximal root of
unity is obtained.

\begin{proposition}\label{prop:tangpresentation}
For a prime $p$ and an integer $m\geq1$, a monic irreducible
polynomial of degree $m$ over $\mathbb F_p$ can be constructed with
$O(m^2(1+\log_2p))$ prime-field factorization calls of degree at
most $m$ and polynomial additional bit cost in $m$ and $\log_2p$.
\end{proposition}

\begin{proof}
The case $m=1$ is immediate. Factor $m$ by trial division, whose
cost is polynomial in $m$, and treat each exact prime-power divisor
$r^a$ separately.

For odd $r\ne p$, factor $\Phi_r$ over $\mathbb F_p$ and choose
one irreducible factor. This constructs a field $E$ of degree at
most $r-1$ containing a primitive $r$th root $u$. Let
$t=v_r(|E|-1)$. An element of order $r^j$, for $j<t$, has an
$r$th root in the cyclic group $E^\times$, and every such root
has order $r^{j+1}$. Starting at $u$, factor $V^r-b$ over $E$
by Proposition~\ref{prop:tangfrobenius} at each step. A linear
factor gives the next root. After $t-1$ steps this produces an
element $c$ of order $r^t$.

The polynomial $V^{r^a}-c$ is irreducible over $E$. To prove this,
any one of its roots has order $r^{t+a}$. Its degree over $E$ is
therefore the least $e$ for which $r^{t+a}$ divides $|E|^e-1$.
The geometric sum formula for the part of $e$ prime to $r$,
followed by the binomial theorem for each factor of $r$, gives
\[
 v_r(|E|^e-1)=t+v_r(e).
\]
The least such $e$ is $r^a$. Thus this binomial constructs a
field containing the degree-$r^a$ extension of $\mathbb F_p$.

For $r=2\ne p$ and $p\equiv1\pmod4$, use $E=\mathbb F_p$,
start with $-1$, and take successive square roots until their
order is $2^{v_2(p-1)}$. The same valuation identity is valid
because $v_2(p-1)\geq2$, and $V^{2^a}-c$ is irreducible.
For $p\equiv3\pmod4$, first use
$E=\mathbb F_p[i]$, $i^2=-1$. This already completes $a=1$.
For $a\geq2$, start at $i$, take square roots to the maximal
$2$-power order in $E^\times$, and use $V^{2^{a-1}}-c$.
Its degree over $\mathbb F_p$ is $2^a$.

For $r=p$, build an Artin--Schreier tower. In each current field
$E$, the $\mathbb F_p$-linear map $z\mapsto z^p-z$ has kernel
$\mathbb F_p$ and image of codimension one. Choose a basis vector
$c$ outside that image. Then $V^p-V-c$ has no root in $E$;
its roots in a splitting field are the translates of one root by
$\mathbb F_p$. Its splitting field is generated by one root and
its Galois group is a nontrivial subgroup of the additive group
$\mathbb F_p$. The group therefore has order $p$, proving
irreducibility. Repeat $a$ times.

In each ambient field constructed above, find the degree-$r^a$
subfield as the kernel of $z\mapsto z^{p^{r^a}}-z$. Choose a basis
element outside the degree-$r^{a-1}$ subfield. Every proper subfield
of the first is contained in the second, so this element has degree
$r^a$. A dependence among its successive powers gives its minimal
polynomial over $\mathbb F_p$.

The degrees $r^a$ for distinct primes are coprime. Adjoining these
generators successively therefore gives fields throughout: a
degree-$r^a$ irreducible polynomial remains irreducible over an
extension of coprime degree, by the length of its Frobenius orbit.
The resulting compositum has degree $m$. The sum of its chosen
generators is primitive. Indeed, if its degree is $e\mid m$, then
for each $r^a\Vert m$ the field of degree
$\operatorname{lcm}(e,m/r^a)$ contains the sum and all the other
generators, hence this generator as well. Thus
$r^a\mid\operatorname{lcm}(e,m/r^a)$, and coprimality gives
$r^a\mid e$. Therefore $e=m$. Compute its minimal polynomial.

The root chain for a prime $r$ has
$O(r(1+\log_2p))$ steps, and each uses at most $r-1$ prime-field
queries of degree at most $r$. The initial cyclotomic query has
degree $r-1$. Since $\sum_{r\mid m}r\leq m$, the sum of the
squares of these primes is at most $m^2$, giving the stated query
bound. The ambient degrees are at most $m^2$. Multiplication tables,
linear systems, minimal polynomials, and binary powering in these
tables consequently take polynomial bit time.
\end{proof}

For a binary integer $q$ promised to be a prime power greater than
one, test exact integer roots with exponents from
$\lfloor\log_2q\rfloor$ down to $1$. The largest successful
exponent gives $q=p^m$ with $p$ prime, by unique factorization.
Integer-root tests are polynomial in $\log_2q$. Since
$m\log_2p=\log_2q$, Proposition~\ref{prop:tangpresentation} and
the prime-field algorithm prove the presentation interface in
polynomial bit time. Proposition~\ref{prop:tangfrobenius} separately
proves the factoring interface. Credit for this polynomial algorithm
belongs to Tang and to the cited prime-field and zero-free results.

\subsection{The two algorithmic interfaces for almost all characteristics}
\label{sec:almost-all-12r}

This section uses \(A\), the Riemann hypothesis for rational Dirichlet
\(L\)-functions, as its only new analytic input. Its algorithmic deductions
also use the purely algebraic auxiliary-prime reduction of OpenAI and
the two algebraic oracle interfaces of Tang. No Hecke or Dedekind
zero-free strip is used.

\begin{proof}[Proof of Theorem~\ref{thm:aa-12r}]
Fix a prime \(r\leq D\), put \(Q=12r\), and set
\(B_r=C_0r^2\log^2(Pr)\).
One has \(\varphi(Q)\leq4r\), also for \(r=2,3\).
Character orthogonality and Lemma~\ref{lem:dirichlet-log-weight} give
\[
 \sum_{\substack{u\leq B_r\\u\equiv1\pmod Q}}
       \Lambda(u)\log(B_r/u)
 =\frac{B_r}{\varphi(Q)}
       +O\bigl(\sqrt{B_r}\log(QB_r)\bigr).
\]
The contribution from proper prime powers is
\(O(\sqrt{B_r}\log B_r)\), by the elementary Chebyshev bound
\(\psi(t)-\vartheta(t)\ll\sqrt t\).
Uniformly for \(r\leq(\log P)^{a_0}\),
\(\log(QB_r)=O_{a_0}(\log\log P)=o(\log P)\).
Consequently the sum over primes alone is at least
\(B_r/(2\varphi(Q))\) for sufficiently large \(P\).
Every prime contributes at most \((\log B_r)^2/4\). If
\[
 \mathcal S_r=\{\ell\leq B_r:\ell\text{ prime},\
                                  \ell\equiv1\pmod Q\},
 \qquad k_r=|\mathcal S_r|,
\]
then
\[
 k_r\geq\frac{2B_r}{\varphi(Q)(\log B_r)^2}
       \geq\frac{B_r}{2r(\log B_r)^2}.
\]
Since \(\log B_r=O_{a_0}(\log\log P)\) and
\(r\log r\geq2\log2\), this implies, uniformly in this range,
\begin{equation}\label{eq:aa-many-primes}
 k_r\log r\geq2\log P
\end{equation}
once \(P\) is large enough. Any fixed sufficiently large absolute
\(C_0\) works, with the threshold enlarged if needed.

For each \(\ell\in\mathcal S_r\), choose a character
\(\chi_\ell\) of exact order \(r\) modulo \(\ell\), and put
\(M_r=\prod_{\ell\in\mathcal S_r}\ell\).
These choices are for the counting proof; the algorithm below
does not compute these characters. Define
\[
 I_r(u)=r^{-k_r}
   \sum_{(j_\ell)\in\{0,\ldots,r-1\}^{\mathcal S_r}}
       \prod_{\ell\in\mathcal S_r}\chi_\ell^{\,j_\ell}(u),
\]
where a zeroth power denotes the principal character of its modulus.
For \((u,M_r)=1\), this is the indicator that \(u\) is an \(r\)-th
power modulo every \(\ell\in\mathcal S_r\); for nonunits it is zero.
The Chinese remainder theorem shows that exactly one character
in the sum is principal. All coefficients are the positive real
number \(r^{-k_r}\). Applying Lemma~\ref{lem:dirichlet-log-weight} at
\(y=2P\) therefore yields
\[
 \sum_{u\leq2P}\Lambda(u)I_r(u)\log(2P/u)
 =\frac{2P}{r^{k_r}}
       +O\bigl(\sqrt P\log(2PM_r)\bigr).
\]
The left side is nonnegative. Each bad prime
\(p\in(P/2,P]\) contributes at least
\(\log(P/2)\log2\), because all candidate primes are smaller than
\(P/2\) for sufficiently large \(P\).
Moreover
\(\log M_r\leq\vartheta(B_r)\ll B_r\), and
\eqref{eq:aa-many-primes} gives \(2P/r^{k_r}\leq2/P\).
Thus the number of bad primes for this \(r\) is
\[
 O\!\left(
 \frac{\sqrt P(B_r+\log P)}{\log P}\right)
 =
 O\!\left(
 \frac{\sqrt P\,r^2\log^2(Pr)}{\log P}\right).
\]
Summing over \(r\leq D\), and using
\(\sum_{r\leq D,\ r\ {\rm prime}}r^2\leq D^3\), proves
\eqref{eq:aa-12r-exceptions}.
Finally \(B_r<P/2\) uniformly in the prescribed range, so
\(\ell\ne p\); the congruence gives \(\ell>12r\), excluding
\(2,3,r\). For \(D<2\) the auxiliary conditions are empty.
\end{proof}

The exceptional proportion among primes in \((P/2,P]\) is
\[
 O\!\left(\frac{D^3\log^2(PD)}{\sqrt P}\right)=o(1).
\]
Thus the simultaneous statement holds for a set of characteristics
of relative density tending to one. No assertion that a specified
prime belongs to that set is needed for the algorithms to check
their auxiliary tables.

\begin{proof}[Proof of Corollary~\ref{cor:aa-finite-field-algorithms}]
First consider a prime-field factorization query of degree \(t\leq n\),
where \(n\) is the outer input degree. With
\(L=\lceil\log_2p\rceil\), write
\(\mathcal B_n=20+(n+1)(L+1)\).
If \(n\leq1\), no table is needed.
If \(p\leq\mathcal B_n^{200000}\), Berlekamp's scalar-enumeration
branch factors all prime-field queries by testing at most \(p\)
scalars per split. Its cost is polynomial in the outer
input parameters, and its auxiliary table is empty.
This branch also covers large degrees relative to \(p\); it does
not require the almost-all estimate.

Otherwise \(p>\mathcal B_n^{200000}>n\). Since \(P<2p\), the
auxiliary primes promised by \eqref{eq:aa-12r-aux}, for every
prime \(r\leq n\), all lie below
\[
 E_n=C_1 n^2\log^2(pn)
\]
for an absolute \(C_1\). Enumerate candidates up to
\(\lceil E_n\rceil\), test primality by trial division, and test
the congruence, exclusions, and modular power in
\eqref{eq:aa-12r-aux}. This requires polynomial time in
\(n+1,L\), and the residue test needs no factorization of
\(\ell-1\). The table works for every query of degree at most \(n\).
The quoted algebraic algorithm and its bound
\[
 O\bigl(\mathcal B_n^{500000}
             +\mathcal B_n^{100}E_n^{20}\bigr)
\]
therefore give a uniform polynomial prime-field procedure.
Any table that is returned is checked directly. On inputs outside
the stated promise, bounded search may report that the required
table was not found; no incorrect factorization is asserted.

For a field of order \(p^e\) presented by an irreducible
quotient or an equivalent multiplication table, Tang's
Proposition~3.1 reduces complete factorization in degree \(n\) to
at most \(\max(0,n-1)\) prime-field factorizations of split
square-free polynomials of degree at most \(n\)
\citep[pp.~4--5]{Tang2026}.
The query-degree bound is independent of \(e\); the remaining
bit cost is polynomial in \(n+1,e\log p\).
This is why the same set of good characteristics works for
every presented extension, with no bound on \(e\).
The input is promised to be a field and its polynomial is nonzero.
A nonzero constant has an empty factor list.

For the presentation problem, exact integer-root tests recover
\(p,m\) from the promised prime power \(q\): among exponents from
\(\lfloor\log_2q\rfloor\) down to \(1\), the largest exponent giving
an integer root at least \(2\) is \(m\), and that root is \(p\).
Binary search and capped powering take polynomial time in
\(\log q\). Tang's Proposition~4.1 constructs a degree-\(m\)
irreducible polynomial using
\(O(m^2(1+\log_2p))\) prime-field factorization queries of
degree at most \(m\), with polynomial additional bit cost
\citep[pp.~5--7]{Tang2026}.
Apply the preceding construction with the parameter \(m\)
in place of \(n\). Since \(m\log p=\log q\), the complete cost
is polynomial in the input length. No factorization of the
integer \(p-1\), primitive-root oracle, or nonabelian analytic
estimate is used.
\end{proof}

This is an average consequence in the characteristic. It leaves the
requested uniform auxiliary-prime assertion for every characteristic
unchanged.


\subsection{The example \texorpdfstring{$p=17$, $r=3$}{p=17, r=3}}

This example shows the exact step provided by each construction.
Here $17\equiv-1\pmod3$ and
$17^2-1=288=2^5\cdot3^2$, so
$\operatorname{ord}_3(17)=2$ and $v_3(17^2-1)=2$.

For the auxiliary-prime construction, $\ell=7$ works:
$17^{(7-1)/3}\equiv2\pmod7$. Moreover
$\operatorname{ord}_7(17)=6$, so $\Phi_7$ is irreducible over
$\mathbb F_{17}$. If $z$ is its root, the involution
$z\mapsto z^{-1}$ fixes a degree-three subfield. Its element
$z+z^{-1}$ has minimal polynomial
\[
 X^3+X^2-2X-1.
\]
The auxiliary prime supplies a cyclotomic algebra from which the
degree-three invariant field is extracted.

Tang instead begins with
$E=\mathbb F_{17}[u]/(u^2+u+1)$. The order of $u$ is three,
whereas the maximal $3$-power order in $E^\times$ is nine.
The root chain has exactly one step. It can choose
\[
 c=2+8u,\qquad c^3=u,\qquad
 c^{-1}=11+9u,\qquad c+c^{-1}=13.
\]
The cube roots of $u$ are $2+8u$, $6+15u$, and $9+11u$.
Their absolute traces are $13$, $14$, and $7$. In the split
algebra $E[V]/(V^3-u)$, the element
$b=V+u^2V^2=V+V^{17}$ is fixed by absolute Frobenius and has
minimal polynomial
\[
 Y^3-3Y+1=(Y-7)(Y-13)(Y-14)
       \quad\hbox{over }\mathbb F_{17}.
\]
Thus one cubic prime-field query separates all three roots and
recovers $c$ by a gcd. This is the step that constructs the
maximal $3$-power unit.

Now adjoin $v$ with $v^3=c$. It has order $27$ and gives a
degree-six field over $\mathbb F_{17}$. Since
$17^3\equiv-1\pmod{27}$, the element
$\alpha=v+v^{-1}$ lies in its degree-three subfield. The identity
$(v+v^{-1})^3-3(v+v^{-1})=c+c^{-1}$ gives
\[
 \alpha^3-3\alpha-13=0.
\]
It has degree three, so $X^3-3X-13$, or
$X^3+14X+4$ in standard residues, is the resulting irreducible
polynomial. For this small prime the prime-field calls use the
empty-table branch: the example requires no zero-free theorem.
If one nevertheless wants a prime meeting Tang's stronger
congruence, $\ell=37$ satisfies $\ell\equiv1\pmod{36}$ and
$17^{12}\equiv26\pmod{37}$. The arithmetic assertions are checked
by the finite computation named in the workbench.

The two finite-field examples are reproduced by the programs \emph{Prime-degree Kummer examples} and \emph{Tang's construction at seventeen}, available through the workbench (\url{\workbenchurl}). Their exact arithmetic checks supplement the displayed calculations.

\subsection{The local revision to Cureton--Thorner}

Tang's review remark \cite[Remark~1.3, p.~2]{Tang2026} distinguishes
the official version of \cite{CT2026} from a separately posted
revision \cite{TangCTReview2026}. Their difference concerns the
residual norm estimate in Lemma~8.1, official pp.~21--22,
equation~(8.4). The original passage appeals to a compact boundary
pairing. The revision replaces that passage by an invariant-form
argument; the theorem statement and the subsequent Sobolev argument
are unchanged.

The comparison begins on the space $V^{\mathrm{fin}}$ of finite
sums of spherical harmonics in the complex local model. For a
unitary central character, Kazhdan--Patterson already realizes its
global residues in $L^2$, before the extension to smooth local
vectors \cite[Theorem~II.1.4, pp.~114--117; Theorem~II.2.1 and the
preceding paragraph, p.~118]{KP1984}. With the other local vectors
fixed, the form
$B(u,v)=\langle\Theta(u),\Theta(v)\rangle_{L^2}$ is therefore
well defined on $V^{\mathrm{fin}}$. This use of the $L^2$ norm does
not assume the continuity estimate being proved.

Let $\mathcal J$ be the complex standard intertwiner. In the
compact picture its eigenvalue on harmonics of degree $\ell$ is
\[
 a_\ell=a_0\prod_{j=1}^{\ell}
       \frac{j-1/3}{j+1/3}>0,
 \qquad a_\ell^{-1}\leq a_0^{-1}(1+\ell).
\]
The latter bound follows by comparing each inverse factor with
$(j+1)/j$. The pairing between normalized inducing representations
with opposite real parameters is invariant, so
$Q(u,v)=\langle\mathcal J^{-1}u,v\rangle_{L^2(\mathbb S^2)}$
is a positive invariant form on the target representation. Both
forms are invariant under the maximal compact subgroup and the
real Lie algebra. The compact types are distinct and have
multiplicity one. Represent $B$ relative to $Q$ on each type and
take the algebraic direct sum of these operators, obtaining $T$.
Invariance gives $T\pi(X)=\pi(X)T$ for each Lie algebra element
$X$. The one-dimensional spherical type is an eigenspace of $T$,
and it generates the irreducible local module
\cite[Theorem~I.6.4(c), p.~104]{KP1984}. Consequently $T=cI$
for a finite $c\geq0$, and
\[
 \|\Theta(u)\|_{L^2}^2
 =c\sum_{\ell\geq0}a_\ell^{-1}\|u_\ell\|_2^2
 \leq ca_0^{-1}\sum_{\ell\geq0}(1+\ell)^2\|u_\ell\|_2^2.
\]
Only finite sums occur at this stage, so no convergence or
unbounded-operator assertion is hidden in the comparison. The
product argument handles simultaneous variation at the complex
places. Differentiation and local Sobolev inequalities then give
uniform bounds for every derivative on compact sets, as required
to extend the realization continuously to smooth vectors.

This verifies the local replacement relative to the cited classical
representation theorems. It does not constitute a new proof of the
remaining analytic arguments in \cite{CT2026}. The official text's
account of its Lean formalization concerns the official version;
it is not a certificate for the separately posted replacement.
The proof of Tang's transfer uses the fixed zero-free conclusion,
and its provenance must be kept separate from a check of that
conclusion's entire proof.


\section{What the abelian quotient does not determine}
This section isolates an algebraic obstruction to a direct transfer of
prime-counting information from the cyclotomic quotient.

\begin{proposition}[Failure of abelian character reconstruction]
\label{prop:abelianobstruction}
Let $G=\F_r\rtimes\F_r^\times$, $r$ odd, and
$N=\F_r$.  The indicator of $N\setminus\{0\}$ is not
a complex linear combination of one-dimensional characters of $G$.
In particular, the condition on the Frobenius element in
Theorem~\ref{thm:kummer} is not a class-function condition
on the abelian quotient alone.
\end{proposition}
\begin{proof}
For $b\ne1$, the commutator of dilation by $b$ and
translation by $c$ is translation by $(b-1)c$.
As $c$ varies these exhaust $N$, and $G/N$ is abelian.
Thus $[G,G]=N$.  Every one-dimensional character is
trivial on $N$, so any linear combination of them has
the same value at the identity and at every nonzero
translation.  The indicated function has values zero
and one there, respectively.  This is a contradiction.
\end{proof}

This proposition is an obstruction to that specific character
reconstruction.  It does not prove a logical independence of
Hecke GRH from Dirichlet GRH, nor a lower bound for field
construction.  A new analytic implication, or a different
algebraic construction, could bypass it.

\subsection{Cyclotomic factor counts and auxiliary primes}

This subsection tests whether the factor-count construction of
Bhargava, Ivanyos, Mittal, and Saxena can bypass the auxiliary-prime search
by using products of complete cyclotomic polynomials.
Their Property~1.1 asks for a squarefree polynomial having, for some degree
divisible by $r$, a number of irreducible factors of that degree not divisible
by $r$; their Theorem~1.4 turns such a polynomial into extensions of any
$r$-power degree \cite[p.~2]{BhargavaEtAl2017}.
In the higher-valuation case, restricting this search to cyclotomic
products leads back to the same prime condition.

\begin{proposition}\label{prop:cyclotomic-count}
Let $p\ne r$ be primes, let $r$ be odd, and put
$s=\operatorname{ord}_r(p)$ and $a=v_r(p^s-1)\ge2$.
Let $M$ be a finite set of distinct positive integers prime to $p$, and put
\[
 f(X)=\prod_{m\in M}\Phi_m(X)\in\mathbb F_p[X],
 \qquad D=\deg f.
\]
Then $f$ is squarefree.  If, for some $d$ divisible by $r$, the number of
irreducible factors of $f$ of degree $d$ is not divisible by $r$, there
exist $m\in M$ and a prime $\ell\mid m$ such that
\[
 \ell\ne p,r,\qquad \ell\equiv1\pmod r,
 \qquad p^{(\ell-1)/r}\not\equiv1\pmod\ell.
 \tag{\(*\)}
\]
Moreover, $\ell\le2D^2$.
Conversely, for every prime $\ell$ satisfying \textup{\((*)\)}, the
polynomial $\Phi_\ell\bmod p$ has the stated factor-count property and
degree $\ell-1$.
If the list $M$ is given, a prime satisfying \textup{\((*)\)} can be
recovered in deterministic time polynomial in $D$ and $\log p$.
\end{proposition}

\begin{proof}
For $m$ prime to $p$, write
$m=r^k\prod_{\ell\ne r}\ell^{e_\ell}$, with positive displayed
$e_\ell$, and set
\[
 B_0=\max(k-1,0),\quad C_0=\max(k-a,0),\quad
 B_\ell=v_r(\ell-1),\quad
 C_\ell=v_r(\operatorname{ord}_\ell(p)).
\]
The relevant valuations are
\begin{align*}
 v_r(\varphi(m))&=B_0+\sum_{\ell\ne r}B_\ell,\\
 v_r(\operatorname{ord}_m(p))
   &=\max\bigl(C_0,\{C_\ell:\ell\mid m,\ \ell\ne r\}\bigr).
\end{align*}
An empty maximum is zero and $\operatorname{ord}_1(p)=1$.
The first equality follows from Euler's product.
For the second, orders over relatively prime moduli combine by least
common multiple.  When $\ell\ne r$, the quotient
$\operatorname{ord}_{\ell^{e}}(p)/\operatorname{ord}_\ell(p)$ is an
$\ell$-power, since the kernel of reduction on the unit groups is an
$\ell$-group.  It does not change the $r$-adic valuation.
For $k\ge1$, the odd-prime lifting-the-exponent identity gives
\[
 v_r(p^{su}-1)=a+v_r(u),\qquad
 \operatorname{ord}_{r^k}(p)=s r^{\max(k-a,0)}.
\]
Indeed, an exponent making $p$ equal to one modulo $r^k$ must be a
multiple of $s$, and the displayed valuation determines the least
possible additional power of $r$.

For $\ell\ne p,r$, condition \textup{\((*)\)} is equivalent to
$C_\ell=B_\ell>0$: writing $\ell-1=r^{B_\ell}h$ with $(r,h)=1$,
the order of $p$ divides $(\ell-1)/r$ exactly when
$C_\ell<B_\ell$.
Suppose no prime divisor of $m$ satisfies \textup{\((*)\)} and write
$C=v_r(\operatorname{ord}_m(p))$.
If $C>0$, take a local factor attaining this maximum.
For a factor $\ell\ne r$ its corresponding $B_\ell$ is at least
$C+1$.
If the $r$-primary factor attains it, then $k>a$ and
$B_0-C_0=a-1\ge1$.
In either case $v_r(\varphi(m))\ge C+1$.  Therefore
\[
 r\mid\operatorname{ord}_m(p)
 \quad\Longrightarrow\quad
 r\mid\frac{\varphi(m)}{\operatorname{ord}_m(p)}.
 \tag{\(**\)}
\]

Every primitive $m$-th root of unity has a Frobenius orbit of length
$\operatorname{ord}_m(p)$.  Thus $\Phi_m$ has
$\varphi(m)/\operatorname{ord}_m(p)$ irreducible factors, all of that
degree.  Different conductors prime to $p$ have disjoint root sets,
which also proves squarefreeness.  For $r\mid d$, the number of
degree-$d$ factors of $f$ is consequently
\[
 \sum_{\substack{m\in M\\\operatorname{ord}_m(p)=d}}
       \frac{\varphi(m)}d.
\]
If none of the conductors has a prime divisor satisfying
\textup{\((*)\)}, every summand is divisible by $r$, by
\textup{\((**)\)}.  This proves the first assertion by contraposition.

For the size bound,
\[
 \frac{\varphi(m)^2}{m}
 =\prod_{\ell^e\Vert m}\ell^{e-2}(\ell-1)^2\ge\frac12.
\]
Every factor is at least one except possibly the factor $1/2$ from
$\ell=2,e=1$.  Hence $\ell\le m\le2\varphi(m)^2\le2D^2$.
Conversely, if \textup{\((*)\)} holds, then
$d=\operatorname{ord}_\ell(p)$ is divisible by $r$, and
$(\ell-1)/d$ is not.  The factorization formula for $\Phi_\ell$
proves the converse.
Finally, trial division of the given conductors and modular
exponentiation test all possible prime divisors in time polynomial in
$D$ and $\log p$, since $|M|\le D$ and each $m\le2D^2$.
\end{proof}

For example, with $p=17$ and $r=3$, one has $s=a=2$, and for every
$k\ge3$ the polynomial $\Phi_{3^k}$ has three irreducible factors,
each of degree $2\cdot3^{k-2}$.  Complete cyclotomic products cannot
remove this divisibility unless a conductor has a prime divisor
satisfying \textup{\((*)\)}.  The prime $7$ does satisfy it:
$17^2\equiv2\pmod7$ and $\operatorname{ord}_7(17)=6$.

Thus a polynomial-size witness within this class is equivalent, up to
a polynomial change of size, to a polynomial-size auxiliary prime.
The proposition does not address proper selected factors, other
polynomial families, or succinct representations of much larger
cyclotomic algebras.  It restricts this construction and makes no claim
that Dirichlet GRH cannot imply the desired field construction by
another argument.


\subsection{Selected factors and invariant subalgebras}

An explicit small factor of a cyclotomic polynomial could give a
different construction, but selecting that factor is a substantive
step.  The following observations distinguish this possibility from
compression by cyclotomic symmetries.

Suppose $c\in D=\mathbb F_p[T]/\Phi_r(T)$ generates the
Teichm\"uller subgroup of order $R=r^a$.
Put $z=c^{R/r}$, an element of order $r$.
The defining equivariance of $c$ gives
$\delta_b(z)=z^b$ for $\delta_b(T)=T^b$ and every
$b\in(\mathbb Z/r\mathbb Z)^*$, since the Teichm\"uller lift of $b$
is congruent to $b$ modulo $r$.
Choose a primitive $r$-th root $\zeta$ in an algebraic closure and
write $z(\zeta)=\zeta^u$.
Evaluating the equivariance gives
$z(\zeta^b)=\zeta^{ub}$ for every $b$.
If $u=0$, every geometric value of $z$ would be one, contradicting
its order.  Hence $(u,r)=1$ and, by equality at all geometric
points of $D$, $c^{r^{a-1}}=T^u$.
Its characteristic polynomial $g(X)$ therefore has degree $r-1$
and divides $\Phi_{r^a}(X)$: its geometric roots are distinct,
because their $r^{a-1}$-th powers are distinct, and each has order
$r^a$.
It follows that
\[
 h(X)=g(X^r)
\]
is a squarefree degree-$r(r-1)$ divisor of $\Phi_{r^{a+1}}(X)$.
Every irreducible factor of $h$ has degree $sr$, and the number of
factors is $(r-1)/s$, which is prime to $r$.
Thus an explicit degree-$r$ field yields a small selected factor by
the Teichm\"uller construction; conversely any explicit divisor with
this factor-count property yields the field by
\cite[Theorem~1.4, p.~2]{BhargavaEtAl2017}.
If the selected factor is irreducible, its quotient is a degree-$sr$
field, whose degree-$r$ subfield is obtained directly as the kernel
of $\operatorname{Frob}_p^r-1$.
More generally, it is enough to select a nonconstant divisor
$g\mid\Phi_{r^a}$ for which $r\nmid\deg(g)/s$.
Every irreducible factor of $g$ has degree $s$; the squarefree
polynomial $g(X^r)$ divides $\Phi_{r^{a+1}}$, has all its
irreducible factors of degree $sr$, and has exactly $\deg(g)/s$
of them.  This proves the same factor-count condition.
Together with the characteristic-polynomial construction above,
this is a polynomial-time equivalence between constructing the
maximal Teichm\"uller unit and constructing such a polynomial-size
selected divisor of $\Phi_{r^a}$.

\begin{proposition}\label{prop:cyclotomic-compression}
Let $p\ne r$ be primes, $r$ odd, put
$s=\operatorname{ord}_r(p)$ and $a=v_r(p^s-1)$, and let $k>a$.
Set
\[
 A_k=\mathbb F_p[T]/\Phi_{r^k}(T),\qquad
 \Delta=(\mathbb Z/r^k\mathbb Z)^*.
\]
The action of $u\in\Delta$ sends $T$ to $T^u$.
For a subgroup $H\le\Delta$, put $B=A_k^H$ and
$b=v_r([\Delta:H])$.
Then $B$ has dimension $[\Delta:H]$, and every simple field factor
has degree $d=\operatorname{ord}_{\Delta/H}(pH)$, with
\[
 v_r(d)=\max(b-a+1,0).
\]
In particular, if $r\mid d$, then $\dim_{\mathbb F_p}B\ge r^a$
and the number of its simple field factors has $r$-adic valuation
exactly $a-1$.
Thus, for $a\ge2$, no such invariant algebra has the factor-count
property, irrespective of its dimension.
Every unital $\mathbb F_p$-subalgebra of $A_k$ stable under all of
$\Delta$ is of the form $A_k^H$ and obeys this conclusion.
\end{proposition}

\begin{proof}
After extension to an algebraic closure, $A_k$ is the algebra of
functions on $\Delta$, with its regular translation action.  Its
$H$-invariants are functions on $\Delta/H$, and Frobenius acts by
multiplication by $p$.  The number of cosets gives the dimension,
and their common orbit length gives the degree of every field factor.

The group $\Delta$ is cyclic, with $r$-primary subgroup of order
$r^{k-1}$.  The $r$-primary component of $p$ has order $r^{k-a}$.
Relative to a generator $\gamma$ it therefore has the form
$\gamma^{r^{a-1}v}$ with $(v,r)=1$.
The $r$-primary quotient of $\Delta/H$ has order $r^b$, so that
element's image has order $r^{\max(b-a+1,0)}$.
If this order is divisible by $r$, then $b\ge a$, which proves the
dimension bound.  The number of factors is $[\Delta:H]/d$, whose
$r$-adic valuation is then $b-(b-a+1)=a-1$.

Finally, a unital subalgebra stable under $\Delta$ becomes, over an
algebraic closure, the functions constant on the blocks of a
translation-invariant partition of $\Delta$.
The block containing the identity is a subgroup $H$, and the blocks
are its cosets: translation invariance gives closure under products
and inverses and identifies every other block as a translate.
The subalgebra is therefore $A_k^H$ after scalar extension, and hence
before scalar extension as well.
\end{proof}

Although $\Phi_{r^{a+1}}(X)=\sum_{j=0}^{r-1}X^{j r^a}$ has a short
sparse description, this description alone does not give arithmetic
in a smaller explicit algebra.  The proposition limits compression
by $\Delta$-stable subalgebras, including fixed-subgroup constructions.
It leaves open a method that computes a small selected factor from
the sparse input and breaks that symmetry.


\subsection{Cyclic fields and the auxiliary prime}
The conductor of a cyclic field gives an exact comparison with the classical auxiliary-prime search.


\begin{proposition}[Cyclic fields and the least auxiliary prime]
\label{prop:cyclic-discriminant-gate}
Let $r$ be an odd prime, let $p\ne r$ be prime, and suppose
\[
 p^{r-1}\equiv1\pmod{r^2}.
\]
Let $\ell_0$ be the least rational prime satisfying
\[
 \ell_0\ne p,\qquad \ell_0\equiv1\pmod r,
 \qquad p^{(\ell_0-1)/r}\not\equiv1\pmod{\ell_0}.
\]
Among cyclic degree-$r$ number fields in which $p$ is unramified
and inert, the least absolute discriminant is exactly
\[
 \ell_0^{\,r-1}.
\]
Equivalently, for every real $B\ge1$, such a cyclic field of
absolute discriminant at most $B$ exists if and only if such an
auxiliary prime at most $B^{1/(r-1)}$ exists.
\end{proposition}

\begin{proof}
We use the classical conductor facts recorded, with a proof, in
\citet[Lemma~8.3.1]{CornelissenHokkenRingeling2026}; the local
character argument needed here is also given explicitly below.
By Kronecker--Weber, a cyclic degree-$r$ field $L$ is cut out by
a primitive Dirichlet character $\chi$ of order $r$, of conductor
$f$.  All the characters $\chi^j$, $1\le j\le r-1$, have the
same kernel, and therefore the same conductor $f$.  The
conductor--discriminant formula gives
\[
 |D_L|=f^{r-1}.
\]

We describe the local factors of $f$.  At a prime $\ell\ne r$,
an order-$r$ character is trivial on the principal units, a
pro-$\ell$ group.  Thus a nontrivial local factor has conductor
$\ell$, and its existence requires $r\mid\ell-1$.
At $r$, the order-$r$ character is trivial on the factor of
order $r-1$ in $\mathbb Z_r^\times$.  The group
$1+r\mathbb Z_r$ is procyclic, and its subgroup of $r$th powers
is $1+r^2\mathbb Z_r$.  Hence the conductor's $r$-part is either
$1$ or $r^2$.  Consequently
\[
 f=r^{2\varepsilon}\prod_{\ell\in S}\ell,
 \qquad \varepsilon\in\{0,1\},
 \qquad \ell\equiv1\pmod r\quad(\ell\in S).
\]

The assumption on $p$ says that its class modulo $r^2$ has
order dividing $r-1$.  Every local order-$r$ character at $r$
therefore takes the value one at $p$.
Since $p$ is unramified and inert in $L$, one has
$p\nmid f$ and $\chi(p)\ne1$.  At least one local character
at a prime $\ell\in S$ must consequently take a nontrivial
value at $p$.  Such a character on
$(\mathbb Z/\ell\mathbb Z)^\times$ has kernel exactly the
subgroup of $r$th powers.  Thus $\ell$ is an auxiliary prime
of the required kind, and
\[
 f\ge\ell\ge\ell_0,
 \qquad |D_L|\ge\ell_0^{r-1}.
\]

Conversely, the unique degree-$r$ subfield of
$\mathbb Q(\zeta_{\ell_0})$ is cut out by an order-$r$
character of conductor $\ell_0$.
The auxiliary-prime condition says precisely that its Frobenius
at $p$ is nonidentity.  Every nonidentity element of its cyclic
Galois group is an $r$-cycle, so $p$ is unramified and inert.
Its discriminant is $\ell_0^{r-1}$, proving equality.
\end{proof}

The least prime in the statement exists unconditionally, for
example by qualitative Chebotarev in the Kummer splitting field;
the proposition imposes no effective bound on it.  In the hard
branch $a=v_r(p^s-1)\ge2$, with $s=\operatorname{ord}_r(p)$,
the displayed congruence on $p$ follows since $s\mid r-1$.
For $r=3$, the least discriminant in this cyclic family is
therefore $\ell_0^2$.  A small complex cubic field with prescribed
inertness can vary its imaginary quadratic resolvent, so its
discriminant condition is a different sufficient condition.

\subsection{An elliptic-curve alternative}

Elliptic curves provide a distinct way to construct irreducible
polynomials; Couveignes and Lercier's algorithm uses isogeny fibres,
with a randomized search for a suitable curve
\cite[\S4]{CouveignesLercier2013}.
The following sufficient condition uses the factorization pattern
of a division polynomial.  It replaces the auxiliary-prime search
by a curve-construction problem, whose uniform deterministic solution
is not established here.

\begin{proposition}\label{prop:elliptic-trace-gate}
Let $p>r$ be primes, with $r$ odd, and let
$E/\mathbb F_{p^2}$ be an explicitly given elliptic curve of
Frobenius trace $t$.  Suppose
\[
 v_r(t^2-4p^2)=1.
\]
Let $d$ be the order of $p$ in
$\mathbb F_r^*/\{1,-1\}$, and put $N=(r-1)/(2d)$.
The monic $r$-division polynomial of $E$ is squarefree and has
exactly $N$ irreducible factors of degree $d$ and $N$ of degree $rd$,
with no other factors.
From $E$ and the given presentation of $\mathbb F_{p^2}$ one can
therefore construct a presentation of $\mathbb F_{p^r}$
deterministically in time polynomial in $r$ and $\log p$.

In particular, the trace $t=2p-r$ satisfies the condition, and an
ordinary elliptic curve with that trace exists over $\mathbb F_{p^2}$.
This existence assertion does not give a polynomial-time
construction of the curve.
\end{proposition}

\begin{proof}
Write $q=p^2$.  The Frobenius and division-polynomial facts used here
are the ones recalled in \cite[\S2, pp.~485--486]{Schoof1985}.
Frobenius acts on the free rank-two $r$-adic Tate
module of $E$ with characteristic polynomial $X^2-tX+q$.
Its reduction $F$ on $E[r]$ has a repeated eigenvalue
$\lambda=t/2\pmod r$, and $\lambda^2=p^2\pmod r$.
Thus $\lambda=\pm p$ in $\mathbb F_r^*$.
The reduction cannot be scalar.  If the integral Frobenius matrix
were $\lambda_0I+rA$, for an integer lift $\lambda_0$ of $\lambda$,
then its characteristic discriminant would be
\[
 t^2-4q=r^2\bigl((\operatorname{tr}A)^2-4\det A\bigr),
\]
contrary to the hypothesis.  Hence
$F=\lambda(I+U)$ in a suitable basis, with $U\ne0$ and $U^2=0$.

The roots of the $r$-division polynomial correspond to nonzero
vectors of $E[r]\simeq\mathbb F_r^2$ modulo the identification
$v\sim-v$.  Since $r\ne p$, these are $(r^2-1)/2$ distinct roots.
On the one-dimensional space $\ker U$, every nonzero vector has
orbit length $d$ modulo sign.  There are $(r-1)/2$ such sign pairs.
For $v\notin\ker U$, the vectors $v,Uv$ are linearly independent,
and
\[
 F^jv=\lambda^j(v+jUv).
\]
The equality $F^jv=\pm v$ holds exactly when $r\mid j$ and
$\lambda^j=\pm1$.  Since $d$ divides $(r-1)/2$, the least such
positive $j$ is $rd$.
There are $r(r-1)/2$ sign pairs outside $\ker U$.
Dividing the two numbers of sign pairs by their respective orbit
lengths gives $N$ factors of degree $d$ and $N$ of degree $rd$.

The integer $N$ is positive and smaller than $r$.
The degree-$rd$ factor count therefore satisfies the property of
\cite[Property~1.1 and Theorem~1.4, p.~2]{BhargavaEtAl2017}.
The division polynomial has degree $(r^2-1)/2$ and is computed by
the usual division-polynomial recurrences in polynomially many
field operations.  The factor-count construction gives a
degree-$r$ extension of $\mathbb F_{p^2}$, hence an explicitly
presented field of order $p^{2r}$, in polynomial time.
Compute the kernel of the $\mathbb F_p$-linear map
$z\mapsto z^{p^r}-z$ in this field.
It is the unique subfield of order $p^r$.
Choose a basis element outside its one-dimensional prime field and
compute its minimal polynomial; since $r$ is prime, that polynomial
has degree $r$.  These are polynomial-time operations in dimension
$2r$ over $\mathbb F_p$.

For the proposed trace,
\[
 (2p-r)^2-4p^2=r(r-4p),
\]
whose $r$-adic valuation is one because $r\ne p$ and $r$ is odd.
Also $0<2p-r<2p$ and $p\nmid2p-r$.
Waterhouse's classification, Theorem~4.1(1), therefore gives an
ordinary elliptic curve over $\mathbb F_{p^2}$ with that trace
\cite[p.~536]{Waterhouse1969}.
\end{proof}

The quadratic base extension is useful here.
Over $\mathbb F_p$, an element of $\operatorname{GL}_2(\mathbb F_r)$
whose order is divisible by $r$ must have a nontrivial unipotent
part and a repeated eigenvalue $\lambda$; its determinant is
$\lambda^2$.
When $p$ is a quadratic nonresidue modulo $r$, Frobenius on
$E[r]$ cannot have this form for any $E/\mathbb F_p$.
In particular, merely requiring trace $2$ modulo $r$ would not
produce a unipotent matrix in the residual case $r\nmid p-1$.

The remaining input in Proposition~\ref{prop:elliptic-trace-gate} is
a deterministic polynomial-time method for finding a curve with
$v_r(t^2-4p^2)=1$.  The trace test itself can be performed by
Schoof's algorithm; an efficient test does not bound the number of
curves that must be examined.
The fixed trace $2p-r$ has characteristic discriminant
$-r(4p-r)$, so a direct CM construction has no established uniform
polynomial bound in $r$ and $\log p$ from this choice.
One could instead seek any trace with the valuation property.
Rational Dirichlet estimates control quadratic splitting conditions,
but a CM construction also needs the appropriate class-polynomial
root and Frobenius behaviour; those requirements are not removed by
a quadratic splitting test.

\paragraph{An explicit curve for $p=17$, $r=3$.}
Let $u^2=3$ over $\mathbb F_{17}$ and set
\[
 E:\quad y^2=x^3+x+(1+6u)
 \qquad\text{over }\mathbb F_{289}.
\]
The element $3$ is a nonsquare modulo $17$.
The nonsingularity expression is
$4+27(1+6u)^2=6+u\ne0$.
For $x=A+Bu$, write
\[
 R=A^3+9AB^2+A+1,\qquad
 I=3A^2B+3B^3+B+6
 \quad\pmod{17}.
\]
Then $x^3+x+(1+6u)=R+Iu$, and this nonzero element is a square
in $\mathbb F_{289}$ precisely when $R^2-3I^2$ belongs to
\[
 \{1,2,4,8,9,13,15,16\}\subset\mathbb F_{17}.
\]
The following table counts the $17$ values of $A$ for each $B$.
The three count columns distinguish zero, nonzero square, and
nonsquare right-hand sides.
\begin{center}
\begin{tabular}{r|rrr}
$B$ & zero & square & nonsquare\\\hline
0&0&7&10\\
1&0&10&7\\
2&0&8&9\\
3&0&8&9\\
4&0&11&6\\
5&0&4&13\\
6&0&7&10\\
7&0&7&10\\
8&0&9&8\\
9&0&7&10\\
10&0&6&11\\
11&0&6&11\\
12&0&9&8\\
13&0&13&4\\
14&0&6&11\\
15&0&8&9\\
16&0&3&14\\\hline
total&0&129&160
\end{tabular}
\end{center}
Including the point at infinity gives
$\#E(\mathbb F_{289})=1+0+2\cdot129=259$.
Thus $t=289+1-259=31$ and $t^2-4\cdot289=-195$, of
$3$-adic valuation one.

The normalized $3$-division polynomial factors as
\begin{align*}
 X^4+2X^2+(4+7u)X+11
  ={}&(X+5+7u)\\[-2pt]
 &\cdot\bigl(X^3+(12+10u)X^2
             +(4+2u)X+10+3u\bigr).
\end{align*}
The cubic factor has no root in $\mathbb F_{289}$, as can be
checked on its $289$ elements; it is therefore irreducible.
Let $\alpha$ be its image in the resulting degree-three extension.
The trace to the degree-three subfield over $\mathbb F_{17}$ is
\[
 \beta=\alpha+\alpha^{17^3}
       =(5+13u)+(5+5u)\alpha+(9+5u)\alpha^2.
\]
Direct reduction using the cubic gives
\[
 \beta^{17^3}=\beta,\qquad
 \prod_{i=0}^2(X-\beta^{17^i})=X^3+7X^2+2X+5.
\]
The latter polynomial has values
\[
 5,15,11,16,2,9,9,8,12,10,8,12,11,11,1,4,9
\]
at $0,1,\ldots,16$, respectively, so it has no root over
$\mathbb F_{17}$ and is irreducible.
This finite computation illustrates the conditional construction;
it gives no uniform bound for locating such a curve.


The numerical elliptic example is checked by \emph{An elliptic cubic extension at seventeen}, available through the same workbench.

\paragraph{A related elliptic-data interface.}
\citet[Definition~1.2 and Theorem~1.3, p.~2]{Spaenlehauer2026}
gives an unconditional deterministic root-finding algorithm when the input includes an elliptic curve, its quadratic twist, their orders and factorizations, and bases of their rational point groups. Its complexity depends on the largest prime factor of these orders. Constructing such data with suitably small prime factors is posed separately in Problem~1.7, p.~3. This conditional-on-data interface does not construct the required elliptic input from the field characteristic alone.

\subsection{Rational three-torsion and cubic extraction}
A rational torsion point makes the remaining extraction step explicit.


\begin{proposition}[The Tate family with rational $3$-torsion]
\label{prop:tate-cubic-gate}
Let $k$ be a field of characteristic different from $3$ containing
a primitive cube root of unity $\zeta$.  Let $a,b\in k$ satisfy
\[
 bR\ne0,\qquad R=a^3-27b.
\]
The elliptic curve
\[
 E_{a,b}: y^2+axy+by=x^3
\]
has the point $(0,0)$ of order $3$, and its $3$-division polynomial is
\[
 \psi_3(X)=Xg(X),\qquad
 g(X)=3X^3+a^2X^2+3abX+3b^2.
\]
The polynomial $g$ is irreducible over $k$ if and only if $R$ is
a noncube in $k^\times$.  If $R$ is a cube, $g$ splits into three
distinct linear factors over $k$.
\end{proposition}

\begin{proof}
The Weierstrass invariants are
$b_2=a^2$, $b_4=ab$, $b_6=b^2$, and $b_8=0$.
They give the discriminant $b^3(a^3-27b)$ and the displayed
division polynomial.  The line $y=0$ is the tangent at $(0,0)$
and has intersection multiplicity three there; hence this
nonidentity point has order three.

Choose $c$ with $c^3=R$ in an algebraic closure.  The three roots
of $g$ are
\[
 x_c=-\frac{a^2+ac+c^2}{9},
 \qquad c^3=R.
\]
Indeed, if $W=a^2+ac+c^2$, then
$27b=(a-c)W$.  Substitution of $X=-W/9$ into $g$ gives zero:
after taking out the factor $W^2/243$, the remaining expression is
\[
 -W+3a^2-3a(a-c)+(a-c)^2=0.
\]
These roots are distinct.  If $c\ne c'$ and $x_c=x_{c'}$, then
$a+c+c'=0$.  The remaining cube root of $R$ is consequently $a$,
so $a^3=R$, contrary to $b\ne0$.

If $R$ is a cube in $k$, all three values of $c$ lie in $k$.
If $R$ is a noncube, $X^3-R$ is irreducible, and its splitting
field over $k$ is a cyclic cubic extension.  Its Galois group
permutes the three values of $c$ transitively, and the injective
formula above shows that it also permutes the three roots of
$g$ transitively.  Thus $g$ is irreducible.
\end{proof}

The root correspondence has the explicit inverse
\[
 c=a+\frac{3b}{x_c};
\]
none of the roots $x_c$ is zero because $g(0)=3b^2\ne0$.
Thus selecting one of the other three cyclic $3$-isogeny kernels
from a curve with the distinguished kernel generated by $(0,0)$
is exactly a cube-root extraction for $R$, even when $R$ is known
to be a cube.  The parent curve is determined by quotienting by
the distinguished kernel, so it adds no extra generic parameter
to these formulas.

For finite $k=\mathbb F_q$ containing $\mu_3$, Frobenius on
$E_{a,b}[3]$ fixes $(0,0)$ and has determinant $q\equiv1\pmod3$.
It is therefore either the identity or a nontrivial unipotent
matrix.  The proposition says that the second case is equivalent
to the noncube condition on $a^3-27b$: its other three pairs
$\{P,-P\}$ form one Frobenius orbit.  In particular, over
$\mathbb F_{p^2}$ with $p\equiv-1\pmod9$, selecting such a
nontrivial unipotent curve from this Tate family retains the
original cubic-nonresidue condition.  Choosing both $a$ and $b$
in $\mathbb F_p$ always gives a cube when $R\ne0$, because cubing
is a bijection on $\mathbb F_p^\times$.

\subsection{A cubic isogeny descent and its exact factor counts}

The following statement concerns the complete collection of descending
isogenies.  It does not rule out a method that selects one branch by using
additional information.

\begin{proposition}\label{prop:cubic-volcano-count}
Let $p>3$ be a prime with $p\equiv-1\pmod9$, and let
$E/\mathbb F_p$ be ordinary, with a rational point of order $3$.
Write $t=p+1-\#E(\mathbb F_p)$ and $v=v_3(t)$, and suppose
$\operatorname{End}_{\overline{\mathbb F}_p}(E)=O$ has discriminant
$D<-4$.  For $k\ge1$, let $C_k(X)$ be the polynomial whose roots
are the $j$-invariants obtained from $E$ by $k$ successive strictly
descending $3$-isogenies, each counted once.  Then
$C_k\in\mathbb F_p[X]$ is squarefree, of degree $2\cdot3^{k-1}$.
Every irreducible factor of $C_k$ has degree
\[
 d_k=2\cdot3^{\max(k-v,0)},
\]
and the number of these factors is
\[
 N_k=3^{\min(k,v)-1}.
\]
Consequently, whenever $3\mid d_k$, one has
$v_3(N_k)=v-1$.  If $v\ge2$, no $C_k$ satisfies the
degree-$3$ factor-count property.
\end{proposition}

\begin{proof}
The rational point implies $3\mid\#E(\mathbb F_p)$, and
$p\equiv-1\pmod3$ therefore gives $3\mid t$.
Ordinarity implies $t\ne0$.  Moreover
\[
 t^2-4p\equiv1\pmod3.
\]
Since $\mathbb Z[\pi]\subseteq O$ for the Frobenius endomorphism
$\pi$, the discriminant relation shows that $3\nmid D$ and
that $3$ splits in $O$.  Thus $E$ is maximal at $3$ and
\[
 O\otimes\mathbb Z_3\simeq\mathbb Z_3\times\mathbb Z_3.
\]
For $O_k=\mathbb Z+3^kO$, the conductor exact sequence, with
$3$ prime to the conductor of $O$, gives
\[
 G_k:=\ker\bigl(\operatorname{Cl}(O_k)\longrightarrow
                         \operatorname{Cl}(O)\bigr)
 \simeq (O/3^kO)^*/(\mathbb Z/3^k\mathbb Z)^*
 \simeq(\mathbb Z/3^k\mathbb Z)^*.
\]
The unit correction is trivial because $D<-4$.
In particular $|G_k|=2\cdot3^{k-1}$; see the exact sequence
and class-number formula in
\cite[\S3, (3) and (5)]{BrokerLauterSutherland2012}.

Here is also why this kernel counts distinct descendants.
In characteristic zero, extension of ideals gives the map from
$O_k$-CM classes to $O$-CM classes.  The fibre over the canonical
CM lift of $E$ is a $G_k$-torsor.  Locally at $3$,
\[
 O_k\otimes\mathbb Z_3
 =\{(x,y)\in\mathbb Z_3^2:x\equiv y\pmod{3^k}\}.
\]
Increasing $k$ gives the unique chain of orders reached by
strict descent.  Equivalently, the standard ideal-lattice
description of a $3$-isogeny gives two descending neighbours
at the split maximal order, and three at each subsequent
level, each with a unique ascending neighbour.
Thus the above fibre is precisely the set of endpoints of the
strictly descending chains, rather than a multiset of paths;
see \cite[\S4.1]{BrokerLauterSutherland2012}.
The ordinary CM reduction theorem applies because $p$ splits
in the CM field and $p\nmid3D$.  The CM values of the same
order remain distinct on reduction, by
\cite[\S2.4, Proposition~4]{CremonaSutherland2024}.
It follows that the $2\cdot3^{k-1}$ descendants remain distinct.

Write $\pi=(\alpha,\beta)$ in the displayed $3$-adic product.
Then $\alpha+\beta=t$ and $\alpha\beta=p$, and both entries
are units.  In the quotient $G_k$ the Frobenius is the class
of $\rho=\alpha/\beta$, up to inversion.  Its image in
$\operatorname{Cl}(O)$ is trivial because the Frobenius prime
ideal is generated by $\pi$.  Hence it acts on the fibre
by translation by this one element.  We have
\[
 \rho\equiv-1\pmod3,
 \qquad v_3(\rho+1)=v_3(t/\beta)=v.
\]
The lifting-the-exponent identity therefore gives
\[
 \operatorname{ord}_{(\mathbb Z/3^k\mathbb Z)^*}(\rho)
 =2\cdot3^{\max(k-v,0)}.
\]
All translation orbits have this length.  Dividing $|G_k|$
by this length gives $N_k$, proves Frobenius stability of the
root set, and proves all the asserted factor counts.
\end{proof}

The same obstruction survives subgroup compression of this
whole fibre.  Indeed, for $H\le G_k$, put
$b=v_3([G_k:H])$.  In $G_k/H$ the $3$-part of the Frobenius
order is $3^{\max(b-v+1,0)}$.  If this order is divisible
by $3$, the number of Frobenius orbits has valuation exactly
$v-1$.  This follows directly by writing $G_k$ as its cyclic
$2$-part times its cyclic $3$-part and observing that the
Frobenius exponent in the latter has valuation $v-1$.
It applies to the invariant algebra
$(\mathbb F_p[X]/C_k)^H$ and does not presume that an
arbitrary selected factor is such an invariant algebra.

When $v=1$, there is an effective short descent.
The two first descendants form an irreducible quadratic
factor of $\Phi_3(j(E),X)$, obtainable by distinct-degree
gcds; horizontal roots can be removed without choosing them.
Use this quadratic as a presentation of $\mathbb F_{p^2}$,
and denote its class of $X$ by $j_1$.  At this first level,
\[
 \Phi_3(j_1,X)/(X-j(E))
\]
is irreducible of degree $3$ over $\mathbb F_{p^2}$.
Its quotient presents $\mathbb F_{p^6}$, and the kernel of
$\operatorname{Frob}_p^3-1$ gives its degree-$3$ subfield by
linear algebra.  All these operations take polynomial time
in $\log p$.  If $v\ge2$, the same quotient at the first
level is a split cubic, and selecting one of its roots is
additional work not performed by the complete-fibre construction.

\subsection{What a known parent supplies in cubic descent}
\label{sec:cubic-parent-choice}
The next calculation separates a known isogeny from a choice of its
continuation.  Suppose $p\equiv-1\pmod9$, and an ordinary curve
$E_0/\mathbb F_p$ has rational three-torsion and nonzero trace divisible
by $9$.  After the first descending isogeny
$\phi:E_0\longrightarrow E_1$ over $k=\mathbb F_{p^2}$, assume the
three-torsion of $E_1$ is rational, as in the intermediate levels of
Proposition~\ref{prop:cubic-volcano-count}.  Its known ascending kernel
can be normalized to $(0,0)$ in
$y^2+axy+by=x^3$.  The normalization uses translation and its tangent
line, and requires no cubic root.

By Proposition~\ref{prop:tate-cubic-gate}, selecting one of the three
other kernels is equivalent, by constant-size algebraic formulas, to
selecting a root of $C^3=a^3-27b$.  The parent is recovered by quotienting
by the known kernel, so these generic parent data do not add an
independent parameter choosing a root.  Over $k(a,b)$ the radicand has
valuation one at the divisor $a^3-27b=0$ and is not a cube.  Consequently
there is no rational function in the generic parameters that makes this
choice.  This assertion does not exclude algorithms using finite-field
Frobenius or the special descent sequence.

There is an explicit Frobenius action, but its effect is a permutation
of the choices.  Write $\pi$ for the $p$-Frobenius of $E_0$ and
$C_0=\ker\phi$.  At the first descent, $\pi(C_0)\ne C_0$.
On $E_1[3]$, the computable endomorphism
$\alpha=\phi\pi\widehat\phi$ is nonzero of rank one: the image of
$\widehat\phi$ is $C_0$, and $\phi$ is nonzero on $\pi(C_0)$.
Moreover $\alpha^2=[3]\phi\pi^2\widehat\phi=0$ there.
Thus its kernel and image are the ascending line, and $1+\alpha$
cyclically permutes the three other lines.  The coordinate algebra
already has this symmetry in the form
\[
 k[X]/(3X^3+a^2X^2+3abX+3b^2)
 \simeq k[C]/(C^3-a^3+27b),\qquad C=a+3b/X,
\]
where $C\mapsto\zeta_3C$ rotates the roots.  Mapping the known points
of $E_0[3]$ through $\phi$ yields only $\ker\widehat\phi$, because
$\widehat\phi\phi=[3]$.  A nonparent continuation instead corresponds
to a cyclic subgroup of order nine above $C_0$; its preimage cannot be
$E_0[3]$, whose image is exactly the excluded ascending line.


\begin{theorem}[Selection on an ordinary descending chain]
\label{thm:restricted-volcano-oracle}
Let $p>3$ be prime with $p\equiv-1\pmod9$, and present
$k=\mathbb F_p[\zeta]$, where $\zeta^2+\zeta+1=0$.
Suppose an ordinary elliptic curve $E_0/\mathbb F_p$ and a
specified point of order three in $E_0(\mathbb F_p)$ are given.
Consider the following restricted oracle: on a descending
$3$-isogeny chain from $E_0$, after the first descent and strictly
above the floor, choose one of the three cyclic kernels other
than the known ascending dual kernel.  The current curve, incoming
isogeny, and a generator of the dual kernel are given.

Relative to this given ordinary starting curve, this oracle and
extraction of a cube root of an arbitrary promised cube in $k$
are polynomial-time oracle equivalent.  More precisely, if $t_0$
is the trace of $E_0$ and $v=v_3(t_0)$, at most $v-1$ restricted
oracle calls produce a maximal $3$-primary unit of $k^\times$,
and hence solve every promised cube-root instance.  All remaining
work is deterministic and polynomial in $\log p$.
\end{theorem}

\begin{proof}
The rational point gives $3\mid p+1-t_0$, hence $3\mid t_0$.
Since $E_0$ is ordinary and $p>3$, $t_0\ne0$.  The Frobenius
polynomial $X^2-t_0X+p$ has discriminant congruent to $1$ modulo
three.  It therefore has distinct roots $\lambda_+,\lambda_-$
in $\mathbb Z_3$, reducing to $1,-1$; its local endomorphism
order is maximal and split.  Moreover
\[
 v_3(\lambda_+^2-\lambda_-^2)
 =v_3((\lambda_+-\lambda_-)t_0)=v.
\]

Here is the local lattice calculation for the descent, including
the termination assertion.  Up to independent scaling in the two
eigenlines, a lattice of local conductor $3^i$ is
\[
 L_i=\mathbb Z_3(1,1)+\mathbb Z_3(0,3^i).
\]
For completeness, normalize both coordinate projections of any
lattice to $\mathbb Z_3$.  Its two coordinate kernels have the
same index $3^i$, and the induced identification modulo $3^i$
is multiplication by a unit.  Rescaling one coordinate gives
exactly the displayed lattice.  In its displayed basis,
the $p^2$-Frobenius matrix is
\[
 \begin{pmatrix}
 \lambda_+^2&0\\
 (\lambda_-^2-\lambda_+^2)/3^i&\lambda_-^2
 \end{pmatrix}.
\]
The reduction is the identity for $i<v$, and is nontrivial
unipotent for $i=v$.

At $i>0$, put $u=(1,1)$ and $w=(0,3^i)$.  Among the four
index-three superlattices, adjoining $w/3$ gives $L_{i-1}$.
The other three adjoin $(u+cw)/3$, $c=0,1,2$; normalizing
their two projections as above gives conductor $3^{i+1}$.
Thus there is exactly one ascending kernel and every other
kernel descends.  At $i=0$ the two Frobenius eigenlines are the
horizontal kernels.  The two remaining kernels are conjugate
over $\mathbb F_p$ and rational over $k$, and either gives the
first descent.  These statements concern kernels, so coincident
$j$-invariants caused by automorphisms do not change the argument.

The first descent is explicit.  Move the given point and its
tangent to $(0,0)$ and $y=0$, obtaining
$y^2+axy+by=x^3$.  Since cubing is bijective on $\mathbb F_p$,
scale by the uniquely computable cube root of $b$ to obtain
$E_K:y^2+Kxy+y=x^3$.  With
$c^3=K^3-27$ in $\mathbb F_p$ and $d=\zeta c$, the point
\[
 Q=\left(-\frac3{K-d},\frac{K-c}{K-d}\right)
\]
generates a descending kernel; the identities are verified in
Proposition~\ref{prop:restricted-mobius-family} below.
Its $x$-coordinate is not in $\mathbb F_p$, so it is not a
horizontal kernel.  Quotienting by this kernel is explicit.
The image of the original rational point generates the dual
kernel: its line is different from the selected kernel, and
the image of $E[3]$ under a $3$-isogeny is its dual kernel.
The same argument tracks an explicit dual generator at every
later step.  Degree-three quotient maps can be computed from
the kernel $x$-coordinate alone.

At every intermediate depth all four kernel lines are rational
over $k$.  At depth $v$, Frobenius fixes only the ascending
line and cyclically permutes the other three.  Hence the cubic
for their $x$-coordinates is irreducible.  No trace computation
is needed: at each stage normalize at the known dual generator
to $y^2+axy+by=x^3$ and test whether $R=a^3-27b$ is a cube.
The elementary Tate formula
\[
 x_z=-\frac{a^2+az+z^2}{9},\qquad z^3=R,
 \qquad z=a+\frac{3b}{x_z}
\]
identifies its three remaining kernel coordinates with the cube
roots of $R$.  Thus a noncube is detected exactly at the floor.
Before the floor the restricted oracle makes the next descent.
There are at most $v-1$ such calls, and Hasse's bound gives
$v\le\log_3(2\sqrt p)$.

Write $p^2-1=3^a m$, with $3\nmid m$.  At the floor the noncube
$R$ gives $\eta=R^m$ of exact order $3^a$.  To extract a root
of a nonzero promised cube $\gamma$, choose $u$ with
$3u\equiv1\pmod m$ and write $3u-1=jm$; for $m=1$ take
$u=0$.  Then $x_0=\gamma^u$ has error
\[
 \gamma/x_0^3=\gamma^{-jm}\in\langle\eta\rangle.
\]
Its discrete logarithm to base $\eta$ is divisible by three.
Successive base-three digits compute this logarithm in
$O(a^2)$ group operations.  If it is $e$ with $0\le e<3^a$,
then $x_0\eta^{e/3}$ is a cube root of $\gamma$.
The zero input is immediate.  Conversely, a general promised
cube-root oracle applied to $R$ in the displayed Tate formula
selects a nonparent kernel at any intermediate stage.
All formulas have bounded degree, and exponentiation and field
arithmetic require polynomial time in $\log p$.
\end{proof}

The ordinary starting curve is an additional input in this
theorem.  The theorem does not construct such an input from $p$
alone and does not establish a smaller sufficient analytic
hypothesis for finite-field construction.

\begin{proposition}[The first continuation instance]
\label{prop:restricted-mobius-family}
Under the same conditions on $p$ and $k$, let
$E_K:y^2+Kxy+y=x^3$, where $K\in\mathbb F_p\setminus\{3\}$.
Put $\delta=\zeta-\zeta^2$, let $c\in\mathbb F_p$ be the
unique cube root of $K^3-27$, and put $d=\zeta c$.
The point $Q$ in the preceding proof has order three.  Moving
$Q$ and its tangent to the origin gives Tate coefficients
\[
 a_0=\frac{K+2d}{\delta},\qquad
 b_0=\frac{\delta d}{K-d},\qquad
 a_0^3-27b_0=\left(\frac{K-d}{\delta}\right)^3.
\]
For $t=K/c$, the resulting first continuation problem is
cube-root extraction for
\[
 b_0(t)=\frac{\delta\zeta}{t-\zeta},
 \qquad t\in\mathbb F_p\setminus\{1\}.
\]
Every such $t$ occurs.  Every nonzero promised cube in $k$ can
be transformed, by deterministic polynomial-time field
operations and a known cube multiplier, to a scalar multiple
of one of these values.  If its associated $E_K$ is ordinary,
one valid restricted selector call extracts its cube root.
Supersingular associated curves are not covered by that
oracle reduction.
\end{proposition}

\begin{proof}
Since $K^3-d^3=27$, the denominators are nonzero, and
\[
 -\frac3{K-d}=-\frac{K^2+Kd+d^2}{9}.
\]
The latter is the Tate $3$-torsion formula.  Substituting
$y=(K-c)/(K-d)$ verifies the curve equation.  Its tangent
slope $m$ satisfies
$K+2m=(K+2d)/\delta$, while
$2y+Kx+1=\delta d/(K-d)$, giving $a_0,b_0$.
Using $\delta^2=-3$ and
$(K-d)(K^2+Kd+d^2)=27$ proves the remaining cube identity.
Conversely, for $t\ne1$ take
\[
 c=\left(\frac{27}{t^3-1}\right)^{1/3}\in\mathbb F_p,
 \qquad K=tc.
\]
The curve's $j$-invariant, useful for identifying the exceptional
supersingular cases, is
\[
 j(E_K)=\frac{27t^3(t^3+8)^3}{(t^3-1)^3}.
\]

For any nonsingular Tate curve $E:y^2+axy+by=x^3$, the quotient
by $(0,0)$ has equation
\[
 E':y^2+axy+by=x^3-5abx-(a^3b+7b^2)
\]
and division polynomial
\[
 \psi_{3,E'}(X)=(3X+a^2)
       \bigl(X^3-9abX-b(a^3+27b)\bigr).
\]
The linear factor is the dual kernel.  The other roots correspond
bijectively to cube roots of $b$ through
\[
 X=az+3z^2,\quad z^3=b,
 \qquad z=\frac{aX+9b}{a^2+3X}.
\]
Substitution verifies both identities; the denominator cannot
vanish because the cubic at $-a^2/3$ is
$-(a^3-27b)^2/27\ne0$.  Applying this with $a_0,b_0$ proves
the assertion about the first continuation.

For the arbitrary-input reduction write $\gamma=A+B\zeta$.
If $B=0$, its root is uniquely computable in $\mathbb F_p$.
Otherwise, when $A\ne2B$ and $A\ne B$, put
\[
 t=\frac{B-2A}{A-2B},\qquad
 s=\frac{A^2-AB+B^2}{A-2B}.
\]
Then $s\in\mathbb F_p^\times$, $t\ne1$, and direct
multiplication gives $\gamma=s b_0(t)$.  All scalars are
cubes in $k$, so $b_0(t)$ is a promised cube, and multiplication
by the prime-field cube root of $s$ transfers roots back.
The two excluded projective classes are those of
$\delta\zeta=-2-\zeta$ and $\zeta^2$.  Multiplication by the
known cube $(1-\zeta)^3=-3\delta$ sends them to the classes
of $\zeta$ and $1-\zeta$, respectively, neither excluded.
After extracting a root, divide by $1-\zeta$.

If $E_K$ is ordinary, a cubic $b_0(t)$ means that the first
descending curve is above its floor, so the corresponding
selector call satisfies the oracle's promise.  Without
ordinarity this inference is unavailable.  The $j=0$ cases
$t=0,-2$ themselves have explicit roots: the latter is scalar,
and for the former $\gamma$ is a scalar multiple of $\delta$,
whose root is $\delta/\sqrt[3]{-3}$ in view of
$\delta^3=-3\delta$.  Other supersingular cases remain outside
the proved reduction.
\end{proof}

The coordinate, tangent, quotient-root, and arbitrary-input
identities in the proposition are checked by exact arithmetic,
for every admissible $K$ and every nonzero cube at
$p=17,53,71,107$, by the dependency-free program
\emph{Exact identities for restricted cubic isogeny descent}, available through the workbench (\url{\workbenchurl}).  These finite checks
supplement the identities; they do not supply an asymptotic
ordinary-start construction.

\subsubsection{An ordinary curve after extending the base}
The ordinary-start condition can be solved over a quartic extension, although the cubic selector remains a separate operation.


\begin{proposition}[An ordinary Tate curve over a quartic extension]
\label{prop:ordinary-quartic-start}
Let $p>3$ be prime with $p\equiv2\pmod3$.
There is an unconditional deterministic algorithm, polynomial in
$\log p$, that constructs $L=\mathbb F_{p^4}$ and an ordinary
curve $E/L$ with a specified $L$-rational point of order three.
After constructing a quadratic presentation
\[
 L=k[U],\qquad k=\mathbb F_{p^2},\qquad U^2=\eta,
 \quad \eta\in k^\times\setminus(k^\times)^2,
\]
one of the two curves
\[
 E_A:y^2+Axy+y=x^3,\qquad A=U\ \text{or}\ 2U,
\]
has $j(E_A)\notin k$ and therefore is ordinary.  In both cases
the point $(0,0)$ has order three.
\end{proposition}

\begin{proof}
The element $d=-3$ is a quadratic nonresidue modulo $p$, so
$k=\mathbb F_p[\delta]$, $\delta^2=d$, is explicit.
Tonelli--Shanks with this known nonresidue extracts every
promised square root in $\mathbb F_p$ in polynomial time.
Promised square roots in $k$ reduce to the same operation.
Indeed, for $z=A+B\delta$ with $B\ne0$, take a square root
$h$ in $\mathbb F_p$ of $A^2-dB^2$.  The two elements
$(A+h)/2,(A-h)/2$ have product $dB^2/4$, a nonsquare, so
exactly one is a square.  If $x^2=(A+h)/2$ is the square
choice, then
\[
 x+\frac{B}{2x}\delta
\]
is a square root of $z$.  Changing the sign of $h$ selects the
other choice when necessary.  For $B=0$, use a prime-field
square root of $A$ if it exists, and otherwise use
$\delta\sqrt{A/d}$.

Starting from $-1\in k$, repeatedly take a square root while
the current element is a square in $k$.  Its order, initially
two, doubles at each step.  After at most $v_2(p^2-1)-1$
steps the result $\eta$ is a nonsquare.  This constructs $L$
and supplies $U^{p^2}=-U$ without any unproved search bound.

For a nonzero $A$ with $A^{p^2}=-A$, neither $A^3-27$ nor
$A^3+27$ vanishes.  The displayed Tate curve is nonsingular,
has the flex $(0,0)$, and has invariant
\[
 j(A)=\frac{A^3(A^3-24)^3}{A^3-27}.
\]
Coefficient conjugation sends $j(A)$ to $j(-A)$.
Writing $u=A^3$, direct expansion gives
\[
 j(A)-j(-A)
 =\frac{2u\bigl(u^4-216u^2-373248\bigr)}{u^2-729}.
\]
Thus $j(A)\in k$ if and only if
\[
 (A^6)^2-216A^6-373248=0.
\]

Put $x=U^6=\eta^3$.  This is a nonsquare of $k$ and hence
does not belong to $\mathbb F_p$, since every nonzero element
of $\mathbb F_p$ is a square in $k$.  If both trial curves had
$j$ in $k$, then $x$ and $64x$ would both be roots of
$X^2-216X-373248$.  They are distinct: $p\ne3,7$, because
$p\equiv2\pmod3$.  Their sum would give
\[
 65x=216.
\]
If $p\mid65$, this is impossible since $p>3$; otherwise it
would put $x$ in $\mathbb F_p$, again impossible.
At least one trial therefore has $j\notin k$.

Finally, a supersingular elliptic curve in characteristic $p$
has purely inseparable multiplication-by-$p$ of degree $p^2$.
It factors as relative Frobenius squared followed by an
isomorphism, so the curve and its $p^2$-Frobenius conjugate
have the same $j$-invariant.  Every supersingular $j$ therefore
lies in $\mathbb F_{p^2}$.  The selected curve is ordinary.
All operations above are field arithmetic, exponentiation,
or promised quadratic extraction in fields of bounded degree.
\end{proof}

This construction supplies an ordinary starting curve over
$\mathbb F_{p^4}$, with explicit three-torsion.  It does not
produce a prime-field ordinary parent for
Theorem~\ref{thm:restricted-volcano-oracle}.  Nor does ordinary
status imply that a given nonparent cubic is irreducible:
the scalar-Frobenius intermediate case can still require a
split cubic selector.  A noncube found over $\mathbb F_{p^4}$
would descend to one over $\mathbb F_{p^2}$ by the relative
norm, but the construction does not assert that either trial
already provides such a noncube.


\begin{example}\label{ex:ordinary-both-cubic}
Ordinarity of the two trial curves does not ensure a noncubic
radicand.  Take $p=53$, write $k=\mathbb F_{53}(\zeta)$ with
$\zeta^2+\zeta+1=0$, and put $\eta=1+17\zeta$.
Then $\eta^{1404}=-1$, so $L=k(U)$, $U^2=\eta$, is quadratic.
For $A=U,2U$ and $B=1$, respectively, the values of $A^6$ are
$19+32\zeta$ and $50+34\zeta$.  Substitution into the ordinarity
polynomial of Proposition~\ref{prop:ordinary-quartic-start} gives
\[
 (A^6)^2-216A^6-373248
 =34+11\zeta,\quad 9+41\zeta.
\]
Both are nonzero, so both $j$-invariants lie outside $k$ and both
curves are ordinary.  On the other hand,
\[
 N_{L/k}(U^3-27)=21+21\zeta,\qquad
 N_{L/k}(8U^3-27)=43+19\zeta.
\]
Each of these elements has $936$th power one in $k$, where
$936=(53^2-1)/3$.  The norm map induces an isomorphism of the
cubic-class groups of $L$ and $k$: for a generator of $L^*$ its
norm is a generator of $k^*$.  Thus both Tate radicands are cubes
in $L$.  All displayed equalities follow by multiplication using
$\zeta^2=-\zeta-1$ and repeated squaring; the workbench program
\emph{Two ordinary Tate trials} records these powers and the
quadratic-field extraction checks.
\end{example}

The fact that supersingular invariants lie in the quadratic field
is classical \citep[Theorem~V.3.1(a)(iii), pp.~145--146]{Silverman2009}.
For ordinary curves over the prime field, Hashimoto and Nuida
\citep[Theorem~3.10 and Corollary~3.11, p.~12]{HashimotoNuida2026}
identify the surface position and first descending Hessian
three-isogeny.  Their supersingularity test takes a cubic
nonresidue in $\mathbb F_{p^2}$ as auxiliary input
\citep[Section~4, pp.~13--14]{HashimotoNuida2026}; Algorithm~2,
step~4(c), subsequently extracts cubic roots there.  These
statements explain the geometric descent but do not construct
the missing nonresidue.  The prescribed-radicand construction
below isolates the remaining root-selection input on an
ordinary curve over a larger field.



\begin{proposition}[Infinitely many failures of the two noncube trials]
\label{prop:quartic-finite-trial-failure}
Fix $a\ge2$.  There are infinitely many primes $p$ satisfying
\[
 p\equiv5\pmod{48},\qquad v_3(p+1)=a,
\]
such that, for every primitive sixteenth root of unity
$U\in\mathbb F_{p^4}$ and every $t\in\{1,2\}$,
\[
 (tU)^3-27\in(\mathbb F_{p^4}^{\times})^3.
\]
For these primes $U^2$ is a nonsquare in $\mathbb F_{p^2}$,
and at least one of the two curves
$y^2+tUxy+y=x^3$ is ordinary by
Proposition~\ref{prop:ordinary-quartic-start}.  Thus neither
trial supplies a noncube, regardless of the square-root choices
in the two-primary tower construction.
\end{proposition}

\begin{proof}
Put $L=\mathbb Q(\zeta_{48})$ and let $\tau$ be the
automorphism $\zeta_{48}\mapsto\zeta_{48}^{5}$.
It has order four.  Let $\mathcal A$ consist of all numbers
\[
 t^3\xi^3-27,
 \qquad t\in\{1,2\},\quad \xi\text{ primitive of order }16.
\]
Every number in $\mathcal A$ is nonzero and integral.
The set is stable under $\operatorname{Gal}(L/\mathbb Q)$.
Consequently
\[
 M=L\bigl(\alpha^{1/3}:\alpha\in\mathcal A\bigr)
\]
is Galois over $\mathbb Q$, and
$\operatorname{Gal}(M/L)$ has exponent dividing three.

Choose any lift $h$ of $\tau$ to $M$.  Since $h^4$ belongs
to that exponent-three kernel, $h^{12}=1$.  The element
$g=h^{-3}$ still restricts to $\tau$, and $g^4=1$.
In particular $g$ has order four.

Let $C=\mathbb Q(\zeta_{3^{a+1}})$.  The extension
$LC/L$ is cyclic of degree $3^a$.  Its intersection with
$M$ has Galois group of exponent at most three, hence
\[
 M\cap LC\subseteq L\mathbb Q(\zeta_9).
\]
Intersecting with $C$ gives
$M\cap C\subseteq\mathbb Q(\zeta_9)$; the reverse
inclusion of $\mathbb Q(\zeta_3)$ is automatic.
On this intersection, $g$ acts by complex conjugation.
Indeed its action has order dividing four and induces
conjugation on $\mathbb Q(\zeta_3)$, while the only
such lift to $\mathbb Q(\zeta_9)$ has order two and
is conjugation.

The automorphism $\sigma$ of $C$ corresponding to
\[
 -1+3^a\pmod{3^{a+1}}
\]
also restricts to conjugation on $\mathbb Q(\zeta_9)$,
since $a\ge2$.  Thus $g$ and $\sigma$ agree on $M\cap C$
and extend to an automorphism $\gamma$ of $MC$.
The qualitative Chebotarev theorem supplies a positive-density
set of rational primes with Frobenius conjugate to $\gamma$;
see \cite[Section~2.1, Theorem~1, p.~131]{Serre1981Chebotarev}.
Exclude the finitely many ramified primes and primes dividing
the norms of numbers in $\mathcal A$.

For each remaining prime the cyclotomic restrictions give
$p\equiv5\pmod{48}$ and
$p\equiv-1+3^a\pmod{3^{a+1}}$, hence $v_3(p+1)=a$.
The residue degree of a prime of $L$ over $p$ is four,
and its Frobenius in $M/L$ is the fourth power of a
conjugate of $g$, hence is trivial.  Thus every element
of $\mathcal A$ has a cube root in the residue field
$\mathbb F_{p^4}$.  Reduction sends the primitive
sixteenth roots in $L$ bijectively to all primitive
sixteenth roots in that field, proving the claim.

Finally $p\equiv5\pmod8$ gives $v_2(p^2-1)=3$.
For every such $U$, the element $U^2$ has order eight
in $\mathbb F_{p^2}^\times$ and is therefore nonsquare.
Repeated square roots of $-1$ in that quadratic field
produce a primitive eighth root, and either square root
of it in the quartic field is primitive of order sixteen.
All possible choices were included in $\mathcal A$.
\end{proof}

This proposition concerns a fixed algebraic family and every fixed
exact valuation $a\ge2$.  It gives no lower bound for a number of
trials growing with $\log p$, and it does not rule out an adaptive
construction using additional information about $p$.

\begin{proposition}[An ordinary Tate curve with prescribed radicand]
\label{prop:ordinary-prescribed-radicand}
Let $p>3$ be prime with $p\equiv2\pmod3$, let
$k=\mathbb F_{p^2}$, and present $L=k[U]$ with $U^2$ a
nonsquare in $k$.  Given any $\gamma\in k^\times$, for each
$A\in\{U,2U\}$ put
\[
 B=\frac{A^3-\gamma}{27},\qquad
 E_A:y^2+Axy+By=x^3.
\]
Both curves are nonsingular, and at least one has
$j(E_A)\notin k$, hence is ordinary.  Its Tate radicand is
exactly $A^3-27B=\gamma$.

Consequently a single oracle call selecting a cyclic order-three kernel other than the distinguished
kernel generated by $(0,0)$ on a presented ordinary Tate curve over $L$, under the promise
that its remaining cubic splits, suffices to extract a cube
root of any promised cube in $k$.  All preprocessing is
unconditional and polynomial in $\log p$.
\end{proposition}

\begin{proof}
Write a bar for the nontrivial automorphism of $L/k$.
Then $\bar A=-A$.  Since $\gamma\in k^\times$, neither
$A^3-\gamma$ nor $A^3+\gamma$ vanishes.  Thus $B\ne0$ and
the discriminant $B^3\gamma$ is nonzero.  The point $(0,0)$
has order three.

Substitution in the Tate invariants, with $u=A^3/\gamma$,
gives
\[
 j(E_A)=J(u):=\frac{27u(u+8)^3}{(u-1)^3},
 \qquad \overline{J(u)}=J(-u).
\]
Expanding the difference gives the identity
\[
 J(u)-J(-u)
 =\frac{54u(u^2+8)(u^4+259u^2+64)}{(u^2-1)^3}.
\]
Here $u\ne0$ and $u^2\ne1$.  Put
\[
 x=\frac{U^6}{\gamma^2}\in k^\times.
\]
Since $U^2$ is nonsquare in $k$, its cube divided by a square
is nonsquare.  In particular $x\notin\mathbb F_p$, and the
same holds for $64x$.  The factor $u^2+8$ cannot vanish at
either trial, because $-8\in\mathbb F_p$.

If both trial invariants were in $k$, then $x$ and $64x$ would
both be roots of
\[
 X^2+259X+64.
\]
They are distinct since $p\nmid63$.  Their sum would imply
$65x=-259$.  If $p\mid65$, this contradicts $p\nmid259$;
otherwise it puts $x$ in $\mathbb F_p$.  Both alternatives
are impossible.  Thus one invariant lies outside $k$, and
the supersingular-invariant theorem makes that curve ordinary.

For a promised cube $\gamma$ in $k$, the remaining Tate
division cubic on this selected curve splits over $L$.
For any root coordinate $X$ selected by the oracle, the
explicit Tate inverse gives
\[
 z=A+\frac{3B}{X},\qquad z^3=\gamma.
\]
The denominator is nonzero because the constant term of the
remaining division cubic is $3B^2\ne0$.  All three cube
roots of $\gamma$ already lie in $k$, since $\mu_3\subset k$.
Therefore the returned $z$ belongs to the embedded copy of $k$;
extracting its coordinates in that subfield is immediate.
The field construction, invariant tests, and formulas have the
claimed deterministic complexity by
Proposition~\ref{prop:ordinary-quartic-start}.
\end{proof}

The oracle here allows ordinary inputs over $\mathbb F_{p^4}$
and is not required to recognize an incoming edge as ascending.
It is therefore a different domain from the restricted prime-field
parent oracle in Theorem~\ref{thm:restricted-volcano-oracle}.
The proposition removes the supersingular exception for this
ordinary-input root-selection formulation; it does not construct
a selector and does not establish a uniform Dirichlet-GRH-only
field algorithm.


\subsection{An ascending selector over the quartic field}
The explicit ordinary curve constructed above removes the starting-curve
input when the selector is allowed over the quartic field. The precise
promise and its use at the volcano floor matter.


\begin{proposition}[The ascending-parent selector over the quartic field]
\label{thm:quartic-ascending-selector}
Let $p>3$ be prime with $p\equiv-1\pmod9$, and put
$k=\mathbb F_{p^2}$ and $L=\mathbb F_{p^4}$.
Suppose a deterministic algorithm selects a cyclic order-three kernel other
than a specified ascending kernel whenever its input is an
ordinary elliptic curve over $L$, strictly above its
$3$-isogeny volcano floor, with all its three-torsion rational
over $L$. The input includes a rational generator of the ascending
kernel, equivalently the dual kernel of a known descending incoming
isogeny. Assume a stated polynomial time bound on these
promised inputs.

Then one can deterministically construct an irreducible cubic
over $\mathbb F_p$, a maximal three-primary unit in $k$, and
cube roots of arbitrary promised cubes in $L$, in time
polynomial in $\log p$ using that algorithm.  No ordinary
starting curve is an additional input.  Conversely, promised
cube extraction in $L$ implements the stated selector.

The reduction runs a polynomial-time total completion of the
selector: each invocation is clocked at its promised time bound,
and every returned kernel is checked.  No behavior off the
geometric promise is assumed.
\end{proposition}

\begin{proof}
Use Proposition~\ref{prop:ordinary-prescribed-radicand} with
$\gamma=1$.  This constructs an ordinary curve
\[
 E:y^2+Axy+By=x^3\quad\text{over }L,
 \qquad A^3-27B=1.
\]
Its four three-isogeny kernels are explicit: one is generated
by $(0,0)$, and the other three have abscissas
\[
 -\frac{A^2+Az+z^2}{9},\qquad z\in\{1,\zeta,\zeta^2\}.
\]
All kernel lines are rational.  Since $\mu_3\subset L$ and
$(0,0)$ is a rational point, Frobenius acts as the identity on
$E[3]$, so their generators are rational as well.  Their
ordinates require only promised square extraction in the
already constructed quadratic tower.  That extraction is
deterministic by the same norm-and-two-square-roots formulas
used to construct the tower.

Start all four outgoing isogeny branches.  At least one first
edge descends: at a maximal-order vertex there are at most
two horizontal kernels, and below it there is one ascending
kernel and three descending ones.  On a branch whose first
edge descends, the dual of the incoming edge is the unique ascending
parent thereafter.  Every nonparent choice continues down.
These general ordinary-volcano facts hold in the split,
inert, and ramified local cases; exact sources are
\cite[Section~2.9, p.~514; Lemma~6, p.~515; Theorem~7 and
Remark~8, pp.~516--517]{Sutherland2013Volcanoes}.
We count cyclic kernel lines, retaining distinct edges whose
target $j$-invariants coincide.  The initial invariant lies
outside $\mathbb F_{p^2}$ and is therefore neither $0$ nor
$1728$.  After a descent the endomorphism order has conductor
divisible by three, so no vertex on that descending branch
has either exceptional invariant.  The multiplicities at
exceptional surface vertices described in the cited Remark~8
therefore do not affect this branch.
Parallel branching to find a descending path is the classical
\textsc{FindShortestPathToFloor} mechanism
\cite[Section~3.1, pp.~519--520]{Sutherland2013Volcanoes};
here the prescribed-radicand preprocessing gives the starting curve and
the stated selector supplies the root selections.

After every quotient, keep a rational generator of the dual
kernel.  At the initial step it is the image of a known
independent three-torsion point.  At later steps it is the
image of the previous parent generator, since the chosen
kernel is different.  A linear coordinate change moving this
point and its tangent to the origin gives a nonsingular Tate
equation $y^2+axy+by=x^3$ with the parent generated by $(0,0)$.

Test $R=a^3-27b$ before calling the selector.  If $R$ is a
noncube in $L$, keep it and stop.  If it is a cube, the
remaining division cubic splits.  Invoke the clocked selector,
check that its output is a three-kernel distinct from the
parent, and continue.  Discard a branch upon timeout or an
invalid output.  On the truly descending branch all calls
strictly before the floor satisfy the promise, so that branch
is never discarded.

The height is $O(\log p)$.  More explicitly, if $t$ is the
$L$-trace, the conductor of $\mathbb Z[\pi]$ has square
dividing $t^2-4p^4$, whose absolute value is at most $4p^4$.
Thus the height is at most
$\tfrac12\log_3(4p^4)$.  A fixed larger integer bound is
used as the clock for every branch.  At a floor reached by
descent there is only one rational kernel
\cite[Proposition~10, p.~518]{Sutherland2013Volcanoes}.
Frobenius fixes the known rational parent point and has
determinant $p^4\equiv1\pmod3$ on three-torsion.  Both
eigenvalues are therefore one.  It cannot be the identity,
which would make all four kernels rational.  Thus it is
nontrivial unipotent, fixing the ascending line and cycling
the other three.
The corresponding remaining cubic is irreducible and $R$
is a noncube.  Thus some branch stops successfully within
the bound, independently of every off-promise answer on the
other branches.  A successful noncube test is itself a
complete algebraic certificate for the subsequent output.

Put $u=N_{L/k}(R)$.  Then
\[
 u^{(p^2-1)/3}=R^{(p^4-1)/3}\ne1,
\]
so $u$ is a noncube in $k$.  If $a=v_3(p+1)$, then
$u^{(p^2-1)/3^a}$ has exact order $3^a$.  Its image in
$L$ has the same order, which is maximal because
$v_3(p^4-1)=a$.  Successive base-three digit extraction
therefore computes roots of every promised cube in $L$.

For a direct cubic output, let
\[
 n=N_{k/\mathbb F_p}(u),\qquad
 v=n^{1/3}\in\mathbb F_p,\qquad
 f(X)=X^3-3vX-\operatorname{Tr}_{k/\mathbb F_p}(u).
\]
The prime-field cube root is unique and computable by
exponentiation.  If $c^3=u$, the three roots of $f$ are
$\zeta^j c+v/(\zeta^j c)$, for $j=0,1,2$.
They are distinct: equality for two roots would imply
$u^2=n$, hence $u=\overline u$, which is impossible for
a noncube in $k$ because every nonzero prime-field scalar
is a cube there.  The cyclic extension $k(c)/k$ acts
transitively on these three roots.  Thus $f$ is irreducible
over $k$ and consequently over $\mathbb F_p$.

Conversely, on any promised selector input move the given
parent point to Tate form.  The split remaining cubic has
the displayed Tate roots with $z^3=a^3-27b$.  One promised
cube extraction supplies such a $z$, and hence a kernel
different from the parent.  All transformations and
verification have polynomial complexity.
\end{proof}

This proposition concerns ordinary inputs over $\mathbb F_{p^4}$
and permits the explicitly specified clocked total completion.
It does not apply the original prime-field-parent oracle to a
new domain without justification.  It proves an equivalence
for the enlarged selector problem and leaves the original
uniform construction question unchanged.

\subsection{Complementary torsion and the local endomorphism action}
The Weil pairing identifies exactly what a second high-order torsion
point would add to the known rational eigenspace.


\begin{proposition}[Complementary rational high torsion]
\label{prop:complementary-torsion-gate}
Let $p>3$ be prime with $p\equiv2\pmod3$, let
$a=v_3(p+1)$, $n=3^a$, and
$k=\mathbb F_p(\sqrt{-3})$.
Suppose an elliptic curve $E/\mathbb F_p$ is given with a
point $P$ of exact order $n$ in $E(\mathbb F_p)$.
Let $E^{(-3)}$ be its quadratic twist by $-3$.
Constructing a point of exact order $n$ in
$E^{(-3)}(\mathbb F_p)$ is deterministically polynomial-time
equivalent to constructing an element of exact order $n$ in
$k^\times$.  The polynomial bound is in $\log p$ and the
presented curve data.
\end{proposition}

\begin{proof}
Write $t$ for the trace of $E$.  The given point implies
$p+1-t\equiv0\pmod n$.  Since $n\mid p+1$, we have
$t\equiv0\pmod n$, so the twist also has group order
$p+1+t$ divisible by $n$.
The rational three-primary subgroup of either curve is
cyclic, because two independent rational three-torsion
points would give $\mu_3\subset\mathbb F_p$ by the
Weil pairing.  Thus a point of the required order exists
on the twist.

Suppose first that such a twist point has been constructed,
and let $Q\in E(k)$ be its image under the standard
quadratic-twist isomorphism.  Prime-field Frobenius fixes
$P$ and sends $Q$ to $-Q$.  If a multiple of $P$ equals
a multiple of $Q$, applying Frobenius shows that this
common point is killed by two.  It is also killed by the
odd integer $n$, and is therefore zero.  The two points
form a basis of $E[n]$, so $e_n(P,Q)$ has exact order $n$.

The pairing requires no auxiliary point search.  In the
shifted Miller formula take $S=2P+Q$:
\[
 e_n(P,Q)=
 \frac{f_P(Q+S)f_Q(-S)}{f_P(S)f_Q(P-S)}.
\]
The two evaluations of $f_P$ are outside $\langle P\rangle$,
because their $Q$-coefficients are one and two.  The two
evaluations of $f_Q$ are outside $\langle Q\rangle$,
because their $P$-coefficients are minus two and minus one.
These coefficients are units modulo $n$.  Hence none of
the intermediate Miller line evaluations meets a zero or
pole in the corresponding cyclic subgroup.  The pairing
is computable in polynomial time by the standard algorithm
\cite[Theorem~XI.8.1 and Remark~XI.8.2, pp.~393--397]{Silverman2009}.

Conversely suppose a maximal three-primary unit in $k$ is
given.  It gives roots of every promised cube in $k$ by
successive base-three digits.  In particular it permits
deterministic factorization of a split separable cubic over
$\mathbb F_p$: depress the cubic, use its square
discriminant and $\sqrt{-3}$ in Cardano's identity, and
extract the resulting promised cube in $k$.  If the depressed cubic has zero linear coefficient, take a cube
root of its constant term directly; if the chosen Cardano radicand
is zero, interchange the two radicands, or use the already known
zero root. These cases avoid division by zero. All three
roots returned by the identity lie in $\mathbb F_p$.
The required prime-field square roots are deterministic
because $-3$ is an explicit nonsquare.

We first obtain a rational point of order three on the
twist.  Put $P_0=[n/3]P$.  Frobenius on $E[3]$ has
distinct eigenvalues $1,-1$.  Its action on the four
kernel lines has two fixed lines and a transposed pair.
Therefore
\[
 \gcd(\psi_{3,E}(X),X^p-X)
\]
has degree two; removing the known line of $P_0$ leaves
the other linear factor.  Its nonzero points have
Frobenius $-1$, so their images on the quadratic twist
are rational.  Their coordinates use only square
extraction with the known nonsquare.  This gives a
rational twist point $Q_1$ of order three.

We describe a constant-degree lifting step on the twist,
now denoted $C$.  Suppose $Q_i\in C(\mathbb F_p)$ has
order $3^i$, where $i<a$.  Cyclicity and divisibility of
the rational group order by $3^a$ imply that $Q_i$ is
three-divisible.  Choose the known rational subgroup
$G=\langle[3^{i-1}]Q_i\rangle$, and compute the
degree-three isogeny
\[
 \phi:C\longrightarrow C'=C/G
\]
and its dual.  These are constant-degree rational maps
computable from the given kernel by the standard explicit
three-isogeny formulas.  The dual kernel is the Cartier
dual of the constant group $G$, hence has Frobenius $-1$
and no nonzero prime-field point.  The dual is therefore
a bijection on rational points, since the two isogenous
curves have the same rational group order.

Its inverse image $R_i$ of $Q_i$ is obtained by a
constant-degree abscissa equation and a linear gcd with
$X^p-X$.  There is exactly one prime-field abscissa:
an additional rational abscissa would give either another
rational preimage or a pair interchanged by negation;
the latter would force $Q_i=-Q_i$, impossible for a
nonzero point of odd order.  Recover the ordinate by
square extraction and verify the map.

Because $Q_i$ is three-divisible and
$\widehat\phi\phi=[3]$, the point $R_i$ lies in
$\phi(C(\mathbb F_p))$.  Its three preimages under
$\phi$ are rational and differ by $G$.  Their abscissas
are distinct: equality of two abscissas would make
the two points negatives, forcing a preimage of order
$3^{i+1}$ to have a double in $G$, which is impossible.
The abscissa equation is therefore a split separable
cubic over $\mathbb F_p$.  Use the preceding cubic
extraction procedure, recover an ordinate, and verify
the image.  This constructs $Q_{i+1}$ with
$[3]Q_{i+1}=Q_i$ and exact order $3^{i+1}$.
There are $a-1=O(\log p)$ such steps, each of constant
algebraic degree and polynomial arithmetic complexity.
\end{proof}

For the exact hard branch $v_3(p+1)=2$, a given rational
nine-torsion point lies in one Frobenius eigenspace; obtaining
a complementary rational nine-torsion point on the twist is
already equivalent to the original maximal-unit gate.  If
only a point of order nine is given when $a>2$, the forward
pairing yields a ninth root, not automatically a maximal unit.
The proposition does not assert that a single given rational
high-torsion point alone solves or is equivalent to the gate.

\begin{proposition}[The full local-order action on three-torsion]
\label{prop:local-order-descending-lines}
Suppose an ordinary three-adic Tate lattice is explicitly identified
with \(I_k=\mathbb Z_3+3^k\mathbb Z_3\omega\), where \(k\geq1\) and
\(\omega^2-T\omega+N=0\). On \(I_k/3I_k\), the image of the whole
local order \(I_k\) is \(\{aI+bJ:a,b\in\mathbb F_3\}\), where
\(J=\left(\begin{smallmatrix}0&0\\1&0\end{smallmatrix}\right)\).
Its unique common eigenline is the ascending line; a unit in this
algebra cyclically permutes the three remaining lines.
\end{proposition}
\begin{proof}
Put \(\vartheta=3^k\omega\). In the basis \((1,\vartheta)\),
multiplication by \(\vartheta\) is
\[
 \begin{pmatrix}0&-3^{2k}N\\1&3^kT\end{pmatrix},
\]
which reduces to \(J\). Thus every local-order action has the asserted
form. The common eigenline is \(\ker J=\operatorname{im}J\), generated
by \(\vartheta\). Enlarging the lattice by \(\vartheta/3\) gives
\(I_{k-1}\), identifying the ascending line in this lattice model.
The other lines have representatives \((1,c)\), \(c\in\mathbb F_3\);
the action of \(1+\vartheta\) sends \(c\) to \(c+1\).
\end{proof}

This is a statement about the explicit lattice model, rather than a
freeness assertion for an unspecified Tate module. A unit in its
reduction modulo three may represent an isogeny of degree prime to
three; it need not represent an automorphism of the elliptic curve.
Dividing Frobenius by further powers of three, whenever the quotient
is an endomorphism, remains in the same algebra and gives no second
eigenline. On the split coordinate algebra of the three descending
kernels, a nontrivial cyclic eigenvector has values
\((c,\zeta c,\zeta^2c)\) and cube \(c^3\). A Fourier projector is a
linear operator on that algebra, rather than a primitive algebra
idempotent. Obtaining an evaluation of the eigenvector still asks for
a cube root. This calculation does not exclude algorithms that use
coordinates or other data to choose a kernel.

\subsection{The first successful Tate model is a simplest-cubic test}

The success condition $v=1$ has an exact elementary
parametrization, so point counting is not needed to test it.
The statement below also includes supersingular curves; it
does not rely on the ordinary volcano description.

\begin{proposition}\label{prop:tate-simplest-cubic}
Let $p>3$, $p\equiv2\pmod3$, and $k\in\mathbb F_p$ with
$k\ne3$.  Set
\[
 E_k:\quad y^2+kxy+y=x^3,
 \qquad h_k(X)=X^3-(k+3)X^2+kX+1.
\]
Then $h_k$ is squarefree and is either split or irreducible
over $\mathbb F_p$.  It is irreducible exactly when
\[
 v_3\bigl(\#E_k(\mathbb F_p)\bigr)=1.
\]
If additionally $p\equiv-1\pmod9$, this is equivalent to
$v_3(t(E_k))=1$.
Every nonsingular Tate form $y^2+axy+by=x^3$ over
$\mathbb F_p$ is isomorphic over $\mathbb F_p$ to one of
these models, with $k=a/c$ and $c^3=b$; this cube root is
unique and computable by exponentiation.
\end{proposition}

\begin{proof}
The discriminant of $E_k$ is $k^3-27$, which is nonzero
precisely when $k\ne3$, because cubing is bijective in
$\mathbb F_p$.  The point $P=(0,0)$ has order $3$.
The quotient isogeny $\phi:E_k\to E'_k=E_k/\langle P\rangle$
has the formulas
\[
 E'_k:\ y^2+kxy+y=x^3-5kx-(k^3+7),
 \qquad \phi_x(x)=x+k/x+1/x^2.
\]
The corresponding second coordinate is
\[
 \phi_y(x,y)=y-
 \frac{(kx+2)y+(kx+1)^2}{x^3}.
\]
For example the $x$-formula follows by adding $P$ and $-P$
in the group law and summing their changes in $x$;
substitution gives the displayed quotient equation.
Its $3$-division polynomial is
\[
 \psi_{3,E'_k}(X)
 =(3X+k^2)\bigl(X^3-9kX-(k^3+27)\bigr).
\]
The point $Q'=(k+3,4)$ lies on $E'_k$ and on its
$3$-division polynomial, so has order $3$.
The equation $\phi_x(x)=k+3$ is exactly $h_k(x)=0$.

A root $x\in\mathbb F_p$ of $h_k$ is nonzero.
Choose a point above it in the fibre of $Q'$ or $-Q'$.
Frobenius fixes its $x$-coordinate, so fixes the point or
negates it.  Negation is impossible because its image is
a rational point of order $3$.  Hence the root gives a
rational preimage of $Q'$ after a possible sign change.
Conversely a rational preimage gives a root.
Thus $h_k$ has a rational root exactly when
$Q'\in\phi(E_k(\mathbb F_p))$.

Both rational groups have the same order and have cyclic
$3$-primary part: full rational $3$-torsion would, by the
Weil pairing, imply $\mu_3\subset\mathbb F_p$, contrary to
$p\equiv2\pmod3$.  The isogeny has rational kernel of
order $3$, so its image has index $3$ in $E'_k(\mathbb F_p)$.
If the common group order has valuation $m\ge1$, that image
contains the unique subgroup of order $3$ exactly when
$m\ge2$.  Since $Q'$ generates it, the asserted equivalence
follows.  Direct calculation gives
\[
 \operatorname{disc}(h_k)=(k^2+3k+9)^2.
\]
The quadratic factor $k^2+3k+9$ never vanishes in
$\mathbb F_p$, since its discriminant $-27$ is a nonsquare.
The square discriminant restricts the Frobenius permutation
to $A_3$, giving the split-or-irreducible alternative.
Finally $p+1\equiv0\pmod9$ gives the trace equivalence,
and the change $x=c^2X$, $y=c^3Y$ gives the normalization.
\end{proof}

The same cubic is an exact one-parameter Eisenstein residue
test.  If $\zeta^2+\zeta+1=0$ in $\mathbb F_{p^2}$, put
$z=k-3\zeta$.  With
$w=(Y+\zeta)/(Y+\zeta^2)$, direct expansion gives
\begin{equation}\label{eq:tate-eisenstein-radicand}
 (w-1)^3h_k\left(\frac{\zeta-w\zeta^2}{w-1}\right)
 =(k-3\zeta)w^3-(k-3\zeta^2).
\end{equation}
Neither excluded denominator value is a root, because
$k-3\zeta$ and $k-3\zeta^2$ are nonzero.  A cubic
over $\mathbb F_p$ is irreducible if and only if it
remains irreducible over its quadratic extension.
The displayed transformation and Kummer theory therefore
give
\[
 h_k\text{ irreducible over }\mathbb F_p
 \quad\Longleftrightarrow\quad
 k-3\zeta\text{ is a noncube in }\mathbb F_{p^2}.
\]
Indeed the Kummer radicand is $z^p/z=N(z)/z^2$,
and $N(z)\in\mathbb F_p^*$ is a cube because
$p\equiv2\pmod3$.  Here $N(z)=k^2+3k+9$.

For integral $k$, the identity
\[
 -h_k(-X)=X^3+(k+3)X^2+kX-1
\]
identifies the family with the simplest cubic polynomial
of parameter $-k-3$; see
\cite[\S1, (1.1)--(1.3)]{HashimotoAoki2024}.
Its irreducibility over $\mathbb Q$ follows directly from
the rational-root test, and the displayed square
discriminant then gives a cyclic cubic field $L_k$.
If $m_k$ is its conductor, the conductor-discriminant
formula and the polynomial-index formula give
\[
 \operatorname{Disc}(L_k)=m_k^2,
 \qquad m_k\mid k^2+3k+9.
\]
Suppose now $p\equiv-1\pmod9$ and $h_k\bmod p$ is
irreducible.  A primitive cubic Dirichlet character
of conductor $m_k$ is nontrivial at $p$.
Its local conductor at $3$ divides $9$, and that local
component is trivial at $p\equiv-1\pmod9$.
At every other prime its nontrivial local component
has conductor a prime $\ell\equiv1\pmod3$ and kernel
the cubes modulo $\ell$.
Consequently some prime
\[
 \ell\le m_k\le k^2+3k+9
 \quad\text{satisfies}\quad
 \ell\equiv1\pmod3,
 \qquad p^{(\ell-1)/3}\not\equiv1\pmod\ell.
\]
Thus a uniformly polylogarithmic bound on an \emph{integer}
parameter $k$ successful in this normalized family would
give the earlier good-auxiliary-prime bound.
Normalization can turn small original coefficients $a,b$
into a large residue $k=a/b^{1/3}$, so this height implication
is specifically about the normalized integer parameter.



\subsection{A supersingular point and an affine-line noncube}
A fixed supersingular curve gives a second geometric description of
the missing cubic witness.  Its point-to-unit map uses the classical
Weil pairing and the order-three automorphism; its reverse map uses
explicit three-isogeny descent.  The proof below identifies both
operations and the precise search that remains.


\begin{proposition}[A supersingular point formulation]
\label{prop:supersingular-point-gate}
Let $p>3$ be prime with $p\equiv2\pmod3$, put
$a=v_3(p+1)$, and present
$k=\mathbb F_p[\delta]$, where $\delta^2=-3$.
For
\[
 E/\mathbb F_p:y^2=x^3-3,
\]
constructing a point of exact order $3^a$ in $E(\mathbb F_p)$
and constructing an element of exact order $3^a$ in $k^\times$
are deterministically polynomial-time equivalent.
Moreover, for every affine point $P=(X,Y)\in E(\mathbb F_p)$,
\[
 P\in3E(\mathbb F_p)
 \quad\Longleftrightarrow\quad
 Y+\delta\in(k^\times)^3.
\]
Thus a nondivisible point can be sought by scanning the single
prime-field parameter $Y$, with
$X=(Y^2+3)^{1/3}$ uniquely computable in $\mathbb F_p$.
\end{proposition}

\begin{proof}
Cubing permutes $\mathbb F_p$, so each $Y$ gives exactly one
affine point and $\#E(\mathbb F_p)=p+1$.  Its $3$-primary
subgroup is cyclic: two independent rational points of order
three would give a primitive cube root in $\mathbb F_p$ by
the Weil pairing.  In particular it has order $3^a$.

First suppose $P$ has that exact order.  With
$\zeta=(-1+\delta)/2$, the automorphism
\[
 \iota:E_k\longrightarrow E_k,
 \qquad (x,y)\longmapsto(\zeta x,y)
\]
has order three.  Write $P_0=[3^{a-1}]P$.
The line generated by $P_0$ cannot be invariant under $\iota$:
an order-three action on a one-dimensional $\mathbb F_3$-space
would be trivial, whereas the nonzero fixed points of $\iota$
have $x=0$ and $y^2=-3$, and hence are not rational over
$\mathbb F_p$.  Thus $P_0,\iota(P_0)$ are independent, so
$P,\iota(P)$ form a basis of $E[3^a]$.
Their Weil pairing has exact order $3^a$ in $k^\times$.

This pairing is computable without a point-search subroutine.
Use the explicit point of order two
\[
 S=(3^{1/3},0)\in E(\mathbb F_p).
\]
It is outside the odd-order subgroup generated by $P$ and
$Q=\iota(P)$.  Miller's algorithm and the shifted evaluation
formula compute
\[
 e_{3^a}(P,Q)
 =\frac{f_P(Q+S)f_Q(-S)}{f_P(S)f_Q(P-S)},
 \qquad
 \operatorname{div}(f_P)=3^a(P)-3^a(O),
\]
in polynomial time.  All evaluation points have a nonzero
two-primary component, so the Miller line functions encounter
neither zeros nor poles from the three-primary subgroups.
See \cite[Chapter~III, Section~8; Theorem~XI.8.1 and
Remark~XI.8.2, pp.~393--397]{Silverman2009}.

We next prove the divisibility criterion and the constructive
reverse implication.  Put $E':y^2=x^3+81$.  The explicit
degree-three isogeny
\[
 \phi:E\longrightarrow E',\qquad
 (x,y)\longmapsto
 \left(\frac{x^3-12}{x^2},
       \frac{y(x^3+24)}{x^3}\right)
\]
has no nonidentity rational kernel point because $-3$ is a
nonsquare.  It is therefore a bijection on prime-field points,
as both curves have $p+1$ points.  Its dual is
\[
 \widehat\phi:E'\longrightarrow E,\qquad
 (z,w)\longmapsto
 \left(\frac{z^3+324}{9z^2},
       \frac{w(z^3-648)}{27z^3}\right),
 \qquad \widehat\phi\phi=[3].
\]
These identities follow by substitution; equivalently they are
the standard $j=0$ three-isogeny and its dual, with the target
scaling made explicit.  The associated map $Y+\delta$ is a
specialization, up to conjugation, of the classical explicit
three-isogeny descent map; see
\cite[Definition~1.3 and Proposition~1.4, p.~372]{CohenPazuki2009}.
We verify the required finite-field case directly below.
Consequently
$\widehat\phi(E'(\mathbb F_p))=3E(\mathbb F_p)$.

For $P=(X,Y)$ with $Y\ne0$, a preimage under the dual exists
over $\mathbb F_p$ exactly when
\[
 f_P(z)=z^3-9Xz^2+324
\]
has a root there.  The ordinate is then
$27Yz^3/(z^3-648)$.  Its denominator cannot vanish for such
a root when $Y\ne0$.  The cubic
\[
 h_P(w)=w^3-3Xw-2Y
\]
has discriminant $324$ and has roots corresponding to those
of $f_P$ through
\[
 z=3(w^2-X),\qquad w=\frac{6Y}{z-6X}.
\]
Substitution proves the identities, and the denominators are
nonzero when $Y\ne0$.

Set $\gamma=Y+\delta$, whose norm is $X^3$.
If $c^3=\gamma$ in $k$, then the norm of $c$ is $X$, because
cubing is bijective in $\mathbb F_p$.  Hence
$w=c+X/c$ belongs to $\mathbb F_p$ and is a root of $h_P$.
Conversely, if $\gamma$ is a noncube, the cyclic cubic
extension $k(c)/k$ acts transitively on the three distinct
values $c+X/c$, so $h_P$ is irreducible over $k$ and over
$\mathbb F_p$.  For distinct cube roots equality of these
values would force their product to be $X$, and cubing would
give $\gamma=\overline\gamma$, impossible since $\delta\ne0$.
This proves the criterion for $Y\ne0$.  For $Y=0$, the point
has order two and is three-divisible; $\delta$ is a cube
because $\delta^2=-3\in\mathbb F_p^\times$ is a cube and
the cube-class group has order three.

Finally, a known maximal three-primary unit gives promised
cube extraction in $k$ by successive base-three digits.
Start with the explicit point
\[
 P_1=(12^{1/3},3)\in E(\mathbb F_p)[3]\setminus\{O\}.
\]
For a point $P_i=(X,Y)$ of order $3^i$, $i<a$, cyclicity
guarantees a rational preimage under $[3]$.  The proved
criterion therefore says that $Y+\delta$ is a promised cube.
Extract $c^3=Y+\delta$ using the maximal unit, and form
$w=c+X/c$ and $z=3(w^2-X)$.  These lie in $\mathbb F_p$
and the displayed formulas give a point of $E'$ mapping to
$P_i$ under the dual.  The resulting $E'$-point has a unique prime-field
preimage under $\phi$; its abscissa is the unique root over
$\mathbb F_p$ of
\[
 Z^3-X' Z^2-12,
\]
found as a linear gcd with $Z^p-Z$.
Indeed the lifted point of $E'$ has odd order, since its image
has odd order and the dual kernel has order three; its ordinate
is therefore nonzero.  For any prime-field root $z$ of the
fiber cubic, $z\ne0$ and $z^3+24\ne0$: the latter equality
would force the target ordinate to vanish.  The ordinate
$Y'z^3/(z^3+24)$ is then rational and gives a preimage.
The bijectivity of $\phi$ on rational points therefore makes
this root unique.  Recover its ordinate by that formula.
This constructs $P_{i+1}$ with $[3]P_{i+1}=P_i$.
After $a-1=O(\log p)$ iterations its order is $3^a$.
All quadratic extraction uses the known nonsquare $-3$.
\end{proof}

For the hard branch $p\equiv-1\pmod9$, if a small integer
$Y$ yields a nondivisible point, the Eisenstein integer
$Y+\sqrt{-3}$ has norm $Y^2+3$.  The established cubic
reciprocity argument then gives a good auxiliary prime
$\ell\le Y^2+3$.  Group operations applied only to points
already in $3E(\mathbb F_p)$ cannot leave that subgroup.
Thus this supersingular formulation makes the remaining
noncube search explicit; it is an equivalence, not a
strictly weaker sufficient condition.

\begin{proposition}[Scalar Frobenius in the supersingular walk]
\label{prop:supersingular-walk-scalar}
Let $p>3$, $p\equiv2\pmod3$, and let $E_0/\mathbb F_p$
be supersingular.  Every curve reached by isogenies defined
over $k=\mathbb F_{p^2}$ has $k$-Frobenius $[-p]$.
Consequently every such curve has all its three-torsion
rational over $k$.  Whenever it is put in Tate form using
a rational point of order three, its remaining cubic splits
and its Tate radicand is a cube.  Such a walk never reaches
an ordinary floor producing a noncube.
\end{proposition}

\begin{proof}
For a supersingular prime-field curve, Frobenius and its
dual are purely inseparable.  Their sum $[t_0]$ has zero
differential, so $p\mid t_0$.  Hasse's bound and $p>3$
give $t_0=0$, hence prime-field Frobenius satisfies
$\pi^2=[-p]$.  If $\phi:E_0\to E_1$ is defined over $k$,
then $\pi_{k,E_1}\phi=\phi\pi_{k,E_0}=[-p]\phi$.
Surjectivity of $\phi$ gives $\pi_{k,E_1}=[-p]$.
Induction handles the whole walk.  Since $-p\equiv1\pmod3$,
Frobenius fixes all three-torsion points.  The explicit Tate
root formula then gives the last assertions.
\end{proof}

The workbench program \emph{Supersingular points and cubic classes} checks the isogeny maps, divisibility criterion, and maximal-order lift on the stated finite examples.

\subsection{A ramified ring-class alternative}

A whole CM class polynomial avoids choosing one CM $j$-invariant.
The resulting sufficient condition still contains a cyclic
$r$-extension splitting condition, which is distinct from quadratic
splitting in the CM field.

\begin{proposition}\label{prop:cm-ring-class-gate}
Let $r$ be an odd prime and $D<-4$ a negative fundamental
discriminant with $r\mid D$ and $r\nmid h(D)$.
Put $K=\mathbb Q(\sqrt D)$, $O=O_K$, and $O_r=\mathbb Z+rO$.
Let $p\nmid rD$ be an odd prime splitting as
$pO=\mathfrak p\overline{\mathfrak p}$.
Let $e$ be the order of $[\mathfrak p]$ in $\operatorname{Cl}(O)$,
and choose $\pi=(t+v\sqrt D)/2$ with
$\mathfrak p^e=(\pi)$.
The reduction of the class polynomial $H_{r^2D}$ modulo $p$ is
squarefree.  If $r\nmid v$, all its irreducible factors have degree
$re$, and their number is $h(D)/e$.
If $r\mid v$, all its factors have degree $e$, and their number is
$r h(D)/e$.
Thus it satisfies the degree-$r$ factor-count property exactly
when $r\nmid v$, equivalently when
\[
 \pi\bmod rO\notin\mathbb F_r^*
 \quad\Longleftrightarrow\quad
 v_r(t^2-4p^e)=1.
\]
In the positive case, a presentation of $\mathbb F_{p^r}$ can be
constructed deterministically in time polynomial in $r$, $|D|$,
and $\log p$.
\end{proposition}

\begin{proof}
Since $D<-4$, both orders have unit group $\{1,-1\}$.
The conductor exact sequence is
\[
1\longrightarrow (O/rO)^*/\mathbb F_r^*
 \longrightarrow\operatorname{Cl}(O_r)
 \longrightarrow\operatorname{Cl}(O)\longrightarrow1;
\]
see \cite[\S3, (3)]{BrokerLauterSutherland2012}.
Here its conductor-coprimality hypothesis holds because $O$ is
maximal.  Since $r$ ramifies and is odd,
$O/rO\simeq\mathbb F_r[\epsilon]/(\epsilon^2)$.
Modulo scalar units, its units are $1+b\epsilon$, with multiplication
given by addition of $b$.  The kernel therefore has order $r$,
and $h(r^2D)=r h(D)$.

Let $g$ be the class of $\mathfrak p\cap O_r$.
Its image has order $e$, which is prime to $r$.
Consequently $\operatorname{ord}(g)$ is either $e$ or $re$.
Its $e$-th power in the kernel is the class of $\pi\bmod rO$,
up to inversion according to the Artin-map convention.
It is trivial exactly when that residue is scalar.
The image of $\sqrt D$ is a nonzero nilpotent generator of $O/rO$,
so this is equivalent to $r\mid v$.
The identity
\[
 t^2-4p^e=Dv^2
\]
and $v_r(D)=1$ give the valuation criterion.

By CM reciprocity, the roots of $H_{r^2D}$ form a regular
$\operatorname{Cl}(O_r)$-torsor.
At $p$ split in $K$ and prime to the conductor and discriminant,
ordinary reduction preserves this torsor and its distinct roots;
Frobenius acts by translation by $g$ or $g^{-1}$.
The factorization statement is the Frobenius computation in
\cite[\S2.3]{CremonaSutherland2024}; its Proposition~4 establishes
squarefreeness for every split prime under consideration.
Every orbit therefore has length $\operatorname{ord}(g)$.
Dividing $r h(D)$ by the two possible orbit lengths gives the stated
factor counts.  In the positive case $h(D)/e$ is prime to $r$, so
the factor-count construction applies.

For the complexity assertion, the elementary complex-evaluation
construction computes $H_\Delta$ deterministically in time
polynomial in $|\Delta|$; a coarse bound suffices here.
Enumerate its primitive reduced quadratic forms, evaluate the
corresponding $j$-values with certified precision, and multiply the
linear factors.  The form bounds give polynomially many roots,
$|q|\le e^{-\pi\sqrt3}$ at every reduced CM point, and
$|j|\le e^{O(\sqrt{|\Delta|})}$.
The exponentially convergent $j$-series therefore gives
polynomial-size coefficient intervals that can be narrowed below
$1/2$ and rounded using CM integrality.
Taking $\Delta=r^2D$ proves the required polynomial bound for
constructing and reducing $H_{r^2D}$.
The remaining algebra has polynomial cost in its degree and
$\log p$.
\end{proof}

The units hypothesis is material: for $D=-3,r=3$ the additional
units remove the proposed order-$r$ kernel.
The ramified-prime calculation above uses the exact sequence
directly, not the unramified-prime hypothesis of
\cite[Lemma~3.1]{BrokerLauterSutherland2012}.

Let $L/K$ be the Hilbert class field and $L_r/K$ the ring class
field of conductor $r$.
The positive condition in Proposition~\ref{prop:cm-ring-class-gate}
says that a prime of $L$ above $\mathfrak p$ has nontrivial
Frobenius in $\operatorname{Gal}(L_r/L)\simeq C_r$.
Rational Dirichlet GRH can control the quadratic splitting test for
$K$, but no polynomial bound for a suitable $D$ with this additional
condition is established here.
The fact that $L_r/K$ is abelian does not identify its characters
with rational Dirichlet characters.
A uniform existence bound $|D|\le\operatorname{poly}(r,\log p)$
for a positive case would suffice: enumerate the discriminants,
count their reduced forms, and apply deterministic distinct-degree
decomposition to $H_{r^2D}\bmod p$ to test the factor-count property.
This search requires neither a generator $\pi$ nor a chosen CM root.

\begin{example}\label{ex:cm-15-135}
Take $D=-15$ and $r=3$.  Here $h(-15)=2$ and $h(-135)=6$.
The primes $17$ and $107$ are both $17\pmod{45}$, but
\[
 4\cdot17^2=14^2+15\cdot8^2,
 \qquad
 4\cdot107^2=106^2+15\cdot48^2.
\]
Both prime ideals have class order $e=2$.
The first coefficient $v$ is prime to $3$, and the second is not.
Accordingly,
\[
 H_{-135}\bmod17
 =X^6+X^5+10X^4+16X^3+15X^2+9X+3
\]
is irreducible, whereas
\[
 H_{-135}\bmod107
 =X^6+46X^5+88X^4+9X^3+70X^2+80X+22
\]
has exactly three distinct irreducible quadratic factors.
For a root $\alpha$ of the degree-six polynomial over $\mathbb F_{17}$,
the element $\beta=\alpha+\alpha^{17^3}$ has minimal polynomial
\[
 X^3+X^2+4X+8\in\mathbb F_{17}[X].
\]
The workbench's \emph{CM class-polynomial and subfield certificate}
derives $H_{-15}$ and $\Phi_3$ from exact modular $q$-series and
certified CM intervals, computes the integer resultant, and verifies
both factor patterns and this subfield polynomial by Frobenius
arithmetic and gcds.
The independently derived input polynomials agree with the
discriminant-$-15$ singular moduli displayed in
\cite[p.~1 of the author manuscript]{Duke2006} and the numerical
$\Phi_3$ coefficients in the official PARI source
\cite{PARIPhi32026}.
\end{example}

\subsection{A class-group period formulation}

The preceding construction can be expressed more generally by taking
periods of CM values over class-group cosets.  Decompositions of this
kind are classical; see \cite[\S2.5]{Sutherland2012CM}, which develops
the constructions of Enge--Morain and Hanrot--Morain.  The proposition
below states the particular finite-field selection condition that is
needed here.  It does not assert a bound for a discriminant satisfying
that condition.

\begin{proposition}\label{prop:cm-period-gate}
Let $\Delta<-4$ be an imaginary quadratic order discriminant,
$A=\operatorname{Cl}(O_\Delta)$, and $h=|A|$.
Let $p\nmid\Delta$ be a prime split in
$K=\mathbb Q(\sqrt\Delta)$, and let $g\in A$ be the class of a
prime above $p$.
Let $r$ be an odd prime.  Write $rA=\{a^r:a\in A\}$ in
multiplicative notation.  If $g\notin rA$, then a presentation of
$\mathbb F_{p^r}$ can be constructed deterministically in time
polynomial in $|\Delta|$ and $\log p$.

More precisely, index the CM $j$-values by $A$ as $(j_b)_{b\in A}$.
For an index-$r$ subgroup $H\le A$ and $1\le k\le h-1$, define
\[
 \theta_{aH,k}=\sum_{b\in aH}j_b^k,
 \qquad
 P_{H,k}(X)=\prod_{aH\in A/H}(X-\theta_{aH,k}).
\]
Every $P_{H,k}$ belongs to $\mathbb Z[X]$ and has degree $r$.
If $g\notin rA$, at least one such polynomial is irreducible
modulo $p$.
All these polynomials can be enumerated and computed in time
polynomial in $|\Delta|$; the class $g$ need not be computed.
Conversely, if $g\in rA$, every $P_{H,k}\bmod p$ is a product
of linear factors, with repetitions allowed.
\end{proposition}

\begin{proof}
Let $L$ be the ring class field.  The CM values form a regular
$A=\operatorname{Gal}(L/K)$-torsor.
Choose the indexing so that complex conjugation sends $j_b$ to
$j_{b^{-1}}$; equivalently, take the principal CM value as the
initial value.
Translations by $A$ and inversion both permute the cosets of $H$.
They therefore permute the $\theta_{aH,k}$.
The coefficients of $P_{H,k}$ are algebraic integers fixed by
$\operatorname{Gal}(L/\mathbb Q)=A\rtimes\langle c\rangle$,
where $c$ acts by inversion.  Hence they are rational integers.
This also explains why passing from coefficients over $K$ to
coefficients over $\mathbb Q$ is valid in this CM setting.

At a prime above $p$, ordinary CM reduction gives $h$ distinct
residues $\bar j_b$; this is the injectivity statement in
\cite[\S\S2.3--2.4, Proposition~4]{CremonaSutherland2024}.
It is stronger than the assertion that $p$ is unramified in $L$.
Frobenius acts on these residues by $b\mapsto gb$, or by its inverse
depending on the Artin-map convention.

The group $A/rA$ is a vector space over $\mathbb F_r$.
If $g\notin rA$, choose a linear functional nonzero on its image,
and let $H$ be the inverse image of its kernel.
Then $[A:H]=r$ and $gH$ generates $A/H$.
Suppose that for every $0\le k\le h-1$ the two sums
\[
 \sum_{b\in H}\bar j_b^k,
 \qquad
 \sum_{b\in gH}\bar j_b^k
\]
were equal.
The difference of the two coset-indicator vectors would then be
annihilated by the Vandermonde matrix
$V=(\bar j_b^k)_{0\le k<h,\,b\in A}$.
Its determinant is a nonzero product of differences of the distinct
$\bar j_b$, so $V$ is invertible.
The indicator difference is nonzero, because $H$ and $gH$ are
disjoint.  This is a contradiction in every characteristic,
including when $p$ divides $|H|$.
The sums for $k=0$ are equal, so some $1\le k\le h-1$
separates the two cosets.

For this $k$, the Frobenius orbit of $\bar\theta_{H,k}$ is
nonconstant and its length divides the prime $r$.
Its length is therefore $r$.
These are exactly the $r$ roots of $P_{H,k}\bmod p$, proving
irreducibility.
If instead $g\in rA$, its image is trivial in every index-$r$
quotient of $A$.  Frobenius then fixes every coset sum, so every
root of every $P_{H,k}\bmod p$ lies in $\mathbb F_p$.

It remains to justify the deterministic complexity assertion.
Put $M=|\Delta|$.
The reduced positive primitive forms have
$1\le a\le\sqrt{M/3}$, $|b|\le a$, and
$c=(b^2+M)/(4a)\ge a$.
Enumerating these integer pairs and applying the usual boundary
conventions gives all $h=O(M)$ classes.
Composition and reduction give the entire class-group table using
polynomially many operations on polynomial-size integers.
From this table one constructs $A/rA$ by exhaustive group operations
and elementary linear algebra.
If its dimension is $d$, the number of index-$r$ subgroups is
$(r^d-1)/(r-1)\le h$.
All of them can thus be enumerated in polynomial time.
This uses an explicitly enumerated group, rather than a discrete
logarithm oracle.

The CM values and the period polynomials can be computed by
certified complex evaluation and integer rounding, with a coarse
polynomial bound sufficient here.
At a reduced CM point,
$q=\exp(2\pi i\tau)$ satisfies
\[
 e^{-\pi\sqrt M}\le |q|\le e^{-\pi\sqrt3}<1/200.
\]
The absolutely convergent identity
\[
 j(q)=\frac{\left(1+240\sum_{n\ge1}\sigma_3(n)q^n\right)^3}
 {q\prod_{n\ge1}(1-q^n)^{24}}
\]
gives $|j(q)|\le 2^{O(\sqrt M)}$.
For example, $\sigma_3(n)\le n^4$ bounds the series tails
geometrically after a polynomial factor; the omitted logarithm
of the product is also geometrically bounded.
Thus $N$ certified absolute bits of each $j$ require only
$O(N+\sqrt M+\log(N+2))$ terms, with polynomial-size rational
interval arithmetic.
For $k<h$, each period has size at most
$h\,2^{O(h\sqrt M)}$, and a coefficient of its degree-$r$
period polynomial has logarithmic size
$O(rh\sqrt M+r\log(h+1))$.
The elementary estimates
$|z^k-w^k|\le k\max(|z|,|w|)^{k-1}|z-w|$
and the product telescoping identity show that a polynomial number
of additional guard bits makes every final coefficient interval
have radius less than $1/4$.
Its unique integer is then recovered exactly.
This establishes a coarse rigorous polynomial bound; it does not
invoke the heuristic near-optimal floating-point error estimate
discussed in \cite[\S1, Corollary~1.3]{Enge2009ClassPolynomials}.

Finally, reduce the resulting integer polynomials modulo $p$ and
test them for irreducibility.
For a monic polynomial $P$ of prime degree $r$, the conditions
\[
 X^{p^r}\equiv X\pmod P,
 \qquad \gcd(P,X^p-X)=1
\]
are equivalent to irreducibility.
Modular powering and gcd computation have polynomial cost in
$r$ and $\log p$, and $r\le h$ whenever an index-$r$ subgroup
exists.  At least one test succeeds in the positive case.
No chosen CM root and no computation of $g$ are required.
\end{proof}

This gives another precise sufficient input for the original field
construction problem: a discriminant
\[
 |\Delta|\le\operatorname{poly}(r,\log p),\qquad
 \left(\frac{\Delta}{p}\right)=1,
 \qquad [\mathfrak p]\notin r\operatorname{Cl}(O_\Delta).
\]
A uniform existence bound of this form would suffice, since one
could enumerate all discriminants in the range and test the period
polynomials.  No such bound is established from rational Dirichlet
GRH here.
In the ramified-conductor situation of
Proposition~\ref{prop:cm-ring-class-gate}, the $r$-Sylow subgroup
of $\operatorname{Cl}(O_r)$ has order $r$.
Consequently its positive case is exactly the condition
$g\notin r\operatorname{Cl}(O_r)$ in the present proposition.


\subsection{An equivalent cubic-field gate when \texorpdfstring{$r=3$}{r=3}}

For $r=3$, the class-group formulation has an exact interpretation
in terms of cubic fields.  This is a specialization of the classical
ring-class correspondence, not a new source of a bounded witness.

\begin{proposition}\label{prop:cubic-cm-equivalent-gate}
Let $p$ be a prime and $B\ge1$ a real number.  The following
conditions are equivalent.
\begin{enumerate}
\item There is an imaginary quadratic order discriminant
      $\Delta<-4$ with $|\Delta|\le B$ and $p\nmid\Delta$, such
      that $p$ splits in $K=\mathbb Q(\sqrt\Delta)$ and the class
      $g$ of a prime above $p$ does not belong to
      $3\operatorname{Cl}(O_\Delta)$.
\item There is a complex cubic number field $M$ with
      $|\operatorname{Disc}(M)|\le B$ and
      $p\nmid\operatorname{Disc}(M)$ in which $p$ is inert.
\end{enumerate}
In the forward direction, $M$ can be chosen with
$\Delta=c^2\operatorname{Disc}(M)$ for an integer $c\ge1$.
In the reverse direction one may take
$\Delta=\operatorname{Disc}(M)$.
Consequently, a uniform bound $B\le\operatorname{poly}(\log p)$
for either existence condition would suffice for deterministic
polynomial-time construction of $\mathbb F_{p^3}$.
\end{proposition}

\begin{proof}
Write $A=\operatorname{Cl}(O_\Delta)$.
The class-field correspondence in
\cite[\S5.1, Lemma~34 and Proposition~35, pp.~256--258]
{BhargavaVarma2016}
identifies index-three subgroups $H\le A$ with isomorphism
classes of cubic fields $M$ satisfying
$\Delta=c^2\operatorname{Disc}(M)$.
In this correspondence, the cyclic cubic extension of $K$
cut out by $H$ is the normal closure of $M$ over $\mathbb Q$.
The discriminant is negative, so $M$ has one real and two complex
embeddings, and its normal closure has group
$C_3\rtimes C_2\simeq S_3$.
The involution acts on $C_3$ by inversion.
The ray-class minus-part in the cited Lemma~34 is essential:
the entire ray class group is not being identified with the ring
class group.

Suppose the first condition holds.
Since $A/3A$ is a vector space over $\mathbb F_3$, there is an
index-three subgroup $H$ not containing $g$.
In the corresponding $S_3$ extension, Frobenius at $p$ lies in
the normal subgroup $C_3$, because $p$ splits in $K$.
Its image is nontrivial because $g\notin H$.
It is therefore a three-cycle on the three embeddings of $M$,
which says precisely that $p$ is inert in $M$.
The discriminant relation gives both
$|\operatorname{Disc}(M)|\le B$ and unramifiedness at $p$.

Conversely, let $M$ satisfy the second condition, and put
$\Delta=\operatorname{Disc}(M)$.
Its quadratic resolvent ring is the imaginary quadratic order of
discriminant $\Delta$.
The same correspondence supplies an index-three subgroup of
$A=\operatorname{Cl}(O_\Delta)$ cutting out the normal closure
of $M$ over its quadratic resolvent field.
Inertness means that Frobenius is a three-cycle.
It is even, so $p$ splits in the quadratic resolvent field, and
it is nontrivial in the quotient $A/H$.
Thus $g\notin H$, and in particular $g\notin3A$.
The correspondence includes nonfundamental discriminants and
ramification at $3$, so neither direction discards these cases.

For the algorithmic consequence, enumerate all negative order
discriminants in the asserted range and apply
Proposition~\ref{prop:cm-period-gate}.
The work is polynomial in $B$ and $\log p$.
The existence formulation does not require a given cubic field,
its defining polynomial, or the ideal class $g$.
\end{proof}

\paragraph{What a uniform counting theorem currently supplies.}
Let $N_-(B;p)$ count complex cubic fields of absolute discriminant
less than $B$ in which $p$ is unramified and inert.
Theorem~1.3 and Table~1 of
\cite[pp.~2--3 of the checked author manuscript]
{BhargavaTaniguchiThorne2024}
give, for every fixed $\epsilon>0$, uniformly in the prime $p$,
\begin{align*}
 N_-(B;p)
 &=\frac{p^2}{12\zeta(3)(p^2+p+1)}B
   +\kappa(p)B^{5/6}
   +O_\epsilon\!\left(p^{2/3}B^{2/3+\epsilon}\right),\\
 \kappa(p)
 &=\frac{4\sqrt3\,\zeta(1/3)}
         {15\Gamma(2/3)^3\zeta(5/3)}
   \frac{1-p^{-1/3}}{1-p^{-5/3}}.
\end{align*}
The constants in the error term depend on $\epsilon$, not on $p$.
The leading coefficient is bounded below by a positive absolute
constant, and $\kappa(p)$ is bounded.
For every fixed $\delta>0$, this estimate therefore proves
$N_-(B;p)>0$ once $B\ge C_\delta p^{2+\delta}$.
For example, take
$\epsilon=\delta/(6(2+\delta))$ and compare the error with $B$.
At $B=(\log p)^C$ for a fixed $C$, this error estimate does not
certify positivity.
This is a limitation of the displayed estimate; it does not rule
out a stronger estimate, an argument using Assumption~A, or another
construction.

\paragraph{The trace still to be controlled.}
For any finite family $\mathcal C$ of complex cubic fields
unramified at $p$, let $D_M$ be the fundamental discriminant of
the quadratic resolvent of $M$.
Let $a_M(p)$ be the trace of the standard two-dimensional
representation of $S_3$ at Frobenius.
The three splitting types $(111),(12),(3)$ give trace values
$2,0,-1$, respectively, while the quadratic character
$\chi_{D_M}(p)$ takes values $1,-1,1$.
Consequently the exact identity
\[
 3\#\{M\in\mathcal C:p\text{ is inert in }M\}
 =|\mathcal C|+\sum_{M\in\mathcal C}\chi_{D_M}(p)
               -\sum_{M\in\mathcal C}a_M(p)
\]
isolates the remaining term.
Class-number and reflection formulas enumerate the available
cubic quotients; they do not, on their own, bound this last
Frobenius-trace sum for a given prime $p$.
No estimate removing the numerical $p^{2/3}$ dependence in the
polylogarithmic discriminant range is established here from
Assumption~A.

\subsection{A quadratic-unit form of the cubic CM gate}

The quadratic-unit formulation makes a residue test and its
output polynomial explicit.  The bounded-parameter selection
condition remains separate.

\begin{lemma}\label{lem:unit-norm-cubic}
Let $p\equiv2\pmod3$ be prime, let
$u\in\mathbb F_{p^2}^*$, and suppose
$u^{p+1}=n\in\{1,-1\}$.  Put $t=u+u^p$.
Then
\[
 u\notin(\mathbb F_{p^2}^*)^3
 \quad\Longleftrightarrow\quad
 X^3-3nX-t\text{ is irreducible over }\mathbb F_p.
\]
The assertion includes scalar $u$ and repeated-discriminant
specializations.
\end{lemma}

\begin{proof}
If $u=w^3$ with $w\in\mathbb F_{p^2}^*$, then
$N(w)^3=n$.  Cubing is bijective on $\mathbb F_p^*$,
so $N(w)=n$.  Thus $x=w+n/w=w+w^p$ belongs to
$\mathbb F_p$ and satisfies $x^3-3nx=t$.
Conversely, if $x\in\mathbb F_p$ satisfies this equation,
choose a root $w\in\mathbb F_{p^2}^*$ of
$Z^2-xZ+n$.  Expanding its third power shows
\[
 w^3+n^3/w^3=t.
\]
Hence $w^3$ is a root of $Z^2-tZ+n$, whose roots
are $u,u^p$.  In either case $u$ is a cube in
$\mathbb F_{p^2}$.  A cubic over a field is reducible
exactly when it has a root there.
\end{proof}

\begin{proposition}\label{prop:quadratic-unit-cm-gate}
Let $d>1$ be squarefree, $3\nmid d$, and put
$C=\mathbb Q(\sqrt d)$ and $F=\mathbb Q(\sqrt{-3d})$.
Write $D_F$ for the fundamental discriminant of $F$ and
assume $3\nmid h(F)$.  Let $\epsilon>1$ be a fundamental
unit of $C$, and set
\[
 n=N_{C/\mathbb Q}(\epsilon)\in\{1,-1\},
 \qquad t=\operatorname{Tr}_{C/\mathbb Q}(\epsilon).
\]
Let $p\nmid6d$ be prime, $p\equiv-1\pmod9$, and
suppose $p$ splits in $F$.  Then $p$ is inert in $C$.
With $u=\epsilon\bmod p\in\mathbb F_{p^2}^*$, the
following conditions are equivalent:
\begin{enumerate}
\item $u$ is a noncube in $\mathbb F_{p^2}$;
\item $X^3-3nX-t$ is irreducible over $\mathbb F_p$;
\item $p$ is inert in the complex cubic field
$M=\mathbb Q(\epsilon^{1/3}+n\epsilon^{-1/3})$;
\item the prime class in the conductor-$3$ ring class
group of $F$ has nontrivial image in its cyclic cubic quotient.
\end{enumerate}
This field $M$ has discriminant $9D_F$.
Given such a $d$, the test in (1) and the output polynomial
in (2) are computable deterministically in time polynomial
in the numerical $d$ and $\log p$.
\end{proposition}

\begin{proof}
Theorems~1.1 and~1.3 of \cite{EvansEtAl2025Hasse}
give the discriminant assertion.  To specify the dependency,
put $v=\epsilon^{1/3}$ and use their ring class field for
$\mathbb Z[\sqrt{-3d}]$.  Inclusion of $v$ there would
give a cubic field of discriminant $D_F$, forcing
$3\mid h(F)$ by Hasse's correspondence.  Thus that
inclusion is excluded, and their Theorem~1.3 gives
$\operatorname{Disc}(M)=9D_F$.

Since $3$ ramifies in $F$,
$h(9D_F)=3h(D_F)$.  As $3\nmid h(D_F)$, the
conductor-$3$ ring class group has a unique quotient of
order $3$.  The cubic-field correspondence identifies
its cubic subfield with $M$.  This proves (3)$\Leftrightarrow$(4).

We have $(-3/p)=-1$ and $(D_F/p)=1$, so $(d/p)=-1$.
This proves inertness in $C$ and gives $N(u)=n$.
The preceding lemma proves (1)$\Leftrightarrow$(2).
For the remaining Frobenius equivalence, let $N$ be the
normal closure of $M$ and put $B=FC=C(\zeta_3)$.
Then $N$ has quadratic resolvent $F$ and
\[
 NB=C(\zeta_3,\epsilon^{1/3}).
\]
The residue field of $B$ at the chosen prime above $p$
is $\mathbb F_{p^2}$.  Its cubic Kummer Frobenius is
nontrivial exactly when $u$ is a noncube.  This Frobenius
restricts to the square of the Frobenius in
$\operatorname{Gal}(N/F)$; squaring preserves nontriviality
in a group of order $3$.  Since $p$ splits in $F$,
nontriviality is precisely the inert splitting type in $M$.
This proves (1)$\Leftrightarrow$(3), without making an
unjustified assumption that $p$ is prime to the index of
the displayed generator of $M$.

Continued fractions compute the fundamental unit in time
polynomial in numerical $d$, with coefficient bit lengths
polynomial in $d$.  Arithmetic in
$\mathbb F_p[Y]/(Y^2-d)$ then tests
$u^{(p^2-1)/3}\ne1$ by exponentiation.  The integers
$n$ and $t$ give the stated output polynomial.
\end{proof}

This gives a sufficient all-input condition: for each hard
prime $p$, find a squarefree $d\le\operatorname{poly}(\log p)$
with $3\nmid d$, $3\nmid h(\mathbb Q(\sqrt{-3d}))$,
$p$ split in that imaginary quadratic field, and the
fundamental unit noncubic modulo its inert prime in
$\mathbb Q(\sqrt d)$.  Quadratic splitting, class-number
indivisibility, and this particular Frobenius condition
are distinct requirements.

Beckwith, Raum, and Richter prove that for every prime
$\ell\ge2$ and finite set $S_+$ of odd primes, some
imaginary quadratic field has class number prime to
$\ell$ and splits at every prime in $S_+$
\cite[Theorem~A]{BeckwithRaumRichter2024}.
Their result includes prescribed split primes
$p\equiv-1\pmod\ell$.  The stated theorem and its
proof do not bound the least discriminant in terms of
$S_+$, do not prescribe ramification at $3$, and do not
test the nontrivial cubic Frobenius required above.
Accordingly that qualitative existence theorem does not
supply the bounded-$d$ selector used here.
Likewise the reciprocity formulas of
\cite[Theorem~1.6]{EvansEtAl2025Hasse} and
\cite[Theorem~1.1]{EvansVanVeen2026Cubic} concern their
specified $p\equiv1\pmod3$ setting and are not used as
an all-input selection bound in the present inert setting.

\subsection{Small cubic nonresidues recover the auxiliary prime}
\label{subsec:cubic-small-candidates}

The following reduction identifies a precise limitation of a search
through small coefficients in the first unresolved prime-degree case.
It is a consequence of classical cubic reciprocity.

\begin{proposition}[Cubic reciprocity and coefficient height]
\label{prop:cubic-height-prime}
Let $p\equiv-1\pmod 9$ be prime, put
$\mathcal O=\mathbb Z[\zeta]$ with $\zeta^2+\zeta+1=0$, and identify
$\mathcal O/(p)$ with $\mathbb F_{p^2}$.
Suppose that the image of $\alpha=A+B\zeta$, where $A,B\in\mathbb Z$,
is a nonzero noncube in $\mathcal O/(p)$.
Then there is a rational prime $\ell$ such that
\[
 \ell\equiv1\pmod3,\qquad
 p^{(\ell-1)/3}\not\equiv1\pmod\ell,
\]
\[
 \ell\le A^2-AB+B^2\le3H^2,
 \quad H=\max\{|A|,|B|\}.
\]
Conversely, every rational prime $\ell\equiv1\pmod3$ satisfying
$p^{(\ell-1)/3}\not\equiv1\pmod\ell$ admits integers $A,B$ such that
\[
 A^2-AB+B^2=\ell,\qquad
 |A|,|B|\le2\sqrt{\ell/3},
\]
and every element of norm $\ell$ has noncube image in
$\mathcal O/(p)$.

Consequently a bound polynomial in $\log p$ for the height of a
cubic nonresidue in this quadratic presentation is equivalent, up to
a change of polynomial, to such a bound for an auxiliary prime
with the displayed property. The corresponding exhaustive searches
are deterministic and polynomial in those numerical bounds and
$\log p$.
\end{proposition}

\begin{proof}
We fix the normalization explicitly. A prime element
$\pi\in\mathcal O$ coprime to $3$ is called primary here when
$\pi\equiv2\pmod{3\mathcal O}$.
Each prime element not above $3$ has a unique primary associate.
For a prime element $\pi$ not above $3$, and an element $\gamma$
coprime to $\pi$, write
\[
 \left(\frac{\gamma}{\pi}\right)_3
   \equiv \gamma^{(N\pi-1)/3}\pmod\pi,
 \qquad
 \left(\frac{\gamma}{\pi}\right)_3\in\{1,\zeta,\zeta^2\}.
\]
This definition uses $N(p)=p^2$ for an inert rational prime $p$.
The cubic reciprocity law in
\citet[equation~(1.1), p.~395; the split--inert proof is in
Section~2, pp.~396--397]{FriesenSpearmanWilliams1986}
says that
\[
 \left(\frac{\pi}{\varpi}\right)_3
 =
 \left(\frac{\varpi}{\pi}\right)_3
\]
for primary primes of distinct norm.

The ring $\mathcal O$ is Euclidean for
$N(A+B\zeta)=A^2-AB+B^2$, hence is a unique factorization domain.
Its prime elements, up to associates, are $1-\zeta$, the rational
primes $q\equiv2\pmod3$, and the elements of rational prime norm
$\ell\equiv1\pmod3$.
We first eliminate all factors except the last type.
Every nonzero rational integer modulo $p$ is a cube already in
$\mathbb F_p$, because $\gcd(3,p-1)=1$.
Also $\zeta$ is a cube in $\mathbb F_{p^2}$:
its cubic character is
$\zeta^{(p^2-1)/3}=1$ since $9\mid p^2-1$.
Thus all six units $\{\pm1,\pm\zeta,\pm\zeta^2\}$ have cube image.
Finally
\[
 (1-\zeta)^2=-3\zeta
\]
is a cube. The quotient
$\mathbb F_{p^2}^{\times}/(\mathbb F_{p^2}^{\times})^3$
has order $3$, so squaring is injective on it;
therefore $1-\zeta$ is itself a cube.
These statements cover units, the ramified prime above $3$,
and every inert rational prime factor other than $p$.
The factor $p$ cannot divide $\alpha$, since its residue is nonzero.

Factor $\alpha$ in $\mathcal O$. Multiplicativity of the cubic
character, together with its nontrivial value on $\alpha$, implies
that some prime factor $\pi$ of norm
$\ell\equiv1\pmod3$ has
$\left(\frac{\pi}{p}\right)_3\ne1$.
Multiplying $\pi$ by a unit makes it primary without changing this
value or its norm. The rational prime $p$ is itself primary, since
$p\equiv2\pmod3$; it is coprime to $\pi$ and has norm
$p^2\ne\ell$. Reciprocity therefore gives
\[
 \left(\frac{p}{\pi}\right)_3
 =
 \left(\frac{\pi}{p}\right)_3\ne1.
\]
The residue field $\mathcal O/(\pi)$ has $\ell$ elements, so this is
exactly $p^{(\ell-1)/3}\not\equiv1\pmod\ell$.
In particular $\ell\notin\{3,p\}$.
Since $\pi\mid\alpha$, its norm divides the positive integer
$N\alpha$, proving $\ell\le N\alpha\le3H^2$.

Conversely, a rational prime $\ell\equiv1\pmod3$ splits in
$\mathcal O$, and unique factorization gives an element $\pi$
of norm $\ell$. Normalize $\pi$ to be primary.
The hypothesis on $p$ means
$\left(\frac{p}{\pi}\right)_3\ne1$ in the residue field
$\mathcal O/(\pi)\cong\mathbb F_\ell$.
Reciprocity gives
$\left(\frac{\pi}{p}\right)_3\ne1$.
The same argument applies to the conjugate prime, and multiplying
by a unit preserves the cube class modulo $p$. Every element of
norm $\ell$ is an associate of one of these two primes, so every
such element is a noncube modulo $p$.
Writing it as $A+B\zeta$, the identities
\[
 N\pi=(A-B/2)^2+3B^2/4
      =(B-A/2)^2+3A^2/4
\]
prove the claimed coordinate bounds.

For the algorithmic assertion, from a height bound $H$ one may
enumerate rational primes at most $3H^2$ and test the displayed
congruence by modular powering; the first part proves success.
From an auxiliary prime $\ell$, enumerate integer pairs in the
displayed box until $A^2-AB+B^2=\ell$.
There are $O(\ell)$ pairs, and the second part proves both existence
and the required noncube property. These running times are
polynomial in the numerical search bounds and $\log p$; no claim
of a running time polynomial in $\log H$ or $\log\ell$ is needed.
\end{proof}

\begin{proposition}[Exact density on an affine line]
\label{prop:quadratic-affine-character}
Let $p$ be prime, $r>1$ an odd divisor of $p+1$, and
$\chi:\mathbb F_{p^2}^{\times}\to\mu_r$ a character of exact
order $r$. For $\alpha\in\mathbb F_{p^2}\setminus\mathbb F_p$,
\[
 \sum_{t\in\mathbb F_p}\chi(\alpha+t)=-1.
\]
More precisely, the value $1$ occurs $(p+1)/r-1$ times, and each
other value in $\mu_r$ occurs $(p+1)/r$ times.
For $r=3$, exactly $2(p+1)/3$ members of this affine line are
noncubes.
\end{proposition}

\begin{proof}
The restriction of $\chi$ to $\mathbb F_p^\times$ is trivial:
its order divides both the odd integer $r\mid p+1$ and $p-1$,
whose greatest common divisor divides $2$.
Thus $\chi$ induces a surjective character on the cyclic quotient
$\mathbb F_{p^2}^\times/\mathbb F_p^\times$, of order $p+1$.
The $p$ elements $\alpha+t$ and the element $1$ are one representative
of each quotient class. Indeed, an equality
$\alpha+t=c(\alpha+u)$ with $c\in\mathbb F_p^\times$
forces $c=1$ and $t=u$, and no $\alpha+t$ is in $\mathbb F_p$.
Every character value appears $(p+1)/r$ times on the full quotient.
Removing the representative $1$ proves both assertions.
\end{proof}

This density statement concerns the entire affine line.
It gives no bound polynomial in $\log p$ on the first successful
integer parameter $t$. In the cubic setting, a successful parameter
$t$ in $\zeta+t$ already yields, by
Proposition~\ref{prop:cubic-height-prime}, a good auxiliary prime
$\ell\le t^2-t+1$.

\begin{proposition}[An explicit cubic candidate family]
\label{prop:cubic-torus-family}
Let $p\equiv2\pmod3$ be an odd prime, put
$\mathbb F_{p^2}=\mathbb F_p(\zeta)$, and let $u\in\mathbb F_p$.
Set
\[
 v=u^2-u+1,\qquad
 t_u=2-\frac3v,\qquad
 f_u(Y)=Y^3-3Y-t_u .
\]
Then $v\ne0$, and $f_u$ is irreducible over $\mathbb F_p$
if and only if $u+\zeta$ is a noncube in $\mathbb F_{p^2}$.
\end{proposition}

\begin{proof}
The quadratic $U^2-U+1$ has no root in $\mathbb F_p$ when
$p\equiv2\pmod3$ and $p\ne2$.
Put $\alpha=u+\zeta$ and
$b=\alpha/\alpha^p$. Its norm is $1$, and a direct calculation gives
\[
 b+b^{-1}=
 \frac{2u^2-2u-1}{u^2-u+1}=t_u .
\]
For the cubic character $\chi$ of $\mathbb F_{p^2}$,
$\chi(\alpha^p)=\chi(\alpha)^p=\chi(\alpha)^{-1}$.
Hence $\chi(b)=\chi(\alpha)^2$, so $b$ is a cube exactly when
$\alpha$ is a cube.

If $b=x^3$ for $x\in\mathbb F_{p^2}$, then
$N(x)^3=N(b)=1$. Since the cube map is bijective on
$\mathbb F_p^\times$, $N(x)=1$.
Thus $y=x+x^{-1}$ is fixed by $p$th-power Frobenius and belongs
to $\mathbb F_p$. The identity
$y^3-3y=x^3+x^{-3}=t_u$ gives a root of $f_u$.
Conversely, if $y\in\mathbb F_p$ is a root of $f_u$, a root $x$
of $X^2-yX+1$ belongs to $\mathbb F_{p^2}$, is nonzero, and
satisfies $x^3+x^{-3}=t_u=b+b^{-1}$.
Therefore $x^3$ is one of $b,b^{-1}$, so $b$ is a cube.
A cubic over a field is irreducible exactly when it has no root
in that field, completing the proof.
\end{proof}



\subsection{A coefficient family and the cost of Dirichlet-character transfer}
\label{subsec:cubic-fourier-transfer}

There is a particularly simple family with a positive density of
successful coefficients.  Its multiplicative Fourier expansion also
quantifies the loss in applying rational Dirichlet estimates separately
to all the resulting characters.

\begin{proposition}[Density and multiplicative Fourier norm]
\label{prop:cubic-fourier-norm}
Let $p\equiv2\pmod3$ be an odd prime, and define
\[
 F(m)=\boldsymbol 1_{\{X^3+X+m\text{ is irreducible over }\mathbb F_p\}}
 \qquad(m\in\mathbb F_p).
\]
For a multiplicative character $\chi$ of $\mathbb F_p^\times$, put
\[
 \widehat F(\chi)=\frac1{p-1}
    \sum_{m\in\mathbb F_p^\times}F(m)\overline{\chi(m)},
 \qquad
 \mu=\frac{p+1}{3(p-1)}.
\]
Then
\[
 \sum_{m\in\mathbb F_p}F(m)=\frac{p+1}{3},
 \qquad \widehat F(1)=\mu.
\]
Every nontrivial character satisfies
\[
 |\widehat F(\chi)|\le\frac{4\sqrt p}{3(p-1)},
\]
and $\widehat F(\chi)=0$ when $\chi(-1)=-1$.
Moreover,
\[
 \frac{(p+1)(p-2)}{6(p-1)\sqrt p}
 \le \sum_{\chi\ne1}|\widehat F(\chi)|
 \le \sqrt{(p-2)\mu(1-\mu)}.
\]
In particular,
\[
 \frac{p-1}{6\sqrt p}
 \le \sum_{\chi\ne1}|\widehat F(\chi)|
 \le \frac{\sqrt{p-2}}2.
\]
Thus the nonprincipal Fourier $\ell^1$ norm has order $\sqrt p$.
\end{proposition}

\begin{proof}
Let $\lambda$ be the quadratic character, extended by zero at zero,
and let $N_m$ be the number of roots of $X^3+X+m$ in $\mathbb F_p$.
Its discriminant is
\[
 D_m=-4-27m^2.
\]
The element $-3$ is a nonsquare: a square root would give a
nontrivial cube root of unity in $\mathbb F_p$, whereas
$3\nmid p-1$.  Consequently $D_m$ never vanishes.  For a separable
cubic, the sign of its Frobenius permutation is the quadratic
character of its discriminant.  Indeed Frobenius acts on the
Vandermonde product of its three roots by this sign, and the square
of that product is the discriminant.  The three possible root
counts $0,1,3$ therefore give the identity
\begin{equation}
 F(m)=\frac{2+\lambda(D_m)-N_m}{3}.
 \label{eq:cubic-indicator}
\end{equation}

Each $x\in\mathbb F_p$ is a root for exactly one parameter,
$m=-x^3-x$, so $\sum_mN_m=p$.
For nonzero $c\in\mathbb F_p$ one has
\[
 \sum_{x\in\mathbb F_p}\lambda(x^2-c)=-1:
\]
the equation $y^2=x^2-c$ has $p-1$ solutions, by the bijection
$(x,y)\mapsto(x-y,x+y)$ with pairs of product $c$.
It follows that
\[
 \sum_m\lambda(-4-27m^2)=-\lambda(-27)=1.
\]
Summing \eqref{eq:cubic-indicator} proves the asserted count.
Also $F(0)=0$, since $X^3+X$ has the root zero, so this count gives
the stated mean on $\mathbb F_p^\times$.

We include the character-sum facts used for the coefficient bound.
Extend all multiplicative characters by zero at zero, including
the trivial character.  For characters $A,B$, set
\[
 J(A,B)=\sum_{z\in\mathbb F_p}A(z)B(1-z).
\]
If $A,B,AB$ are nontrivial, then $|J(A,B)|=\sqrt p$.
For completeness, let $e_p(z)=\exp(2\pi i z/p)$ and
$G(A)=\sum_zA(z)e_p(z)$.  Multiplicative orthogonality gives
\[
 |G(A)|^2
 =\sum_{u\ne0}A(u)\sum_{y\ne0}e_p((u-1)y)=p
\]
for nontrivial $A$.  Sorting the terms of $G(A)G(B)$ by
$x+y$ gives
\[
 G(A)G(B)=J(A,B)G(AB).
\]
The terms with $x+y=0$ sum to zero because $AB$ is nontrivial.
The claimed magnitude follows.

The function $F$ is even, since replacing $X$ by $-X$ identifies
the roots at $m$ and $-m$.  Hence all its odd-character coefficients
vanish.  Fix a nontrivial even character $\chi$, and write
$\psi=\overline\chi$.  The character group is cyclic, and
$\psi(-1)=1$, so there is a character $\eta$ with
$\eta^2=\psi$.  Neither $\eta$ nor $\eta\lambda$ is trivial.

Put $A_0=-4$ and $B_0=-27$.  Changing variables $y=m^2$ gives
\begin{align*}
 S_D&:=\sum_{m\ne0}\psi(m)\lambda(A_0+B_0m^2)\\
 &=\sum_y\eta(y)\lambda(A_0+B_0y)
   +\sum_y(\eta\lambda)(y)\lambda(A_0+B_0y).
\end{align*}
The two sums are, up to factors of modulus one,
$J(\eta,\lambda)$ and $J(\eta\lambda,\lambda)$.
Their two characters and their products are nontrivial, so
\[
 |S_D|\le2\sqrt p.
\]
Similarly,
\begin{align*}
 S_N&:=\sum_{m\ne0}\psi(m)N_m
     =\sum_x\psi\bigl(-x(x^2+1)\bigr)\\
 &=\sum_y\eta(y)\psi(1+y)
   +\sum_y(\eta\lambda)(y)\psi(1+y).
\end{align*}
Here the two Jacobi sums, again up to factors of modulus one,
are $J(\eta,\psi)$ and $J(\eta\lambda,\psi)$.
Their products are $\eta^3$ and $\eta^3\lambda$.
Since $3\nmid p-1$, cubing is an automorphism of the character
group.  Thus $\eta^3=1$ would force $\eta=1$, and
$\eta^3=\lambda$ would force $\eta=\lambda$; both are excluded.
All the required characters are therefore nontrivial, and
\[
 |S_N|\le2\sqrt p.
\]
The constant term in \eqref{eq:cubic-indicator} has zero sum
against $\psi$.  Consequently
\[
 |\widehat F(\chi)|
   =\frac{|S_D-S_N|}{3(p-1)}
   \le\frac{4\sqrt p}{3(p-1)}.
\]

Parseval's identity and $F^2=F$ give
\[
 \sum_{\chi\ne1}|\widehat F(\chi)|^2
   =\mu-\mu^2
   =\frac{2(p+1)(p-2)}{9(p-1)^2}.
\]
Dividing by the upper bound for the largest summand's modulus
proves the lower bound for the $\ell^1$ norm.  The further
simplification uses
$(p+1)(p-2)\ge(p-1)^2$, valid for $p\ge3$.
Cauchy--Schwarz with $p-2$ nonprincipal characters proves the
upper bound; finally $\mu(1-\mu)\le1/4$.
\end{proof}

To quantify a direct transfer from rational Dirichlet estimates, use the
fixed nonnegative smooth function $w$ supported in $(1,2)$ from
Lemma~\ref{lem:smooth}, with $c_w=\int_1^2w(t)\,dt>0$.
For a Dirichlet character modulo $p$, put
\[
 \Psi_{\chi,w}(X)=\sum_{n\ge1}\Lambda(n)\chi(n)w(n/X).
\]
Under $A$, uniformly in these characters and $X\ge2$,
\begin{equation}\label{eq:complex-dirichlet-smooth}
 \Psi_{\chi,w}(X)=\boldsymbol1_{\chi=1}c_wX
       +O_w\!\left(\sqrt X\log(2p)\right).
\end{equation}
Here the error is an absolute bound on a complex number, as needed
for the possibly complex Fourier coefficients. To verify it, apply
Mellin inversion to $-L'/L$ for the primitive character inducing
$\chi$ and shift the contour past its pole and zeros. The zero term is
$-\sum_\rho X^\rho\widehat w(\rho)$. Three integrations by parts
give $|\widehat w(\sigma+it)|\ll_w(1+|t|)^{-3}$ on
$0\le\sigma\le1$. The unit-height zero count is
$O(\log(2p)+\log(|t|+3))$, so the absolute sum of these transforms
is $O_w(\log(2p))$. Under $A$, $|X^\rho|=\sqrt X$.
The pole at one gives the stated main term; the gamma factor and
trivial zeros contribute $O_w(\log(2p))$, exactly as in the proof
of Lemma~\ref{lem:smooth}. For the principal character modulo $p$,
the omitted Euler factor contributes only prime powers of $p$ in
$(X,2X)$; there is at most one, and its contribution is
$O_w(\log p)$. This proves \eqref{eq:complex-dirichlet-smooth}.

Fourier inversion and a termwise triangle inequality now give
\begin{equation}\label{eq:cubic-dirichlet-transfer}
 \sum_{n\ge1}\Lambda(n)F(n\bmod p)w(n/X)
   =\mu c_wX+O_w\!\left(\sqrt{pX}\log(2p)\right).
\end{equation}
Proper prime powers contribute $O_w(\sqrt X)$ by Chebyshev's
bound. Thus this estimate guarantees a prime parameter in $(X,2X)$
when $X=C_w p\log^2(2p)$ and $C_w$ is sufficiently large.
This still depends on the numerical characteristic $p$.
Proposition~\ref{prop:cubic-fourier-norm} shows that the factor
of order $\sqrt p$ is present in the $\ell^1$ norm used by this
particular termwise transfer. It does not give a lower bound for
the error in \eqref{eq:cubic-dirichlet-transfer}, and it does not
exclude cancellation between characters or a different family.


\subsection{Logarithmic normalized Tate parameters for almost all inputs}

The quadratic norms in the normalized Tate family provide linearly many
independent rational cubic characters.  This improves the parameter
range in the average construction to a constant times $\log P$.
The product estimate is the classical quadratic-value mechanism;
compare \citet[Lemmas~1, 3, and 4 and \S2.4, pp.~1131--1134]{Cilleruelo2011}.
We give its specialization in full to track the private conductor
primes and the leading constant.  

\begin{lemma}[Large primes in the quadratic norms]
\label{lem:tate-large-prime-norms}
For an integer $K\ge4$, put
\[
 N(k)=k^2+3k+9,\qquad
 \mathcal S_K=\{q>K+3:q\text{ prime},\ q\mid N(k)
                      \text{ for some }1\le k\le K\}.
\]
Then
\begin{equation}\label{eq:tate-large-prime-count}
 |\mathcal S_K|\ge\left(\frac12-o(1)\right)K
 \qquad(K\longrightarrow\infty).
\end{equation}
Each $N(k)$ has at most one prime divisor exceeding $K+3$, and
that divisor has valuation one.  A prime $q\in\mathcal S_K$
divides at most two of the values; if it divides two, then
$q\le2K+3$.
\end{lemma}

\begin{proof}
Stirling's formula gives
\begin{equation}\label{eq:tate-norm-product-size}
 \sum_{k=1}^K\log N(k)
 =2\log K!+\sum_{k=1}^K\log(1+3/k+9/k^2)
 =2K\log K-2K+O(\log K).
\end{equation}
For a prime $q\ne3$, the congruence $N(k)\equiv0\pmod q$
has two simple roots if $q\equiv1\pmod3$, and no roots
otherwise.  Indeed, division by nine transforms it into
$u^2+u+1=0$; its roots are the nonidentity cube roots of unity.
For $q=2$ there are no roots.  Hensel lifting therefore gives
exactly two roots modulo every $q^a$ in the split case.

Set $L_q=\lfloor\log N(K)/\log q\rfloor$.  For
$q\equiv1\pmod3$,
\begin{equation}\label{eq:tate-small-prime-valuation}
 \sum_{k=1}^K v_q(N(k))
 \le2\sum_{a=1}^{L_q}\left(\frac K{q^a}+1\right)
 \le\frac{2K}{q-1}+2L_q.
\end{equation}
The accumulated floor error after multiplication by $\log q$
and summation over $q\le K+3$ is at most
\[
 2\pi(K+3)\log N(K)=O(K).
\]
The prime number theorem in the fixed progression $1\pmod3$,
followed by partial summation, yields
\[
 \sum_{\substack{q\le x\\q\equiv1\bmod3}}
                  \frac{\log q}{q-1}
 =\frac12\log x+o(\log x).
\]
The prime three contributes only $O(K)$: if $3\nmid k$ its
valuation is zero, while for $k=3u$ one has
$N(k)=9(u^2+u+1)$ and $v_3(u^2+u+1)\le1$.
Thus the total logarithmic contribution from primes at most
$K+3$ is at most $K\log K+o(K\log K)$.  Subtracting this
from \eqref{eq:tate-norm-product-size} gives
\begin{equation}\label{eq:tate-large-prime-mass}
 \sum_{q>K+3}\log q\sum_{k=1}^K v_q(N(k))
 \ge K\log K+o(K\log K).
\end{equation}

Since
\[
 N(k)\le K^2+3K+9<(K+3)^2,
\]
no $N(k)$ contains two prime factors greater than $K+3$,
and every such factor occurs once.  A prime $q>K+3$ divides
at most two values, because the parameters are distinct modulo
$q$ and the quadratic has two roots.  If it divides $N(a)$
and $N(b)$ with $a\ne b$, then
$a+b\equiv-3\pmod q$.  The bounds $1\le a,b\le K<q-3$
force $a+b=q-3$, hence $q\le2K+3$.

A prime occurring once consequently contributes at most
$\log N(K)<2\log(K+3)$ to
\eqref{eq:tate-large-prime-mass}; one occurring twice contributes
at most $2\log(2K+3)$.  Therefore
\[
 2|\mathcal S_K|\log(2K+3)
 \ge K\log K+o(K\log K),
\]
which proves \eqref{eq:tate-large-prime-count} and the remaining
assertions.
\end{proof}

\begin{lemma}[Cyclic cubic characters and ramification]
\label{lem:tate-cyclic-character}
For every positive integer $k$, there is a primitive cubic
Dirichlet character $\theta_k$ of conductor $f_k\mid N(k)$
such that, for every prime $p\nmid N(k)$,
\begin{equation}\label{eq:tate-global-character-test}
 h_k(Y):=Y^3-(k+3)Y^2+kY+1
 \text{ is irreducible over }\mathbb F_p
 \quad\Longleftrightarrow\quad\theta_k(p)\ne1.
\end{equation}
If $q\ne3$ and $v_q(N(k))=1$, then $q\mid f_k$.
\end{lemma}

\begin{proof}
The rational-root test proves that $h_k$ is irreducible over
$\mathbb Q$, since
\[
 h_k(1)=-1,\qquad h_k(-1)=-2k-3.
\]
Its discriminant is $N(k)^2$, so a root generates a cyclic
cubic splitting field $L_k/\mathbb Q$.
The Kronecker--Weber theorem and the conductor--discriminant
formula give a primitive cubic Dirichlet character of conductor
$f_k$, with
\[
 \operatorname{disc}(L_k)=f_k^2;
\]
see \citet[p.~472]{SilvesterSpearmanWilliams2006} for these
cyclic-cubic facts.  If $\alpha_k$ is a root and
$I_k=[\mathcal O_{L_k}:\mathbb Z[\alpha_k]]$, the polynomial
discriminant identity gives
\[
 N(k)^2=f_k^2I_k^2,
\]
and hence $f_k\mid N(k)$.  No equality between $f_k$ and
$N(k)$ is assumed.  For $p\nmid N(k)$, Frobenius has the
cycle type of the factorization of $h_k$ modulo $p$.
The only possibilities in the cyclic group of order three are
the identity and a three-cycle.  A nontrivial character is
faithful on that group, proving
\eqref{eq:tate-global-character-test}.

Finally, direct expansion gives
\begin{equation}\label{eq:tate-eisenstein-translation}
 h_k\!\left(Z+\frac{k+3}{3}\right)
 =Z^3-\frac{N(k)}3Z-\frac{(2k+3)N(k)}{27}.
\end{equation}
If $q\ne3$ and $v_q(N(k))=1$, the identity
$(2k+3)^2=4N(k)-27$ implies $q\nmid2k+3$.
All denominators in \eqref{eq:tate-eisenstein-translation}
are $q$-adic units, and the translated polynomial is Eisenstein
over $\mathbb Q_q$.  Thus $q$ is totally ramified in $L_k$,
so $q\mid\operatorname{disc}(L_k)$ and $q\mid f_k$.
\end{proof}

\begin{corollary}[A linear-size independent parameter set]
\label{cor:tate-linear-independent-set}
For every sufficiently large $K$, there is a set
$\mathcal K_K\subset\{1,\ldots,K\}\setminus\{3\}$ of
size at least $(1/2-o(1))K$ whose associated characters
$\theta_k$ are linearly independent over $\mathbb F_3$.
Each such character has a conductor prime greater than $K+3$
which divides none of the other selected conductors.
For any $m$ selected parameters, their conductors satisfy
\[
 \log\prod_{i=1}^m f_{k_i}
 \le m\log N(K)<2m\log(K+3).
\]
\end{corollary}

\begin{proof}
For each $q\in\mathcal S_K$, choose one $k(q)\le K$ with
$q\mid N(k(q))$.  These parameters are distinct, because a
value has at most one prime divisor greater than $K+3$.
The valuation at $q$ is one, so
Lemma~\ref{lem:tate-cyclic-character} gives $q\mid f_{k(q)}$.
For a different selected parameter $k(q')$, the value
$N(k(q'))$ contains $q'>K+3$ and therefore cannot also
contain $q$.  Consequently $q\nmid f_{k(q')}$.
Restriction to the $q$-component of the common modulus forces
the exponent of $\theta_{k(q)}$ in any principal character
product to be zero modulo three.  Doing this for every selected
$q$ proves independence.  The parameter three cannot occur,
since $N(3)=27$ has no prime factor exceeding $K+3$.
The size assertion follows from
Lemma~\ref{lem:tate-large-prime-norms}, and the conductor
bound follows from $f_k\mid N(k)\le N(K)$.
\end{proof}

\begin{proof}[Proof of Theorem~\ref{thm:almost-all-small-tate}]
Put $K=\lfloor C\log P\rfloor$ with $C>1/\log3$ and let
\[
 m=\left\lceil\frac{\log P}{2\log3}\right\rceil.
\]
By Corollary~\ref{cor:tate-linear-independent-set}, there are at
least $(C/2-o(1))\log P$ independent characters from parameters
at most $K$.  The strict condition on $C$ therefore allows us
to select exactly $m$ of them for all sufficiently large $P$.
Write them as $\theta_{k_1},\ldots,\theta_{k_m}$ and put
\[
 Q=\prod_{i=1}^m f_{k_i}.
\]
Then
\begin{equation}\label{eq:tate-logarithmic-conductor-budget}
 \log Q\le m\log N(K)
 =O_C(\log P\log\log P),\qquad 3^m\ge\sqrt P.
\end{equation}
Also $N(K)<P/2$ eventually, so primes in the input interval
divide none of the polynomial discriminants or conductors.

Regard the characters as characters modulo $Q$, extended by
zero on nonunits, including each exponent-zero character.
The indicator $G(n)$ of common triviality is
\[
 G(n)=\frac1{3^m}
       \sum_{\boldsymbol e\in\{0,1,2\}^m}
         \prod_{i=1}^m\theta_{k_i}(n)^{e_i}.
\]
Here $G(n)=0$ when $(n,Q)>1$.  Independence implies that
the displayed sum contains $3^m$ distinct characters, exactly
one principal.  Its coefficients are real and have absolute
values summing to one.  Lemma~\ref{lem:dirichlet-log-weight},
applied termwise and followed by taking real parts, consequently
gives
\[
 \sum_{n\le x}\Lambda(n)G(n)\log(x/n)
 =\frac{x}{3^m}+O\!\left(\sqrt x\log(2Q)\right).
\]
If every admissible polynomial fails at a prime
$p\in(P/2,P]$, then all the selected polynomials fail and
\eqref{eq:tate-global-character-test} gives $G(p)=1$.
At $x=2P$ each such prime contributes at least
$\log(P/2)\log2$ to this nonnegative sum.  Thus the number
of exceptional primes is
\[
 O\!\left(
 \frac{P/3^m+\sqrt P\log(2Q)}{\log P}\right)
 =O_C(\sqrt P\log\log P)
\]
by \eqref{eq:tate-logarithmic-conductor-budget}.  This counts
all primes in the interval, so it also bounds the exceptions
among $p\equiv-1\pmod9$.

Every selected parameter is different from three.  Since
$K<p$ eventually, it also satisfies $k\not\equiv3\pmod p$.
The deterministic algorithm simply scans all $1\le k\le K$,
skipping that residue, and tests
\[
 \gcd\bigl(h_k(Y),Y^p-Y\bigr)=1.
\]
For a cubic this is equivalent to irreducibility, since a
reducible cubic has a linear factor.  Repeated squaring computes
$Y^p\bmod h_k$ in polynomial bit time in $\log p$.
There are $O_C(\log p)$ tests.  A successful test gives the
presentation $\mathbb F_p[Y]/(h_k)$ of $\mathbb F_{p^3}$;
its correctness is independent of A.  The independent character
set and its conductors are used only in the proof and need not
be computed by the algorithm.
\end{proof}

The conclusion is an average over the varying input prime.
It retains the displayed exceptional set and supplies no
pointwise bound for a prescribed exceptional prime.


\paragraph{A smaller explicit prime-interval menu.}
The lower comparison below uses a different family: one parameter for
each split prime in a prescribed interval.  Its size is asymptotic to
$K/(2\log K)$, while Corollary~\ref{cor:tate-linear-independent-set}
uses large prime divisors of the full product and gives linear size.
The following statement fixes the smaller menu exactly.

\begin{lemma}\label{lem:tate-private-prime-menu}
Let $K\ge4$ be an integer and put $N(k)=k^2+3k+9$.
For each prime
\[
 \ell\in(K+3,2K+3],\qquad \ell\equiv1\pmod3,
\]
let $k_\ell$ be the smaller least nonnegative representative of a
root of $N(k)\equiv0\pmod\ell$.  Then $1\le k_\ell\le K$,
$k_\ell\ne3$, and $v_\ell(N(k_\ell))=1$.
No other prime in the displayed interval divides $N(k_\ell)$.

For each such parameter there is a cubic Dirichlet character
$\theta_\ell$ of conductor dividing $N(k_\ell)$ such that, for
every prime $p\equiv-1\pmod9$ not dividing $N(k_\ell)$,
\[
 h_{k_\ell}(Y)=Y^3-(k_\ell+3)Y^2+k_\ell Y+1
 \text{ is irreducible over }\mathbb F_p
 \quad\Longleftrightarrow\quad \theta_\ell(p)\ne1.
\]
These characters are linearly independent over $\mathbb F_3$:
every product $\prod_\ell\theta_\ell^{e_\ell}$ with some
$e_\ell\not\equiv0\pmod3$ is nonprincipal.  More strongly,
it has a nontrivial component at a prime $\ell\ne3$.
\end{lemma}

\begin{proof}
The discriminant of $N$ is $-27$, so it has two distinct roots
modulo every prime $\ell\equiv1\pmod3$.  Here $\ell>7$.
Neither $\ell-1$ nor $\ell-2$ is a root, since $N(-1)=N(-2)=7$.
Thus the sum of the two least representatives is $\ell-3$,
and the smaller is at most $(\ell-3)/2\le K$.
It cannot be zero or three because $N(0)=9$ and $N(3)=27$.
Since
\[
 N(k_\ell)\le K^2+3K+9<(K+3)^2<\ell^2,
\]
its $\ell$-valuation is one.  The product of any two primes
larger than $K+3$ exceeds the same upper bound, proving the
other divisibility assertion.  The exclusion of $\ell=7$ is
necessary: its two least roots are five and six.

We spell out the reciprocity character to control its conductor.
In the Eisenstein integers factor
$\alpha=k_\ell-3\zeta_3$ into units, the ramified prime,
inert rational primes, and primary split prime factors.  For a
hard prime $p$, the units and ramified prime are cubes modulo $p$:
$\zeta_3$ is a cube because $9\mid p+1$, and
$(1-\zeta_3)^2=-3\zeta_3$ is a cube.  Inert rational factors
are scalars, hence also cubes.  Cubic reciprocity for the remaining
primary factors, in the split/inert case of
\cite[p.~395, equation~(1.1), and Section~2, pp.~396--397]{FriesenSpearmanWilliams1986},
therefore expresses the cubic residue symbol of
$\alpha$ modulo $p$ as a product of rational cubic Dirichlet
characters, each of conductor a split prime dividing $N(k_\ell)$.
At conjugate prime factors their characters are inverses.  This
defines $\theta_\ell$ with conductor dividing $N(k_\ell)$.
The established Tate--Kummer identity proves the displayed
irreducibility criterion.

The designated prime $\ell$ divides the Eisenstein norm exactly
once, so precisely one of its two prime factors divides $\alpha$,
with exponent one.  Consequently $\theta_\ell$ has a nontrivial
local component at $\ell$.  No other selected character has a
component there, by the norm-divisibility assertion above.  Hence
a product with $e_\ell\ne0$ retains that nontrivial component.
This proves independence.  A character may have additional
conductor primes; only the designated prime is private to it.
\end{proof}

\subsection{Classical character-product lower comparisons}

The character product also produces exceptional primes at controlled
scales.  This is the classical mechanism for least-quadratic-nonresidue
lower bounds: \cite[p.~432, Exercise~9(a)--(e)]{MontgomeryVaughan2007}
explicitly attributes that argument to Montgomery (1971, p.~122).
We record the specializations to the present searches.

\begin{lemma}\label{lem:compatible-bad-coset}
Assume A.  Let $Q\ge2$, let $H\le(\mathbb Z/Q\mathbb Z)^*$
have index $I$, and let $J\ge I$.  For an absolute sufficiently
large $C$, the coset $(-1)H$ contains a prime in $(x/4,x]$, where
\[
 x=CJ^2\log^2(2Q).
\]
\end{lemma}

\begin{proof}
Extended by zero on nonunits, the coset indicator is
\[
 1_{(-1)H}(n)=I^{-1}\sum_{\chi\in H^\perp}\chi(-1)\chi(n).
\]
Its coefficients are real, equal to $1/I$ or $-1/I$, and have
absolute-value sum one.  Lemma~\ref{lem:dirichlet-log-weight} gives
\[
 S(y):=\sum_{n\le y}\Lambda(n)1_{(-1)H}(n)\log(y/n)
       =y/I+O(\sqrt y\log(2Q)).
\]
The nonnegative logarithmic tent
$S(x)-2S(x/2)+S(x/4)$, supported on $(x/4,x]$ and bounded
by $\log2$, has value
\[
 x/(4I)+O(\sqrt x\log(2Q)).
\]
Proper prime powers contribute $O(\sqrt x)$ by Chebyshev's bound.
At the stated $x$, the main term is at least
$CJ\log^2(2Q)/4$, while the error is
$O(\sqrt C J\log^2(2Q))$.  An absolute choice of $C$ makes
the prime contribution positive.
\end{proof}

\begin{proposition}\label{prop:auxiliary-lower-comparison}
Assume A and fix an odd prime $r$.  There are arbitrarily large
primes $p\equiv-1\pmod{r^2}$ with no good auxiliary prime at
most $B$, where
\[
 B=\left(\frac{r-1}{2\log r}+o(1)\right)
                   \log p\log\log p.
\]
These are hard inputs: $\operatorname{ord}_r(p)=2$ and
$v_r(p^2-1)\ge2$.  The $o(1)$ is along the constructed sequence
with $r$ fixed.
\end{proposition}

\begin{proof}
Let $B\to\infty$.  Take all primes $\ell\le B$ with
$\ell\equiv1\pmod r$, let $m$ be their number, and put
$M=\prod\ell$, $Q=r^2M$.  The conditions
$n\equiv-1\pmod{r^2}$ and that $n$ be an $r$th power at
every selected $\ell$ describe a compatible coset $(-1)H$:
every odd-order residue character takes value one at $-1$.
The Chinese remainder theorem gives the exact index
\[
 I=\varphi(r^2)r^m.
\]
Lemma~\ref{lem:compatible-bad-coset}, with $J=I$, supplies a
prime $p\in(x/4,x]$ satisfying these conditions, for
$x=CI^2\log^2(2Q)$.  The fixed-progression prime number theorem
and Chebyshev's bound give
\[
 m=(1+o(1))\frac B{(r-1)\log B},\qquad \log Q\ll_r B.
\]
Consequently
\[
 \log p=2m\log r+O_r(\log B)
   =\left(\frac{2\log r}{r-1}+o(1)\right)\frac B{\log B}.
\]
Inverting this relation gives the asserted scale and shows
$p\to\infty$.  The hard-case assertions follow from the chosen
congruence and oddness of $r$.
\end{proof}

\begin{proposition}\label{prop:private-tate-lower-comparison}
Assume A.  For every sufficiently large integer $K$, the private-prime
menu of Lemma~\ref{lem:tate-private-prime-menu} has a hard prime
$p\equiv-1\pmod9$ at which all its cubics split, with
\[
 \log p=(\log3+o(1))K/\log K,
 \qquad
 K=(1/\log3+o(1))\log p\log\log p.
\]
\end{proposition}

\begin{proof}
Let $m$ be the menu size and $Q$ nine times the least common
multiple of its character conductors.  The private conductor primes
give the bad coset exact index $I=6\cdot3^m$.  It is compatible
with $p\equiv-1\pmod9$ because cubic characters take value one
at $-1$.  Apply Lemma~\ref{lem:compatible-bad-coset} with $J=I$.
Since $m=(1/2+o(1))K/\log K$ and $\log Q\ll K$, its prime
satisfies $\log p=2m\log3+O(\log K)$, giving the formulas.
For large $K$, $p>K^2+3K+9$, so the cubics are unramified.
Their residue characters are trivial.  A squarefree cubic with
square discriminant has only the split and irreducible possibilities
here; the nontrivial-character alternative is precisely the latter.
Thus all menu cubics split.
\end{proof}

The private menu is not nested as $K$ varies.  This proposition
concerns the displayed pairs $(p,K)$; it does not by itself refute
a different prescribed rule $K=K(p)$ whose moving prime interval
is not the same one.

\begin{proposition}\label{prop:full-tate-lower-comparison}
Assume A.  There are arbitrarily large hard primes $p$ such that
every $h_k$ with $1\le k\le K$ splits over $\mathbb F_p$, where
\[
 K=\left(\frac1{2\log3}+o(1)\right)\log p.
\]
Thus scanning only $1\le k\le c\log p$ fails for infinitely
many hard inputs whenever $c<1/(2\log3)$.  This does not exclude
larger polynomial windows in $\log p$.
\end{proposition}

\begin{proof}
For each $1\le k\le K$, take the cubic reciprocity character
of $k-3\zeta_3$, allowing the principal character.  Its conductor
divides $N(k)$ by the same factorization argument used above.
Common triviality and $p\equiv-1\pmod9$ form a compatible coset
$(-1)H$.  Its index satisfies $I\le6\cdot3^K$, with no
independence assumption.  For $Q$ nine times the least common
multiple of the conductors,
\[
 \log Q\le\log9+\sum_{k\le K}\log N(k)\ll K\log(2K).
\]
Choose the explicit upper index $J=6\cdot3^K$ in
Lemma~\ref{lem:compatible-bad-coset}, whether or not $I=J$.
The resulting prime satisfies
$\log p=2K\log3+O(\log K)$, and all these cubics split by
the unramified residue-character criterion.  The parameter intervals
are nested, so the final scanning assertion follows.
When interpreted via elliptic curves, the singular parameter $k=3$
can be omitted; it causes no difficulty for the polynomial statement.
\end{proof}

The first proposition compares directly with the almost-all upper
coefficient $(r-1)/(2\log r)$.  The existence of exceptional
primes is consistent with that upper theorem and with a possible
uniform bound having a larger constant or a higher logarithmic
power.  These specializations do not establish independence of the
desired all-input construction from A.


\subsection{Restrictions on all-parameter rational noncube families}

\begin{proposition}\label{prop:rational-noncube-family}
Let $p>3$ be prime with $p\equiv2\pmod3$, and put $q=p^2$.
Let $C=A/B\in\mathbb F_q(T)^*$, where $A,B$ are relatively
prime polynomials of degrees at most $D$.  Suppose
\[
 p>(8D+4)^2
\]
and $C(t)$ is a noncube for every $t\in\mathbb F_p$ at which it
is finite and nonzero.  Then
\[
 C(T)=c\,b(T)\,g(T)^3,
 \qquad b\in\mathbb F_p(T)^*,\quad
 g\in\mathbb F_q(T)^*,
\]
where
\[
 c=\operatorname{lc}(A)/\operatorname{lc}(B)
       \in\mathbb F_q^*\setminus(\mathbb F_q^*)^3.
\]
Conversely every function of this form has noncubic value at every
base-field argument at which it is finite and nonzero.
\end{proposition}

\begin{proof}
Write a bar for coefficientwise $p$-Frobenius, and put
$R=\overline C/C$.  Suppose first that $R$ is not a cube over
$\overline{\mathbb F}_p(T)$.  The smooth projective model of
\[
 W^3=R(T)
\]
is geometrically connected.  The semilinear involution that conjugates
constants, fixes $T$, and sends $W$ to $W^{-1}$ descends this
curve to $\mathbb F_p$.  The reduced numerator and denominator of
$R$ have degrees at most $2D$; moreover their degrees are equal,
so infinity is unramified.  There are at most $4D$ branch points.
The tame degree-three Riemann--Hurwitz formula therefore gives
$g\le4D-2$.

For a valid $t\in\mathbb F_p$, put $z=C(t)$.
The fibre equation is $w^3=z^{p-1}$.  It has a solution in
$\mathbb F_q$ exactly when $z$ is a cube, since $p-1\equiv1\pmod3$.
Every such solution has norm one: its norm has cube one, and cubing
is bijective on $\mathbb F_p^*$.  These are precisely the
$\mathbb F_p$-rational points of the descended fibre.  By hypothesis
there are none.  All rational points must therefore lie above the
at most $2D$ finite zeros and poles of $C$, or above infinity.
There are at most three points above each, giving
\[
 \#X(\mathbb F_p)\le3(2D+1).
\]
On the other hand, Hasse--Weil gives
\[
 \#X(\mathbb F_p)\ge p+1-2g\sqrt p
                    \ge p+1-(8D-4)\sqrt p.
\]
For $\sqrt p>8D+4$ this exceeds $3(2D+1)$, a contradiction.
The Kummer and Hasse--Weil inputs are standard; see, for example,
\cite[Section~1.1, Theorems~1.1--1.2]{BartoliEtAl2024}.

It follows that the divisor of $R$ is divisible by three.
Factor $C$ into monic irreducibles over $\mathbb F_q$:
\[
 C=c\prod_f f^{e_f},\qquad
 c=\operatorname{lc}(A)/\operatorname{lc}(B).
\]
The divisor condition says $e_f\equiv e_{\bar f}\pmod3$.
Choose a common representative in $\{0,1,2\}$ on each conjugation
orbit.  The products over these orbits, with those representative
exponents, give a monic polynomial $b\in\mathbb F_p[T]$.
The remaining exponents are divisible by three and give a monic
rational function $g$ with $C=cbg^3$.

There is at least one valid base-field argument because $p>2D$.
At any such argument the local orders of $b$ and $g$ can cancel,
but their leading units still give
$C(t)=c\,b_0g_0^3$ with $b_0\in\mathbb F_p^*$ and
$g_0\in\mathbb F_q^*$.  Every element of $\mathbb F_p^*$ is a
cube in $\mathbb F_q^*$.  Thus the cubic class of $C(t)$ is exactly
that of $c$.  This proves that $c$ is a noncube and also proves
the converse, including possible cancelling zeros or poles in the
displayed factorization.
\end{proof}

The base-field factor is essential.  If $a$ is any noncube, then
$C(T)=aT$ has noncubic values at every nonzero base-field argument,
although $C$ is not a constant times a rational cube.  The proposition
therefore factors out both base-field rational functions and cubes.
It says that an all-parameter family of the stated small degree already
exposes a noncube as its leading-coefficient ratio.  It does not bound
the least successful parameter in a family taking both cubic and
noncubic values.

\begin{proposition}\label{prop:prime-parameter-congruence}
Let $p>3$ be prime with $p\equiv-1\pmod9$ and put
\[
 h_k(Y)=Y^3-(k+3)Y^2+kY+1.
\]
Let $m\ge1$ be prime to $p$, and let $(a,m)=1$.
There are infinitely many rational primes $k\equiv a\pmod m$
for which $k\not\equiv3\pmod p$ and $h_k$ is irreducible
over $\mathbb F_p$.  There are also infinitely many such prime
parameters for which $h_k$ splits completely over $\mathbb F_p$.
\end{proposition}

\begin{proof}
Choose $\zeta_3\in\mathbb F_{p^2}$.  The established Tate-family
identity says that $h_k$ is irreducible exactly when
$k-3\zeta_3$ is a noncube.  The lines represented by
$k-3\zeta_3$, for $k\in\mathbb F_p$, are precisely the
$p$ nonscalar lines of
$\mathbb F_{p^2}^*/\mathbb F_p^*$.
This quotient is cyclic of order $p+1$; its cube subgroup has
index three, and the omitted scalar line belongs to it.
Hence exactly $(p-2)/3$ parameters are cubic and
$2(p+1)/3$ are noncubic.

The parameter $k=0$ is cubic because $\zeta_3$ is a cube when
$9\mid p+1$.  Thus some noncubic parameter is different from
both $0$ and $3$.  A nonzero nonsingular split parameter is
$k=-3$, since $-3-3\zeta_3=3\zeta_3^2$ is a cube.
Fix one residue of each kind.  The Chinese remainder theorem combines
it with $k\equiv a\pmod m$ to give a reduced class modulo $mp$.
Dirichlet's theorem gives infinitely many primes in each class.
\end{proof}

Thus a single rational congruence condition of modulus prime to $p$
cannot force the status of all prime parameters in this family.
The conclusion has no least-prime bound and does not exclude a
bounded-window selector.  Cubic reciprocity instead makes a fixed
integer parameter into a cubic Dirichlet character of the varying
input prime; that orientation is compatible with the preceding
almost-all theorem but supplies no pointwise selection theorem.


\subsection{Two possible transfers that do not identify the desired witness}
The following checks distinguish a rational character at a selected
input from a character that detects the required finite-field output.

\begin{proposition}\label{prop:parameter-character-counterexample}
A small prime detected by a cubic character associated with
$p-3\zeta$ need not supply an irreducible polynomial $h_k$ modulo $p$.
For $p=17$, the norm of $17-3\zeta$ is the prime $349$.
The prime $7$ is a cubic nonresidue modulo $349$, since
\[
 7^{116}\equiv226\pmod{349}.
\]
Nevertheless
\[
 h_7(Y)\equiv(Y-2)(Y-9)(Y-16)\pmod{17}.
\]
\end{proposition}
\begin{proof}
The norm is $17^2+3\cdot17+9=349$; trial division by
$2,3,5,7,11,13,17$ proves primality.  Repeated squaring gives
the displayed power, and multiplying the three linear factors
gives $Y^3-10Y^2+7Y+1$ modulo $17$.  A nontrivial cubic
character modulo $349$ therefore detects $7$, but its detection
does not imply that $h_7$ is irreducible at the prescribed input.
\end{proof}

\begin{proposition}\label{prop:split-index-obstruction}
Let a number field $L$ of degree $N$ split completely at a prime
$p$, and suppose $N>p$.  For every integral primitive element
$\alpha$ of $L$, the index $[\mathcal O_L:\mathbb Z[\alpha]]$
is divisible by $p$.
\end{proposition}
\begin{proof}
Complete splitting gives
$\mathcal O_L/p\mathcal O_L\cong\mathbb F_p^N$.
The subalgebra generated by a single element of this product
has dimension equal to the number of its distinct coordinates,
which is at most $p$.  It cannot be the whole algebra when
$N>p$.  If the displayed index were prime to $p$, reduction
of $\mathbb Z[\alpha]\subset\mathcal O_L$ modulo $p$ would be
surjective, contradicting this dimension bound.
\end{proof}
Thus a large compositum splitting completely at $p$ is not
contradicted by counting roots of an arbitrary primitive-element
polynomial.  A use of that polynomial must also control its index.


\subsection{A finite comparison of cubic witnesses}
The following computation illustrates the different sufficient conditions;
it provides no bound beyond its stated range.
For all $1{,}599$ primes $p\le100{,}000$ with $p\equiv-1\pmod9$,
we compared exact modular tests for the six quantities below.
The cubic menu consists of monic polynomials $X^3+aX^2+bX+c$ with
$a\in\{-1,0,1\}$, $-12\le b\le12$, $1\le c\le24$,
negative polynomial discriminant, and no rational root.
These restrictions leave $1{,}451$ entries, ordered by absolute
polynomial discriminant and then lexicographically by coefficients.
A successful entry must be unramified and irreducible modulo $p$.
\begin{table}[htbp]
\centering
\small
\begin{tabular}{@{}p{.60\textwidth}rr@{}}
\toprule
Quantity minimized for each prime & Largest minimum & First prime\\
\midrule
Good auxiliary prime $\ell$ & $43$ & $5{,}669$\\
Eisenstein height $\max(|A|,|B|)$ & $5$ & $5{,}669$\\
Affine parameter $t\ge0$ with $t+\zeta_3$ noncubic & $11$ & $38{,}501$\\
Absolute polynomial discriminant in the stated menu & $204$ & $41{,}579$\\
Positive $m$ with $X^3+X+m$ irreducible & $23$ & $96{,}857$\\
Positive $m$ with $X^3-X+m$ irreducible & $26$ & $35{,}099$\\
\bottomrule
\end{tabular}
\caption{Exact cubic searches for the stated finite range.}
\label{tab:cubic-finite-probe}
\end{table}
At $p=41{,}579$, the first successful menu entry is
$X^3+X^2+X+3$, of discriminant $-204$.
Its field discriminant divides the polynomial discriminant, but
maximal orders and least field discriminants were not computed.
The program \emph{Finite comparison of cubic construction witnesses},
available through the \href{\workbenchurl}{workbench}, records every
prime and witness.  Its irreducibility test uses $\gcd(f,X^p-X)$;
three hundred small cases were also compared with direct root enumeration.
The fixed-menu result below explains why success of any one finite menu
on this interval cannot establish success at every prime.

\subsection{Restrictions on fixed candidates and global units}
\label{subsec:candidate-barriers}

\begin{proposition}[A finite menu of cyclotomic candidates]
\label{prop:fixed-cyclotomic-menu}
Fix an odd prime $r$, an integer $a\ge2$, and finitely many nonzero
elements $\alpha_1,\ldots,\alpha_m$ of $K=\mathbb Q(\zeta_r)$.
There is a positive-density set of rational primes $p$ such that
\[
 \ord_r(p)=2,\qquad v_r(p^2-1)=a,
\]
and every $\alpha_j$ reduces to an $r$th power in every residue
field of $K$ above $p$. The finitely many primes at which a chosen
element has zero, pole, or ramified reduction are omitted.
\end{proposition}

\begin{proof}
Let $M$ be obtained from $K$ by adjoining $r$th roots of all
$\mathbb Q$-conjugates of every $\alpha_j$.
Then $M/\mathbb Q$ is Galois, and $M/K$ is abelian of exponent
dividing $r$ by elementary Kummer theory.
Write $C_b=\mathbb Q(\zeta_{r^b})$ and $L=MC_a$.
The group $\operatorname{Gal}(L/K)$ embeds in
\[
 \operatorname{Gal}(M/K)\times\operatorname{Gal}(C_a/K);
\]
its exponent therefore divides $r^{a-1}$.
Since $C_{a+1}/K$ is cyclic of degree $r^a$, this implies
\[
 L\cap C_{a+1}=C_a.
\]
Indeed, the intersection contains $C_a$, and any larger
intersection would be all of $C_{a+1}$, giving a cyclic quotient
of exponent $r^a$.

Let $c$ denote complex conjugation on $L$.
Choose any $t\in\{1,\ldots,r-1\}$, and let $\gamma$ be the
automorphism of $C_{a+1}$ with cyclotomic exponent
$-1+tr^a\pmod{r^{a+1}}$.
On the intersection $C_a$, both $c$ and $\gamma$ act by
inversion. They consequently extend to a common automorphism
$g$ of $LC_{a+1}$.
By the qualitative Chebotarev density theorem, a positive-density
set of unramified rational primes has Frobenius conjugate to $g$.
Every such prime satisfies
\[
 p\equiv-1+tr^a\pmod{r^{a+1}},
\]
because the cyclotomic Galois group is abelian.
Thus $p\equiv-1\pmod r$ and $v_r(p+1)=a$.

The restriction of each Frobenius representative to $M$ is a
conjugate of complex conjugation and has square $1$.
Each prime of $K$ above $p$ has residue degree $2$.
Its relative Frobenius in $M/K$ is therefore trivial.
Consequently its residue field extension in $M$ is trivial:
the $r$th roots adjoined to $M$ reduce into that same
$\mathbb F_{p^2}$ residue field. This proves the asserted
$r$th-power property at every prime of $K$ above $p$.
\end{proof}

\begin{proposition}[A finite menu of fixed polynomials]
\label{prop:fixed-polynomial-menu}
Fix an odd prime $r$, an integer $a\ge2$, and finitely many
separable polynomials in $\mathbb Q[X]$, each of positive odd
degree at most $r$.
There is a positive-density set of primes $p$ with
$\ord_r(p)=2$ and $v_r(p^2-1)=a$ such that every polynomial
in the menu has a root modulo $p$.
In particular, a fixed finite menu of degree-$r$ polynomials
cannot cover all primes in this exact-$a$ branch by irreducible
reductions.
\end{proposition}

\begin{proof}
Let $M$ be the compositum of $K=\mathbb Q(\zeta_r)$ and the
splitting fields of the polynomials.
The group $\operatorname{Gal}(M/K)$ embeds in a product of
subgroups of $S_r$. The exponent of $S_r$ is
$\operatorname{lcm}(1,\ldots,r)$, whose $r$-adic valuation is
$1$. Hence the $r$-primary part of the exponent of
$\operatorname{Gal}(M/K)$ divides $r$.
For $L=M\mathbb Q(\zeta_{r^a})$, the $r$-primary part of the
group exponent divides $r^{a-1}$.
The intersection and compatibility argument in the preceding
proof therefore applies without requiring $M/K$ to be abelian.
It gives positive-density primes of the prescribed exact
valuation whose Frobenius on $M$ is a conjugate of complex
conjugation.

Complex conjugation is an involution on the distinct roots of
each polynomial. An involution on an odd finite set has a fixed
point. After omitting the primes dividing denominators, leading
coefficients, and discriminants, Frobenius fixed points are
exactly the linear factors of the reduction.
Thus every polynomial in the menu has a root.
\end{proof}

These statements concern fixed finite menus. They do not rule out
an adaptive deterministic algorithm, a menu that grows with
$\log p$, or an element whose coefficients depend on previously
computed data.

\begin{proposition}[Global $r$-units have small Teichm\"uller projection]
\label{prop:global-unit-teich-barrier}
Let $r$ be an odd prime, $K=\mathbb Q(\zeta_r)$, and
$u\in\mathcal O_K[1/r]^\times$.
Let $p\ne r$ be prime, put
\[
 s=\ord_r(p),\quad q=p^s,\quad R=r^a,\quad
 a=v_r(q-1),
\]
\[
 w=(q-1)/R,\qquad D=\mathbb F_p[T]/\Phi_r.
\]
Reduce $u$ to a unit of $D$, and define
\[
 \mathcal P(u)=
 \prod_{b\in(\mathbb Z/r\mathbb Z)^\times}
 \sigma_b(u^w)^{\omega_a(b)^{-1}},
 \qquad \sigma_b(T)=T^b,
\]
with inverse exponents taken modulo $R$.
Then $\mathcal P(u)$ has order dividing $r$.
It cannot be a maximal Teichm\"uller unit when $a\ge2$.
If in addition $s$ is even and $a\ge2$, then the residue of $u$
is itself an $r$th power in every field component of $D$.
\end{proposition}

\begin{proof}
The unique prime of $K$ above $r$ is stable under complex
conjugation. Since the divisor of $u$ is supported there, all
finite valuations of $u/\bar u$ vanish. Thus $u/\bar u$ is an
algebraic unit. Every complex conjugate has absolute value $1$,
because $K$ is a CM field and its embeddings commute with
complex conjugation.
An algebraic integer all of whose conjugates have absolute
value $1$ is a root of unity: the monic polynomials whose roots
are the conjugates of its positive integral powers have bounded
integer coefficients and bounded degree, so only finitely many
such powers occur. Applied here, this gives
\[
 u/\bar u\in\mu(K)=\mu_{2r}.
\]
For completeness, the last equality uses the familiar ramification
and degree facts for prime-conductor cyclotomic fields.
A root of unity of order divisible by an odd prime other than
$r$, or by $4$, would give a subfield ramified at a prime other
than $r$. A root of order divisible by $r^2$ would give a
subfield of degree at least $r(r-1)>[K:\mathbb Q]$.
Neither is possible, and $K$ already contains $\mu_{2r}$.

The map $\mathcal P$ is a group homomorphism into the
$r$-primary part of $D^\times$ and has Teichm\"uller
equivariance. In particular, since
$\omega_a(-1)=-1\pmod R$,
\[
 \mathcal P(\bar u)=\mathcal P(u)^{-1},\qquad
 \mathcal P(u)^2=\mathcal P(u/\bar u).
\]
The last term has order dividing $r$: the $r$-primary part of
the reduction of a root in $\mu_{2r}$ has order dividing $r$.
Squaring is an automorphism of any finite group of odd
$r$-power order, so $\mathcal P(u)$ also has order dividing $r$.

For the final assertion, suppose $s$ is even.
The unique order-two element of
$(\mathbb Z/r\mathbb Z)^\times$ is $-1$, whence
$p^{s/2}\equiv-1\pmod r$.
Consequently, in every component $\mathbb F_q$,
complex conjugation acts as $p^{s/2}$th-power Frobenius.
The equation $u/\bar u\in\mu_{2r}$ then implies that the
$r$-primary order of
$u^{p^{s/2}-1}$ is at most $r$.
Since $r\nmid p^{s/2}-1$, the same is true of $u$.
The $r$-primary subgroup of $\mathbb F_q^\times$ is cyclic of
order $r^a$. When $a\ge2$, all of its elements of order at most
$r$ are $r$th powers, while the $r$th-power map is bijective on
the prime-to-$r$ part. Hence $u$ is an $r$th power in that
component.
\end{proof}

The proposition applies to the entire global unit group and to
the units obtained by allowing denominators at $r$.
For Gauss or Jacobi sums it applies only when the sum in question
lies in $K$ and has divisor supported above $r$; arbitrary such
sums need not satisfy those hypotheses.

\subsection{Jacobi sums and the semiprimitive case}
Point counting can sometimes construct algebraic constants, but the
splitting conditions in that use of point counting are consequential.
For a fixed odd prime $\ell$, Pila gives an unconditional algorithm
polynomial in $\log p$ for the $\ell$th roots of unity in
$\mathbb F_p$ when $p\equiv1\pmod\ell$
\citep[Theorem~D, author p.~2]{Pila1990}.
His argument uses Jacobi sums to distinguish the prime ideals of
$\mathbb Q(\zeta_\ell)$ above the completely split prime $p$
\citep[\S5, especially Propositions~5.1--5.3, author pp.~14--18]{Pila1990}.
The exponent and constants in his general point-counting theorem
also depend on the defining equations and their degrees
\citep[Theorem~A, author p.~1]{Pila1990}; the statement is not a
uniform polynomial bound in a growing $\ell$.

The following elementary Frobenius symmetry explains a further
restriction on this particular Jacobi-sum strategy.  It is the
standard semiprimitive phenomenon; the same symmetry appears in
\citet[Lemma~1, author p.~3]{Akiyama1996} for a quadratic second
character.  We give the two-character form needed here.

\begin{proposition}\label{prop:semiprimitive-jacobi}
Let $q=p^{2j}$, where $p$ is prime and $j\ge1$, and let
$m>2$ divide $p^j+1$.  Let $\chi$ be a multiplicative character
of exact order $m$ on $\mathbb F_q^\times$.
If $a,b,a+b$ are all nonzero modulo $m$, then
\[
 J(\chi^a,\chi^b)
 :=\sum_{x\in\mathbb F_q}\chi^a(x)\chi^b(1-x)
 \quad\hbox{belongs to }\{p^j,-p^j\},
\]
where all characters are extended by zero at zero.
Consequently these values generate no nontrivial number field,
and their principal ideals do not distinguish primes above $p$.
\end{proposition}

\begin{proof}
Put $Q=p^j$.  The substitution $x\mapsto x^Q$ permutes
$\mathbb F_q$ and satisfies $(1-x)^Q=1-x^Q$.  Therefore
\[
 J(\chi^a,\chi^b)=J(\chi^{Qa},\chi^{Qb})
   =J(\chi^{-a},\chi^{-b})
   =\overline{J(\chi^a,\chi^b)}.
\]
The sum is real.  To recall its absolute value, take any nontrivial
additive character $e$ of $\mathbb F_q$ and put
$G(\alpha)=\sum_x\alpha(x)e(x)$.
For a nontrivial multiplicative character $\alpha$, substituting
$x=uy$ and using additive and multiplicative orthogonality gives
\[
 |G(\alpha)|^2
 =\sum_{u\ne0}\alpha(u)\sum_{y\ne0}e((u-1)y)=q.
\]
If $\alpha,\beta,\alpha\beta$ are nontrivial, sorting the
terms in $G(\alpha)G(\beta)$ according to $x+y$ gives
$G(\alpha)G(\beta)=J(\alpha,\beta)G(\alpha\beta)$;
the terms with $x+y=0$ sum to zero.
Thus $|J(\alpha,\beta)|=\sqrt q$.
Apply this with $\alpha=\chi^a$, $\beta=\chi^b$.
The real number in question has absolute value $Q$, so it is
$Q$ or $-Q$.  Its ideal in $\mathbb Z[\zeta_m]$ is
$(p)^j$ in either case.
\end{proof}

In the remaining Kummer case, suppose $s=\operatorname{ord}_r(p)$
is even and $a=v_r(p^s-1)$.
Then $r\nmid p^{s/2}-1$, whereas $r\mid p^{s/2}+1$;
hence $r^a\mid p^{s/2}+1$.
Proposition~\ref{prop:semiprimitive-jacobi} applies with
$q=p^s$, $j=s/2$, and $m=r^a$.
Thus extending the base field in Pila's argument does not preserve
its Jacobi-sum separation mechanism in this case.
This observation concerns those Jacobi sums; it does not exclude
other curves, twists, or uses of point counting.

\section{An auxiliary cyclotomic conductor for Jacobi sums}
\label{sec:auxiliary-cm-conductor}

An enlargement of the cyclotomic conductor can remove the
semiprimitive phenomenon and can make Jacobi ideals distinguish
an entire \(r\)-power subfield. The statements below establish
that algebraic distinction. They do not assert a uniform
polynomial-time method for computing the required Jacobi generator.

\begin{lemma}\label{lem:quadratic-cm-auxiliary}
Assume rational Dirichlet GRH. Let \(p,r\) be distinct odd primes,
let \(d=\operatorname{ord}_r(p)\) be even, and put \(t=v_2(d)\).
There is a prime \(q\notin\{2,p,r\}\) such that
\[
 q\equiv1\pmod{2^{t+1}},\qquad
 \left(\frac pq\right)=-1,\qquad
 q\ll d^2\log^2(2pd),
\]
with an absolute effective implied constant. Such a prime can be
found deterministically in time polynomial in \(r\) and \(\log p\).
If \(a=v_r(p^d-1)\), then for every \(1\leq b\leq a\) there is
no integer \(j\) satisfying \(p^j\equiv-1\pmod{r^bq}\).
\end{lemma}
\begin{proof}
Put \(M_2=2^{t+1}\), and let \(\kappa\) be the primitive quadratic
Dirichlet character of \(\mathbb Q(\sqrt p)\). On odd primes
different from \(p\), it takes the value \((p/q)\).
Expand the indicator of \(1\bmod M_2\) by the characters modulo
\(M_2\), and multiply it by
\((\kappa_0-\kappa)/2\), where \(\kappa_0\) is the principal
character of the same modulus as \(\kappa\).
All terms have modulus dividing \(4pM_2\).
Exactly one is principal, with coefficient
\(1/(2\varphi(M_2))\); the sum of absolute coefficients is at most
\(1\). The nontrivial character at the odd prime \(p\) ensures
that no term involving \(\kappa\) is principal.

The complex smoothed Dirichlet estimate of
Lemma~\ref{lem:complex-dirichlet-tent} gives
\[
 \sum_{m\leq x}\Lambda(m)\log(x/m)\,
 1_{m\equiv1\bmod M_2}
 \frac{\kappa_0(m)-\kappa(m)}2
 =
 \frac{x}{2\varphi(M_2)}
 +O\bigl(\sqrt x\log(2pM_2)\bigr).
\]
Even prime powers contribute zero. Higher odd prime powers
contribute at most
\(O(x^{1/3}(\log x)^2)=O(\sqrt x)\).
Taking \(x\) a sufficiently large absolute multiple of
\(\varphi(M_2)^2\log^2(2pM_2)\) therefore gives a prime of the
required kind. Since \(M_2\leq2d\), this is the displayed bound.
Testing all integers up to the bound, using primality testing and
the Jacobi symbol, has the stated polynomial cost.

The quadratic condition gives
\[
 v_2(\operatorname{ord}_q p)=v_2(q-1)\geq t+1.
\]
For \(b\leq a\), the order of \(p\bmod r^b\) is still \(d\).
An exponent giving \(-1\) modulo \(r^b\) must have
\(2\)-adic valuation \(t-1\), whereas one giving \(-1\) modulo
\(q\) has valuation at least \(t\). They cannot coincide.
The same order comparison excludes \(q=r\), the quadratic
condition excludes \(q=p\), and the progression excludes \(q=2\).
\end{proof}

The assumption that \(p\) is odd is necessary for this construction:
when \(p=2\), a sufficiently strong \(2\)-power progression forces
\((2/q)=1\). That characteristic already belongs to the separate
small-characteristic algorithmic case.

\begin{proposition}\label{prop:jacobi-wild-subfield}
Let \(p,r,d,a,q\) satisfy Lemma~\ref{lem:quadratic-cm-auxiliary},
let \(2\leq b\leq a\), and put
\[
 M=r^bq,\qquad K_M=\mathbb Q(\zeta_M),\qquad
 G=(\mathbb Z/M\mathbb Z)^\times,\qquad H=\langle p\rangle.
\]
Let \(R_b\) be the unique subfield of \(\mathbb Q(\zeta_{r^b})\)
of degree \(r^{b-1}\) over \(\mathbb Q\).
Fix a prime \(\mathfrak P\) of \(K_M\) above \(p\), and let
\(\chi_{\mathfrak P}\) be the order-\(M\) power-residue character
of its residue field. The Jacobi sum
\[
 J=-\sum_{z\in\mathcal O_{K_M}/\mathfrak P}
       \chi_{\mathfrak P}(z)\chi_{\mathfrak P}(1-z)
 \quad\text{satisfies}\quad
 R_b\subseteq\mathbb Q(J).
\]
Moreover the tuple of valuations of all Galois conjugates of \(J\) distinguishes
any two primes above \(p\) whose restrictions to \(R_b\) are
different.
\end{proposition}
\begin{proof}
Write \(r-1=2^u m\) with \(m\) odd. Choose a character \(\tau\)
modulo \(r\) of order \(2^u\); it is odd, and \(\tau(p)\) has
order \(2^t\). Choose a character \(\eta\) modulo \(q\) of
\(2\)-power order satisfying
\(\eta(p)=\tau(p)^{-1}\).
This is possible because \(p\) is a quadratic nonresidue modulo
\(q\), so evaluation at \(p\) maps the \(2\)-primary character
group modulo \(q\) onto its full \(2\)-power roots of unity.
The character \(\eta\) has order \(2^t\), is nonprincipal, and
is even: the \(2\)-part of \(\operatorname{ord}_q p\) is strictly
larger than \(2^t\).

Let \(\nu\) be a primitive character modulo \(r^b\) of order
\(r^{b-1}\), trivial on the Teichm\"uller subgroup.
Because \(p\) has order \(d\) modulo \(r^b\), we have \(\nu(p)=1\).
Thus
\[
 \psi=\tau\nu\eta
\]
is an odd primitive character of conductor exactly \(M\), and
\(\psi(p)=1\). The factors \(\tau\eta\) and \(\nu\) have coprime
orders. In particular,
\begin{equation}\label{eq:cm-character-kernel}
 \ker\psi\subseteq\ker\nu,
\end{equation}
where \(\nu\) is inflated to \(G\).

Use the convention
\[
 \theta_M(c)=
 \sum_{s\in G}\left\langle-\frac{cs}{M}\right\rangle\sigma_s^{-1},
 \qquad \alpha=2\theta_M(1)-\theta_M(2).
\]
The composite-conductor Stickelberger formula gives
\[
 (J)=\mathfrak P^\alpha.
\]
The identity holds for arbitrary residue degree; it is stated in
\citet[Proposition~5.1, author p.~17, proof p.~18]
{BernardKucera2021}, with precisely this convention for
\(\theta_M\) and for the sign of \(J\).
Here \(\alpha\) has coefficients \(0\) or \(1\).

Evaluating at \(\psi\), a substitution \(v=cs\bmod M\) gives
\[
 \psi(\theta_M(c))
   =\psi(c)L(0,\overline\psi)
 \quad ((c,M)=1).
\]
Indeed the constant term in
\(\langle-s/M\rangle=1-s/M\) sums to zero for a nonprincipal
character. Consequently
\[
 \psi(\alpha)=(2-\psi(2))L(0,\overline\psi)\ne0.
\]
The first factor is nonzero because \(|\psi(2)|=1\); the second
is nonzero for an odd primitive Dirichlet character, by its
functional equation and nonvanishing at \(1\).
Equivalently, the nonvanishing follows from the classical
relative class-number product recorded in
\citet[equation~(31), author p.~18]{BernardKucera2021}.

The primes above \(p\) form the \(G\)-set \(G/H\).
Let \(\overline\alpha\in\mathbb Z[G/H]\) be the valuation vector
of \((J)\). Since \(\psi|_H=1\), the evaluation just computed
descends to this permutation module and remains nonzero.
If \(g\in G\) stabilizes \((J)\), then
\[
 \psi(\overline\alpha)
 =\psi(g\overline\alpha)
 =\psi(g)\psi(\overline\alpha),
\]
so \(g\in\ker\psi\subseteq\ker\nu\).
Every automorphism fixing \(J\) fixes its ideal. Hence
\[
 \operatorname{Gal}(K_M/\mathbb Q(J))
 \subseteq\ker\nu.
\]
The fixed field of \(\ker\nu\) is \(R_b\), proving the inclusion.

For the final assertion, if two primes have different
restrictions to \(R_b\), the group element taking one to the
other has \(\nu(g)\ne1\). It cannot stabilize the entire
valuation vector. Translating a coordinate where the valuations
differ amounts to replacing \(J\) by a Galois conjugate.
That conjugate distinguishes the given pair.
\end{proof}

\paragraph{An exact small example.}
For \(p=17\), \(r=3\), \(q=5\), and \(M=45\), the decomposition
subgroup is \(H=\{1,17,19,8\}\). The kernel of reduction from
conductor \(45\) to \(15\) is \(\{1,16,31\}\).
The three cosets are
\[
 \{1,8,17,19\},\quad
 \{2,16,34,38\},\quad
 \{4,23,31,32\}.
\]
Label primes by inverse cosets, consistently with the coefficient
of \(\sigma_s^{-1}\) in \(\alpha\). That coefficient is
\(1\) exactly when \(s<M/2\), their valuation sums are
\(4,2,1\). Thus the enlarged conductor gives
ideal distinction, whereas conductor \(9\) alone is semiprimitive.
This calculation gives an ideal-valuation certificate; it does
not compute the element \(J\).

\paragraph{The conductor and computation requirements.}
Taking \(b=2\) gives
\(M\ll r^2d^2\log^2(2pd)\), polynomial in \(r\) and \(\log p\).
The proposition then distinguishes the first degree-\(r\)
cyclotomic layer. It does not supply a maximal \(r\)-primary unit
when \(a>2\).
Taking \(b=a\) distinguishes all the required \(r\)-power layers,
but its numerical conductor contains \(r^a\), which need not be
polynomial in \(r\) and \(\log p\).
There is also a degree requirement intrinsic to this particular
Jacobi generator:
\[
 [\mathbb Q(J):\mathbb Q]\geq r^{b-1}.
\]
If \(J\) occurs as a Frobenius eigenvalue of an abelian variety
of dimension \(g\), then \(2g\geq r^{b-1}\).

Even at \(b=2\), the point-counting theorem of
\citet[Theorem~A, author p.~1]{Pila1990} has an exponent depending
on the defining equations and dimension. It does not give a
uniform polynomial bound when the auxiliary conductor \(q\)
and the Fermat genus grow with \(\log p\).
Knowing the ideal exponent \(\alpha\), or its nonzero Fourier
coefficient, also does not give a polynomial-time computation
of a generator of \(\mathfrak P^\alpha\).
The positive distinction above therefore leaves a concrete
computational requirement: construct a suitably separating
Jacobi generator uniformly within the allowed input budget,
and then give the required maximal-unit extraction.
Neither requirement is supplied by the ideal calculation alone.

More explicitly, even when \(a=b=2\), the output needed for the
finite-field problem is an explicitly represented element in the
given target field whose \(r\)-primary part has order \(r^2\).
The proposition does not produce that element. A route through its
Jacobi sum must still compute a representation of the sum, or a
separating Frobenius polynomial with sufficiently explicit CM
endomorphisms, and justify the reduction and factor-selection step
in the presented finite field. The valuation assertion concerns
the entire tuple of conjugates; it does not assert that the single
valuation of one chosen sum distinguishes every prime.

\begin{corollary}\label{cor:sextic-cm-field}
For \(r=3\) and \(b=2\), the field cut out by \(\psi\) above is
\[
 L=\mathbb Q(\zeta_9+\zeta_9^{-1})
       \mathbb Q(\sqrt{-3q}).
\]
It is cyclic of degree \(6\), \(p\) splits completely in it, and
its absolute discriminant is \(3^9q^3\).
\end{corollary}
\begin{proof}
Indeed, the order-\(3\) part is the real cubic field of conductor
\(9\), and the order-\(2\) part is the product of the quadratic
characters of conductors \(3\) and \(q\). Both characters take
the value \(1\) at \(p\). The real cubic field has discriminant
\(81\), while \(\mathbb Q(\sqrt{-3q})\) has discriminant
\(-3q\), since \(q\equiv1\pmod4\). The tensor product of integral
bases gives an order in \(L\), so its discriminant yields
\[
 |D_L|\leq81^2(3q)^3=3^{11}q^3.
\]
The five nonprincipal characters have conductors
\(9,9,3q,9q,9q\). The conductor-discriminant identity therefore
gives the exact value \(3^9q^3\). At \(3\), the wild cubic
part cannot cancel the tame quadratic part; at \(q\), the
quadratic part is nontrivial. This proves the conductor list.
\end{proof}
Thus a suitable CM construction in dimension \(3\) would avoid a
dimension growing with \(q\). This field calculation supplies
neither such a construction with a polynomial cost in \(q\), nor
a reduction procedure that avoids the original finite-field
factor-selection problem.

The distinction is material in the available CM factorization
literature. \citet[Chapter~III, printed p.~50/PDF p.~62]{Boyer2014}
fixes the conductor, measures complexity only in \(\log p\), and
allows arbitrary finite preprocessing independent of \(p\).
His Proposition~2.6, printed p.~67/PDF p.~79, proves separation
for the displayed CM curves with field
\(\mathbb Q(\zeta_l+\zeta_l^{-1},i)\), for prime \(l>5\),
\(p\equiv1\pmod4\), and \(p\equiv-1\pmod l\).
The broader even-order claims are Conjecture~2.7
(printed p.~68/PDF p.~80) and Conjecture~2.11
(printed p.~72/PDF p.~84). Those results do not provide a uniform
varying-\(q\) model for the sextic field just displayed.

\section{Explicit primitive CM periods for the sextic construction}
\label{sec:sextic-cm-periods}

The particular sextic fields arising above admit explicit
principally polarized periods. Their computation does not need
a general period-reduction algorithm.

\begin{proposition}\label{prop:explicit-sextic-cm-periods}
Let \(q\ne3\) be an odd prime with \(q\equiv1\pmod4\), put
\(D=3q\), and let
\[
 f(X)=X^3-3X+1,\quad R=\mathbb Q(\theta),\quad f(\theta)=0,
 \quad s=\sqrt{-D},\quad L=R(s).
\]
Let \(\theta_1,\theta_2,\theta_3\) be the real roots of \(f\),
\(S=(\theta_i^{j-1})_{1\leq i,j\leq3}\), and
\(\epsilon_i=\operatorname{sign}f'(\theta_i)\).
The matrix
\[
 B=\begin{pmatrix}0&0&1\\0&1&0\\1&0&3\end{pmatrix},
 \qquad
 \tau=\frac12B^{-1}
 +\frac{i\sqrt D}{2}
 S^{-1}\operatorname{diag}(|f'(\theta_i)|)S^{-T}
\]
belongs to \(\mathcal H_3\). The principally polarized torus
\(\mathbb C^3/(\mathbb Z^3+\tau\mathbb Z^3)\)
has primitive CM by the order
\[
 \mathcal O=\mathbb Z[\theta,(1+s)/2]\subset L.
\]
There are absolute positive constants \(c,C\) such that
\[
 c\sqrt D\,I_3\preceq\operatorname{Im}\tau
                   \preceq C\sqrt D\,I_3.
\]
The periods and all genus-three theta constants at \(\tau\)
with characteristic denominator \(2\) or \(4\)
can be approximated with certified absolute error \(2^{-N}\)
in deterministic time polynomial in \(D\) and \(N\).
Moreover,
\[
 |\operatorname{disc}(L)|=3^9q^3,
 \qquad [\mathcal O_L:\mathcal O]=3.
\]
\end{proposition}
\begin{proof}
The polynomial \(f\) has discriminant \(81\) and three distinct
real roots; it defines the real cubic subfield of
\(\mathbb Q(\zeta_9)\). Its power-basis trace dual is
\(f'(\theta)^{-1}\mathbb Z[\theta]\).
The quadratic ring of integers is
\(\mathbb Z[\omega]\), where \(\omega=(1+s)/2\), with trace
dual \(s^{-1}\mathbb Z[\omega]\).
The two fields are linearly disjoint. The product basis is an
integral basis of the order \(\mathcal O\), and its trace dual
is
\[
 \mathcal O^\vee=\xi\mathcal O,
 \qquad \xi=(f'(\theta)s)^{-1}.
\]
This assertion concerns the product order and does not require
that it be the maximal order of \(L\).
In fact, the six characters of the cyclic extension \(L/\mathbb Q\)
are \(1,\nu,\nu^2,\chi,\chi\nu,\chi\nu^2\), where \(\nu\) is a
cubic character of conductor \(9\) and \(\chi\) is the quadratic
character of conductor \(3q\). Their conductors are respectively
\(1,9,9,3q,9q,9q\). At \(3\), multiplication by the tame quadratic
character cannot cancel the nontrivial wild cubic character; at \(q\)
the quadratic character remains ramified. The conductor-discriminant
identity therefore gives \(|\operatorname{disc}(L)|=3^9q^3\).
The product basis has discriminant
\(81^2(3q)^3=3^{11}q^3\), proving the index assertion.

Choose the CM embeddings \(\phi_i\) by
\(\phi_i(\theta)=\theta_i\) and
\(\phi_i(s)=\epsilon_i i\sqrt D\).
On \(\mathcal O\), define
\[
 E(x,y)=\operatorname{Tr}_{L/\mathbb Q}(\xi\overline{x}y).
\]
The trace-dual identity shows that \(E\) is an integral
unimodular alternating form. With this sign convention,
\(\phi_i(\xi)=-i/(|f'(\theta_i)|\sqrt D)\), so
\(E(z,iz)>0\) for every nonzero complex tangent vector.
It is therefore a principal Riemann form.

For \(0\leq j,k\leq2\),
\[
 E(\theta^j,\omega\theta^k)
 =\operatorname{Tr}_{R/\mathbb Q}
             \left(\frac{\theta^{j+k}}{f'(\theta)}\right)
 =B_{j+1,k+1}.
\]
Here the traces for exponents \(0,1,2,3,4\) are
\(0,0,1,0,3\), by the power-basis dual identity and the
relation \(\theta^3=3\theta-1\).
Both remaining diagonal blocks of \(E\) are zero.
Since \(\det B=-1\), the basis with first block
\((1,\theta,\theta^2)\) and second block
\(\omega(1,\theta,\theta^2)B^{-1}\) is symplectic.
Its complex period blocks are
\[
 \Omega_1=S,\qquad
 \Omega_2=
 \left(\frac12 I_3+\frac{i\sqrt D}{2}
                        \operatorname{diag}(\epsilon_i)\right)
 S B^{-1}.
\]
The identity
\(B=S^T\operatorname{diag}(1/f'(\theta_i))S\)
now gives exactly \(\tau=\Omega_1^{-1}\Omega_2\) as displayed.
Its imaginary part is positive definite. Its eigenvalues are
\(\sqrt D\) times those of one fixed positive definite matrix,
which proves the two uniform inequalities.

The field \(L/\mathbb Q\) is cyclic of degree \(6\).
Its only proper CM subfield is \(\mathbb Q(s)\).
An induced CM type would choose the same sign of \(s\) at
all three real embeddings. The signs of \(f'\) at the three
ordered real roots are \(+,-,+\), so the chosen type is
primitive. The Riemann form construction gives a principally
polarized abelian variety; primitivity implies simplicity, by
the standard CM criterion recalled in
\citet[author/PDF p.~3]{BILV2016}.

For completeness, certified theta evaluation here can use direct
summation. For \(r\in\{2,4\}\) and a characteristic
\(a/r,b/r\), where \(a,b\in\{0,\ldots,r-1\}^3\), the
absolute value of its term indexed by \(n\in\mathbb Z^3\) is
at most
\[
 \exp\bigl(-\pi c\sqrt D\,\|n+a/r\|^2\bigr).
\]
The sum outside \(\|n\|_\infty\leq K\), for \(K\geq2\),
is bounded by \(C_1\exp(-c_1\sqrt D\,K^2)\), with absolute
effective positive constants. This follows by separating the
three coordinates, bounding the one-dimensional Gaussian tails
by an integral, and taking a union over the three coordinates.
It suffices to take
\(K=O(1+\sqrt{(N+1)/\sqrt D})\), and the number of terms is
polynomial in \(N\).
All entries of the period matrix are fixed-degree algebraic
numbers given explicitly above. Certified algebraic arithmetic
and interval evaluation of exponentials therefore calculate the
finite sum and its tail bound with a polynomial cost in
\(D\) and \(N\); guard precision polynomial in these quantities
suffices. No assumption about a nonzero theta constant is used.
\end{proof}

The next proposition gives effective degree and height bounds for these theta coordinates and recovers them exactly. Subsequent sections construct the finite-field model and its addition atlas.

\section{Effective algebraization of the explicit sextic theta point}
\label{sec:sextic-cm-algebraization}

The exact period description above gives more than a numerical candidate:
under rational Dirichlet GRH its theta-null point can be recovered as an
exact algebraic point with a polynomial cost in the numerical auxiliary
prime. The statement here concerns that point, rather than a smooth
integral model and an addition atlas.

\begin{proposition}\label{prop:sextic-theta-recognition}
Assume rational Dirichlet GRH, and use the notation of
Proposition~\ref{prop:explicit-sextic-cm-periods}.
Fix \(r\in\{2,4\}\). Let
\(T_r\in\mathbb P^{r^6-1}(\overline{\mathbb Q})\) be the theta-null point
with coordinates
\(\vartheta[a/r,b/r](\tau,0)\), for
\(a,b\in\{0,\ldots,r-1\}^3\), and write \(T=T_r\).
There are effectively computable absolute constants \(C_0,C_1\) such that
\[
 [\mathbb Q(T):\mathbb Q]
 \leq B:=C_0q^{3/2}(1+\log q)^5,
 \qquad h(T)\leq C_1\log(2q).
\]
One can choose a nonzero coordinate, recover all ratios to that
coordinate exactly as algebraic numbers, and construct a common
number-field presentation of those ratios in deterministic time
polynomial in the numerical value of \(q\).
The polynomial's exponent is absolute; it does not depend on \(q\).
\end{proposition}
\begin{proof}
Write \(A\) for the principally polarized complex abelian variety
defined by the periods. It is algebraic over \(\overline{\mathbb Q}\),
by CM algebraicity; see \citet[Proposition~7.9, p.~54]{MilneCM2020}.
The algebraic theta embedding and its identification with these
analytic coordinates are given in
\citet[Sections~2.3.1--2.3.4, pp.~27--31]{Pazuki2012}.

First we bound the height. Saturating its lattice from \(\mathcal O\)
to \(\mathcal O_L\) gives a degree-\(3\) isogeny \(A\to A'\), where
\(A'\) has full CM by \(\mathcal O_L\). Multiplication by \(3\) gives
an isogeny \(A'\to A\) of degree \(3^5\), because
\(3\mathcal O_L\subset\mathcal O\) and
\([\mathcal O_L:\mathcal O]=3\). The isogeny inequality therefore gives
\[
 h_F(A)\leq h_F(A')+\tfrac52\log3;
\]
see \citet[Remark~2.3, p.~27]{Pazuki2012}.
The abelian Colmez formula applies to \(A'\). In the original source,
Theorem~0.3(ii), p.~631, combines with Theorem~III.2.9(i), p.~680;
the conductor here is odd, so the relevant function is invariant under
inertia at \(2\) and belongs to the subspace defined on p.~676
\citep{Colmez1993}. Thus no undetermined multiple of \(\log2\) occurs.

For an explicit upper bound one may use
\citet[Corollary~4.3, author/PDF pp.~14--15;
Theorem~4.8, p.~18]{Wei2021}. The reflex field is a subfield of
the abelian field \(L\), so its Galois closure is itself. Every
nontrivial odd Artin character in that corollary is consequently
a one-dimensional Dirichlet character over \(\mathbb Q\), of conductor at most
\(9q\). Rational Dirichlet GRH excludes its zeros from the short
region in Corollary~4.3 with \(c=1/4\). That corollary and the Colmez
formula imply
\[
 h_F(A')\leq C\log|\operatorname{disc}(L)|+C
           \leq C'\log(2q),
\]
with effective absolute constants; the dimension is always \(3\).
Pazuki's Corollary~1.3(2), pp.~21--22, now gives
\(h(T)=h_\Theta(A)\leq C_1\log(2q)\).
The Euclidean-norm height used there is at least the usual maximum-norm
projective height, so it also gives the latter bound claimed here.

Next we count conjugates. An \(\mathcal O\)-stable lattice \(I\subset L\)
satisfies
\[
 3\mathcal O_L I\subset I\subset\mathcal O_L I.
\]
Up to multiplication by \(L^\times\), the ideals \(\mathcal O_LI\)
have \(h_L\) possible classes. For each class the intermediate
lattices correspond to subgroups of a group of order \(3^6\), hence
their number is bounded by an absolute constant. There are at most
\(2^3\) CM types. On a fixed simple CM variety, two principal
polarizations differ by a totally positive unit of the real subfield;
equivalence under automorphisms identifies units differing by
\(u\overline u\). This quotient is a quotient of
\(\mathcal O_R^{\times,+}/(\mathcal O_R^\times)^2\), because
\(\mathcal O_R\subset\mathcal O\), and has at most \(2^3\) elements.
These lattice and isomorphism classifications are the analytic CM
classification of \citet[Section~3, pp.~29--31]{MilneCM2020}.
Choosing a symmetric representative of the polarization and a fixed
theta structure introduces only an absolute finite number of choices.
In particular, Pazuki's construction on pp.~28--29 changes its
coordinates only by the fixed finite group of permutations and
root-of-unity diagonal multipliers described there. Since the
construction commutes with field automorphisms, the Galois orbit of
\(T\) has size at most \(C h_L\).

Minkowski's ideal-class bound gives an integral ideal of norm at most
\(C\sqrt{|\operatorname{disc}(L)|}\) in every class. The number of
ideals of a given norm \(n\) is at most the six-fold divisor function
\(d_6(n)\): this follows coefficientwise from the local Euler factors.
The elementary bound
\(\sum_{n\leq X}d_6(n)\leq X(1+\log X)^5\)
follows by summing over five of the six factors. Since
\(|\operatorname{disc}(L)|=3^9q^3\), the asserted degree bound follows.

There is also an effective analytic normalization. Since
\(\operatorname{Re}\tau=B_0^{-1}/2\), where \(B_0\) denotes the
fixed trace matrix in the period formula, and \(B_0^{-1}\) is integral,
summing over the \(r^3\) characteristics with
\(a=0\) gives
\[
 \sum_{b\in\{0,\ldots,r-1\}^3}\vartheta[0,b/r](\tau,0)
 =r^3\sum_{n\in(r\mathbb Z)^3}
       \exp(-\pi n^T\operatorname{Im}\tau\,n)\geq r^3.
\]
The Fourier sum over \(b\) kills every other lattice vector, and the
real phases on the remaining vectors equal \(1\), since \(r\) is even.
Therefore some coordinate has absolute value at least \(1\).
Evaluating these \(r^3\)
coordinates to a fixed certified accuracy chooses one with absolute
value at least \(1/2\). All theta constants have an absolute upper
bound independent of \(q\), from the Gaussian estimate in the period
proposition. Their ratios can thus be approximated to \(2^{-N}\)
with polynomial cost in \(q,N\).

For any normalized coordinate \(\alpha\), its degree is at most
\(B\) and \(h(\alpha)\leq h(T)\). Its primitive integer minimal
polynomial therefore has coefficient height at most
\[
 H=2^B\exp\bigl(B C_1\log(2q)\bigr),
\]
by the relation between Mahler measure and absolute logarithmic height.
The reconstruction theorem of
\citet[Theorem~1.19, p.~241]{KannanLenstraLovasz1988}
recovers that polynomial from a certified approximation using
polynomially many bit operations in \(B,\log H\). For example,
\(N=O(B^2+B\log H+B\log B)\) suffices. Root separation then gives
an isolating complex description of the intended root. There are only
\(r^6\leq4096\) coordinates. Standard characteristic-zero factorization and
primitive-element constructions combine their presentations in a
common number field; \citet[Theorem~4.5, pp.~253--254]{Lenstra1983}
provides deterministic polynomial-time number-field factorization.
Even the crude compositum degree bound \(B^{4096}\) is polynomial in
\(q\). Intermediate heights, separation precision, and primitive-element
coefficient sizes are polynomial in these degree and height bounds.
This proves the cost assertion.
\end{proof}

This proposition constructs the characteristic-zero marked theta point. Simultaneous projective reduction and the marked good-reduction theorem below identify its reduction at every prime over the target characteristic.


\begin{proposition}\label{prop:unconditional-cm-recognition}
For the sextic family of Proposition~\ref{prop:explicit-sextic-cm-periods},
exact algebraic recognition of its fixed-level theta point has
deterministic cost polynomial in the numerical auxiliary prime \(q\)
without any Riemann-hypothesis assumption. A sufficient unconditional
height bound is \(h(T_4)\le Cq^3\). The sharper logarithmic height
bound in Proposition~\ref{prop:sextic-theta-recognition} uses \(A\).
\end{proposition}
\begin{proof}
For a maximal-order CM variety, apply
\citet[Lemma~7.10, printed p.~63; equation~(7.7), p.~62]{VonKanelKret2023}
with dimension three and \(\epsilon=1/2\). Its real field is fixed,
and its relative discriminant norm is
\(c=|D_L|/D_R=3^5q^3\). The real class number is fixed, the CM
class number is at least one, the Hasse unit index is at most two,
and the number of roots of unity is bounded in degree six. The explicit
bound is therefore \(h_F\le Cc+C\log c+C\le C'q^3\).
That lemma uses the proved averaged Colmez formula and effective
analytic estimates; it assumes no zero hypothesis. The degree-three
lattice isogeny and its reverse change the height by a constant.
Pazuki's fixed-level comparison gives \(h(T_4)\le Cq^3\).
The orbit degree and certified theta approximation in
Proposition~\ref{prop:sextic-theta-recognition} are unconditional.
Replacing its logarithmic height bound by \(Cq^3\) leaves the
minimal-polynomial coefficient length, recognition precision, and
all algebraic-number operations polynomial in \(q\). This proves
the assertion.
\end{proof}

\section{Simultaneous projective reduction at a given prime}
\label{sec:kummer-projective-reduction}

Reduction of projective coordinates can require division by different
common factors at different primes. The coordinate ideal retains
these factors without choosing any residue-field component.

\begin{lemma}\label{lem:kummer-rank-one-patching}
Let \(B\) be an explicitly presented finite \'etale algebra over
\(\mathbb F_p\), and let \(M\) be an explicitly presented
\(B\)-module of rank one, with specified generators
\(x_0,\ldots,x_m\). A basis \(w\) of \(M\) over \(B\), and
the coordinates \(c_i\in B\) defined by \(x_i=c_iw\), can be
computed deterministically using polynomially many operations in
\(\mathbb F_p\). The \(c_i\) generate the unit ideal of \(B\).
No decomposition of \(B\) into fields is required.
\end{lemma}
\begin{proof}
For \(w\in M\), compute the kernel
\(A_w=\operatorname{Ann}_B(w)\) of the linear map
\(B\longrightarrow M\), \(b\longmapsto bw\).
Every ideal of a product of fields has an identity idempotent.
Find the identity \(f_w\) of \(A_w\) by solving
\[
 f_w\in A_w,\qquad f_wa_j=a_j
\]
for a vector-space basis \(a_j\) of \(A_w\). These are linear
equations, with a unique solution; if \(A_w=0\), take \(f_w=0\).
In a field component, \(f_w\) is \(1\) where \(w=0\) and
\(0\) where \(w\ne0\).

Starting with \(w=0\), successively replace
\[
 w\longleftarrow w+f_wx_i\qquad(0\leq i\leq m).
\]
An already nonzero component is preserved, and a zero component
acquires the component of \(x_i\). Since the \(x_i\) generate
\(M\), the final \(w\) is nonzero in every component and hence
is a basis. Solve \(c_iw=x_i\) by linear algebra. The isomorphism
\(B\longrightarrow M\), \(b\mapsto bw\), takes the ideal
generated by the \(c_i\) onto the submodule generated by the
\(x_i\), so that ideal is \(B\).
\end{proof}

\begin{proposition}\label{prop:kummer-simultaneous-projective-reduction}
Let a number field \(K\) be given by an irreducible rational
polynomial or an explicit rational multiplication table. Given a
rational prime \(p\) and explicit coordinates of a point
\[
 [\alpha_0:\cdots:\alpha_m]\in\mathbb P^m(K),
\]
there is an unconditional deterministic algorithm, polynomial in
the input bit length and \(\log p\), that outputs structure
constants for a finite \'etale algebra
\[
 B\simeq\prod_{v\mid p}\kappa(v)
\]
and elements \(c_0,\ldots,c_m\in B\) generating its unit ideal.
Under the natural residue-field identification, the component
\([c_0(v):\cdots:c_m(v)]\) is the projective reduction of the
given point at each \(v\mid p\). The algorithm neither factors
a polynomial over \(\mathbb F_p\) nor selects a residue field.
\end{proposition}
\begin{proof}
Construct an order \(O_0\subset K\) of polynomial bit size.
For a defining polynomial
\(a_nX^n+\cdots+a_0\in\mathbb Z[X]\), one may use
\(\mathbb Z[a_n\theta]\); the scaled generator is integral.
For a rational multiplication table on a basis \(e_i\), choose
\(D\) clearing its multiplication constants and use the lattice
\(\mathbb Z1+\sum_i\mathbb ZDe_i\), which is a subring.

Compute the overorder \(O\) of \(O_0\) that is maximal at
\(p\) and has \(p\)-power index over \(O_0\). This is the
unconditional deterministic polynomial-time algorithm of
\citet[Algorithm~6.1 and Proposition~6.2, pp.~246--247]
{BuchmannLenstra1994}. Its operations are Frobenius kernels and
ideal arithmetic; it requires no finite-field factorization.
Thus
\[
 O\otimes\mathbb Z_p=O_K\otimes\mathbb Z_p.
\]

Set \(A=O/pO\). If \(p^h\geq n=[K:\mathbb Q]\), the
nilradical \(N\) of \(A\) is the kernel of the
\(\mathbb F_p\)-linear map \(x\mapsto x^{p^h}\).
Indeed a nilpotent element of an \(n\)-dimensional algebra has
nilpotence exponent at most \(n\). Let \(J\subset O\) be the
inverse image of \(N\), and compute
\[
 B=A/N=O/J.
\]
Since \(O\) is maximal at \(p\), its localizations above
\(p\) are the corresponding discrete valuation rings of
\(O_K\). Hence \(B\) is naturally the product of their residue
fields, including when \(p\) ramifies in \(K\).

Clear the rational denominators of all coordinates by a common
positive integer \(C\), so that \(\beta_i=C\alpha_i\in O\),
and form the nonzero ideal
\[
 I=\sum_{i=0}^{m}\beta_iO.
\]
Compute a lattice basis by Hermite normal form. The quotient
\[
 M=I/JI
\]
is a \(B\)-module computable by linear algebra over
\(\mathbb F_p\): since \(pI\subseteq JI\), first form
\(I/pI\), then quotient by the span of the products of
generators of \(J\) and \(I\) modulo \(pI\).
The action of \(B\) is computed from multiplication in \(K\).

At each \(v\mid p\), the ideal \(I_v\) is principal in the
discrete valuation ring \(O_v\). Therefore
\[
 M\otimes_B\kappa(v)=I_v/\mathfrak m_vI_v
\]
has dimension one over \(\kappa(v)\). Thus \(M\) has rank
one over \(B\), and the images \(x_i\) of \(\beta_i\)
generate it. Apply Lemma~\ref{lem:kummer-rank-one-patching} to
obtain \(w\) and coordinates \(x_i=c_iw\).

To verify the reduction, choose any lift \(\widetilde w\in I\)
of \(w\). Write
\[
 t_v=\min_i v(\beta_i),\qquad
 I_v=\mathfrak m_v^{t_v}.
\]
The component of \(w\) in
\(I_v/\mathfrak m_vI_v\) is nonzero, so
\(v(\widetilde w)=t_v\). The coordinate equation then gives
\[
 c_i(v)=\overline{\beta_i/\widetilde w}\in\kappa(v).
\]
At least one coordinate is nonzero in each component, and this
is exactly the projective reduction of the original point.
Different module bases multiply all coordinates by a unit in
\(B\), which preserves the projective point.

For complexity, the initial order and its discriminant have
polynomial bit size. The overorder's index squared divides the
initial discriminant, giving a polynomial bound for its lattice
encoding; the source also supplies this bound in Section~5.9,
p.~246. Ideal sums, products and bases are polynomial-time
integer lattice operations. The coordinate ideal and all its
product generators have polynomial bit length. Both \(B\) and
\(M\) have dimension at most \(n\). Frobenius powering uses an
exponent with polynomially many bits, and support patching takes
\(m+1\) rounds of finite-field linear algebra. There is no loop
of length \(p\), and no loop proportional to the numerical size
of a common local factor. The output has at most \(n^3\)
multiplication constants and \((m+1)n\) coordinate entries.
\end{proof}

\paragraph{The role of common factors.}
Reducing the original coordinates directly in \(O/J\) can erase
every coordinate at a component. The quotient \(I/JI\) instead
removes the minimum valuation at each prime simultaneously.
The proposition outputs the full product of residue fields; it
provides no rule for choosing an individual field factor.

\section{Integral reduction of a marked CM theta point}
\label{sec:kummer-cm-theta-reduction}

Fix the theta type \(\delta=(16,16,16)\).  Let
\(U_\delta\) denote Mumford's moduli scheme of abelian threefolds
with a normalized, relatively ample, totally symmetric invertible
sheaf of type \(\delta\) and a symmetric \(\delta\)-theta
structure.  Theorem and Further Theorem of
\citet[Section~6, pp.~83--84]{Mumford1967} realize
\(U_\delta\) as an open subscheme of the explicit projective
theta-null scheme over \(\mathbb Z[1/2]\), with its universal
abelian scheme.  The hypothesis that every entry of \(\delta\)
is divisible by \(8\) is satisfied.  The fixed type here avoids
any ambiguity between different conventions for a ``level-eight''
marking.

\begin{proposition}\label{prop:kummer-integral-cm-theta}
Let \(H\) be a number field and let \(u\in U_\delta(H)\).
Assume that its associated abelian variety has geometric complex
multiplication.  At every finite place \(v\) of \(H\) of odd
residue characteristic, the normalized projective reduction of
the theta-null point \(u\) belongs to \(U_\delta\).
In particular, if the coordinates of \(u\) are given explicitly,
Proposition~\ref{prop:kummer-simultaneous-projective-reduction}
computes, deterministically in polynomial bit time, a point of
\(U_\delta\) over
\[
 B=\prod_{v\mid p}\kappa(v)
\]
for every odd rational prime \(p\).  Pulling back the fixed
universal family gives an abelian scheme over \(B\), without
selecting a residue-field factor.
\end{proposition}
\begin{proof}
Write \((A,L,\alpha)\) for the marked variety.  The theta
structure identifies the polarization kernel with
\[
 H(L)\simeq(\mathbb Z/16\mathbb Z)^3\times\mu_{16}^3.
\]
This group is killed by \(16\) and has order \(16^6\), so
\(H(L)=A[16]\).  The marking therefore makes \(A[16]\)
unramified at every odd place: both the constant factor and
\(\mu_{16}\) are unramified there.  Adjoining \(\zeta_{16}\),
if desired, makes the full torsion rational by a fixed-degree
extension unramified at these places.

The CM action is defined after a finite extension of \(H\).
Theorem~6(a) of \citet[p.~503]{SerreTate1968} implies potential
good reduction, which consequently also holds over \(H_v\).
Corollary~2(b) to their Theorem~2, printed p.~497, now applies
with \(m=16\): potential good reduction and unramified
\(16\)-torsion imply good reduction.  The potential-good
hypothesis is part of this implication.  Equivalently, inertia
has finite image on \(T_2(A)\), while the principal congruence
kernel modulo \(4\) in \(\operatorname{GL}_6(\mathbb Z_2)\)
has no nonidentity finite-order element.  Thus inertia acts
trivially, and the N\'eron--Ogg--Shafarevich criterion applies.

It remains to extend the specified marking, rather than only the
underlying abelian variety.  Work over the strict henselization
\(R\) of the local valuation ring, with fraction field \(F\).
Let \(\mathcal A/R\) be the good model.  Since
\(\ker\lambda_L=A[16]\), the polarization factors uniquely as
\(\lambda_L=16\lambda_0\), with \(\lambda_0\) principal.
The homomorphism \(\lambda_0\), its inverse, and its polarization
property extend over \(R\); see
\citet[Propositions~2.11 and 2.14, author pp.~17--18]{Milne1992}.

Let \(\mathcal P\) be the rigidified Poincar\'e bundle on
\(\mathcal A\times\mathcal A^\vee\).  Set
\[
 \mathcal M=(\operatorname{id},\lambda_0)^*\mathcal P,
 \qquad \mathcal L=\mathcal M^{\otimes8}.
\]
The normalized symmetric bundle \(\mathcal M\) induces
\(2\lambda_0\).  Hence \(\mathcal L\) is normalized, totally
symmetric, relatively ample, and induces \(16\lambda_0\).
Its generic fibre is the specified \(L\).  Indeed two totally
symmetric bundles with the same polarization differ by an
element of \(\operatorname{Pic}^0[2]\) whose character on
\(A[2]\) is trivial.  Perfect \(2\)-Weil duality makes that
element zero, and rigidification makes the isomorphism unique.
This is the uniqueness observation of
\citet[Section~2, pp.~304--305 and 307]{Mumford1966}.

The theta group \(G(\mathcal L)\) is an extension of
\(\mathcal A[16]\) by \(\mathbb G_m\).  Over the strictly
henselian ring \(R\), its finite quotient and \(\mu_{16}\)
are constant.  Choose a symplectic basis of the quotient and
lift its basis elements to the theta group.  Such lifts exist
because their fibres are \(\mathbb G_m\)-torsors over the local
ring \(R\).  Their sixteenth powers are units.  Hensel's lemma
permits rescaling the lifts so that these powers equal one.
Their commutators are the prescribed pairing, and the resulting
generators define an integral theta structure.

Compare this structure with the specified generic structure
\(\alpha\).  A central-fixing automorphism of the standard
Heisenberg group induces a symplectic automorphism of its finite
quotient.  For each chosen quotient generator, its central
coefficient satisfies a sixteenth-power equation with a unit
right side.  A solution in \(F\) is an integral unit.  Thus
this generic automorphism extends over \(R\), and so does
\(\alpha\).  Its symmetry condition is an equality of
morphisms true on the generic fibre and hence over \(R\).
We have therefore obtained an \(R\)-point of \(U_\delta\).

Its image in projective space is the unique extension of the
generic theta-null point, since projective space is separated.
Its special point consequently equals normalized projective
reduction.  Membership in an open subscheme can be checked after
strict henselization, so the assertion holds at the original
place \(v\).  Applying the simultaneous projective-reduction
algorithm gives coordinates over \(B\) whose component at every
\(v\mid p\) lies in \(U_\delta\).  The resulting
\(B\)-point factors through this open, proving the final claim.
\end{proof}

\begin{remark}\label{rem:kummer-fixed-theta-denominators}
All geometric extension arguments above certify the location
of the calculated point; the algorithm does not compute local
uniformizers or local theta structures.  Evaluating a given
fixed universal presentation and a fixed complete morphism atlas
has polynomial cost.  One may also spread additional fixed
characteristic-zero coordinate formulas over
\(\mathbb Z[1/M]\), for an effectively computable fixed integer
\(M\) independent of the input.  Its finitely many prime divisors
are constant exceptional cases.  This does not establish the
degree or height bounds needed to recognize the initial algebraic
theta point, nor any Frobenius-separation or field-construction
claim.
\end{remark}

\section{A complete addition atlas from one fixed theta level}
\label{sec:theta-model-atlas}

The theta constants used for recognition can be converted to an
explicit projective abelian variety by fixed-size linear algebra.
Using level sixteen supplies both the equations and the addition law
for a level-eight embedding. In particular, this step does not require
selecting a root of a CM class polynomial or searching for an addition
law over the finite field. The complete biquadratic addition-law
mechanism is classical; compare
\citet[Proposition~2.2 and Lemma~2.4, pp.~310--312]{AreneKohelRitzenthaler2012}.
The construction below makes its fixed-size input requirements explicit.

\begin{lemma}\label{lem:theta-fixed-fourier}
Fix an even positive integer \(r\), put \(n=r^2\), and let
\(\tau\) be a point of the Siegel upper half space of degree \(g\).
For \(A,B\in(\mathbb Z/r\mathbb Z)^g\), represented by vectors with
entries in \(\{0,\ldots,r-1\}\), put
\[
 T_{A,B}(z)=\theta[A/r,B/r](\tau,rz),\qquad
 S_c(z)=\theta[c/n,0](n\tau,nz),\quad
 c\in(\mathbb Z/n\mathbb Z)^g.
\]
Then, with \(\zeta_n=e^{2\pi i/n}\),
\begin{align}
 T_{A,B}(z)&=\sum_{j\bmod r}
       \zeta_n^{(A+rj)\cdot B}S_{A+rj}(z),\label{eq:theta-fourier-forward}\\
 S_{A+rj}(z)&=r^{-g}\sum_{B\bmod r}
       \zeta_n^{-(A+rj)\cdot B}T_{A,B}(z).
       \label{eq:theta-fourier-inverse}
\end{align}
Thus the \(r^{2g}\) recognized theta constants at zero determine the
standard level-\(n\) theta-null point by a fixed invertible linear
transformation over \(\mathbb Q(\zeta_n)\). For \(g=3,r=4\), this has
fixed degree and height cost and introduces only powers of two in its
denominators.
\end{lemma}
\begin{proof}
In the defining, absolutely convergent theta series write its summation
index as \(k=rt+j\), with \(j\bmod r\). Then
\(k+A/r=r(t+(A+rj)/r^2)\). The factor involving \(B\) is
\(\exp(2\pi i(k+A/r)\cdot B/r)
 =\zeta_n^{(A+rj)\cdot B}\).
The remaining terms are precisely the series defining \(S_{A+rj}\).
The inverse follows from orthogonality of the characters of
\((\mathbb Z/r\mathbb Z)^g\). This identity also identifies the
line-bundle sections: all the displayed functions of \(z\) have the
common automorphy factor of the \(n\)-th power of the principal theta
line bundle. In the \(S_c\) basis its standard theta-group action is
translation and multiplication by characters. Consequently this is
Mumford's type-\((n,\ldots,n)\) marking, rather than merely a projective
coordinate change with an unspecified moduli interpretation.
\end{proof}

\begin{proposition}\label{prop:fixed-theta-complete-atlas}
Fix \(g=3\) and \(n=8\). Let \(k\) be a field of characteristic
different from two. Suppose a projective vector
\[
 v=(v_a)_{a\in(\mathbb Z/2n\mathbb Z)^g}
\]
is a valid point of Mumford's theta-null moduli problem of type
\((2n,\ldots,2n)\), including a totally symmetric line bundle and
a symmetric marking.
There is a deterministic construction, using a fixed number of field
arithmetic operations and equality tests, of the following data for the
same underlying abelian variety in its compatible level-\(n\)
embedding in \(\mathbb P^{n^g-1}\): its identity, its inverse, a set of
quadratic equations defining it scheme-theoretically in projective
space, and a complete set of \(n^g\) addition laws of bidegree
\((2,2)\).

The construction also works over a finite \'etale \(\mathbb F_p\)-algebra
of dimension \(N\), given by structure constants, with a projective
vector \(v\) whose coordinates generate the unit ideal and with the
same validity promise on every geometric component. It returns the
data on a partition into direct factors in deterministic time
polynomial in \(N,\log p\) and the input length, without finding an
irreducible factor or selecting a component field.
\end{proposition}
\begin{proof}
Write \(G=(\mathbb Z/n\mathbb Z)^g\),
\(H=(\mathbb Z/2n\mathbb Z)^g\), and
\(E=\{0,1\}^g\). Use integer representatives for \(i,j\in G\)
in formulas with indices in \(H\); changing a representative only
permutes the summation over \(E\).
In compatible theta bases let \(s_i\) be the level-\(n\) sections
and \(t_a\) the level-\(2n\) sections. Mumford's fundamental
addition formula and multiplication formula give, up to a common
nonzero scalar arising from choices of trivialization,
\begin{align}
 s_i(x+y)s_j(x-y)
   &=\sum_{\eta\in E}
       t_{i+j+n\eta}(x)t_{i-j+n\eta}(y),
       \label{eq:theta-product-addition}\\
 s_i(z)s_j(z)
   &=\sum_{\eta\in E}
       v_{i-j+n\eta}t_{i+j+n\eta}(z).
       \label{eq:theta-multiplication-matrix}
\end{align}
The scalars do not affect any of the kernels or projective maps below.
For the analytic standard bases these formulas follow directly by
replacing the two theta summation variables by their sum and difference
and sorting their common parity. Their algebraic validity is
\citet[\S3, fundamental addition formula, p.~324;
multiplication formula, p.~330]{Mumford1966}; the proof extends to
abelian schemes and to residue characteristics different from two,
as checked in \citet[\S6, Step~III, pp.~92--94]{Mumford1967}.

First form the symmetric matrix
\[
 P_{ij}=\sum_{\eta\in E}
     v_{i+j+n\eta}v_{i-j+n\eta}.
\]
It is a nonzero scalar multiple of \((e_i e_j)_{i,j\in G}\), where
\((e_i)\) are the compatible level-\(n\) coordinates of zero.
Thus it has rank one and is nonzero. Choose \(j\) with
\(P_{jj}\ne0\); its \(j\)-th column gives the identity \([e_i]\).
This is also the explicit level-forgetting morphism of
\citet[\S6, Step~II, p.~85]{Mumford1967}. No square root is needed.
The inverse is \([X_i]\mapsto[X_{-i}]\).

Next use \eqref{eq:theta-multiplication-matrix} to form the matrix
\[
 M:\operatorname{Sym}^2(k^G)\longrightarrow k^H,
 \qquad
 X_iX_j\longmapsto
 \sum_{\eta\in E}v_{i-j+n\eta}[i+j+n\eta].
\]
Here \([a]\) denotes the standard basis vector indexed by \(a\in H\),
and one uses one monomial for each unordered pair \(\{i,j\}\).
Symmetry is consistent because \(v_a=v_{-a}\).
This matrix is surjective: the polarization kernel contains the full
four-torsion, so \citet[\S4, Theorem~1, pp.~338--340]{Mumford1966}
applies to the multiplication of sections.
Gaussian elimination gives a basis of \(\ker M\), hence all the
quadratic relations on the \(s_i\). These relations cut out the abelian
variety scheme-theoretically in projective space. This is the
projective consequence of
\citet[\S4, the graded-ring lemmas on pp.~336--338 and
Theorem~2, pp.~341--348]{Mumford1966}, used explicitly in
\citet[\S6, Step~III, p.~94]{Mumford1967}; equivalently they span the
quadratic equations of the universal family in the Further Theorem
on p.~84 of the latter paper. Only the saturated projective ideal is
asserted here: Theorem~2 itself states a relation in degree four,
whereas the preceding graded-ring argument supplies equality after
saturation. On each standard affine chart the
resulting ideal is radical, since this scheme is an abelian variety.

Choose by the same elimination a linear right inverse of \(M\).
Its columns define quadratic polynomials \(Q_a(X)\), for \(a\in H\),
representing the sections \(t_a\), up to the same common nonzero
normalization. The map \([X_i]\mapsto[Q_a(X)]\) identifies this
level-eight model with the given level-sixteen model; hence the
construction preserves its form over the ground field, and not only
its geometric isomorphism class. Form the \(n^g\) by \(n^g\) matrix of biquadratic
polynomials
\[
 F_{ij}(X,Y)=\sum_{\eta\in E}
    Q_{i+j+n\eta}(X)Q_{i-j+n\eta}(Y).
    \label{eq:theta-biquadratic-atlas}
\]
At every geometric pair \((x,y)\) on the variety this matrix is a
common nonzero scalar multiple of the outer product
\((s_i(x+y))_i(s_j(x-y))_j^{\mathsf T}\). Both factors are nonzero,
because the complete level-\(n\) linear system gives an embedding.
Consequently at least one column is nonzero, and every nonzero column
represents \(x+y\). The \(n^g\) columns therefore give a complete
addition atlas, including over the algebraic closure; this is stronger
than completeness merely on \(k\)-rational pairs.

All matrices have sizes depending only on \(g,n\): here
\(|G|=512\), \(|H|=4096\). The kernel, right inverse, and a nonzero
diagonal entry of \(P\) require only arithmetic and zero tests.
Over a finite \'etale algebra, perform these operations with the
structure-constant simulation of
Proposition~\ref{prop:structure-etale-simulation}. Split when a pivot
is zero on some components and nonzero on others, and invert it only
on the latter factors. Surjectivity and the rank-one identity hold on
every geometric component, so all branches terminate with the stated
data. The total number of leaves is at most \(N\), and all fixed-size
arithmetic has polynomial bit cost. No search through the elements of
a component field occurs.
\end{proof}

\begin{remark}\label{rem:theta-atlas-reduction-scope}
The input promise in Proposition~\ref{prop:fixed-theta-complete-atlas}
is essential: satisfying a subset of theta relations does not establish
that an arbitrary projective vector parametrizes an abelian variety.
For a CM theta-null point, one must separately prove that its
projective reduction is the point of a good-reduction abelian scheme
with its marking. Once that statement has been established, the
proposition supplies all the geometric input required by
Corollary~\ref{cor:etale-pila}, with ambient dimension, equation degrees,
and addition-law degrees fixed independently of the auxiliary
prime and the characteristic. There is no growing-degree CM
endomorphism among these inputs.
\end{remark}

\section{Extracting a ninth root from sextic CM Frobenius}
\label{sec:sextic-cm-frobenius-output}

\begin{proposition}\label{prop:sextic-cm-output}
Let \(p>3\) be prime with \(p\equiv-1\pmod9\), and let
\(q\notin\{2,3,p\}\) be prime with
\(q\equiv1\pmod4\) and \((p/q)=-1\).
Put \(R=\mathbb Q(\theta)\), \(\theta^3-3\theta+1=0\),
and \(L=R(\sqrt{-3q})\).
Suppose a principally polarized abelian threefold of the primitive
CM type in Proposition~\ref{prop:explicit-sextic-cm-periods}
has been reduced with good reduction to a field \(\mathbb F_{p^e}\),
and that its Frobenius characteristic polynomial is given.
Then one can deterministically find a root of
\(X^3-3X+1\) in \(\mathbb F_p\), and an element of order \(9\)
in any presented quadratic extension of \(\mathbb F_p\), in time
polynomial in \(e,\log p,\log q\) and the given polynomial's length.
This implication uses no unproved analytic hypothesis.
If \(v_3(p+1)=2\), that element generates the entire
\(3\)-primary subgroup of \(\mathbb F_{p^2}^\times\).
\end{proposition}
\begin{proof}
The prime \(p\) splits completely in \(R\), since its class modulo
\(9\) is \(-1\). It also splits in \(\mathbb Q(\sqrt{-3q})\):
\((-3/p)=-1\), while quadratic reciprocity and \(q\equiv1\pmod4\)
give \((q/p)=(p/q)=-1\). Thus it splits completely in \(L\).

Let \(F\) denote the \(p^e\)-power Frobenius of the reduction.
The characteristic-zero endomorphism algebra is \(L\), since the
CM type is primitive. Galois acts on it through
\(\operatorname{Aut}(L/\mathbb Q)\), of order \(6\). Compatibility
of good reduction with endomorphisms implies that \(F^6\) centralizes
the embedded \(L\). Its centralizer in the reduction's rational
endomorphism algebra is \(L\): on an auxiliary prime-adic Tate
module the embedded field has degree \(6\), acts faithfully on a
space of dimension \(6\), and has rank-one regular representation.
Consequently \(F^6\) is an algebraic integer \(\pi\in\mathcal O_L\).

The Shimura--Taniyama formula, in
\citet[Theorem~8.1 and Corollaries~8.2--8.3, pp.~66--68]{MilneCM2020},
gives
\[
 (\pi)=\prod_{\mathfrak p\in\Phi_p}\mathfrak p^{6e},
\]
where \(\Phi_p\) is a conjugate of the primitive CM type, viewed as
three of the six primes over \(p\), with one prime from each
complex-conjugate pair. To apply the stated maximal-order version,
one may first make the degree-\(3\) lattice isogeny to full CM and
extend the local field until the isogeny and its endomorphisms are
defined. Good reduction is preserved, and the formula for a power of
\(\pi\) gives the displayed valuations after division by that power.

This support has trivial stabilizer in the cyclic group of order
\(6\). A subgroup of order \(2\) would be complex conjugation and
cannot preserve a CM type. Invariance under the subgroup of order
\(3\) would make the type induced from the quadratic subfield,
contrary to primitivity. Hence \(\mathbb Q(\pi)=L\).
The characteristic polynomial of \(F^6\), computed by a resultant
from that of \(F\), is therefore the degree-six minimal polynomial
of \(\pi\). Factoring it over the explicit degree-six field \(L\)
produces its six roots in deterministic polynomial time
\citep[Theorem~4.5, pp.~253--254]{Lenstra1983}.
Any one of those roots has the same support property.

Because \([\mathcal O_L:\mathbb Z[\theta,(1+\sqrt{-3q})/2]]=3\),
express the chosen root in this product basis with denominators
dividing \(3\). These denominators and \(2\) are invertible modulo
\(p\). The known nonsquare \(-3\) permits deterministic square-root
extraction in \(\mathbb F_p\); apply it to the square \(-3q\) to
obtain \(t\). Reduce the chosen expression for \(\pi\) after
substituting \(\sqrt{-3q}=t\), obtaining a polynomial \(g(\theta)\)
of degree at most \(2\) modulo \(p\).

Among the three embeddings restricting to either chosen sign of
the quadratic subfield, a primitive CM type contains exactly one
or two. Thus
\[
 1\leq\deg\gcd(g(X),X^3-3X+1)\leq2.
\]
A linear factor, either this greatest common divisor or its
complement, yields a root \(u\in\mathbb F_p\).

Finally solve \(Z^2-uZ+1=0\) in the presented quadratic extension.
Its discriminant is a nonsquare in \(\mathbb F_p\), since otherwise
that field would contain a ninth root of unity. More explicitly,
any root \(z\) satisfies \(z^6+z^3+1=0\), hence has order \(9\)
when \(p\ne3\); and \(9\nmid p-1\).
Quadratic extraction here is deterministic: choose a nonzero
trace-zero element \(v\) in the presented extension by linear
algebra. Then \(v^2\) is a nonsquare in \(\mathbb F_p\), so
\((u^2-4)/v^2\) is a square there. Extract its square root using
the same known nonsquare \(-3\), multiply by \(v\), and use the
quadratic formula. The polynomial and field arithmetic and all
square-root computations have the claimed polynomial cost.
If \(v_3(p+1)=2\), then \(v_3(p^2-1)=2\), proving the last claim.
\end{proof}

The input promise of good CM reduction is essential to this
proposition. In the proposed construction it must follow from
an integral model with a complete addition law. For
\(v_3(p+1)>2\), this procedure gives order \(9\), not the maximal
\(3\)-power order; that lifting problem is also a separate step.

\section{A fixed-dimensional CM construction of ninth roots}
\label{sec:composed-cubic-cm}

The following construction combines classical CM height theorems,
algebraic-number recognition, Mumford's theta equations, and Pila's
Frobenius algorithm. The use of an auxiliary quadratic conductor
allows a primitive sextic CM type, while computation over the full
finite residue algebra avoids selecting a root of a CM invariant
polynomial. The constituent geometric and computational tools are
classical; Theorem~\ref{thm:dirichlet-ninth-root-cubic} is the specific composition proved
here. Pila's own root-of-unity
Theorem~D treats a fixed odd prime order and primes congruent to
one modulo that order \citep[Theorem~D, author p.~2; Section~5, pp.~14--18]{Pila1990}; it is not the
ninth-root statement below.

\subsection*{Proof of Theorem~\ref{thm:dirichlet-ninth-root-cubic}}
\begin{proof}
Put \(\lambda=2+\lceil\log_2p\rceil\). The congruence on \(p\)
implies \(p>3\), \(\operatorname{ord}_3(p)=2\), and
\(({-3}/p)=-1\).

\paragraph{Auxiliary prime and exact theta point.}
Apply Lemma~\ref{lem:quadratic-cm-auxiliary} with \(r=3\) and
\(d=2\). By testing primes in increasing order, find
\[
 q\equiv1\pmod4,\qquad (p/q)=-1,\qquad
 q\notin\{2,3,p\},\qquad q\le C\lambda^2.
\]
For completeness, this specialization needs only the four
Dirichlet characters obtained by multiplying the indicator of
\(1\pmod4\) by the quadratic nonresidue indicator for \(p\).
Their smoothed prime count has main term \(x/4\), and its error
under \(\mathrm A\) is \(O(\sqrt x\log(2p))\). Even powers
have zero nonresidue indicator; the remaining higher prime powers
contribute \(O(x^{1/3}(\log x)^2)\). Thus a sufficiently large
absolute multiple of \(\lambda^2\) contains a suitable prime.
This search can be implemented in \(O(\lambda^6)\) bit operations
using trial division for its small candidates and integer
arithmetic; the bound is intentionally coarse.

Let
\[
 R=\mathbb Q(\theta),\quad \theta^3-3\theta+1=0,
 \qquad L=R(\sqrt{-3q}).
\]
Proposition~\ref{prop:explicit-sextic-cm-periods} constructs the
period matrix of a principally polarized simple abelian threefold
with primitive CM type and multiplication by
\(\mathbb Z[\theta,(1+\sqrt{-3q})/2]\). Its lattice has index
\(3\) in \(\mathcal O_L\), and
\(|\operatorname{disc}(L)|=3^9q^3\).
Apply Proposition~\ref{prop:sextic-theta-recognition} with
\(r=4\), obtaining all \(4096\) theta ratios exactly in a
number field. This step has deterministic polynomial cost in
the numerical \(q\), rather than in an unspecified CM
preprocessing oracle.

Apply Lemma~\ref{lem:theta-fixed-fourier}, adjoining
\(\zeta_{16}\) by characteristic-zero algebraic computation.
It transforms the recovered point into the standard
level-sixteen theta-null point \(v\), defined over an explicit
number field \(K\). It is a point of Mumford's moduli problem
of type \((16,16,16)\), including its totally symmetric line
bundle and symmetric marking. All number-field operations here
are in characteristic zero.

\paragraph{Reduction without choosing a residue field.}
Apply Proposition~\ref{prop:kummer-simultaneous-projective-reduction}
to the homogeneous coordinates of \(v\) at \(p\). It computes
structure constants for
\[
 \mathcal B=\prod_{\mathfrak p\mid p}\kappa(\mathfrak p)
\]
and a coordinate vector generating its unit ideal. This
normalization removes the separate common coordinate valuations
at all primes above \(p\); reducing one arbitrary affine chart
would not suffice. Proposition~\ref{prop:kummer-integral-cm-theta}
shows that every component is still a valid marked abelian
threefold. Its proof uses potential good reduction of CM
varieties and unramified full sixteen-torsion, followed by
extension of the specified marking.

Proposition~\ref{prop:fixed-theta-complete-atlas} now constructs,
over a partition of \(\mathcal B\), the same abelian threefold
in a level-eight embedding in \(\mathbb P^{511}\). It gives
quadratic defining equations with radical affine-chart ideals,
the identity, the inverse, and \(512\) complete biquadratic
addition laws. The construction uses multiplication-matrix
kernels and right inverses; its matrix dimensions are absolute
constants. In particular, the equation degrees and the addition
atlas do not grow with \(q\).

\paragraph{Frobenius and prime separation.}
Separate \(\mathcal B\) into blocks whose field components have
a common residue degree \(e\). This uses the Frobenius-linear
construction in Proposition~\ref{prop:structure-etale-simulation},
not irreducible factor selection. Run the deterministic Pila
algorithm with field size \(p^e\), using
Corollary~\ref{cor:etale-pila}. At every equality test split a
block only as required by the tested element. At termination
there is a nonempty list of blocks, each with an integer
Frobenius characteristic polynomial valid on all its components.
Choose any polynomial on this list. Choosing this integer
output does not choose a residue-field embedding.

All such components are reductions of conjugates of the same
primitive CM construction. The field \(L\) is Galois of degree
six, and \(p\) splits completely in it: it splits in \(R\)
because \(p\equiv-1\pmod9\), and
\[
 (-3q/p)=(-3/p)(q/p)=(-1)(-1)=1
\]
by quadratic reciprocity and \(q\equiv1\pmod4\).
Consequently Proposition~\ref{prop:sextic-cm-output} applies to
any selected integer Frobenius polynomial. Briefly, the sixth
power of Frobenius is a primitive element \(\pi\in L\), whose
ideal is supported at the three primes in a primitive CM type.
Compute its degree-six polynomial and factor it over the
explicit characteristic-zero field \(L\). Express a root in
the product order with denominators dividing \(3\), extract
\(\sqrt{-3q}\) modulo \(p\) using the known nonsquare \(-3\),
and substitute either quadratic sign. The resulting greatest
common divisor with \(f\) has degree one or two. A linear
factor, or the complementary linear factor, yields \(u\).

The same proposition solves \(Z^2-uZ+1\) in the presented
quadratic field and yields an element \(z\) of exact order nine.
The canonical quadratic presentation is valid because
\(({-3}/p)=-1\). For another presentation, a nonzero trace-zero
vector is found by linear algebra, and its square supplies the
quadratic coordinate conversion using only square roots in
\(\mathbb F_p\) with a known nonsquare. This completes the
first two outputs.

\paragraph{The cubic extension when the valuation is two.}
Now suppose \(v_3(p+1)=2\). Then \(v_3(p^2-1)=2\), so an
element of exact order nine in \(\mathbb F_{p^2}\) is not a
cube. We have \(z+z^{-1}=u\). If \(H\) had a root
\(\beta\in\mathbb F_p\), choose a root \(w\) of
\(W^2-\beta W+1\) in \(\mathbb F_{p^2}\). Then
\[
 w^3+w^{-3}=\beta^3-3\beta=u,
\]
so \(w^3\) is one of the roots \(z,z^{-1}\) of
\(Z^2-uZ+1\). This contradicts their being noncubes.
Hence \(H\) has no root in \(\mathbb F_p\), and a cubic with
no root is irreducible. Equivalently, a cube root of \(z\) has
order twenty-seven and cannot lie in \(\mathbb F_{p^2}\).
The displayed polynomial is the promised explicit presentation.

\paragraph{Uniform bit bound.}
Here are parameter bounds sufficient to make the composition's
bit complexity explicit. Put
\[
 B_0=C_0q^{3/2}(1+\log q)^5,
 \qquad N=8\lceil B_0\rceil^{4096}.
\]
The factor eight allows \(\mathbb Q(\zeta_{16})\). The number
field degree and the dimension of \(\mathcal B\) are at most
\(N\). Let \(S\) be the total bit length of its
characteristic-zero presentation and theta coordinates. The
recognition proof, including its fixed number of composita,
gives
\[
 S\le C(B_0\log(2q)+2)^\kappa
\]
for an absolute effectively computable \(\kappa\ge4096\).
The degree and coefficient bounds used by recognition are
specified in Proposition~\ref{prop:sextic-theta-recognition};
in particular its requested precision is polynomial in
\(B_0\) and \(B_0\log(2q)\).

Let \(c_0\ge8\) be a fixed exponent bounding the deterministic
recognition, characteristic-zero, lattice, projective-reduction, and final
fixed-degree operations by \(C(S+\lambda+2)^{c_0}\).
Let \(\Delta_0\ge2\) be a fixed exponent for Pila's algorithm
on the model just constructed, counting its coefficient-field
operations and its non-field bit operations. These are effective
absolute constants supplied by the cited algorithms; neither
is an exponent depending on the input. The Pila parameters may
be taken as ambient dimension \(511\), at most
\(512\cdot513/2\) quadratic equations, at most \(512\)
addition charts, total addition degree \(4\), and dimension
\(3\). The explicit parameter-dependent bound for its exponent
is in \citet[Section~4, author pp.~12--14]{Pila1990}.

Arithmetic, inversion, support splitting, changes of basis, and
reconstruction of structure constants in an algebra of dimension at
most \(N\) have the conservative bound \(O(N^6\lambda^3)\)
per simulated step with school arithmetic and Gaussian elimination. There are at most
\(N\) branches, and each field computation has
\(e\le N\). After increasing \(c_0\) for elementary
preprocessing, a bit bound for the entire construction is
\[
 C\left\{\lambda^6+(S+\lambda+2)^{c_0}
             +N^7\lambda^3(N\lambda+2)^{\Delta_0}\right\}.
 \tag{\(\ast\)}\label{eq:composed-cubic-bit-bound}
\]
Since \(q\le C\lambda^2\), we may bound
\(B_0\le C\lambda^8\),
\(N\le C\lambda^{32768}\), and
\(S\le C\lambda^{9\kappa}\).
For example the absolute exponent
\[
 C_* = 10+\max\{6,\,9\kappa c_0,\,
                     229379+32769\Delta_0\}
\]
suffices in the theorem. Formula \((\ast)\) records the actual
source of the bound; it is not a claim that these deliberately
large exponents are practical.

Finally, the two conditional uses of \(\mathrm A\) were the
quadratic auxiliary-prime estimate and the effective height
bound for an abelian CM field. In the latter, the odd Artin
characters in the Colmez expression are Dirichlet characters over
\(\mathbb Q\), because \(L/\mathbb Q\) is abelian. The algebraic
recognition, marked reduction, complete atlas, Frobenius
computation, and extraction from \(\pi\) are unconditional
once their stated inputs have been obtained. Proposition~\ref{prop:unconditional-cm-recognition}
allows the height step to use an unconditional polynomial bound instead;
then only the auxiliary-prime step needs \(\mathrm A\).
\end{proof}

\paragraph{Scope of the output.}
For \(v_3(p+1)>2\), the construction still produces an element
of exact order nine and a root of \(X^3-3X+1\). It does not
produce an element of maximal three-power order in
\(\mathbb F_{p^2}^{\times}\), and the displayed cubic is then
not an irreducible extension polynomial. The construction has
fixed dimension three; it gives no polynomial-in-degree
extension of the method to unbounded cyclotomic conductor.


\begin{example}\label{ex:sextic-cm-seventeen}
At \(p=17\), the auxiliary quadratic prime can be \(q=5\):
\(5\equiv1\pmod4\) and \((17/5)=-1\).
The roots of \(X^3-3X+1\) modulo \(17\) are \(7,13,14\).
To display the final conversion, take \(u=7\) and write
\(s^2=-3\) in \(\mathbb F_{17^2}\). Since
\(u^2-4=11=2s^2\) and \(6^2=2\pmod{17}\),
\[
 z=\frac{7+6s}{2}=12+3s,\qquad
 z^3=8+8s,\quad z^6=8+9s,\quad z^9=1.
\]
Thus \(z\) has exact order nine. The irreducible cubic in
Theorem~\ref{thm:dirichlet-ninth-root-cubic} is
\(Y^3-3Y-7\). Direct substitution of the seventeen field elements
also verifies that it has no root. These calculations illustrate the
exact final conversion; they are not an execution of the theta
recognition and Frobenius algorithm.
\end{example}

\section{Residue maps for a fixed totally real Galois field}
\label{sec:general-real-residue}
The same geometric interfaces construct a residue map for any fixed totally real Galois field at its completely split primes.


\begin{proposition}\label{prop:fixed-real-residue}
Fix a totally real Galois number field $K/\mathbb Q$ of degree $m\geq3$,
with an integral basis and its multiplication table.  Assume the Riemann
hypothesis for all Dirichlet $L$-functions over $\mathbb Q$.  There is a
deterministic algorithm which, on a prime $p$ splitting completely in $K$,
constructs a ring homomorphism
\[
 \lambda:\mathcal O_K\longrightarrow\mathbb F_p
\]
in time polynomial in $\log p$.  The polynomial and all fixed precomputation
may depend on $K$.  The Galois group of $K/\mathbb Q$ may be nonabelian and
may have even order.  No hypothesis about Hecke, Dedekind, or nonabelian
Artin $L$-functions is required.
\end{proposition}

\begin{proof}
We give the geometric and arithmetic interfaces explicitly.  In particular,
the construction of a model includes a complete addition-law atlas, and the
reduction step does not select a prime by factoring a polynomial over a finite
field.  Finitely many primes depending only on the fixed data may be handled
by exhaustive arithmetic; this has constant cost in the present complexity
statement.

\paragraph{A small imaginary quadratic field.}
Write $D_K=\operatorname{disc}(K)>0$ and let $R_K$ be the product of its
distinct prime divisors.  Suppose first that $p$ is odd and outside the fixed
exceptional set.  On $(\mathbb Z/Q\mathbb Z)^*$, where $Q=4pR_K$, the
quadratic characters $(-1/\cdot)$ and $(p/\cdot)$ are independent.  The latter
has conductor dividing $4p$, with nontrivial $p$-part; the former detects the
modulo-four factor.  Consequently
\[
 q\equiv3\pmod4,\qquad (p/q)=1,\qquad q\nmid pD_K
\]
specifies a coset of an index-four subgroup.  The least-prime-in-a-coset
theorem of Lamzouri--Li--Soundararajan
\cite[Theorem 1.4, p.~3]{LLS2015} gives, under the stated Dirichlet hypothesis,
\[
 q=O_K((\log p)^2).
\]
The theorem's fixed alternative bound and its finitely many small moduli are
absorbed into the constant.  Search and character tests up to this bound are
deterministic polynomial time.  Set
\[
 L=\mathbb Q(\sqrt{-q}),\qquad E=KL,
 \qquad G=\operatorname{Gal}(K/\mathbb Q),
 \qquad H=G\times\langle c\rangle,
\]
where $c$ is complex conjugation.  The discriminant of $L$ is $-q$, so it is
coprime to $D_K$, including when $D_K$ is even.  Reciprocity gives
$(-q/p)=(p/q)=1$.  Thus $p$ splits completely in $E$, and
\[
 \mathcal O_E=\mathcal O_K\mathcal O_L,
 \qquad |\operatorname{disc}(E)|=D_K^2q^m.
\]
The integral-basis assertion follows either from the coprime-discriminant
criterion or locally from normality of the tensor order when at least one
factor is etale.  Equation (27) is the sole conditional step below.

\paragraph{A primitive type with two trivial stabilizers.}
Identify the embeddings of $E$ with $H$ and choose
\[
 \Phi=\{(g,1):g\ne1\}\cup\{(1,c)\}.
\]
This type has one minus sign over the $m$ real embeddings of $K$.
Both its left and its right stabilizer in $H$ are trivial.  A translation
without $c$ must fix the unique location of the minus sign, hence is the
identity.  A translation with $c$ exchanges the numbers $1$ and $m-1$ of
minus and plus signs, so cannot preserve the type when $m\geq3$.
An induced type from a proper CM subfield would be invariant under right
multiplication by the nontrivial subgroup fixing that subfield.  Thus
$\Phi$ is primitive.  The CM classification
\cite[Proposition 2.1, p.~54]{Milne1999} implies that the associated abelian
variety is simple with rational endomorphism algebra $E$.  This classification
does not assume that $H$ is abelian.

\paragraph{Principalization by a fixed-index lattice enlargement.}
Let $b_1,\ldots,b_m$ be an integral basis of $K$, let
$F=(\rho_i(b_j))$ be the real embedding matrix, and let
$S=\operatorname{diag}(s_i)$ record the signs in (29).  Weak approximation,
followed by clearing denominators, gives a fixed $\alpha\in\mathcal O_K$
with $\operatorname{sign}(\rho_i(\alpha))=-s_i$.  No unit-signature assertion
is needed.  Define
\[
 M_{ij}=-\operatorname{Tr}_{K/\mathbb Q}(\alpha b_i b_j),
 \quad h=|\det M|,
 \quad V=\operatorname{diag}(-s_i\rho_i(\alpha)).
\]
The matrix $M$ is symmetric, integral, nonsingular, and fixed with $K$;
$V$ is positive and $M=F^tSVF$.  Write
$\omega=(1+\sqrt{-q})/2$ and $\beta=\sqrt q/2$.
The positive integral alternating form
\[
 \mathcal E(x,y)=\operatorname{Tr}_{E/\mathbb Q}
 \left(\frac{\alpha}{\sqrt{-q}}x\overline y\right)
\]
on $\Phi(\mathcal O_E)$ has block matrix
$\left(\begin{smallmatrix}0&M\\-M^t&0\end{smallmatrix}\right)$.
Enlarge its second half to the lattice
\[
 \Lambda'=F\mathbb Z^m+
 \left(\tfrac12 I+i\beta S\right)FM^{-1}\mathbb Z^m.
\]
The index $[\Lambda':\Phi(\mathcal O_E)]$ is $h$, and the form becomes
unimodular on the displayed basis.  Consequently
$A'=\mathbb C^m/\Lambda'$ is principally polarized, and the natural isogeny
from $A=\mathbb C^m/\Phi(\mathcal O_E)$ has fixed degree $h$.  Since
\[
 \Phi(\mathcal O_E)\subseteq\Lambda'
 \subseteq h^{-1}\Phi(\mathcal O_E),
\]
the endomorphism order contains $\mathbb Z+h\mathcal O_E$, and
$\operatorname{End}^0(A')=E$.  A normalized period matrix is
\[
 \tau=\tfrac12 M^{-1}+i\beta C,
 \qquad C=F^{-1}V^{-1}F^{-t}>0.
\]
Indeed $F^{-1}SF M^{-1}=C$.  Its real part and positive imaginary direction
are fixed; only the scalar $\beta$ varies.

\paragraph{Polynomial degree and effective unconditional height.}
A lattice stable under $\mathbb Z+h\mathcal O_E$, with maximal-order hull
$J$, lies between $hJ$ and $J$.  The number of such intermediate lattices
per ideal class is bounded by a constant depending on $h,m$.  There are
$h_E$ maximal-order ideal classes.  For a fixed simple CM torus, principal
polarizations are classified by a quotient of totally positive units of its
real endomorphism order by norms of units.  Those norms contain squares of
real units, so at most $2^m$ classes occur.  Types, symmetric line-bundle
choices, and a fixed theta marking contribute further constants.  Hence the
degree of the marked theta point is at most $C_Kh_E=q^{O_K(1)}$.
Minkowski's bound and elementary ideal counting in fixed degree give the
last estimate; a class group is not computed.

The unconditional CM-height estimate of von Kanel--Kret applies to every
CM field and every CM type, including the present nonabelian Galois field
\cite[Section 7.2.1, (7.7), Lemma 7.10, pp.~62--63]{VonKanelKret2023}.
Taking their $\epsilon=1/2$ and
$c=|\operatorname{disc}(E)|/D_K=D_Kq^m$ gives
\[
 h_F(A)\leq C_K(1+q^m).
\]
The class-number ratio in the displayed bound is bounded by the fixed
$h_K$, and the number of roots of unity is bounded in fixed degree.
The isogeny-height inequality
\cite[(7.11), p.~67]{VonKanelKret2023}, also recorded in
\cite[Remark 2.3, p.~27]{Pazuki2012}, yields
$|h_F(A')-h_F(A)|\leq\tfrac12\log h$.
Pazuki's comparison \cite[Corollary 1.3, pp.~21--22]{Pazuki2012} therefore
bounds fixed-level theta height by $q^{O_K(1)}$.
The coarse effective bound (32) uses the proved averaged Colmez formula
and unconditional estimates; no individual nonabelian Colmez conjecture
or conditional Artin estimate is needed.

\paragraph{Exact theta coordinates and fixed geometric data.}
Use the fixed level $N=16$.  Mumford's theta moduli scheme of type
$(N,\ldots,N)$ and its embedded universal abelian scheme are fixed
finite-type objects over the ring obtained by inverting two and adjoining
the fixed necessary roots of unity
\cite[Section 6, pp.~76--84]{Mumford1967}.  Include a finite affine cover,
projective scheme equations, and the identity section in the fixed data.
Include a complete finite addition atlas as well.  To see this last
assertion, on each fixed affine base chart let
$\mu:\mathcal A\times\mathcal A\to\mathcal A$ be addition.
For a sufficiently large fixed $t$, the line bundle
$\mathcal O(t,t)\otimes\mu^*\mathcal O(1)^{-1}$ is relatively generated by
finitely many sections.  Multiplying each by the pulled-back output
coordinates gives a tuple representing $\mu$ wherever that section is
nonzero.  Increasing $t$ if necessary makes these sections restrictions of
ambient bihomogeneous polynomials.  Their nonvanishing domains cover every
geometric point.  Fixed symbolic projective algebra computes such data;
any finite characteristics lost on spreading out are fixed exceptions.
The equations of each smooth fiber have radical affine dehomogenizations.
This is the bounded presentation, including complete group-operation data,
required by Pila \cite[Section 3, pp.~752--753]{Pila1990}.

For clarity, the analytic coordinate convention is matched exactly.
For $a,b,j\in(\mathbb Z/4\mathbb Z)^m$, put
\[
 P_{a,b}=\theta[\,a/4\,;\,b/4\,](\tau,0),\qquad
 Q_k=\theta[\,k/16\,;\,0\,](16\tau,0).
\]
Splitting the theta-series index into classes modulo four gives
\[
 P_{a,b}=\zeta_{16}^{a\cdot b}
 \sum_j\zeta_4^{j\cdot b}Q_{4j+a},\qquad
 Q_{4j+a}=4^{-m}\sum_b
 \zeta_4^{-j\cdot b}\zeta_{16}^{-a\cdot b}P_{a,b}.
\]
This fixed invertible Fourier transform relates the theta-null conventions
in \cite[Section 2.3.4, p.~30]{Pazuki2012} and is integral away from two.
It changes heights by a fixed additive constant.

The Gaussian tails at (31) allow certified evaluation to $P$ bits in time
polynomial in $P,q$, since $m,C$ are fixed.  The zero-characteristic
coordinate at $16\tau$ tends to one as $q$ increases, so it supplies a
uniformly bounded-away-from-zero normalization for all sufficiently large
$q$.  The finitely many smaller cases are fixed separately.  Polynomial
degree and height bounds give polynomial coefficient-bit bounds for the
normalized algebraic coordinates.  Algebraic-number recognition
\cite[Theorem 1.19, p.~241]{KannanLenstraLovasz1988}, followed by deterministic factoring over
characteristic-zero number fields, constructs a coefficient field $B$ and
the exact theta point with degrees and encoding lengths $q^{O_K(1)}$.
There are only $16^m$ coordinates, a fixed number.  Adjoin $E$ and
$\zeta_{16}$ at bounded additional degree.

\paragraph{Defined endomorphisms and extension of the specified marking.}
The marking makes $A'[16]$ rational over $B$.  The finite Galois action on
the lattice $\operatorname{End}(A'_{\overline B})$ is congruent to the identity
modulo $16$: an endomorphism vanishing on $A'[16]$ factors through $[16]$.
A finite subgroup of a principal integral congruence group of level at
least three is trivial.  Thus every geometric endomorphism is defined over
$B$, with no assumption about the commutativity of $H$.
CM gives potential good reduction \citep[Theorem~6(a), p.~503]{SerreTate1968}.  At an odd place,
the resulting finite inertia image on the two-adic Tate module lies in the
torsion-free group $1+4M_{2m}(\mathbb Z_2)$.  It is trivial, so the
good-reduction criterion gives an abelian scheme over the valuation ring.

We also need the \emph{specified} theta-marked object to extend.
Let $R$ be this valuation ring, with fraction field $B_v$, and let
$\mathcal A/R$ be the good abelian scheme.  Regularity and the integral
special fiber make $\operatorname{Pic}(\mathcal A)\to
\operatorname{Pic}(\mathcal A_{B_v})$ an isomorphism: divisors extend by
closure, while a vertical divisor is a multiple of the principal special
fiber.  Extend the marked line bundle and normalize its rigidification at
zero by rescaling the generic identification.  If its polarization is
$N\lambda_0$ generically, the principal homomorphism $\lambda_0$ and its
inverse extend, so the special polarization homomorphism has finite kernel.
Semicontinuity gives a nonzero special-fiber section.  An effective line
bundle with finite translation kernel on an abelian variety is ample
\cite[Chapter II, Section 6, Application 1, pp.~60--61]{MumfordAV1970}.
Thus the extended line bundle is relatively ample.

The full torsion marking extends since $N$ is a unit in $R$.  Compare a
specified generic theta lift of a standard order-$N$ torsion generator with
an arbitrary integral lift.  Their ratio is a scalar $u\in B_v^*$.
The integral lift's $N$th iterate is a unit and the specified lift's iterate
is one; therefore $Nv(u)=0$, so $u$ is a unit.  The exact specified lift
extends.  Relations extend by equality on the generic fiber, giving the
specified theta-group isomorphism.  The symmetry involution extends by the
same argument with exponent two; its restriction to the two-torsion, and
compatibility with the theta structure, extend as well.  These are precisely
Mumford's normalized, totally symmetric, theta-marked moduli data
\cite[Section 6, pp.~76--79, 83]{Mumford1967}.
The marked moduli point therefore extends into the moduli open.  Its
projective extension is the unique projective specialization, so the
coordinate reduction used below cannot land on an unwanted boundary point.

\paragraph{Reduction without choosing a residue field.}
Compute a $p$-maximal order $\mathcal O$ in $B$ by
\cite[Algorithm~6.1 and Proposition~6.2, pp.~246--247]{BuchmannLenstra1994}; this algorithm does not factor $p$ into
individual prime ideals.  Let
$\mathfrak r=\sqrt{p\mathcal O}$ and $R_0=\mathcal O/\mathfrak r$.
Frobenius kernels compute the radical, so $R_0$ is an explicit product of
finite fields.  Let $I$ be the fractional ideal generated by the homogeneous
theta coordinates.  The module $I/\mathfrak r I$ is rank-one projective over
$R_0$, since the localizations of $\mathcal O$ above $p$ are DVRs.
Choose a coordinate with nonzero image and quotient $R_0$ by its
annihilator.  On the retained nonempty support it generates the module;
linear algebra computes every affine coordinate ratio.  The same
annihilator operation selects a nonempty chart and makes its denominators
invertible.  Validity of the marked specializations guarantees that such a
chart exists.

For $i=1,\ldots,\dim_{\mathbb F_p}R_0$, quotient by the ideal generated by
$x^{p^i}-x$ for an $\mathbb F_p$-basis of $R_0$.  A field factor survives
exactly when its degree divides $i$.  The first nonzero quotient consists of
copies of a single field $\mathbb F_{p^i}$, and
$i\leq[B:\mathbb Q]$.  No least common multiple of residue degrees is used.
Run Pila's algorithm over that product.  Field arithmetic and zero tests
are simulated by multiplication matrices.  If a test has both zero and
nonzero support, retain either nonempty support and restart; each restart
strictly decreases dimension.  After at most $\dim R_0-1$ such decreases,
the computation returns the Frobenius polynomial of at least one surviving
fiber over $\mathbb F_{p^i}$.  The complete fixed geometric atlas above
supplies exactly the required input.  The running time is polynomial in
$q,\log p$.

\paragraph{Frobenius in a possibly nonabelian Galois CM field.}
Write $\pi$ for Frobenius on a surviving fiber.  All CM endomorphisms are
defined over its residue field, hence commute with $\pi$.  For a prime
$\ell\ne p$, the Tate module has rank one over $E\otimes\mathbb Q_\ell$.
The centralizer of that action on the Tate module is
$E\otimes\mathbb Q_\ell$.  Faithfulness implies that the rational
centralizer of $E$ inside the endomorphism algebra has dimension at most
$2m$; it already contains $E$, so equals $E$.  Thus $\pi\in E$.
It is integral because its Frobenius polynomial is monic integral.
This uses the commutativity of the field $E$, not of its Galois group.

The Shimura--Taniyama formula
\cite[Section 5, p.~64, before Proposition 5.3]{Milne1999} applies with
$E=\operatorname{End}^0(A')$ and does not require a maximal integral
endomorphism order.  As $p$ splits completely in $E$, the valuations of
$\pi$ at its $2m$ primes over $p$ are $0$ or $i$.  Choose a base prime
$P_1$ and label primes by $P_x=x(P_1)$, $x\in H$.  In a compatible
embedding convention the positive support is $T=\Phi^{-1}$; a different
convention translates this set.  Both its stabilizers are trivial:
the left stabilizer of $\Phi^{-1}$ is the inverse of the right stabilizer
of $\Phi$, and conversely.  If $\tau\pi=\pi$, then $\tau T=T$, hence
$\tau=1$.  Therefore $\pi$ generates $E$ and its Frobenius polynomial is
its degree-$2m$ minimal polynomial.  Factor that polynomial over the explicit
characteristic-zero field $E$ and choose any root.  Its conjugates have the
same trivial-stabilizer properties.

\paragraph{Separating supports and reading the residue map.}
In the explicit split algebra
$R_E=\mathcal O_E/p\mathcal O_E\simeq\mathbb F_p^{2m}$ compute
\[
 e_\tau=(\tau\pi)^{p-1}\pmod p\qquad(\tau\in H).
\]
The values of this idempotent are zero on $\tau T$ and one elsewhere.
These masks separate all prime components, including when $H$ is
nonabelian.  Indeed, equality of all masks at $P_x,P_y$ would imply
\[
 \tau^{-1}x\in T\iff\tau^{-1}y\in T\quad\hbox{for every }\tau\in H.
\]
Putting $u=\tau^{-1}x$ shows
$u\in T\iff u(x^{-1}y)\in T$ for all $u$.  The trivial \emph{right}
stabilizer of $T$ then forces $x=y$.
Enumerate the fixed number $2^{2m}$ of products choosing $e_\tau$ or
$1-e_\tau$ at each $\tau$.  Every nonzero such product is a primitive
idempotent.  Choose one, say $e$.  For each integral basis element $b_j$ of
$K$, linear algebra computes the unique $t_j\in\mathbb F_p$ with
$b_je=t_je$.  The assignment $b_j\mapsto t_j$ is the desired ring map.

All dimensions, geometric degrees, numbers of coordinates, lattice indices,
and mask counts in this proof are fixed with $K$.  All varying
characteristic-zero data have size polynomial in $q$, and the residue
degree is polynomially bounded as well.  Equation (27) therefore makes the
total cost polynomial in $\log p$.  At the finitely many excluded primes,
complete splitting guarantees a residue map and exhaustive search in the
fixed finite multiplication-table algebra finds one.  This completes the
proof.
\end{proof}

\paragraph{Scope.}
The proposition concerns a fixed field and completely split input primes.
It does not give a bound polynomial in a varying field's degree or
discriminant, nor does it replace complete splitting by the existence of a
single degree-one prime.  In particular, the nonabelian assertion is a
fixed-field algorithmic conclusion, not a conditional estimate for
nonabelian Artin $L$-functions.


\section{The CM construction for fixed wild cyclotomic parameters}
\label{sec:fixed-wild-cm}

The same construction extends to every fixed odd prime and every
fixed wild layer. The real field needed here is the wild subfield
of the cyclotomic field, rather than its full maximal real
subfield. This distinction allows every tame residue degree.
The constants in this section depend on the fixed parameters;
their dependence is displayed sufficiently to distinguish this
result from a polynomial bound uniform in the cyclotomic degree.

\subsection*{Proof of Theorem~\ref{thm:fixed-wild-cm}}

\begin{proof}
All characteristic-zero data depending only on \(r,b\) may be
computed once by finite exact algorithms. In particular, realize
\(R_b\) as the fixed field of the order-\(r-1\) subgroup of
\(\operatorname{Gal}(K_b/\mathbb Q)\), compute an integral basis,
and choose an integral primitive element by enumeration.
Put \(I=[\mathcal O_{R_b}:\mathbb Z[\theta]]\).
Both \(I\) and the discriminant of \(f\) are now fixed integers.
The finitely many input primes dividing \(2rI\) can be handled
by exhaustive finite-ring arithmetic. This has a constant cost
depending on \(r,b\). The existence of the required outputs in
these cases follows from the cyclotomic reduction at \(p\ne r\)
used at the end of the proof. We henceforth assume \(p\nmid2rI\).

The order of \(p\) modulo \(r^b\) divides \(d\), by the input
valuation condition. Its Frobenius therefore lies in the tame
subgroup, so \(p\) splits completely in \(R_b\). The field is
totally real, cyclic, and of odd degree \(g\). Its discriminant
is
\begin{equation}
 D_R=r^{E_{r,b}},\qquad
 E_{r,b}=\sum_{j=1}^{b-1}(j+1)\varphi(r^j)
        =bg-\frac{g-1}{r-1}-1.
 \label{eq:wild-real-discriminant}
\end{equation}
Indeed its characters of exact order \(r^j\) have conductor
\(r^{j+1}\), and there are \(\varphi(r^j)\) of them; the
conductor--discriminant identity gives \eqref{eq:wild-real-discriminant}.

\paragraph{An auxiliary imaginary quadratic field.}
Let \(\lambda=2+\lceil\log_2p\rceil\).
Under \(\mathrm A\), a prime \(q\) can be found with
\begin{equation}
 q\equiv3\pmod4,\qquad (p/q)=1,\qquad
 q\notin\{p,r\},\qquad q\leq C\lambda^2,
 \label{eq:wild-cm-auxiliary}
\end{equation}
where \(C\) is absolute. Here is the precise character content.
On integers coprime to \(4p\), the desired indicator is
\[
 \tfrac14(1-\chi_{-4})(1+\chi_{D_p}),\qquad
 D_p=\begin{cases}p,&p\equiv1\pmod4,\\4p,&p\equiv3\pmod4,
 \end{cases}
\]
and it is zero on other integers. The four characters are
Dirichlet characters of modulus \(4p\), with only the first
principal. The log-tent explicit formula in
Lemma~\ref{lem:complex-dirichlet-tent} under \(\mathrm A\)
gives main term \(x/4\) and error
\(O(\sqrt x\log(2p))\). Even prime powers have zero indicator;
higher odd powers contribute
\(O(x^{1/3}(\log x)^2)\). Removing the one possible prime
\(r\), including its powers, costs \(O((\log x)^2)\).
Taking a sufficiently large absolute multiple of \(\lambda^2\)
proves \eqref{eq:wild-cm-auxiliary}. Trial division and quadratic-symbol tests suffice
for the small-prime search, at bit cost \(O(\lambda^6)\).
Quadratic reciprocity gives
\(({-q}/p)=(p/q)=1\), so \(p\) splits in
\(E=\mathbb Q(s)\), where \(s^2=-q\).

Set \(L=R_bE\) and
\(\mathcal O=\mathbb Z[\theta,(1+s)/2]\).
The discriminants of \(R_b\) and \(E\) are coprime. Hence
\begin{equation}
 \mathcal O_L=\mathcal O_{R_b}\mathcal O_E,\qquad
 |D_L|=D_R^2q^g,\qquad [\mathcal O_L:\mathcal O]=I^2.
 \label{eq:wild-cm-order}
\end{equation}
The usual coprime-discriminant assertion can also be checked
locally: at every prime at least one of the two integer algebras
is unramified, and their tensor product is the maximal order of
the linearly disjoint compositum. The extension \(L/\mathbb Q\)
is abelian of degree \(2g\), has odd conductor dividing
\(r^bq\), and \(p\) splits completely in it.

\paragraph{Explicit principal periods and a primitive type.}
Order the real roots \(\theta_1<\cdots<\theta_g\) of \(f\),
put \(S=(\theta_i^j)_{1\leq i\leq g,\,0\leq j<g}\), and set
\[
 B_{jk}=\operatorname{Tr}_{R_b/\mathbb Q}
             \frac{\theta^{j+k}}{f'(\theta)}
 \quad(0\leq j,k<g).
\]
The trace-dual formula for a monogenic order shows that \(B\)
is integral. Its entries with \(j+k<g-1\) vanish, and its
anti-diagonal entries are one. Thus \(\det B=\pm1\).
For \(\epsilon_i=\operatorname{sign}f'(\theta_i)\), use the
embeddings \(\phi_i(s)=\epsilon_i i\sqrt q\) and the form
\[
 E(x,y)=\operatorname{Tr}_{L/\mathbb Q}
       \left(\frac{\overline x y}{f'(\theta)s}\right)
                  \quad(x,y\in\mathcal O).
\]
The product-order trace dual is
\((f'(\theta)s)^{-1}\mathcal O\). The form is consequently
integral, alternating, and unimodular. At the chosen embeddings
its Riemann positivity follows from
\(\phi_i((f'(\theta)s)^{-1})=-i/(|f'(\theta_i)|\sqrt q)\).
The symplectic period computation therefore gives
\begin{equation}
 \tau=\tfrac12B^{-1}
       +\tfrac{i\sqrt q}{2}
          S^{-1}\operatorname{diag}(|f'(\theta_i)|)S^{-T}
                    \in\mathcal H_g.
 \label{eq:wild-cm-period}
\end{equation}
In particular \(\operatorname{Im}\tau\) is \(\sqrt q\) times
one fixed positive definite matrix.

The CM type is primitive. To see this without assuming any
particular ordering of the Galois action on the real roots,
write \(\operatorname{Gal}(L/\mathbb Q)=G\times\{1,\iota\}\),
where \(|G|=g\) is odd and \(\iota\) is complex conjugation.
A stabilizer element with negative second component has an odd
power equal to \(\iota\), which cannot preserve a CM type.
The stabilizer is therefore a subgroup of \(G\).
Derivative signs at ordered simple roots alternate, and \(g\)
is odd; the number of positive signs is \((g+1)/2\).
An invariant subset of the regular \(G\)-set is a union of
orbits of the stabilizer. Its order must thus divide both
\(g\) and \((g+1)/2\), which are coprime. The stabilizer is
trivial. This is the primitive CM criterion, so the principally
polarized variety defined by \eqref{eq:wild-cm-period} is simple and has rational
endomorphism algebra \(L\).

\paragraph{Algebraization with fixed dimension.}
We give the parameter bounds needed to apply the preceding
theta construction in genus \(g\). Its ingredients are valid
in every fixed genus. The analytic theta sums at denominator
four have \(16^g\) coordinates. They satisfy the tail bound
\[
 \operatorname{Tail}(K)\leq
 C_{r,b}\exp(-c_{r,b}\sqrt q\,K^2)\qquad(K\geq2).
\]
Thus certified precision \(2^{-M}\) costs a polynomial in
\(q,M\) with exponent depending only on \(r,b\). Moreover,
\begin{equation}
 \sum_{a\in(\mathbb Z/4\mathbb Z)^g}
       \vartheta[0,a/4](\tau,0)
 =4^g\sum_{n\in(4\mathbb Z)^g}
       \exp(-\pi n^T\operatorname{Im}\tau\,n)\geq4^g.
 \label{eq:wild-theta-normalization}
\end{equation}
The phases in \eqref{eq:wild-theta-normalization} equal one because \(B^{-1}\) is integral.
Thus a nonzero normalization coordinate can be selected by
certified approximation with a fixed positive lower bound.

Saturating the lattice from \(\mathcal O\) to
\(\mathcal O_L\) changes the Faltings height by at most a
constant depending on \(I,g\). The abelian Colmez formula
applies without a conductor-at-two ambiguity
\citep[Theorem~0.3(ii), p.~631; Theorem~III.2.9(i), p.~680]{Colmez1993}.
The bound in \citet[Corollary~4.3, pp.~14--15; Theorem~4.8,
p.~18]{Wei2021}, followed by
\citet[Corollary~1.3(2), pp.~21--22]{Pazuki2012}, gives
theta height at most \(C_{r,b}\log(2q)\).
Every conditional zero estimate here concerns a Dirichlet
character over \(\mathbb Q\) of conductor at most \(r^bq\),
because the CM field and its reflex field are abelian.

For clarity, the orbit-degree bound also remains polynomial.
Any \(\mathcal O\)-stable lattice \(J\) lies between
\(I^2\mathcal O_LJ\) and \(\mathcal O_LJ\), giving a fixed
number of choices for each ideal class. The real units
congruent to one modulo \(I^2\) lie in the order. Their squares
are norms of automorphism units, so the number of inequivalent
principal polarizations is bounded by a constant depending on
the fixed field and order. Fixed-level markings and CM types
likewise introduce only fixed factors. Minkowski's class bound,
together with
\(\sum_{n\leq Y}d_{2g}(n)\leq Y(1+\log Y)^{2g-1}\),
therefore gives the theta-point degree bound
\begin{equation}
 B_0=C_{r,b}q^{g/2}(1+\log q)^{2g-1}.
 \label{eq:wild-theta-degree}
\end{equation}
This is the CM lattice classification of
\citet[Section~3, pp.~29--31]{MilneCM2020}, with the fixed
order index accounted for explicitly.

The degree bound \eqref{eq:wild-theta-degree}, the logarithmic height bound, and \eqref{eq:wild-theta-normalization}
permit exact algebraic recognition of every coordinate ratio
by \citet[Theorem~1.19, p.~241]{KannanLenstraLovasz1988}.
For example coefficient height
\(2^{B_0}\exp(C_{r,b}B_0\log(2q))\) and precision polynomial
in its logarithm and \(B_0\) suffice. Characteristic-zero
factorization and compositum construction are deterministic
polynomial-time operations
\citep[Theorem~4.5, pp.~253--254]{Lenstra1983}.
Adjoin \(\zeta_{16}\) and apply the fixed Fourier identity
of Lemma~\ref{lem:theta-fixed-fourier}. The resulting exact
level-sixteen theta point has coefficient-field degree at most
\begin{equation}
                         N=8\lceil B_0\rceil^{16^g}.
 \label{eq:wild-coefficient-degree}
\end{equation}

The simultaneous projective reduction and integral marked-CM
reduction arguments apply in any fixed genus: full
sixteen-torsion, potential good reduction, and the torsion-free
principal congruence kernel give good reduction at odd \(p\);
the chosen symmetric marking extends. Mumford's moduli theorem
and universal family apply to type \((16,\ldots,16)\)
\citep[pp.~83--84,92--94]{Mumford1967}.
The multiplication-matrix construction in
Proposition~\ref{prop:fixed-theta-complete-atlas}, with \(g\)
in place of three, gives a level-eight embedding with
\(8^g\) homogeneous coordinates, quadratic equations with
radical affine-chart ideals, and \(8^g\) complete biquadratic
addition charts. Surjectivity of the multiplication map and
projective equality are the same dimension-independent
theorems of \citet[pp.~324,330,338--341]{Mumford1966} and
\citet[p.~94]{Mumford1967}. All matrix sizes are fixed here.

Compute over the full finite residue algebra, separate equal
residue-degree blocks by Frobenius linear algebra, and use the
arithmetic/equality simulation of Pila's algorithm. This gives
at least one integer Frobenius characteristic polynomial of a
good CM reduction over some \(\mathbb F_{p^e}\), with
\(e\leq N\). It does not choose a field factor of the residue
algebra. Pila's exponent depends on the genus and model sizes,
which are fixed, rather than on \(q\)
\citep[Theorem~A, p.~1; Section~4, pp.~8--14]{Pila1990}.

\paragraph{Separating all real embeddings by conjugate supports.}
Let \(F\) be the selected Frobenius. Its \(2g\)-th power
centralizes the CM field, since the Galois action on \(L\)
has order dividing \(2g\). The centralizer is \(L\): on a
Tate module, \(L\) acts with rank one and degree \(2g\).
Thus \(F^{2g}=\pi\in\mathcal O_L\).
The Shimura--Taniyama formula gives
\begin{equation}
                 (\pi)=\prod_{\mathfrak p\in\Phi_p}
                                      \mathfrak p^{2ge},
 \label{eq:wild-cm-frobenius-support}
\end{equation}
where \(\Phi_p\) is a translate of the primitive CM type
\citep[Theorem~8.1 and Corollaries~8.2--8.3, pp.~66--68]{MilneCM2020}.
One can pass to the fixed-degree isogeny with full CM and a
finite extension to apply the maximal-order formulation, then
divide valuations by the extension degree. The type's trivial
stabilizer implies \(\mathbb Q(\pi)=L\).

Compute the characteristic polynomial of \(F^{2g}\) from that
of \(F\), and factor it over the explicit number field \(L\)
in characteristic zero. It has exactly the \(2g\) conjugate
roots \(\pi_1,\ldots,\pi_{2g}\) in \(L\). Since
\(I^2\mathcal O_L\subset\mathcal O\), their product-basis
expressions have denominators dividing \(I^2\).
Use a quadratic nonresidue of \(\mathbb F_p\), obtainable
under \(\mathrm A\) by testing \(O(\lambda^2)\) small integers,
to extract a square root \(t\) of \(-q\) in \(\mathbb F_p\).
Substituting \(s=t\) in every \(\pi_j\) yields
\(G_j(X)\in\mathbb F_p[X]/(f)\); all denominators are units
because \(p\nmid2I\).

The zero sets of the \(G_j\) separate the roots of \(f\).
Indeed, identify the primes above \(p\) with the regular
Galois set of \(L\). For distinct primes \(P,Q\) restricting
to the chosen quadratic sign, equality of membership in every
translate of \(\Phi_p\) would say that translation by
\(QP^{-1}\ne1\) preserves \(\Phi_p\), a contradiction.
By \eqref{eq:wild-cm-frobenius-support}, membership is exactly the vanishing of the appropriate
conjugate of \(\pi\). Repeatedly replace each current factor
\(h\) of \(f\) by the nonconstant members of
\[
                      \gcd(h,G_j),\qquad
                      h/\gcd(h,G_j).
\]
After all \(2g\) conjugates, every factor is linear. Thus a
root \(u\in\mathbb F_p\) and the corresponding residue
homomorphism on \(\mathcal O_{R_b}/p\mathcal O_{R_b}\) are computed using
only greatest common divisors. Integrality of the basis
conversion follows from \(p\nmid I\).

\paragraph{The formal Teichm\"uller unit.}
The fields \(R_b\) and \(T_r\) have coprime degrees, so
\(K_b=R_bT_r\). Their integer rings generate the full integer
ring locally away from \(r\): both extensions are unramified
there, and their tensor product is finite \'etale with fraction
field \(K_b\). Consequently a fixed product-basis expression
\begin{equation}
                         \zeta_{r^b}=A(\theta,\zeta_r)
 \label{eq:wild-product-expression}
\end{equation}
has denominators with prime factors only among those of \(rI\).
It can be precomputed in characteristic zero. Evaluate \eqref{eq:wild-product-expression} at
\((u,Z)\) in \(\mathcal D_p\) to obtain \(U\).
The compatible choice \(\zeta_r=\zeta_{r^b}^{r^{b-1}}\)
also gives \(U^{r^{b-1}}=Z\). Every component is the
reduction of a primitive \(r^b\)-th root, and remains
primitive because \(p\ne r\).
The automorphism of \(K_b\) fixing \(R_b\) and sending
\(\zeta_r\) to \(\zeta_r^c\) sends \(\zeta_{r^b}\) to
\(\zeta_{r^b}^{\omega_b(c)}\). Reducing this identity proves
\eqref{eq:fixed-wild-equivariance}. All components of
\(\mathcal D_p\) have degree \(d\), which proves the final
maximal-order assertion when the valuation equals \(b\).

\paragraph{Bit bound and analytic scope.}
Let \(S_0\) denote the characteristic-zero presentation length.
Recognition and the fixed number \(16^g\) of composita give
\[
 S_0\leq C_{r,b}(B_0\log(2q)+2)^{\kappa_{r,b}}
\]
for an effective fixed exponent \(\kappa_{r,b}\).
Let \(c_0\) bound the polynomial-time characteristic-zero,
lattice, and reduction algorithms, and let \(\Delta_g\)
bound Pila's arithmetic and non-field operations for the model
above. Including support splitting and changes of structure
constants, a conservative total bit bound is
\begin{equation}
 C_{r,b}\{\lambda^6+(S_0+\lambda+2)^{c_0}
                   +N^7\lambda^3(N\lambda+2)^{\Delta_g}\}.
 \label{eq:wild-fixed-bit-bound}
\end{equation}
For \(\nu=(3g-1)16^g\), \eqref{eq:wild-theta-degree}--\eqref{eq:wild-coefficient-degree} imply
\(N\leq C_{r,b}\lambda^\nu\) and
\(S_0\leq C_{r,b}\lambda^{3g\kappa_{r,b}}\).
For example
\[
 10+\max\{6,3g\kappa_{r,b}c_0,
                  7\nu+3+(\nu+1)\Delta_g\}
\]
is an effective sufficient exponent in the theorem.
The fixed precomputation and finitely many excluded primes
are absorbed in \(C_{r,b}\).
The only conditional analytic steps used Dirichlet characters
over \(\mathbb Q\): conductors dividing \(4p\) for \eqref{eq:wild-cm-auxiliary},
conductor \(p\) for quadratic nonresidue selection, and
conductors at most \(r^bq\) for the abelian CM height bound.
\end{proof}

\subsection*{Proof of Corollary~\ref{cor:fixed-wild-prime-degree-field}}
\begin{proof}
The unit \(U\) of Theorem~\ref{thm:fixed-wild-cm} has the
maximal \(r\)-power order and the equivariance required by
Theorem~\ref{thm:teich}. The unconditional reduction in that
theorem therefore supplies the degree-\(r\) field. Explicitly, form
\(\mathcal D_p[Y]/(Y^r-U)\), let \(c\) act by
\(Z\mapsto Z^c\) and \(Y\mapsto Y^{\omega_{b+1}(c)}\),
and compute the fixed algebra by linear algebra. The theorem
proves that this algebra is a field of dimension \(r\).
Any basis element outside its scalar subspace generates it,
since \(r\) is prime. Its minimal polynomial, computed from
linear dependences among its powers, is the desired monic
irreducible polynomial. These operations have bit cost
polynomial in \(r,\log p\) and preserve the displayed bound.
\end{proof}

The theorem outputs an element of the formal algebra
\(\mathcal D_p\); it does not implicitly select an embedding of
a component into an independently presented extension field.
It applies to all tame residue degrees \(d\), including odd
ones. Its dimension is \(r^{b-1}\), and \eqref{eq:wild-fixed-bit-bound} does not supply a
polynomial bound uniform in \(r,b\).

\begin{proposition}\label{prop:bounded-degree-wild-containment}
Let \(R_b\) be as above. If \(R_b\) is contained in the
compositum of finitely many number fields, each of degree at
most \(D\), then \(D\geq r^{b-1}\).
If all the fields are abelian CM fields of degree at most
\(2h\), then \(h\geq r^{b-1}\).
\end{proposition}
\begin{proof}
The Galois group of the normal closure of a degree-\(d\)
number field embeds in the symmetric group \(S_d\).
The \(r\)-part of its exponent is therefore at most
\(r^{\lfloor\log_rD\rfloor}\): the exponent of \(S_D\)
is \(\operatorname{lcm}(1,\ldots,D)\).
The Galois group of the compositum of the normal closures
embeds in their direct product and has the same upper bound
on the \(r\)-part of its exponent. Since \(R_b/\mathbb Q\)
is cyclic of degree \(r^{b-1}\), containment would give a
quotient of that order and exponent. Hence
\(r^{b-1}\leq r^{\lfloor\log_rD\rfloor}\leq D\).

In the abelian CM case the fields are already Galois. The
\(r\)-part of the exponent of each Galois group divides its
degree. A quotient of exponent \(r^{b-1}\) of a subgroup of
their product consequently requires some factor's degree to
be divisible by \(r^{b-1}\). That degree is even, and thus
at least \(2r^{b-1}\), proving the second assertion.
\end{proof}

This proposition restricts the characteristic-zero containment
mechanism used in the construction. In particular, arbitrarily
many bounded-dimensional abelian CM fields cannot replace its
growing wild field while retaining that mechanism. It does not
rule out finite-field computations using such varieties that
avoid characteristic-zero field containment.



\begin{example}\label{ex:fixed-wild-r5-f43}
The fixed-parameter construction also has a noncubic endpoint that
can be checked entirely by polynomial arithmetic. Let \(r=5\),
\(b=2\), \(p=43\), and \(\zeta=\zeta_{25}\). The period
\(\theta=\zeta+\zeta^7+\zeta^{18}+\zeta^{24}\) satisfies
\[
 f(T)=T^5-10T^3+5T^2+10T+1.
\]
Its roots modulo \(43\) are \(10,19,24,36,40\). Taking
\(\theta\mapsto10\), the resulting element of
\(\mathbb F_{43}[Z]/\Phi_5(Z)\) is
\[
 U=12+31Z+18Z^2+32Z^3,
 \qquad U^5=Z,\qquad U^{25}=1.
\]
For \(c=1,2,3,4\), the lifts \(\omega_2(c)\) are
\(1,7,18,24\), respectively, and substitution verifies
\(U(Z^c)=U(Z)^{\omega_2(c)}\).

Here is an exact characteristic-zero certificate for the transfer.
Write \(\zeta_5=\zeta^5\) and define
\[
\begin{aligned}
A_0(T)&=25+25T-23T^2-T^3+2T^4,\\
A_1(T)&=19+12T+T^2-3T^3-T^4,\\
A_2(T)&=34-T-31T^2+2T^3+3T^4,\\
A_3(T)&=47+54T-62T^2-3T^3+6T^4.
\end{aligned}
\]
Reduction modulo \(\Phi_{25}(X)\) proves both
\(f(\theta)=0\) and
\[
 35\zeta=\sum_{j=0}^3 A_j(\theta)\zeta_5^j.
\]
The period is nonrational, so its minimal polynomial has degree
five, the prime degree of the tame fixed field. The displayed
polynomial is therefore its minimal polynomial. Its discriminant
is \(49\cdot5^8\), whereas the field discriminant is
\(5^8\) by the conductor--discriminant formula; the period order
has index seven. The identity can consequently be reduced at
\(43\), where \(35\) is invertible. Substituting \(T=10\)
and reducing its four coefficients modulo \(43\) gives the stated
\(U\). Multiplication modulo \(\Phi_5\) verifies all its power
and equivariance identities. This checks the exact algebraic
conversion from a period root; it does not execute the CM search.
\end{example}

\section{Fixed higher three-power strata}
\label{sec:kummer-higher-three-powers}

The order-nine construction extends to every fixed three-adic
valuation. Its dimension dependence must remain explicit. Write
\(C_m\in\mathbb Z[X]\) for the monic Chebyshev polynomial
defined by
\[
 C_m(Z+Z^{-1})=Z^m+Z^{-m};\qquad C_3(X)=X^3-3X.
\]

\begin{proposition}\label{prop:kummer-trace-cubic-lift}
Let \(p>3\), \(p\equiv-1\pmod9\), and
\(a=v_3(p+1)\). Suppose \(z_b\in\mathbb F_{p^2}\) has
exact order \(3^b\), with \(2\le b\le a\), and put
\(u_b=z_b+z_b^{-1}\in\mathbb F_p\).
If \(b<a\), the polynomial \(X^3-3X-u_b\) splits into
three distinct linear factors over \(\mathbb F_p\).
Choosing any root, followed by quadratic extraction, yields a
cube root of \(z_b\) and hence an element of order
\(3^{b+1}\). Conversely a cube root yields a root of the
displayed polynomial. If \(b=a\), that polynomial is
irreducible over \(\mathbb F_p\).
\end{proposition}
\begin{proof}
A cube root \(w\) of \(z_b\) has order \(3^{b+1}\).
When \(b<a\), all three such roots lie in the norm-one
subgroup of \(\mathbb F_{p^2}^\times\), so their traces
\(w+w^{-1}\) belong to \(\mathbb F_p\) and satisfy the
displayed cubic. Equal traces make the corresponding roots
equal or inverse. The inverse case would imply
\(z_b=z_b^{-1}\), which is impossible. Conversely, from a
trace root \(v\), solve \(W^2-vW+1=0\) and choose the root
whose cube is \(z_b\) by direct testing. The known nonsquare
\(-3\) supplies deterministic quadratic extraction. If
\(b=a\), a root of the cubic would produce an element of
order \(3^{a+1}\) in \(\mathbb F_{p^2}\), a contradiction.
\end{proof}

\begin{theorem}[Every fixed three-power stratum]\label{thm:kummer-fixed-three-power-stratum}
Fix \(b\ge2\). Under rational Dirichlet GRH, there is an
effective deterministic algorithm which, given a prime
\(p\equiv-1\pmod{3^b}\), constructs an element of exact order
\(3^b\) in any presented quadratic extension of
\(\mathbb F_p\), in bit time
\[
 O_b\bigl((1+\log p)^{C_b}\bigr).
\]
Here \(C_b\) is an effective constant depending on \(b\).
If \(v_3(p+1)=b\), the algorithm also constructs an irreducible
cubic over \(\mathbb F_p\). No uniform bound on \(C_b\)
as \(b\) varies is asserted.
\end{theorem}

\subsection*{Proof of Theorem~\ref{thm:kummer-fixed-three-power-stratum}}
\begin{proof}
Put \(g=3^{b-1}\), \(f_b=C_g+1\), and
\[
 R_b=\mathbb Q(\theta),\qquad
 \theta=\zeta_{3^b}+\zeta_{3^b}^{-1}.
\]
The identity
\[
 Z^g f_b(Z+Z^{-1})=Z^{2g}+Z^g+1=\Phi_{3^b}(Z)
\]
shows that \(f_b\) is the degree-\(g\) minimal polynomial
of \(\theta\). Moreover \(\mathcal O_{R_b}=\mathbb Z[\theta]\):
every element of \(\mathbb Z[\zeta_{3^b}]\) has a unique
expression \(A(\theta)+B(\theta)\zeta_{3^b}\), and being
fixed by conjugation forces \(B=0\).

The quadratic auxiliary-prime argument gives
\[
 q\equiv1\pmod4,\quad (p/q)=-1,\quad
 q\notin\{2,3,p\},\quad q=O((1+\log p)^2).
\]
It is exactly the same rational-character estimate as in
Theorem~\ref{thm:dirichlet-ninth-root-cubic}. Define
\[
 L_b=R_b(\delta),\quad \delta^2=-3q,\qquad
 O=\mathbb Z[\theta,(1+\delta)/2].
\]
Then \(L_b/\mathbb Q\) is cyclic of degree \(2g\), and
\[
 |D_{L_b}|=3^{(2b-1)g}q^g,\qquad
 [\mathcal O_{L_b}:O]=3^{(g-1)/2}.                 \label{eq:fixedb-discriminant}
\]
Indeed the cyclotomic derivative formula gives
\(D_{\mathbb Q(\zeta_{3^b})}\) of absolute value
\(3^{(2b-1)g}\). The relative basis \((1,\zeta_{3^b})\)
over \(R_b\) has discriminant of absolute norm \(3\), so
\(D_{R_b}=3^{((2b-1)g-1)/2}\). The characters of \(L_b\)
have the same three-primary conductors as those of the
cyclotomic field; its \(g\) odd characters acquire the
additional quadratic conductor \(q\). Conductor--discriminant
gives the first assertion of \eqref{eq:fixedb-discriminant}, and comparison with
\(\operatorname{disc}(O)=D_{R_b}^2(3q)^g\) gives the second.

Let \(\theta_i\) be the real roots of \(f_b\), and set
\[
 S_{ij}=\theta_i^{j-1},\qquad
 B_{ij}=\operatorname{Tr}_{R_b/\mathbb Q}
       \left(\frac{\theta^{i+j-2}}{f_b'(\theta)}\right).
\]
The coefficient-of-\(X^{g-1}\) trace identity makes \(B\)
integral, symmetric, and anti-triangular with anti-diagonal
entries one. Thus \(\det B=\pm1\). On \(O\), the form
\[
 E(x,y)=\operatorname{Tr}_{L_b/\mathbb Q}
       \left(\frac{\overline{x}y}{\delta f_b'(\theta)}\right)
\]
has matrix \(\left(\begin{smallmatrix}0&B\\-B&0\end{smallmatrix}\right)\)
in the two blocks \(\theta^j\) and
\((1+\delta)\theta^j/2\). Choose the CM embedding above
\(\theta_i\) by
\(\delta\mapsto i\sqrt{3q}\,\operatorname{sign}f_b'(\theta_i)\).
Then \(E(x,ix)>0\), giving a principal polarization and the
explicit period matrix
\[
 \tau=\frac12 B^{-1}
   +\frac{i\sqrt{3q}}2 S^{-1}
      \operatorname{diag}(|f_b'(\theta_i)|)S^{-T}.       \label{eq:fixedb-periods}
\]
For fixed \(b\), the imaginary part is bounded above and
below by positive constants depending on \(b\), times
\(\sqrt q I_g\).

This CM type is primitive. Its derivative signs contain
\((g+1)/2\) plus signs. A stabilizer in
\(\operatorname{Gal}(R_b/\mathbb Q)\times C_2\) cannot have
negative quadratic sign, since an odd power would then be
complex conjugation. Its order divides both \(g\) and
\((g+1)/2\), and therefore is one. The resulting abelian
variety is simple with rational endomorphism algebra \(L_b\).

We now apply the general classical interfaces used in the
order-nine construction, retaining their dependence on \(g\).
The abelian odd-conductor Colmez formula and the Dirichlet-GRH
estimate for its characters, whose conductors divide \(3^bq\),
give Faltings height \(O_b(\log(2q))\). The fixed lattice
index in \eqref{eq:fixedb-discriminant}, saturation to full CM, and the reverse isogeny
change this bound by \(O_b(1)\). Pazuki's theta-height
comparison gives the same bound for the fixed-level theta
point. These are precisely the sources and deductions in
Proposition~\ref{prop:sextic-theta-recognition}, with their
dimension-dependent constants retained
\citep{Colmez1993,Wei2021,Pazuki2012}.

Lattice and polarization classification bound the orbit degree
of this point by \(C_bh_{L_b}\). Minkowski's bound and the
\(2g\)-fold divisor-function estimate give
\[
 h_{L_b}\le C_bq^{g/2}(1+\log q)^{2g-1}.             \label{eq:fixedb-class-number}
\]
There are \(16^g\) quarter-characteristic theta constants,
a fixed number when \(b\) is fixed. In \eqref{eq:fixedb-periods}, \(B^{-1}\) is
integral, so the sum over the \(4^g\) second characteristics
with first characteristic zero is real and at least \(4^g\).
One coordinate has absolute value at least one. Gaussian
summation gives certified approximations, and the same
algebraic-number recognition and characteristic-zero
composita as before cost \(q^{O_b(1)}\)
\citep{KannanLenstraLovasz1988,Lenstra1983}.

The Fourier conversion gives a type-\((16,\ldots,16)\)
theta point. Simultaneous projective reduction has no fixed
dimension requirement. The proof of
Proposition~\ref{prop:kummer-integral-cm-theta} also works in
dimension \(g\): potential good reduction, unramified
sixteen-torsion, and extension of the marking are unchanged.
Mumford's moduli theorem and the complete addition construction
apply in every fixed dimension \citep{Mumford1966,Mumford1967}.
The level-eight model now has \(8^g\) coordinates. Its shape
is fixed for fixed \(b\), so Pila's theorem and the finite
etale simulation compute the Frobenius polynomials in
\(q^{O_b(1)}(1+\log p)^{O_b(1)}\) time
\citep[Section~4, author pp.~12--14]{Pila1990}.

The prime \(p\) splits completely in \(L_b\). Let \(F\)
be the \(p^e\)-Frobenius on any resulting field component.
The power \(F^{2g}\) centralizes \(L_b\); its centralizer
on the rank-one \(L_b\)-Tate module is \(L_b\). It is
therefore an integer \(\pi\in\mathcal O_{L_b}\). The CM
formula gives valuations \(0\) and \(2ge\), supported on
a translate of the primitive CM type
\citep[Theorem~8.1 and Corollaries~8.2--8.3, pp.~66--68]{MilneCM2020}.
Its trivial stabilizer gives \(\mathbb Q(\pi)=L_b\).
Compute the polynomial of \(F^{2g}\) and factor it over the
explicit characteristic-zero field \(L_b\).

Choose any resulting root \(\pi\). Express all its conjugates
in the product basis of \(O\), with denominators supported
at three. Compute a square root \(t\) of \(-3q\bmod p\)
using the known nonsquare \(-3\). After substituting
\(\delta=t\), each conjugate gives a polynomial
\(P_\sigma(X)\bmod f_b\) over \(\mathbb F_p\).
Starting with the partition \(\{f_b\}\), successively split
every current factor \(h\) by
\[
 \gcd(h,P_\sigma),\qquad h/\gcd(h,P_\sigma),
\]
omitting constants. All final factors are linear: two surviving
roots would give two primes with the same membership vector
in every translate of the CM support, and their quotient would
stabilize that support. This contradicts primitivity. Thus
polynomial gcds, without a finite-field factorization oracle,
give a root \(u_b\) of \(f_b\).

Solving \(Z^2-u_bZ+1=0\) in the presented quadratic field
gives an element of exact order \(3^b\). Indeed
\(f_b(u_b)=0\) implies that \(z^g\) is a nontrivial cube
root of unity. Proposition~\ref{prop:kummer-trace-cubic-lift}
gives the irreducible cubic when \(v_3(p+1)=b\). Finally
\(q=O((1+\log p)^2)\) proves the stated bound.
\end{proof}

\paragraph{The limit of this representation.}
This is a polynomial-time statement on each fixed stratum,
not a uniform polynomial-time bound for variable \(b\).
To see the representation growth on actual prime inputs,
choose a prime
\[
 p\equiv 2\cdot3^b-1\pmod{3^{b+1}}.
\]
It has \(v_3(p+1)=b\). Under \(\mathrm A\), the
least-prime-in-a-progression bound gives
\(p\ll (\varphi(3^{b+1})\log(3^{b+1}))^2=3^{O(b)}\)
\citep[Corollary~1.2]{LLS2015}. Also \(p\ge2\cdot3^b-1\),
so \(\log p=\Theta(b)\). The dimension
\(g=3^{b-1}\), and especially the explicit \(8^g\) or
\(16^g\) coordinates, grow faster than every polynomial
in \(\log p\). This is a limitation of this written CM
representation, not a lower bound for other algorithms.

\begin{proposition}\label{prop:kummer-cyclotomic-three-unit-bound}
Let \(2\le b\le a=v_3(p+1)\), with \(p>3\), and set
\(K_b=\mathbb Q(\zeta_{3^b})\). The reduction of every
\(v\in\mathcal O_{K_b}[1/3]^\times\) at a prime over
\(p\) has three-primary part of order at most \(3^b\).
\end{proposition}
\begin{proof}
The unique prime above three is fixed by conjugation, so
\(v/\overline v\) is a global unit. All its complex
absolute values equal one; Kronecker's theorem makes it a
root of unity. The roots of unity of \(K_b\) have order
dividing \(2\cdot3^b\). Since \(p\equiv-1\pmod{3^b}\),
conjugation on the residue field is \(p\)-Frobenius. On
its three-primary multiplicative subgroup Frobenius is
inversion. The three-primary part of \(v/\overline v\)
is therefore the square of that of \(v\); squaring preserves
its order, which is at most \(3^b\).
\end{proof}

The proposition restricts the whole cyclotomic three-unit group.
It does not restrict arbitrary sums, which may acquire other
prime divisors, or varying Frobenius integers supported at \(p\).
Also, adjoining a finite extension cannot make the order-nine
seed a noncube when \(a\ge3\):
\(v_3(p^{2m}-1)=a+v_3(m)\), and norms cannot increase its
order. The uniform lifting problem remains the selection of
roots in the successive promised split trace cubics.

\section{A fixed collection of CM torsion data can have only residue degrees that are powers of two}
\label{sec:kummer-fixed-cm-data}

\begin{proposition}\label{prop:kummer-fixed-cm-data-collapse}
Fix distinct primes \(q_i\equiv1\pmod4\), \(q_i\ne3\),
and put \(R=\mathbb Q(\zeta_9+\zeta_9^{-1})\) and
\(L_i=R(\sqrt{-3q_i})\). For each \(i\), fix algebraic
coefficients for a CM theta model, finitely many torsion levels,
and finitely many algebraic endomorphism data. There is an
effectively computable constant \(a_0\), depending on this
finite data, such that, for every \(a\ge a_0\), a positive-density
set of rational primes \(p\) satisfies
\[
 v_3(p+1)=a,\qquad (p/q_i)=-1\quad\hbox{for every }i,
\]
and all residue fields generated by the fixed coefficients and
torsion coordinates have degree a power of two over
\(\mathbb F_p\). In particular, their finite tensor products,
subalgebras and finite etale factor quotients have no field
component of degree divisible by three.
\end{proposition}
\begin{proof}
Let
\[
 F_0=R(\sqrt{-3},\sqrt{q_1},\ldots,\sqrt{q_m}).
\]
The independent quadratic square classes define an involution
\(c_0\) fixing \(R\) and negating every displayed square root.
It fixes every \(L_i\). Choose a finite Galois extension
\(E/\mathbb Q\) containing \(F_0\) and all fixed data fields.
In the preimage of \(\langle c_0\rangle\) in
\(\operatorname{Gal}(E/\mathbb Q)\), a Sylow two-subgroup
maps onto \(\langle c_0\rangle\). Hence it contains an
element \(s\) restricting to \(c_0\); its order is a power
of two.

The intersection of \(E\) with the three-power cyclotomic
tower is a fixed finite field contained in
\(\mathbb Q(\zeta_{3^c})\). The integer \(c\) can be
bounded effectively from \([E:\mathbb Q]\). This intersection
contains \(\mathbb Q(\zeta_9)\), since \(F_0\) contains
\(R\) and \(\sqrt{-3}\). The restriction of \(s\) to
the intersection is complex conjugation: its order is a power
of two, and it is nontrivial on \(\zeta_3\), while the
three-power cyclotomic Galois groups have a unique involution.

For \(a\ge c\), the residue class
\[
 \tau_a=-1+3^a\pmod{3^{a+1}}
\]
restricts to complex conjugation on that intersection. Thus
\((s,\tau_a)\) defines an element of
\(\operatorname{Gal}(E\mathbb Q(\zeta_{3^{a+1}})/\mathbb Q)\).
The qualitative Chebotarev theorem supplies a positive-density
set of primes with this conjugacy class
\citep[Section~2.1, Theorem~1, p.~131]{Serre1981Chebotarev}.
The cyclotomic restriction gives \(v_3(p+1)=a\).
The action on \(\sqrt{q_i}\) gives \((q_i/p)=-1\), and
quadratic reciprocity gives \((p/q_i)=-1\). The action on
each \(L_i\) is trivial, so \(p\) splits in each CM field.

After omitting finitely many ramified primes and denominators,
the residue degree in \(E\) is the order of \(s\), a power
of two. Residue degrees in subfields divide that number.
Composita of finite fields of two-power degrees again have
two-power degree, and the same holds for subfields. This proves
the assertion for the stated finite-algebra operations.
\end{proof}

This proposition concerns fixed coordinate-algebra witnesses.
It does not restrict an auxiliary-prime menu or torsion level
growing with \(p\), or adaptive calculations with the varying
Frobenius integers. It is not an impossibility theorem for
finite-field construction.


\section{Further information from a fixed CM Frobenius}
\label{sec:kummer-fixed-frobenius}

\begin{lemma}\label{lem:kummer-unpowered-cm-frobenius}
In Proposition~\ref{prop:sextic-cm-output}, the original
\(p^e\)-Frobenius \(F\) already belongs to the embedded
CM field \(L\). Its degree over \(\mathbb Q\) is six,
and \(e\) is even.
\end{lemma}
\begin{proof}
The support argument there proves \(F^6=\pi\in L\) and
\(\mathbb Q(\pi)=L\). Since \(F\) commutes with its
power \(\pi\), it centralizes \(L\); the rank-one
\(L\)-Tate representation has centralizer \(L\).
Thus \(F\in L\), and its primitive support gives
\(\mathbb Q(F)=L\). The prime three is totally ramified
in \(L\), with residue field \(\mathbb F_3\).
Reducing \(F\overline F=p^e\) at that prime gives
\(1=(-1)^e\), so \(e\) is even.
\end{proof}

In particular, the tempting normalized sixth-power radical is
already trivial in characteristic zero:
\[
 \frac{\pi}{p^{3e}}=\left(\frac{F^2}{p^e}\right)^3.
\]
The same support argument applies to every positive power of \(F\).
Every geometric endomorphism is defined over a finite extension of
the residue field and therefore commutes with some \(F^m\).
Since \(\mathbb Q(F^m)=L\), the centralizer argument gives
\(\operatorname{End}^0_{\overline{\mathbb F}_p}(A)=L\).
Thus this ordinary reduction does not acquire additional
noncommuting endomorphisms after extending its constant field.
There is a broader local restriction on radicals of elements
of a field splitting at the target prime.

\begin{lemma}\label{lem:kummer-split-field-radicals}
Let a number field \(L\) split completely at a prime
\(p\equiv-1\pmod{3^a}\). If \(1\le k\le a\), then
adjoining \(3^k\)-th roots of finitely many elements of
\(L^\times\), together with their conjugates and
\(\mu_{3^k}\), produces only residue degrees dividing two
at primes over \(p\).
\end{lemma}
\begin{proof}
Each completion of \(L\) at \(p\) is \(\mathbb Q_p\).
Write \(\alpha=p^v u\), where \(u\in\mathbb Z_p^\times\).
The map \(x\mapsto x^{3^k}\) is an automorphism on
\(\mathbb F_p^\times\), and Hensel's lemma lifts its
unique inverse to \(\mathbb Z_p^\times\). Thus every local
radical lies in
\[
 \mathbb Q_p\bigl(p^{1/3^k},\mu_{3^k}\bigr).
\]
The cyclotomic part is unramified quadratic, because
\(\operatorname{ord}_{3^k}(p)=2\); over it the displayed
radical of \(p\) is totally ramified. The containing field
has residue degree two, proving the claim for all subfields
and finite composita in question.
\end{proof}

This excludes a residue-degree witness from these radicals
alone. It does not exclude new combinations with \(\zeta_3\)
before radical extraction, or arbitrary arithmetic searches
inside the quadratic residue field.

\subsection{Exact orbit counts on a full local CM module}

For \(L=\mathbb Q(\theta,\delta)\), with
\(\theta^3-3\theta+1=0\) and \(\delta^2=-3q\), put
\(\lambda=\delta/(\theta-2)\). Direct elimination gives
\[
 \lambda^6+15q\lambda^4+18q^2\lambda^2+3q^3=0.
\]
This polynomial is Eisenstein at three. Thus
\(O=\mathcal O_L\otimes\mathbb Z_3=\mathbb Z_3[\lambda]\),
\(v_\lambda(3)=6\), \(\overline\lambda=-\lambda\), and
\(O/3=\mathbb F_3[\lambda]/(\lambda^6)\).
The next proposition assumes a Tate lattice free of rank one
over this full local order; the original index-three product
order does not automatically give this hypothesis.

\begin{proposition}\label{prop:kummer-cm-torsion-counts}
Suppose Frobenius acts on \(O/(\lambda^h)\) by multiplication
by \(\epsilon u\), where \(\epsilon\in\{1,-1\}\) and
\(u\equiv1\pmod\lambda\). Put
\(t_j=v_\lambda(u^{3^j}-1)\), allowing infinity.
For \(j\ge1\), the orbit count at length \(3^j\) when
\(\epsilon=1\), or length \(2\cdot3^j\) when
\(\epsilon=-1\), is respectively
\[
 \frac{3^{\min(t_j,h)}-3^{\min(t_{j-1},h)}}{3^j},
 \qquad
 \frac{3^{\min(t_j,h)}-3^{\min(t_{j-1},h)}}{2\cdot3^j}.
\]
A nonzero such count is prime to three only if
\(t_0=1\) and \(j=1\). Conversely that case has a count
prime to three whenever \(h>1\).
\end{proposition}
\begin{proof}
The kernel of multiplication by \(u^{3^j}-1\) has
\(3^{\min(t_j,h)}\) elements. Subtract successive fixed
sets and divide by the orbit length. For \(\epsilon=-1\),
every nonzero orbit has even length, and multiplication by
\(u^{2\cdot3^j}-1\) has the same kernel as multiplication
by \(u^{3^j}-1\).
When the numerator is nonzero, the count has three-adic
valuation \(t_{j-1}-j\). The binomial expansion and
\(v_\lambda(3)=6\) give
\[
 t_j\ge\min(3t_{j-1},t_{j-1}+6)\ge t_{j-1}+2.
\]
Consequently \(t_{j-1}-j\ge t_0+j-2\), proving the
necessity and the asserted converse.
\end{proof}

For \(h=6\), depth one gives counts \(N_3=8,N_9=78\)
when \(\epsilon=1\), or \(N_6=4,N_{18}=39\) when
\(\epsilon=-1\). The norm relation modulo three permits
depths \(1,3,5\), as well as scalar reduction: its first
nonzero term must have odd degree in \(\lambda\).
Depth three instead gives \(N_3=234\) or \(N_6=117\),
both divisible by three. Nonsemisimplicity alone therefore
does not give the required factor count.

\subsection{A stronger quotient test with explicit CM action}

\begin{proposition}\label{prop:kummer-cm-unit-quotient}
Let a presented finite field \(k_0\), a fixed cyclic CM torsion
module \(O/(\lambda^h)\), and its integral CM action be
explicitly given. Put \(U_h=(O/(\lambda^h))^\times\),
and let \(u_F\in U_h\) be its Frobenius unit. An invariant
algebra of degree three, obtained from the primitive torsion
point algebra using an index-three subgroup of \(U_h\),
is a field for at least one such subgroup if and only if
\(u_F\notin U_h^3\). For fixed \(h\), all such invariant
algebras are computable by linear algebra from the given
endomorphism action.
\end{proposition}
\begin{proof}
Primitive torsion points form a torsor under the finite abelian
group \(U_h\). For an index-three subgroup \(H\), the
quotient is a three-element torsor for \(U_h/H\), and its
Frobenius is translation by \(u_FH\). Its coordinate algebra
is a field exactly when that translation is nontrivial. The
intersection of all index-three subgroups is \(U_h^3\).
The invariant coordinate algebra is the simultaneous kernel
of \(g-1\) for the given permutations \(g\in H\).
No torsion point or residue-field component is selected.
\end{proof}

For example \(1+\lambda^5\in O/3\) is a norm-one
noncube, whereas its full point-algebra counts above are
divisible by three. Thus the quotient test is stronger than
the full point count. The endomorphism units used here act on
the finite torsion scheme; they need not be automorphisms of
the abelian variety.

Over a coefficient product with \(t\) components of degree
\(e\), identical successful quotient outcomes give \(t\)
factors of absolute degree \(3e\). The existing factor-count
transfer applies when \(3\nmid t\); a component cannot be
chosen when that count fails. No uniform bound is currently
proved for finding an auxiliary \(q\) satisfying this test
and the component-count condition. Consequently this is a
sufficient conditional construction, not a smaller established
uniform gate.

\subsection{The canonical CM three-isogeny preserves the local order}

Under the full local-order hypothesis above, the rank-one
\(O\)-module has a unique \(O\)-stable subgroup of order three.
Quotienting by it replaces the local Tate lattice by its
\(\lambda^{-1}\)-multiple. This has the same multiplier order
and the same Frobenius multiplier. Consequently the canonical
CM ideal isogeny preserves the depth and cube-class tests above;
it does not provide a descent to a smaller local order.
Moreover the unique prime \(\mathfrak p_3\) of \(L\) over
three satisfies \(\mathfrak p_3^6=(3)\). Repeating this ideal
isogeny six times returns to a homothetic global lattice, hence
to the same unpolarized CM isomorphism class. This observation
concerns the \(O\)-stable ideal-isogeny chain. It does not
exclude choosing a different, non-\(O\)-stable subgroup by a
separate coordinate algorithm.


\section{The original CM lattice removes the depth-one case}
\label{sec:kummer-original-cm-lattice}
The integral lattice of the original period model imposes a stronger restriction than its rational endomorphism field alone.

\begin{proposition}\label{prop:kummer-original-cm-depth}
Consider the specific sextic CM model with period lattice
\[
 \mathcal O=\mathbb Z[\theta,(1+\delta)/2],\qquad
 \theta^3-3\theta+1=0,\quad\delta^2=-3q,
\]
and a good ordinary reduction at the target prime \(p\).
Let \(F\in L=\mathbb Q(\theta,\delta)\) be its
\(p^e\)-Frobenius as in
Lemma~\ref{lem:kummer-unpowered-cm-frobenius}.
With \(\lambda=\delta/(\theta-2)\) and
\(O=\mathcal O_L\otimes\mathbb Z_3\), one has
\[
 \mathcal O\otimes\mathbb Z_3=\mathbb Z_3+\lambda^2O,
 \qquad F\in(\mathcal O\otimes\mathbb Z_3)^\times.
\]
In particular, for the sign \(\epsilon\in\{1,-1\}\)
with \(F\equiv\epsilon\pmod\lambda\),
\(v_\lambda(\epsilon F-1)\ge2\).
After an isogeny saturates the local Tate lattice to a rank-one
\(O\)-lattice, no full torsion-point algebra at any fixed
three-power level has a nonzero count prime to three at a
factor degree divisible by three, measured over its field of
definition. Modulo three one has exactly
\[
 F=\epsilon(1+c_3\lambda^3+c_5\lambda^5),
 \qquad c_3,c_5\in\mathbb F_3.
\]
The invariant-unit quotient test on the saturated three-torsion
module succeeds if and only if \(c_5\ne0\).
\end{proposition}
\begin{proof}
The full local ring is \(O=\mathbb Z_3[\lambda]\), with
\(v_\lambda(3)=6\). The elements \(\theta-2\) and
\(\delta\) have valuations two and three. The product order
is therefore contained in \(\mathbb Z_3+\lambda^2O\).
The former has index three in the maximal order, by the
discriminant calculation for the period model. The latter
also has index three, since reduction modulo \(\lambda^2\)
leaves precisely one additional \(\mathbb F_3\)-coordinate.
The inclusion is consequently an equality.

The argument uses the integral Tate lattice. The characteristic-zero uniformization with
lattice \(\mathcal O\) gives
\(T_3(A)=\mathcal O\otimes\mathbb Z_3\), compatibly with
the CM action. To compare with the good reduction, pass to a
strictly henselian valuation ring. For every \(n\), the kernel
of multiplication by \(3^n\) on the abelian scheme is finite
etale: multiplication is finite flat and its differential is
invertible because \(p\ne3\). A finite etale scheme over a
strictly henselian local ring is constant, so its geometric
generic and special fibres are identified compatibly with
all CM endomorphisms. Taking inverse limits preserves the
specified \(\mathcal O\)-lattice.

Since \(F\in L\) and acts integrally on this lattice, it
belongs to its multiplier ring
\((\mathcal O_3:\mathcal O_3)=\mathcal O_3\).
It is a unit because its norm relation is
\(F\overline F=p^e\) and \(p\ne3\). Its constant residue
is \(\epsilon\), and a scalar in \(\mathbb Z_3\) differs
from that sign by a multiple of three, of valuation at least
six. Thus its normalized depth is at least two.

An isogeny over the residue field does not change the element \(F\) of the rational
CM algebra. If its field of definition is enlarged, Frobenius becomes
a positive power of \(F\), whose normalized depth is still at least two. On a saturated local lattice the orbit-count
formula in Proposition~\ref{prop:kummer-cm-torsion-counts}
therefore excludes a count prime to three at every
three-divisible orbit length. This assertion is about that
full-order module; the original module
\(\mathcal O_3/3\mathcal O_3\) is a different integral
module and is not identified with \(O/3O\).

For the final assertion write
\(\epsilon F=1+c_2\lambda^2+c_3\lambda^3+
c_4\lambda^4+c_5\lambda^5\) in
\(O/3O=\mathbb F_3[\lambda]/(\lambda^6)\).
The relation \(F\overline F=1\) there uses the already
proved evenness of \(e\). As
\(\overline\lambda=-\lambda\), its coefficients of
\(\lambda^2\) and then \(\lambda^4\) give
\(c_2=c_4=0\). Finally, the cubes in the unit group of
this ring are precisely
\[
 \{\epsilon'(1+c\lambda^3):\epsilon'\in\{1,-1\},
                                   \ c\in\mathbb F_3\},
\]
by the characteristic-three binomial identity. The unit is
therefore a noncube exactly when \(c_5\ne0\), and
Proposition~\ref{prop:kummer-cm-unit-quotient} applies.
\end{proof}

This identifies a concrete remaining condition on the varying
auxiliary prime: its original CM Frobenius must have nonzero
\(\lambda^5\)-coefficient, and the coefficient algebra must
also satisfy the required component-count condition. Neither
condition has been forced uniformly by rational Dirichlet GRH.

\section{A class-number test for varying the auxiliary quadratic field}
\label{sec:kummer-varying-q-classnumber}
The quadratic class number gives a computable auxiliary-prime test, but the integral order contributes a separate factor of three.

\begin{proposition}\label{prop:kummer-varying-q-classnumber}
Let \(q\equiv1\pmod4\) be prime, \(q\ne3\), and put
\(E=\mathbb Q(\sqrt{-3q})\),
\(R=\mathbb Q(\zeta_9+\zeta_9^{-1})\), and \(L=ER\).
Then
\[
 3\nmid h(L)\quad\Longleftrightarrow\quad3\nmid h(E).
\]
For the product order
\(\mathcal O=\mathbb Z[\theta,(1+\sqrt{-3q})/2]\)
used in the period construction, one nevertheless has
\[
 |\operatorname{Pic}(\mathcal O)|=3h(L).
\]
Thus making the maximal-order class number prime to three
does not remove the three-part introduced by the product order.
\end{proposition}
\begin{proof}
The extension \(L/E\) is cyclic cubic, totally ramified at
the unique prime of \(E\) over three, and unramified at
every other finite prime. Both fields are totally imaginary.
The units of \(E\) are \(\{1,-1\}\), and both are norms
from \(L\), since the extension degree is odd. The ambiguous
class-number formula therefore gives
\[
 |\operatorname{Cl}(L)^{\operatorname{Gal}(L/E)}|=h(E).
\]
For an exact primary proof and formula, see
\citep[Theorem~1.1, p.~2]{YuAmbiguous2016}.
A nonzero finite three-primary group acted on by a group of
order three has a nonzero fixed subgroup: count its orbits,
whose nontrivial sizes are three. Applying this to the
three-primary part of \(\operatorname{Cl}(L)\) proves the
first equivalence.

The product order has index three and local conductor
\(\mathfrak p_3^2\), with no conductor away from three.
The standard exact sequence for extension of invertible ideals
from an order to its normalization is
\[
 \mathcal O_L^\times\longrightarrow
 (\mathcal O_L/\mathfrak p_3^2)^\times/
 (\mathcal O/\mathfrak p_3^2)^\times
 \longrightarrow\operatorname{Pic}(\mathcal O)
 \longrightarrow\operatorname{Cl}(L)\longrightarrow1.
\]
The last map is surjective because every ideal class has an ideal representative coprime to the conductor. An invertible ideal whose extension is principal is obtained by gluing the trivial modules off the conductor and on its local support; their gluing ratio is a unit of the normalization modulo the conductor. Changing either trivialization divides by a residue unit of the order or by a global unit of the normalization. This proves exactness in the displayed form.
The middle residue-unit quotient has order three, since the
local rings are \(\mathbb F_3[\lambda]/(\lambda^2)\)
and \(\mathbb F_3\), respectively. Its global-unit image
is trivial. Indeed \(L\) has only the roots of unity
\(\{1,-1\}\). For every unit \(u\), Kronecker's theorem
gives \(u/\overline u\in\{1,-1\}\). Reduction at the
unique prime above three gives \(u/\overline u=1\) in
\(\mathbb F_3\), excluding the negative sign. Thus every
unit lies in the real subfield \(R\), whose full integer
ring is \(\mathbb Z[\theta]\subset\mathcal O\).
The exact sequence now gives the asserted factor three.
\end{proof}

The quadratic class number \(h(-3q)\) is computable in
time polynomial in \(q\), for example by enumerating reduced
positive binary quadratic forms. This supplies a concrete
test during a small-\(q\) search. No uniform bound under
rational Dirichlet GRH has been proved here for an auxiliary
\(q\) passing this class-number test together with the
Frobenius \(\lambda^5\)-coefficient and component-count
tests. In particular, an ideal-class cardinality has not been
identified with a chosen theta coefficient-field degree
without the additional reflex-norm and marking calculation.

\section{What sparse cyclotomic compression leaves to select}
\label{sec:sparse-cyclotomic-selector}

The conductor \(r^{a+1}\) has a short binary description even
when its cyclotomic field has large degree. The following
calculation isolates the coefficient selection left by using
this sparse presentation directly in characteristic \(p\).
It is a statement about this proposed compression method,
rather than a lower bound for finite-field construction.

\begin{proposition}\label{prop:sparse-wild-coefficient-algebra}
Let \(p\ne r\) be primes with \(r\) odd, and put
\(d=\operatorname{ord}_r(p)\), \(a=v_r(p^d-1)\), and
\(g=r^{a-1}\). For \(j=a,a+1\), write
\[
 C_j=\mathbb F_p[X_j]/\Phi_{r^j}(X_j),
\]
and let the tame group \(H\cong(\mathbb Z/r\mathbb Z)^\times\)
act by \(X_j\mapsto X_j^{\omega_j(c)}\).
Set \(A=C_a^H\) and \(B=C_{a+1}^H\).
Then
\[
 A\cong\mathbb F_p^{\,g},\qquad
 B\cong(\mathbb F_{p^r})^{g}.
\]
Under the embedding \(X_a\mapsto X_{a+1}^r\), the algebra
\(A\) is exactly the subalgebra of \(B\) fixed by the
absolute Frobenius. Consequently, for \(b\in B\), the
polynomial
\[
                   \prod_{i=0}^{r-1}(T-b^{p^i})
\]
has coefficients in \(A\), which need not be scalar in
\(\mathbb F_p\).
\end{proposition}
\begin{proof}
The order of \(p\) modulo \(r^a\) is \(d\), whereas its
order modulo \(r^{a+1}\) is \(dr\). The first assertion
follows from the valuation definition and the second from
lifting the exponent, since \(r\) is odd and \(r\nmid d\).
The geometric points of \(C_j\) form a regular set for
\((\mathbb Z/r^j\mathbb Z)^\times\). Taking \(H\)-invariants
passes to its orbits, a regular set for a cyclic wild group
of order \(r^{j-1}\). On that quotient, the Frobenius is
trivial for \(j=a\) and has exact order \(r\) for \(j=a+1\).
The finite-field product descriptions follow. The image of
\(A\) in \(B\) is Frobenius-fixed and has dimension \(g\),
which is the dimension of the fixed algebra just computed;
therefore the two algebras coincide. Frobenius cyclically
permutes the factors of the displayed polynomial, proving
the coefficient assertion.
\end{proof}

There is an exact correspondence between selecting values
for these coefficient functions and supplying a maximal
normalized Teichm\"uller unit. The correspondence can itself
be implemented with short arithmetic circuits; it is not an
artifact of a large dense power basis.

\begin{proposition}\label{prop:succinct-selector-teich-correspondence}
Use \(A,C_a,g\) as above and put
\(D=\mathbb F_p[Z]/\Phi_r(Z)\). There is a natural bijection
between unital \(\mathbb F_p\)-algebra homomorphisms
\(\rho:A\to\mathbb F_p\) and elements \(U\in D\) satisfying
\[
 U^g=Z,\qquad U^{r^a}=1,
 \qquad \sigma_c(U)=U^{\omega_a(c)}\quad(c\bmod r).
\]
From such a \(U\), evaluation of an invariant arithmetic
circuit in \(C_a\) under its corresponding \(\rho\) costs
polynomial time in the circuit size, \(r\), and \(\log p\).
Conversely, a consistent evaluation oracle for one such
\(\rho\), on \(r-1\) explicitly constructible invariant
circuits of size polynomial in \(r,\log p\), gives \(U\)
within the same polynomial bound apart from oracle time.
\end{proposition}
\begin{proof}
There is an isomorphism
\[
                  A\otimes_{\mathbb F_p}D\ \longrightarrow C_a,
                  \qquad Z\longmapsto X_a^g.
\]
Over an algebraic closure, each tame orbit of primitive
\(r^a\)-th roots maps bijectively to the primitive
\(r\)-th roots under \(x\mapsto x^g\). This proves the
isomorphism of finite \'etale algebras, without selecting
any component.

An evaluation \(\rho\) consequently sends \(X_a\) to an
element \(U\in D\) with all the stated identities. Conversely,
such a \(U\) defines the equivariant homomorphism
\(C_a\to D\), \(X_a\mapsto U\). Its restriction to
\(A\) lands in \(D^H=\mathbb F_p\), giving \(\rho\).
The two constructions are inverse.

For the computational assertion, use the basis
\(1,Z,\ldots,Z^{r-2}\) of \(D\), and let
\(b_0(Z),\ldots,b_{r-2}(Z)\) be its dual basis for the
nondegenerate trace pairing of the \'etale algebra \(D\).
This basis is computed by linear algebra of dimension \(r-1\).
Under the displayed isomorphism, write
\[
 X_a=\sum_{j=0}^{r-2}A_jZ^j,
 \qquad
 A_j=\sum_{c\in H}
       \sigma_c\bigl(X_a b_j(X_a^g)\bigr)\in A.
\]
These are the usual dual-trace coefficient formulas.
The powers have binary exponents of length
\(O(a\log r)=O(r\log p)\); the tame lifts are computed
by Hensel lifting. Thus the \(r-1\) displayed invariant
circuits have polynomial size. Their consistent evaluations
give
\(U=\sum_j\rho(A_j)Z^j\).
In the reverse direction, substitute the known \(U\) into
an invariant circuit. Its value in \(D\) is scalar by
equivariance, so its scalar coefficient gives the value of
\(\rho\). The asserted bounds follow.
\end{proof}

The word ``consistent'' is essential: independently choosing
a root for each coordinate polynomial need not preserve
products or define a homomorphism. The proposition identifies
the remaining operation in this specific sparse-algebra
proposal with the maximal-unit input. It provides no new
algorithm for that operation.

For example, when \(r=3\) and \(p\equiv-1\pmod{3^a}\)
with exact valuation \(a\), put \(x=X_{a+1}\) and
\(b=x+x^{-1}\). Then \(b\in B\), and its Frobenius
polynomial is
\[
               T^3-3T-(x^3+x^{-3}).
\]
Its apparently single scalar parameter is actually an
element of \(A\cong\mathbb F_p^{3^{a-1}}\).
A straight-line formula for that element does not select
one of its \(\mathbb F_p\)-values.


\section{Questions remaining}
The remaining problem is a bound uniform in the prime degree and
the valuation for a prescribed hard input. The separate polynomial-time
output algorithms of Tang, using the cited OpenAI and Cureton--Thorner
theorems, are already available under those accepted results.
Theorem~\ref{thm:fixed-wild-cm} and
Corollary~\ref{cor:fixed-wild-prime-degree-field} prove the
construction using \(A\) alone for every fixed odd prime degree
and fixed valuation. The unsettled uniform case is
\[
 r\nmid p-1,\qquad a=v_r(p^{\operatorname{ord}_r(p)}-1)\ge2,
\]
with \(r\) and \(a\) allowed to grow and a requested bound
polynomial in \(r,\log p\) with an absolute exponent.
Theorem~\ref{cor:two-primary-cyclotomic-cofactor} also handles arbitrary
valuation whenever its numerical cofactor is polynomially bounded
in \(r,\log p\), and Corollary~\ref{cor:fermat-supported-cubic}
handles the stated smooth cubic family unconditionally. The uniform
problem therefore remains outside these additional input classes.
The CM dimension \(r^{a-1}\) and its explicit theta representation
prevent the displayed fixed-parameter proof from giving that bound.
Proposition~\ref{prop:bounded-degree-wild-containment} restricts
one attempt to replace its characteristic-zero field by a compositum
of bounded-degree fields; it imposes no lower bound on other
finite-field algorithms.

The classical short-auxiliary-prime estimate remains a separate open
derivation even in the solved output strata. The distinction is about
the ingredients of a derivation. It is not a claim of logical
independence from \(A\), since the separate unconditional zero-free
theorem already implies the output and auxiliary-prime conclusions.

\paragraph{The analytic question.}
The classical sufficient RH hypothesis is the single induced Hecke function in \eqref{eq:onehecke}, equivalently the pure-field zeta function under $A$. A weaker sufficient input is the one-sided window of Proposition~\ref{prop:window}: a fixed gap below twice the principal term for the smoothed pure-field prime sum at a scale polynomial in $r$ and $\log p$. The residual-zero identity in Section~\ref{sec:kummer-residual-zero} makes strict positivity equivalent to a good prime in its support. It therefore describes precisely what must be proved, but does not establish that positivity. The full tensor cone permits equality at the wrong endpoint, and arbitrary rational twists preserve a conductor contribution of size at least $p^{r-1}$ when $\operatorname{ord}_r(p)>1$. The quoted twist-density bound does not certify a fixed zero-free width at a polynomial-logarithmic twist scale. A successful analytic argument must add information beyond these displayed constraints.

\paragraph{The algebraic question.}
Theorem~\ref{thm:teich} identifies an exact alternative input: a maximal $r$-primary unit with the stated Teichm\"uller equivariance in the explicit cyclotomic algebra. A field of degree divisible by $r$ already provides it, but this reverse reduction does not construct that field from its order. An explicit factor-count witness with component count prime to $r$ also suffices. Complete cyclotomic products, their symmetric compressions, fixed candidate menus, and the semiprimitive Jacobi sums considered here do not produce such a witness uniformly. These are limitations of the specified constructions, with their hypotheses; none is a general impossibility theorem.

\paragraph{Elliptic and class-field inputs.}
The elliptic route still needs an effectively constructed curve whose Frobenius discriminant has $r$-adic valuation one. In the cubic prime-field Tate family, its success is exactly the noncubic class of $k-3\zeta_3$. Given an ordinary prime-field parent, the restricted nonparent selector is polynomial-time equivalent to general promised-cube extraction. The two-trial construction over $\mathbb F_{p^4}$ supplies an ordinary curve with rational three-torsion, but changes that parent domain and does not force the necessary noncube. The CM route needs a polynomially bounded order with prime class nonzero modulo $r$; for $r=3$ its stated budget is equivalently an inert complex cubic field within the same discriminant bound. Neither quadratic splitting nor existence of the relevant isogeny class supplies these effective choices.

\paragraph{What the average theorems settle.}
For fixed $r$, the short auxiliary search succeeds outside $O_{r,C_r}(\sqrt P\log\log P)$ primes in a dyadic interval. The simultaneous $12r$ table gives both algebraic algorithmic interfaces for almost all characteristics in the displayed degree range, and the cubic Tate theorem gives a direct output after $O(\log p)$ tests outside the displayed exceptional set. The supersingular point formulation is equivalent to finding a maximal three-primary unit; the two unconditional ordinary quartic-field trials do not by themselves supply a noncube. These results retain exceptional input primes. The character-product lower comparisons show that exceptional primes with long unsuccessful searches do occur; their lower scales still allow a uniform bound with a larger logarithmic power. The fixed-parameter CM construction removes exceptional inputs from every fixed prime-degree and valuation stratum. A bound uniform in both parameters, or a different construction of the maximal equivariant unit with such a bound, remains the central question.

\paragraph{Method and records.}
The mathematics and text were produced by GPT-6 Astra in ChatGPT Pro mode.
Independent statement-alone re-proofs of the cofactor construction, the fixed-parameter construction and its ninth-root case, principal algebraic reductions and average theorems used the same model; the configured effort was not exposed
in their records. These were not different-model or formal checks.  Earlier project records describe targeted checks
by Claude Opus 5.5 at effort high; those checks concern the earlier
consequences, not Tang's later algorithm.  No human has checked this paper.
Our text is dedicated to the public domain under
\href{https://creativecommons.org/publicdomain/zero/1.0/}{CC0~1.0}.
The accompanying programs and records are in the
\href{\workbenchurl}{zeta and $L$-functions workbench}; the permanent
record address is \url{\recordurl}.

\begin{thebibliography}{99}

\bibitem[Adleman and Lenstra(1986)]{AdlemanLenstra1986}
Adleman, L. M., \& Lenstra, H. W., Jr. (1986).
Finding irreducible polynomials over finite fields.
In \emph{Proceedings of the eighteenth annual ACM symposium on theory of computing}
(pp.~350--355). ACM. \url{https://doi.org/10.1145/12130.12166}

\bibitem[Akiyama(1996)]{Akiyama1996}
Akiyama, S. (1996).
On the pure Jacobi sums.
\emph{Acta Arithmetica, 75}(2), 97--104.
Checked author manuscript, Lemma~1, p.~3,
\url{https://www.math.tsukuba.ac.jp/~akiyama/papers/Jacobi_priv.pdf}.

\bibitem[Alp\"oge(2026)]{Alpoge2026}
Alp\"oge, L. (2026, November 20).
\emph{Announcement concerning the Riemann hypothesis} [Author's four-page note].
\url{https://alpo.ge/remark.pdf}

\bibitem[Arthur and Clozel(1989)]{ArthurClozel1989}
Arthur, J., \& Clozel, L. (1989).
\emph{Simple algebras, base change, and the advanced theory of the trace formula}
(Annals of Mathematics Studies, Vol.~120). Princeton University Press.
\url{https://www.claymath.org/library/cw/arthur/pdf/30.pdf}.

\bibitem[Bach and Sorenson(1996)]{BachSorenson1996}
Bach, E., \& Sorenson, J. (1996).
Explicit bounds for primes in residue classes.
\emph{Mathematics of Computation, 65}(216), 1717--1735.
\url{https://doi.org/10.1090/S0025-5718-96-00763-6}

\bibitem[Bach et~al.(2001)]{BGL2001}
Bach, E., von zur Gathen, J., \& Lenstra, H. W., Jr. (2001).
Factoring polynomials over special finite fields.
\emph{Finite Fields and Their Applications, 7}(1), 5--28.
\url{https://doi.org/10.1006/ffta.2000.0306}.

\bibitem[Balakrishnan et~al.(2016)]{BILV2016}
Balakrishnan, J.~S., Ionica, S., Lauter, K., \& Vincent, C. (2016).
Constructing genus-3 hyperelliptic Jacobians with CM.
\emph{LMS Journal of Computation and Mathematics, 19}(Suppl.~A), 283--300.
\url{https://doi.org/10.1112/S1461157016000322}.
Author version: \url{https://arxiv.org/pdf/1603.03832v2}.

\bibitem[Bartoli et~al.(2024)]{BartoliEtAl2024}
Bartoli, D., Marino, G., Neri, A., \& Vicino, L. (2024).
Exceptional scattered sequences.
\emph{Algebraic Combinatorics, 7}(5), 1405--1431.
\url{https://doi.org/10.5802/alco.377}.

\bibitem[Beckwith et al.(2024)]{BeckwithRaumRichter2024}
Beckwith, O., Raum, M., \& Richter, O. K. (2024).
Imaginary quadratic fields with $\ell$-torsion-free class groups and
specified split primes.
\emph{International Mathematics Research Notices, 2024}(16),
11582--11596. \url{https://doi.org/10.1093/imrn/rnae130}

\bibitem[Belabas et~al.(2008)]{BDF2008}
Belabas, K., Diaz y Diaz, F., \& Friedman, E. (2008).
Small generators of the ideal class group.
\emph{Mathematics of Computation, 77}(262), 1185--1197.
\url{https://doi.org/10.1090/S0025-5718-07-02003-0}.
Author version consulted: \url{https://www.math.u-bordeaux.fr/~kbelabas/research/OnBach.pdf}.

\bibitem[Bernard and Ku\v{c}era(2021)]{BernardKucera2021}
Bernard, O., \& Ku\v{c}era, R. (2021).
\emph{A short basis of the Stickelberger ideal of a cyclotomic field}
[Preprint, arXiv:2109.13329v1].
\url{https://arxiv.org/pdf/2109.13329v1}.

\bibitem[Bhargava et~al.(2024)]{BhargavaTaniguchiThorne2024}
Bhargava, M., Taniguchi, T., \& Thorne, F. (2024).
Improved error estimates for the Davenport--Heilbronn theorems.
\emph{Mathematische Annalen, 389}(4), 3471--3512.
\url{https://doi.org/10.1007/s00208-023-02684-w}.
Checked author manuscript,
\url{https://www.math.kobe-u.ac.jp/HOME/tani/dh-improved-public.pdf},
Theorem~1.3 and Table~1, pp.~2--3; Theorem~1.4 and Corollary~1.5,
p.~4.

\bibitem[Bhargava and Varma(2016)]{BhargavaVarma2016}
Bhargava, M., \& Varma, I. (2016).
The mean number of 3-torsion elements in the class groups and ideal
groups of quadratic orders.
\emph{Proceedings of the London Mathematical Society, 112}(2),
235--266.
\url{https://doi.org/10.1112/plms/pdv062}.
Checked author-hosted journal offprint,
\url{https://www.math.toronto.edu/~ila/BV.pdf},
\S5.1, Lemmas~33--34 and Proposition~35, pp.~256--258.

\bibitem[Bhargava et al.(2017)]{BhargavaEtAl2017}
Bhargava, V., Ivanyos, G., Mittal, R., \& Saxena, N. (2017).
Irreducibility and deterministic $r$-th root finding over finite fields.
In \emph{Proceedings of the 2017 ACM International Symposium on Symbolic
and Algebraic Computation}.
\url{https://doi.org/10.1145/3087604.3087620}
Author manuscript:
\url{https://www.cse.iitk.ac.in/users/nitin/papers/irred-and-rth-root.pdf}.

\bibitem[Biasse and Erukulangara(2023)]{BiasseErukulangara2023}
Biasse, J.-F., \& Erukulangara, M.~R. (2023).
A proof of the conjectured run time of the Hafner--McCurley class group algorithm.
\emph{Advances in Mathematics of Communications, 17}(6), 1388--1405.
\url{https://doi.org/10.3934/amc.2021055}.
Author version consulted: \url{https://par.nsf.gov/servlets/purl/10324557}.

\bibitem[Boyer(2014)]{Boyer2014}
Boyer, I. (2014).
\emph{Vari\'et\'es ab\'eliennes et jacobiennes de courbes hyperelliptiques,
en particulier \`a multiplication r\'eelle ou complexe}
[Doctoral dissertation, Universit\'e Paris Diderot--Paris VII].
\url{https://www.normalesup.org/~iboyer/files/TheseIBoyer.pdf}.

\bibitem[Br\"oker et~al.(2012)]{BrokerLauterSutherland2012}
Br\"oker, R., Lauter, K., \& Sutherland, A. V. (2012).
Modular polynomials via isogeny volcanoes.
\emph{Mathematics of Computation, 81}(278), 1201--1231.
\url{https://doi.org/10.1090/S0025-5718-2011-02508-1}.
Checked author manuscript: arXiv:1001.0402v2, \S3, equations (3)
and (5), pp.~4--5; \S4.1, pp.~6--7.

\bibitem[Buchmann and Lenstra(1994)]{BuchmannLenstra1994}
Buchmann, J.~A., \& Lenstra, H.~W., Jr. (1994).
Approximating rings of integers in number fields.
\emph{Journal de Th\'eorie des Nombres de Bordeaux, 6}(2), 221--260.
\url{https://doi.org/10.5802/jtnb.113}.
\url{https://www.numdam.org/item/JTNB_1994__6_2_221_0.pdf}.

\bibitem[Buchmann and Shoup(1996)]{BuchmannShoup1996}
Buchmann, J., \& Shoup, V. (1996).
Constructing nonresidues in finite fields and the extended Riemann hypothesis.
\emph{Mathematics of Computation, 65}(215), 1311--1326.
\url{https://doi.org/10.1090/S0025-5718-96-00751-X}.
Checked author manuscript, \url{https://www.shoup.net/papers/nonresidues.pdf},
pp.~2--5 and \S4, Fact~6, p.~9.

\bibitem[Cilleruelo(2011)]{Cilleruelo2011}
Cilleruelo, J. (2011).
The least common multiple of a quadratic sequence.
\emph{Compositio Mathematica, 147}(4), 1129--1150.
\url{https://doi.org/10.1112/S0010437X10005191}.
Lemmas~1, 3, and 4 and \S2.4, pp.~1131--1134.

\bibitem[Cohen and Pazuki(2009)]{CohenPazuki2009}
Cohen, H., \& Pazuki, F. (2009).
Elementary 3-descent with a 3-isogeny.
\emph{Acta Arithmetica, 140}(4), 369--404.
\url{https://doi.org/10.4064/aa140-4-6}.

\bibitem[Cornelissen et~al.(2026)]{CornelissenHokkenRingeling2026}
Cornelissen, G., Hokken, D., \& Ringeling, B. (2026).
\emph{The asymptotic Mahler measure of Gaussian periods}
(arXiv:2507.09303v3), \S8.3, Lemma~8.3.1.
\url{https://arxiv.org/html/2507.09303v3}.

\bibitem[Couveignes and Lercier(2013)]{CouveignesLercier2013}
Couveignes, J.-M., \& Lercier, R. (2013).
Fast construction of irreducible polynomials over finite fields.
\emph{Israel Journal of Mathematics, 194}, 77--105.
\url{https://doi.org/10.1007/s11856-012-0070-8}
Author manuscript, September 30, 2009:
\url{https://www.math.u-bordeaux.fr/~jcouveig/publi/couveignes-lercier.pdf}.

\bibitem[Cremona and Sutherland(2024)]{CremonaSutherland2024}
Cremona, J. E., \& Sutherland, A. V. (2024).
Computing the endomorphism ring of an elliptic curve over a number field.
In \emph{LuCaNT: LMFDB, computation, and number theory}
(\emph{Contemporary Mathematics, 796}, pp.~75--102).
American Mathematical Society.
\url{https://doi.org/10.1090/conm/796/15998}.
Checked author manuscript dated 30 June 2023, \S2.3, pp.~4--5;
\S2.4, Proposition~4, p.~5.

\bibitem[Cureton and Thorner(2026)]{CT2026}
Cureton, D., \& Thorner, J. (2026).
\emph{A quasi-Riemann hypothesis for Dedekind zeta functions}
(Version 2). Hexagon, 2610.00200v2.
\url{https://hexagonmath.org/2610.00200v2}.
Theorem~1.1, p.~3; Lemma~8.1, pp.~21--22.

\bibitem[Duke(2006)]{Duke2006}
Duke, W. (2006).
Modular functions and the uniform distribution of CM points.
\emph{Mathematische Annalen, 334}(2), 241--252.
\url{https://doi.org/10.1007/s00208-005-0706-7}.
The $j$-series and discriminant $-15$ singular modulus are displayed
on p.~1 of \url{https://math.ucla.edu/~wdduke/preprints/modud.pdf}.

\bibitem[Enge(2009)]{Enge2009ClassPolynomials}
Enge, A. (2009).
The complexity of class polynomial computation via floating point
approximations.
\emph{Mathematics of Computation, 78}(266), 1089--1107.
Checked author manuscript arXiv:cs/0601104v2, 25 July 2008,
\S1, Theorems~1.1--1.2 and heuristic Corollary~1.3, pp.~2--3.
\url{https://arxiv.org/abs/cs/0601104}.

\bibitem[Evans et~al.(2025)]{EvansEtAl2025Hasse}
Evans, R., Lemmermeyer, F., Sun, Z.-H., \& Van Veen, M. (2025).
Ring class fields and a result of Hasse.
\emph{Journal of Number Theory, 266}, 33--61.
\url{https://doi.org/10.1016/j.jnt.2024.07.001}.
Checked April 2024 author manuscript,
\url{https://mathweb.ucsd.edu/~revans/hasse8.pdf},
Theorems~1.1--1.4, pp.~1--2, and Theorem~1.6, p.~2.

\bibitem[Evans and Van Veen(2026)]{EvansVanVeen2026Cubic}
Evans, R., \& Van Veen, M. (2026).
Cubic residuacity of real quadratic integers.
\emph{Journal of Number Theory, 289}, 35--45.
\url{https://doi.org/10.1016/j.jnt.2026.05.005}.
Available online; assigned to the December 2026 issue.
Checked November 2025 author manuscript,
\url{https://mathweb.ucsd.edu/~revans/cubicresidue4.pdf},
Theorem~1.1, p.~2, which explicitly assumes $p\equiv1\pmod3$.

\bibitem[Friesen et~al.(1986)]{FriesenSpearmanWilliams1986}
Friesen, C., Spearman, B. K., \& Williams, K. S. (1986).
Another proof of Eisenstein's law of cubic reciprocity and its supplement.
\emph{Rocky Mountain Journal of Mathematics, 16}(2), 395--402.
\url{https://people.math.carleton.ca/~williams/papers/pdf/141.pdf}.
Equation (1.1), p.~395; \S2, pp.~396--397.

\bibitem[Guo(2017)]{Guo2017Thesis}
Guo, Z. (2017).
\emph{$\mathcal P$-schemes and deterministic polynomial factoring over finite fields}
[Doctoral dissertation, California Institute of Technology].
Checked version arXiv:1706.10028v2, \S1.1, Hypothesis~$(*)$, p.~3.
\url{https://arxiv.org/abs/1706.10028}.

\bibitem[Hasanalizade et~al.(2022)]{HSW2022}
Hasanalizade, E., Shen, Q., \& Wong, P.-J. (2022).
Counting zeros of Dedekind zeta functions.
\emph{Mathematics of Computation, 91}(333), 277--293.
\url{https://doi.org/10.1090/mcom/3665}.
Author version consulted: arXiv:2102.04663v3,
\url{https://arxiv.org/pdf/2102.04663v3}.

\bibitem[Hashimoto and Aoki(2024)]{HashimotoAoki2024}
Hashimoto, Y., \& Aoki, M. (2024).
Normal integral bases and Gaussian periods in the simplest cubic fields.
\emph{Annales math\'ematiques du Qu\'ebec, 48}, 157--173.
\url{https://doi.org/10.1007/s40316-022-00204-x}

\bibitem[Hashimoto and Nuida(2026)]{HashimotoNuida2026}
Hashimoto, Y., \& Nuida, K. (2026).
\emph{Ordinary 3-isogeny graphs and improvement of supersingularity
testing for twisted Hessian curves over prime fields} [Preprint].
arXiv:2609.11037v1.
\url{https://doi.org/10.48550/arXiv.2609.11037}.

\bibitem[Hooley(1967)]{Hooley1967}
Hooley, C. (1967). On Artin's conjecture.
\emph{Journal f\"ur die reine und angewandte Mathematik, 225}, 209--220.
Published translation consulted: Hooley, C. (1968).
O gipoteze Artina [On Artin's conjecture] (A.~I. Vinogradov, Trans.).
\emph{Matematika, 12}(1), 37--47.
\url{https://www.mathnet.ru/eng/mat455}.

\bibitem[Humphries and Brumley(2019)]{HumphriesBrumley2019}
Humphries, P., \& Brumley, F. (2019).
Standard zero-free regions for Rankin--Selberg $L$-functions via sieve theory.
\emph{Mathematische Zeitschrift, 292}, 1105--1122.
\url{https://doi.org/10.1007/s00209-018-2136-8}.

\bibitem[Ivanyos et~al.(2012)]{IKRS2012}
Ivanyos, G., Karpinski, M., R\'onyai, L., \& Saxena, N. (2012).
Trading GRH for algebra: Algorithms for factoring polynomials and related structures.
\emph{Mathematics of Computation, 81}(277), 493--531.
\url{https://doi.org/10.1090/S0025-5718-2011-02505-6}.
Author pagination consulted:
\url{https://www.cse.iitk.ac.in/users/nitin/papers/nogrh.pdf}.

\bibitem[Jao et~al.(2009)]{JMV2009}
Jao, D., Miller, S.~D., \& Venkatesan, R. (2009).
Expander graphs based on GRH with an application to elliptic curve cryptography.
\emph{Journal of Number Theory, 129}(6), 1491--1504.
\url{https://doi.org/10.1016/j.jnt.2008.11.006}.
Author version consulted: \url{https://arxiv.org/pdf/0811.0647v2}.

\bibitem[Kazhdan and Patterson(1984)]{KP1984}
Kazhdan, D. A., \& Patterson, S. J. (1984).
Metaplectic forms.
\emph{Publications Math\'ematiques de l'IH\'ES, 59}, 35--142.
\url{https://doi.org/10.1007/BF02698770}.
Theorem~I.6.4(c), p.~104; Theorem~II.1.4, pp.~114--117;
Theorem~II.2.1 and the preceding paragraph, p.~118.

\bibitem[Koiran(1996a)]{Koiran1996}
Koiran, P. (1996a). Hilbert's Nullstellensatz is in the polynomial hierarchy.
\emph{Journal of Complexity, 12}(4), 273--286.
\url{https://doi.org/10.1006/jcom.1996.0019}.

\bibitem[Koiran(1996b)]{KoiranExtended1996}
Koiran, P. (1996b). \emph{Hilbert's Nullstellensatz is in the polynomial hierarchy}
(DIMACS Technical Report 96-27, extended version).
\url{https://perso.ens-lyon.fr/pascal.koiran/Publis/DIMACS.96-27.ps}.

\bibitem[Koiran(1997)]{KoiranErratum1997}
Koiran, P. (1997, August 8).
\emph{Erratum to Hilbert's Nullstellensatz is in the polynomial hierarchy}.
\url{https://perso.ens-lyon.fr/pascal.koiran/Publis/hn.erratum.ps}.

\bibitem[Kolountzakis and R\'ev\'esz(2006)]{KR2006Turan}
Kolountzakis, M. N., \& R\'ev\'esz, S. Gy. (2006).
Tur\'an's extremal problem for positive definite functions on groups.
\emph{Journal of the London Mathematical Society, 74}(2), 475--496.
\url{https://doi.org/10.1112/S0024610706023234}.
Author version:
\url{https://eigen-space.org/mk/ps/turan-groups.pdf}.

\bibitem[Lamzouri et~al.(2015)]{LLS2015}
Lamzouri, Y., Li, X., \& Soundararajan, K. (2015).
Conditional bounds for the least quadratic non-residue and related problems.
\emph{Mathematics of Computation, 84}(295), 2391--2412.
\url{https://doi.org/10.1090/S0025-5718-2015-02925-1}.
Corrected author version consulted: arXiv:1309.3595v3, December 2016,
\url{https://arxiv.org/pdf/1309.3595v3}.

\bibitem[Lemke Oliver and Thorner(2019)]{LemkeOliverThorner2019}
Lemke Oliver, R.~J., \& Thorner, J. (2019).
Effective log-free zero density estimates for automorphic \(L\)-functions
and the Sato--Tate conjecture.
\emph{International Mathematics Research Notices, 2019}(22), 6988--7036.
\url{https://doi.org/10.1093/imrn/rnx309}.
Author version consulted: 19~August~2017,
\url{https://lemkeoliver.github.io/papers/15-AutoLFZDE.pdf}.

\bibitem[Lenstra(1991)]{Lenstra1991}
Lenstra, H. W., Jr. (1991). Finding isomorphisms between finite fields.
\emph{Mathematics of Computation, 56}(193), 329--347.
\url{https://pub.math.leidenuniv.nl/~lenstrahw/PUBLICATIONS/1991b/art.pdf}

\bibitem[Miller(1976)]{Miller1976}
Miller, G.~L. (1976). Riemann's hypothesis and tests for primality.
\emph{Journal of Computer and System Sciences, 13}(3), 300--317.
\url{https://www.cs.cmu.edu/~glmiller/Publications/Papers/Mi76.pdf}.

\bibitem[Montgomery and Vaughan(2007)]{MontgomeryVaughan2007}
Montgomery, H.~L., \& Vaughan, R.~C. (2007).
\emph{Multiplicative number theory I: Classical theory}.
Cambridge University Press.
\url{https://doi.org/10.1017/CBO9780511618314}.

\bibitem[OpenAI(2026a)]{OpenAI2026}
OpenAI. (2026a, October 4).
\emph{Deterministic polynomial factorization over prime fields}
[Preprint].
\url{https://github.com/openai/math/tree/main/preprints/Deterministic-Polynomial-Factorization-over-Prime-Fields-October-4-2026}.
Theorem~1.2, p.~3; Proposition~10.1, pp.~41--44;
Lemma~11.2, p.~45.

\bibitem[OpenAI(2026b)]{OpenAIArtin2026}
OpenAI. (2026b, October~4).
\emph{\href{https://raw.githubusercontent.com/openai/math/main/preprints/Primitive-roots-for-every-admissible-integer-base-October-4-2026/primitive-roots-all-integer-bases.pdf}{Primitive roots for every admissible integer base}}
[Preprint].

\bibitem[PARI Group(2026)]{PARIPhi32026}
PARI Group. (2026).
\emph{polmodular.c source coverage snapshot: PARI/GP v2.19.0,
development 31092-e6893b0017} [Numerical coefficient data].
\url{https://pari.math.u-bordeaux.fr/lcov-report/basemath/polmodular.c.gcov.html}.
Snapshot dated 21 November 2026; $\Phi_3$ coefficient list before
\texttt{phi3\_ZV}, source lines 2081--2093.

\bibitem[Pila(1990)]{Pila1990}
Pila, J. (1990).
Frobenius maps of abelian varieties and finding roots of unity in finite fields.
\emph{Mathematics of Computation, 55}(192), 745--763.
\url{https://doi.org/10.1090/S0025-5718-1990-1035941-X}.
Checked author manuscript, \url{https://people.maths.ox.ac.uk/pila/Frobenius.pdf},
Theorems~A and D, pp.~1--2; \S5, pp.~14--18.

\bibitem[Rubinstein and Sarnak(1994)]{RubinsteinSarnak1994}
Rubinstein, M., \& Sarnak, P. (1994). Chebyshev's bias.
\emph{Experimental Mathematics, 3}(3), 173--197.
\url{https://www.math.uwaterloo.ca/~mrubinst/publications/Chebyshev.pdf}.

\bibitem[Schoof(1985)]{Schoof1985}
Schoof, R. (1985).
Elliptic curves over finite fields and the computation of square roots mod $p$.
\emph{Mathematics of Computation, 44}(170), 483--494.
\url{https://doi.org/10.1090/S0025-5718-1985-0777280-6}.

\bibitem[Schoof(2008)]{Schoof2008}
Schoof, R. (2008). Four primality testing algorithms.
In J.~P. Buhler \& P. Stevenhagen (Eds.),
\emph{Algorithmic number theory: Lattices, number fields, curves and cryptography}
(MSRI Publications, Vol.~44, pp.~101--126). Cambridge University Press.
\url{https://pub.math.leidenuniv.nl/~stevenhagenp/ANTproc/05rene.pdf}.

\bibitem[Serre(1981)]{Serre1981Chebotarev}
Serre, J.-P. (1981).
Quelques applications du th\'eor\`eme de densit\'e de Chebotarev.
\emph{Publications Math\'ematiques de l'IH\'ES, 54}, 123--201.
\url{https://doi.org/10.1007/BF02698692}.

\bibitem[Shoup(1992)]{Shoup1992}
Shoup, V. (1992). Searching for primitive roots in finite fields.
\emph{Mathematics of Computation, 58}(197), 369--380.
\url{https://shoup.net/papers/primroots.pdf}.

\bibitem[Silverman(2009)]{Silverman2009}
Silverman, J. H. (2009).
\emph{The arithmetic of elliptic curves} (2nd ed.).
Graduate Texts in Mathematics, 106. Springer.
\url{https://doi.org/10.1007/978-0-387-09494-6}.

\bibitem[Silvester et~al.(2006)]{SilvesterSpearmanWilliams2006}
Silvester, A. K., Spearman, B. K., \& Williams, K. S. (2006).
Cyclic cubic fields of given conductor and given index.
\emph{Canadian Mathematical Bulletin, 49}(3), 472--480.
\url{https://people.math.carleton.ca/~williams/papers/pdf/292.pdf}.
The cyclic-cubic conductor and discriminant facts are on p.~472.

\bibitem[Solovay and Strassen(1977)]{SolovayStrassen1977}
Solovay, R., \& Strassen, V. (1977). A fast Monte-Carlo test for primality.
\emph{SIAM Journal on Computing, 6}(1), 84--85.
\url{https://doi.org/10.1137/0206006}.
Correction (1978), \emph{7}(1), 118,
\url{https://doi.org/10.1137/0207009}.

\bibitem[Spaenlehauer(2026)]{Spaenlehauer2026}
Spaenlehauer, P.-J. (2026).
\emph{Deterministic root-finding over finite fields with elliptic curves}
(arXiv:2610.32508v1).
\url{https://arxiv.org/abs/2610.32508v1}.
Theorem~1.3, p.~2; Conjecture~1.5 and Problem~1.7, p.~3;
Proposition~6.4, p.~35 (printed pagination).

\bibitem[Sutherland(2011)]{Sutherland2011}
Sutherland, A.~V. (2011).
Computing Hilbert class polynomials with the Chinese remainder theorem.
\emph{Mathematics of Computation, 80}(273), 501--538.
\url{https://doi.org/10.1090/S0025-5718-2010-02373-7}.
Corrected author version consulted: arXiv:0903.2785v4, November 2013,
\url{https://arxiv.org/pdf/0903.2785v4}.

\bibitem[Sutherland(2012)]{Sutherland2012CM}
Sutherland, A. V. (2012).
Accelerating the CM method.
\emph{LMS Journal of Computation and Mathematics, 15}, 172--204.
\url{https://doi.org/10.1112/S1461157012001015}.
Checked author manuscript arXiv:1009.1082v4, 3 September 2012,
\S2.5, pp.~5--6; \S2.4, pp.~4--5 distinguishes its Dedekind-GRH
assumption for the CRT algorithm.

\bibitem[Sutherland(2013)]{Sutherland2013Volcanoes}
Sutherland, A.~V. (2013). Isogeny volcanoes.
In E.~W. Howe \& K.~S. Kedlaya (Eds.),
\emph{Algorithmic number theory: Tenth international symposium}
(The Open Book Series, Vol.~1, pp.~507--530). Mathematical Sciences Publishers.
\url{https://doi.org/10.2140/obs.2013.1.507}.
\url{https://msp.org/obs/2013/1-1/obs-v1-n1-p25-s.pdf}.

\bibitem[Tang(2026a)]{Tang2026}
Tang, Q. (2026a).
\emph{Deterministic polynomial-time factorization over finite fields}
(Version 2). Hexagon, 2611.00135v2.
\url{https://hexagonmath.org/2611.00135v2}.
Theorem~1.1 and Remarks~1.2--1.3, p.~2; Lemma~2.1, pp.~2--4;
Proposition~3.1, pp.~4--5; Proposition~4.1 and the final proof,
pp.~5--7.

\bibitem[Tang(2026b)]{TangCTReview2026}
Tang, Q. (2026b).
\emph{Cureton--Thorner v2 review}
[Review repository and separate revised manuscript].
\url{https://github.com/tqyqiqi1509/cureton-thorner-v2-review}.
Replacement proof of the residual norm estimate in Lemma~8.1.

\bibitem[Thorner and Zaman(2019)]{ThornerZaman2019}
Thorner, J., \& Zaman, A. (2019).
A unified and improved Chebotarev density theorem.
\emph{Algebra \& Number Theory, 13}(5), 1039--1068.
\url{https://doi.org/10.2140/ant.2019.13.1039}.
Publisher version:
\url{https://msp.org/ant/2019/13-5/ant-v13-n5-p02-p.pdf}.

\bibitem[Waterhouse(1969)]{Waterhouse1969}
Waterhouse, W. C. (1969). Abelian varieties over finite fields.
\emph{Annales scientifiques de l'\'Ecole Normale Sup\'erieure, 2}(4),
521--560. \url{https://doi.org/10.24033/asens.1183}

\bibitem[Zaman(2018)]{Zaman2018}
Zaman, A. (2018). The least unramified prime which does not split completely.
\emph{Forum Mathematicum, 30}(3), 651--661.
\url{https://doi.org/10.1515/forum-2017-0081}.
Author version consulted: arXiv:1704.03451v1,
\url{https://arxiv.org/pdf/1704.03451v1}.

\bibitem[Milne(2020)]{MilneCM2020}
Milne, J.~S. (2020, July~14).
\emph{Complex multiplication} (Version~0.10) [Author's notes].
\url{https://www.jmilne.org/math/CourseNotes/CM.pdf}.

\bibitem[Pazuki(2012)]{Pazuki2012}
Pazuki, F. (2012).
Theta height and Faltings height.
\emph{Bulletin de la Soci\'et\'e Math\'ematique de France, 140}(1), 19--49.
\url{https://doi.org/10.24033/bsmf.2623}.
\url{https://www.numdam.org/item/10.24033/bsmf.2623.pdf}.

\bibitem[Colmez(1993)]{Colmez1993}
Colmez, P. (1993).
P\'eriodes des vari\'et\'es ab\'eliennes \`a multiplication complexe.
\emph{Annals of Mathematics, 138}(3), 625--683.
\url{https://doi.org/10.2307/2946559}.
Author-hosted scan: \url{https://webusers.imj-prg.fr/~pierre.colmez/artin.pdf}.

\bibitem[Wei(2023)]{Wei2021}
Wei, X. (2023). Smallness of Faltings heights of CM abelian varieties.
\emph{Journal of Number Theory, 242}, 73--101.
\url{https://doi.org/10.1016/j.jnt.2022.07.007}.
Author version consulted: arXiv:2111.00617v1 (2021),
\url{https://arxiv.org/pdf/2111.00617v1}.

\bibitem[Kannan et~al.(1988)]{KannanLenstraLovasz1988}
Kannan, R., Lenstra, A.~K., \& Lov\'asz, L. (1988).
Polynomial factorization and nonrandomness of bits of algebraic and
some transcendental numbers.
\emph{Mathematics of Computation, 50}(181), 235--250.
\url{https://www.jstor.org/stable/2007927}.
Original article in the consulted course-reader PDF, pp.~7--22:
\url{https://www.math.cmu.edu/users/af1p/Teaching/AdditiveCombinatorics/LLLL.pdf}.

\bibitem[Lenstra(1983)]{Lenstra1983}
Lenstra, A.~K. (1983).
Factoring polynomials over algebraic number fields.
In J.~A. van Hulzen (Ed.), \emph{Computer algebra: EUROCAL '83}
(Lecture Notes in Computer Science, Vol.~162, pp.~245--254). Springer.
\url{https://ir.cwi.nl/pub/2541/2541D.pdf}.

\bibitem[Mumford(1966)]{Mumford1966}
Mumford, D. (1966). On the equations defining abelian varieties. I.
\emph{Inventiones Mathematicae, 1}, 287--354.
\url{https://doi.org/10.1007/BF01389737}

\bibitem[Mumford(1967)]{Mumford1967}
Mumford, D. (1967). On the equations defining abelian varieties. II.
\emph{Inventiones Mathematicae, 3}, 75--135.
\url{https://doi.org/10.1007/BF01389741}

\bibitem[Arene et~al.(2012)]{AreneKohelRitzenthaler2012}
Arene, C., Kohel, D., \& Ritzenthaler, C. (2012).
Complete addition laws on abelian varieties.
\emph{LMS Journal of Computation and Mathematics, 15}, 308--316.
\url{https://doi.org/10.1112/S1461157012001027}


\bibitem[Serre and Tate(1968)]{SerreTate1968}
Serre, J.-P., \& Tate, J. (1968). Good reduction of abelian varieties.
\emph{Annals of Mathematics, 88}(3), 492--517.
\url{https://doi.org/10.2307/1970722}.

\bibitem[Milne(1992)]{Milne1992}
Milne, J. S. (1992). The points on a Shimura variety modulo a prime of
good reduction. In R. P. Langlands \& D. Ramakrishnan (Eds.),
\emph{The zeta function of Picard modular surfaces} (pp.~151--253).
Les Publications CRM. Author manuscript:
\url{https://jmilne.org/math/articles/1992aP.pdf}.


\bibitem[von K\"anel and Kret(2023)]{VonKanelKret2023}
von K\"anel, R., \& Kret, A. (2023).
\emph{Integral points on coarse Hilbert moduli schemes}
[Preprint, arXiv:2307.06944v1].
\url{https://arxiv.org/pdf/2307.06944}.


\bibitem[Sze(2011)]{Sze2011}
Sze, T.-W. (2011).
\emph{On taking r-th roots without r-th nonresidues over finite
fields and its applications} (arXiv:1105.5852v1).
\url{https://arxiv.org/abs/1105.5852v1}.
The arXiv header is dated 30 May 2011; the retrieved version's
internal title page says ``Preliminary version, November 21, 2018.''


\bibitem[Milne(1999)]{Milne1999}
Milne, J. S. (1999). Lefschetz motives and the Tate conjecture.
\emph{Compositio Mathematica}, 117, 45--76.
\url{https://doi.org/10.1023/A:1000776613765}.

\bibitem[Mumford(1970)]{MumfordAV1970}
Mumford, D. (1970). \emph{Abelian varieties}.
Oxford University Press. Chapter~II, Section~6, Application~1,
pp.~60--61.


\bibitem[Yu(2016)]{YuAmbiguous2016}
Yu, C.-F. (2016). \emph{A remark on Chevalley's ambiguous class
number formulas} (arXiv:1412.1458v2).
\url{https://arxiv.org/abs/1412.1458v2}.
The stable version is dated 18 May 2016; its internal manuscript
also displays 18 October 2021.

\end{thebibliography}
\end{document}
